%%% Hlavní soubor. Zde se definují základní parametry a odkazuje se na ostatní části. %%%

%% Verze pro jednostranný tisk:
% Okraje: levý 40mm, pravý 25mm, horní a dolní 25mm
% (ale pozor, LaTeX si sám přidává 1in)
\documentclass[12pt,a4paper]{report}
\setlength\textwidth{145mm}
\setlength\textheight{247mm}
\setlength\oddsidemargin{15mm}
\setlength\evensidemargin{15mm}
\setlength\topmargin{0mm}
\setlength\headsep{0mm}
\setlength\headheight{0mm}
% \openright zařídí, aby následující text začínal na pravé straně knihy
\let\openright=\clearpage

%% Pokud tiskneme oboustranně:
% \documentclass[12pt,a4paper,twoside,openright]{report}
% \setlength\textwidth{145mm}
% \setlength\textheight{247mm}
% \setlength\oddsidemargin{15mm}
% \setlength\evensidemargin{0mm}
% \setlength\topmargin{0mm}
% \setlength\headsep{0mm}
% \setlength\headheight{0mm}
% \let\openright=\cleardoublepage

%% Pokud používáte csLaTeX (doporučeno):
%\usepackage{czech}
% Pokud nikoliv:
\usepackage[czech]{babel}
\usepackage[T1]{fontenc}
\usepackage[scaled]{beramono}

%% Použité kódování znaků: obvykle latin2, cp1250 nebo utf8:
\usepackage[utf8]{inputenc}

%% Ostatní balíčky
\usepackage{fancyvrb}
\usepackage{color}
\usepackage{relsize}
\usepackage{listings}
\usepackage{tikz}
\usepackage{todonotes}
\usepackage{graphicx}
\usepackage{amsthm}
\usepackage{amsmath}
\usepackage{amsfonts}
\usepackage{epstopdf}
\usepackage{bm}
\usepackage{algpseudocode}
\usepackage{algorithm}
\usepackage{cite}
\usepackage{setspace}
\usetikzlibrary{bayesnet}
\usetikzlibrary{arrows}

%% Balíček hyperref, kterým jdou vyrábět klikací odkazy v PDF,
%% ale hlavně ho používáme k uložení metadat do PDF (včetně obsahu).
%% POZOR, nezapomeňte vyplnit jméno práce a autora.
\usepackage[unicode]{hyperref}   % Musí být za všemi ostatními balíčky
\hypersetup{pdftitle=Název práce}
\hypersetup{pdfauthor=David Marek}
%\usepackage{showlabels}
%\showlabels{hypertarget}

%%% Drobné úpravy stylu

% Tato makra přesvědčují mírně ošklivým trikem LaTeX, aby hlavičky kapitol
% sázel příčetněji a nevynechával nad nimi spoustu místa. Směle ignorujte.
\makeatletter
\def\@makechapterhead#1{
  {\parindent \z@ \raggedright \normalfont
   \Huge\bfseries \thechapter. #1
   \par\nobreak
   \vskip 20\p@
}}
\def\@makeschapterhead#1{
  {\parindent \z@ \raggedright \normalfont
   \Huge\bfseries #1
   \par\nobreak
   \vskip 20\p@
}}
\makeatother

\makeatletter
\def\PY@reset{\let\PY@it=\relax \let\PY@bf=\relax%
    \let\PY@ul=\relax \let\PY@tc=\relax%
    \let\PY@bc=\relax \let\PY@ff=\relax}
\def\PY@tok#1{\csname PY@tok@#1\endcsname}
\def\PY@toks#1+{\ifx\relax#1\empty\else%
    \PY@tok{#1}\expandafter\PY@toks\fi}
\def\PY@do#1{\PY@bc{\PY@tc{\PY@ul{%
    \PY@it{\PY@bf{\PY@ff{#1}}}}}}}
\def\PY#1#2{\PY@reset\PY@toks#1+\relax+\PY@do{#2}}

\expandafter\def\csname PY@tok@gu\endcsname{\let\PY@bf=\textbf}
\expandafter\def\csname PY@tok@gs\endcsname{\let\PY@bf=\textbf}
\expandafter\def\csname PY@tok@cm\endcsname{\let\PY@it=\textit}
\expandafter\def\csname PY@tok@gp\endcsname{\let\PY@bf=\textbf}
\expandafter\def\csname PY@tok@ge\endcsname{\let\PY@it=\textit}
\expandafter\def\csname PY@tok@cs\endcsname{\let\PY@it=\textit}
\expandafter\def\csname PY@tok@gh\endcsname{\let\PY@bf=\textbf}
\expandafter\def\csname PY@tok@ni\endcsname{\let\PY@bf=\textbf}
\expandafter\def\csname PY@tok@nn\endcsname{\let\PY@bf=\textbf}
\expandafter\def\csname PY@tok@s2\endcsname{\let\PY@it=\textit}
\expandafter\def\csname PY@tok@s1\endcsname{\let\PY@it=\textit}
\expandafter\def\csname PY@tok@nc\endcsname{\let\PY@bf=\textbf}
\expandafter\def\csname PY@tok@ne\endcsname{\let\PY@bf=\textbf}
\expandafter\def\csname PY@tok@si\endcsname{\let\PY@bf=\textbf\let\PY@it=\textit}
\expandafter\def\csname PY@tok@nt\endcsname{\let\PY@bf=\textbf}
\expandafter\def\csname PY@tok@ow\endcsname{\let\PY@bf=\textbf}
\expandafter\def\csname PY@tok@c1\endcsname{\let\PY@it=\textit}
\expandafter\def\csname PY@tok@kc\endcsname{\let\PY@bf=\textbf}
\expandafter\def\csname PY@tok@c\endcsname{\let\PY@it=\textit}
\expandafter\def\csname PY@tok@sx\endcsname{\let\PY@it=\textit}
\expandafter\def\csname PY@tok@err\endcsname{\def\PY@bc##1{\setlength{\fboxsep}{0pt}\fcolorbox[rgb]{1.00,0.00,0.00}{1,1,1}{\strut ##1}}}
\expandafter\def\csname PY@tok@kd\endcsname{\let\PY@bf=\textbf}
\expandafter\def\csname PY@tok@ss\endcsname{\let\PY@it=\textit}
\expandafter\def\csname PY@tok@sr\endcsname{\let\PY@it=\textit}
\expandafter\def\csname PY@tok@k\endcsname{\let\PY@bf=\textbf}
\expandafter\def\csname PY@tok@kn\endcsname{\let\PY@bf=\textbf}
\expandafter\def\csname PY@tok@kr\endcsname{\let\PY@bf=\textbf}
\expandafter\def\csname PY@tok@s\endcsname{\let\PY@it=\textit}
\expandafter\def\csname PY@tok@sh\endcsname{\let\PY@it=\textit}
\expandafter\def\csname PY@tok@sc\endcsname{\let\PY@it=\textit}
\expandafter\def\csname PY@tok@sb\endcsname{\let\PY@it=\textit}
\expandafter\def\csname PY@tok@se\endcsname{\let\PY@bf=\textbf\let\PY@it=\textit}
\expandafter\def\csname PY@tok@sd\endcsname{\let\PY@it=\textit}

\def\PYZbs{\char`\\}
\def\PYZus{\char`\_}
\def\PYZob{\char`\{}
\def\PYZcb{\char`\}}
\def\PYZca{\char`\^}
\def\PYZam{\char`\&}
\def\PYZlt{\char`\<}
\def\PYZgt{\char`\>}
\def\PYZsh{\char`\#}
\def\PYZpc{\char`\%}
\def\PYZdl{\char`\$}
\def\PYZhy{\char`\-}
\def\PYZsq{\char`\'}
\def\PYZdq{\char`\"}
\def\PYZti{\char`\~}
% for compatibility with earlier versions
\def\PYZat{@}
\def\PYZlb{[}
\def\PYZrb{]}
\makeatother

\makeatletter
% code by Andrew Stacey: http://tex.stackexchange.com/a/27996/3954
\tikzset{
  through point/.style={
    to path={%
      \pgfextra{%
        \tikz@scan@one@point\pgfutil@firstofone(\tikztostart)\relax
        \pgfmathsetmacro{\pt@sx}{\pgf@x * 0.03514598035}%
        \pgfmathsetmacro{\pt@sy}{\pgf@y * 0.03514598035}%
        \tikz@scan@one@point\pgfutil@firstofone#1\relax
        \pgfmathsetmacro{\pt@ax}{\pgf@x * 0.03514598035 - \pt@sx}%
        \pgfmathsetmacro{\pt@ay}{\pgf@y * 0.03514598035 - \pt@sy}%
        \tikz@scan@one@point\pgfutil@firstofone(\tikztotarget)\relax
        \pgfmathsetmacro{\pt@ex}{\pgf@x * 0.03514598035 - \pt@sx}%
        \pgfmathsetmacro{\pt@ey}{\pgf@y * 0.03514598035 - \pt@sy}%
        \pgfmathsetmacro{\pt@len}{\pt@ex * \pt@ex + \pt@ey * \pt@ey}%
        \pgfmathsetmacro{\pt@t}{(\pt@ax * \pt@ex + \pt@ay * \pt@ey)/\pt@len}%
        \pgfmathsetmacro{\pt@t}{(\pt@t > .5 ? 1 - \pt@t : \pt@t)}%
        \pgfmathsetmacro{\pt@h}{(\pt@ax * \pt@ey - \pt@ay * \pt@ex)/\pt@len}%
        \pgfmathsetmacro{\pt@y}{\pt@h/(3 * \pt@t * (1 - \pt@t))}%
      }
      (\tikztostart) .. controls +(\pt@t * \pt@ex + \pt@y * \pt@ey, \pt@t * \pt@ey - \pt@y * \pt@ex) and +(-\pt@t * \pt@ex + \pt@y * \pt@ey, -\pt@t * \pt@ey - \pt@y * \pt@ex) .. (\tikztotarget)
    }
  }
}
\makeatother

% Toto makro definuje kapitolu, která není očíslovaná, ale je uvedena v obsahu.
\def\chapwithtoc#1{
\chapter*{#1}
\addcontentsline{toc}{chapter}{#1}
}

\renewcommand{\vec}[1]{\bm{#1}}
\hyphenation{prav-dě-po-do-bno-sti}
\hyphenation{prav-dě-po-do-bno-stní}
\hyphenation{prav-dě-po-do-bno-stních}
\hyphenation{pra-vdě-po-do-bnost-ní-mi}
\hyphenation{roz-dě-le-ní}
\hyphenation{sdru-že-né}
\hyphenation{dis-krét-ní-mi}
\hyphenation{pro-měn-nou}
\hyphenation{pro-měn-né}
\hyphenation{při-řa-ze-ní-mi}
\hyphenation{po-zo-ro-va-ný-mi}
\hyphenation{hod-not}
\hyphenation{mar-gi-nál-ních}
\hyphenation{mar-gi-na-li-zo-vat}
\hyphenation{pod-mí-ně-ných}
\hyphenation{Bayes-ov-ská}

\floatname{algorithm}{Algoritmus}
\algblockdefx[INIT]{Init}{EndInit}{\textbf{init}}{\textbf{end init}}

\newfloat{example}{thp}{loe}
\floatname{example}{Příklad}
\fvset{fontsize=\small}

\newtheorem{definice}{Definice}
\newtheorem{theorem}{Věta}

\usepackage{mysphinx}

\begin{document}

\lstset{
  language=Python,
  showstringspaces=false,
  formfeed=\newpage,
  tabsize=4,
  commentstyle=\itshape,
  basicstyle=\footnotesize\ttfamily,
  morekeywords={lambda,>>>,...},
  frame=single,
}

% Trochu volnější nastavení dělení slov, než je default.
\lefthyphenmin=2
\righthyphenmin=2

%%% Titulní strana práce

\pagestyle{empty}
\begin{center}

\large

Univerzita Karlova v~Praze

\medskip

Matematicko-fyzikální fakulta

\vfill

{\bf\Large DIPLOMOVÁ PRÁCE}

\vfill

\centerline{\mbox{\includegraphics[width=60mm]{../img/logo.eps}}}

\vfill
\vspace{5mm}

{\LARGE David Marek}

\vspace{15mm}

% Název práce přesně podle zadání
{\LARGE\bfseries Implementace aproximativních Bayesovských metod pro odhad stavu v~dialogových systémech}

\vfill

% Název katedry nebo ústavu, kde byla práce oficiálně zadána
% (dle Organizační struktury MFF UK)
Ústav formální a aplikované lingvistiky

\vfill

\begin{tabular}{rl}

Vedoucí diplomové práce: & Ing. Mgr. Filip Jurčíček, Ph.D. \\
\noalign{\vspace{2mm}}
Studijní program: & Informatika \\
\noalign{\vspace{2mm}}
Studijní obor: & Teoretická Informatika \\
\end{tabular}

\vfill

% Zde doplňte rok
2013

\end{center}

\newpage

%%% Následuje vevázaný list -- kopie podepsaného "Zadání diplomové práce".
%%% Toto zadání NENÍ součástí elektronické verze práce, nescanovat.

%%% Na tomto místě mohou být napsána případná poděkování (vedoucímu práce,
%%% konzultantovi, tomu, kdo zapůjčil software, literaturu apod.)

\openright

\noindent
Chtěl bych poděkovat svému vedoucímu Ing. Mgr. Filipovi Jurčíčkovi, Ph.D. za cenné rady a čas, který mi věnoval. 
Dále bych chtěl poděkovat svým kolegům Lukáši Žilkovi, Matěji Korvasovi, Ondřeji Plátkovi a Ondřeji Duškovi, s~kterými jsem mohl spolupracovat na projektu Vystadial.
Nakonec bych chtěl poděkovat svým nejbližším za podporu, kterou mi poskytli po celou dobu studia.

Tato práce byla částečně financována z~grantu Ministerstva školství, mládeže a tělovýchovy pod identifikačním kódem LK11221, Vystadial~-- Vývoj metod pro návrh statistických mluvených dialogových systémů.

\newpage

%%% Strana s čestným prohlášením k diplomové práci

\vglue 0pt plus 1fill

\noindent
Prohlašuji, že jsem tuto diplomovou práci vypracoval(a) samostatně a výhradně
s~použitím citovaných pramenů, literatury a dalších odborných zdrojů.

\medskip\noindent
Beru na~vědomí, že se na moji práci vztahují práva a povinnosti vyplývající
ze zákona č. 121/2000 Sb., autorského zákona v~platném znění, zejména skutečnost,
že Univerzita Karlova v~Praze má právo na~uzavření licenční smlouvy o~užití této
práce jako školního díla podle §60 odst. 1 autorského zákona.

\vspace{10mm}

\hbox{\hbox to 0.5\hsize{%
V~........ dne ............
\hss}\hbox to 0.5\hsize{%
Podpis autora
\hss}}

\vspace{20mm}
\newpage

%%% Povinná informační strana diplomové práce

\vbox to 0.5\vsize{
\setlength\parindent{0mm}
\setlength\parskip{5mm}

\textbf{Název práce:}
Implementace aproximativních Bayesovských metod pro odhad stavu v~dialogových systémech
% přesně dle zadání

\textbf{Autor}:
David Marek

\textbf{Ústav}:  % Případně Ústav:
Ústav formální a aplikované lingvistiky
% dle Organizační struktury MFF UK

\textbf{Vedoucí diplomové práce}:
Ing. Mgr. Filip Jurčíček, Ph.D., Ústav formální a aplikované lingvistiky
% dle Organizační struktury MFF UK, případně plný název pracoviště mimo MFF UK

\textbf{Abstrakt}:
% abstrakt v rozsahu 80-200 slov; nejedná se však o opis zadání diplomové práce
Každý dialogový systém obsahuje komponentu řízení dialogu, která vykonává akce dialogového systému na základě dialogového stavu.
Dialogový stav reprezentuje všechny dostupné informace.
Systém pro odhad stavu slouží k~udržování stavu dialogu a integraci nových informací ze systému porozumění.
Pro efektivní práci komponenty řízení dialogu je důležité co nejpřesněji modelovat pravděpodobnostní rozdělení definované nad všemi stavy dialogu. 
Těchto stavů může být velké množství, proto se většinou používají aproximativní techniky.
V~této práci se budeme zabývat implementací aproximativních metod pro odhad stavu dialogu pomocí inference v~Bayesovských sítích.
Výsledkem je knihovna použitelná v~systému pro odhad stavu reálného dialogového systému.


\textbf{Klíčová slova:} dialogový systém, odhad dialogového stavu, aproximativní Bayesovské metody
% 3 až 5 klíčových slov

\vss}\nobreak\vbox to 0.49\vsize{
\setlength\parindent{0mm}
\setlength\parskip{5mm}

\textbf{Title:} Approximative Bayes methods for belief monitoring in spoken dialogue systems
% přesný překlad názvu práce v angličtině

\textbf{Author:}
David Marek

\textbf{Department:}
Institute of Formal and Applied Linguistics
% dle Organizační struktury MFF UK v angličtině

\textbf{Supervisor:}
Ing. Mgr. Filip Jurčíček, Ph.D., Institute of Formal and Applied Linguistics
% dle Organizační struktury MFF UK, případně plný název pracoviště
% mimo MFF UK v angličtině

\textbf{Abstract:}
% abstrakt v rozsahu 80-200 slov v angličtině; nejedná se však o překlad
% zadání diplomové práce
The most important component of virtually any dialog system is a dialogue manager.
The aim of the dialog manager is to propose an action (a continuation of the dialogue) given the last dialog state.
The dialog state summarises all the past user input and the system input and ideally it includes all information necessary for natural progress in the dialog.
For the dialog manager to work efficiently, it is important to model the probability distribution over all dialog states as precisely as possible.
It is possible that the set of dialog states will be very large, so approximative methods usually must be used.
In this thesis we will discuss an implementation of approximate Bayes methods for belief state monitoring.
The result is a library for dialog state monitoring in real dialog systems.

\textbf{Keywords:} dialog system, belief state monitoring, approximative Bayes methods
% 3 až 5 klíčových slov v angličtině

\vss}

\newpage

%%% Strana s automaticky generovaným obsahem diplomové práce. U matematických
%%% prací je přípustné, aby seznam tabulek a zkratek, existují-li, byl umístěn
%%% na začátku práce, místo na jejím konci.

\openright
\pagestyle{plain}
\setcounter{page}{1}
\tableofcontents

%%% Jednotlivé kapitoly práce jsou pro přehlednost uloženy v samostatných souborech
\onehalfspacing
% !TEX root = prace.tex
\chapter*{Úvod}
\addcontentsline{toc}{chapter}{Úvod}

Dialog je přirozený způsob dorozumívání a sdělování informací mezi lidmi.
Počítač, který by dokázal vést dialog s uživatelem, byl vždy snem nejen
příznivců vědecko-fantastické literatury.
Už pro první počítače vnikaly programy, které se snažily využívat přirozenou
řeč pro interakci s uživatelem.
Jedním z takových programů byl například Eliza, program, který předstíral, že
jej zajímá, co mu uživatel říká.
Fungoval na principu rozpoznání textu pomocí gramatiky a následné transformace
textu do promluv dle pravidel.
Avšak gramatiky a pravidlové systémy se ukázaly nedostačné pro praktické
aplikace a tak se vývoj přesunul do statistických metod.
S využitím statistickým metod a metod strojového učení bylo možné začít s
porozumíváním mluveného slova.
Přijetí bylo zpočátku chladné a veřejnost byla 

\section{Cíle}

Dialogové systémy často využívají pouze nejlepší hypotézu ze systému porozumění přirozené řeči.
Většina ovšem umí vytvořit seznam $n$ nejlepších hypotéz.
Tato práce si klade za cíl představit metodu pro inferenci dialogového stavu v dialogovém systému s využitím více hypotéz.
Představená metoda bude založena na reprezentaci dialogového stavu pomocí dynamických bayesovských sítí.

Bude představen algoritmus Loopy Belief Propagation, který bude implementován pro využití v reálných systémech pro odhad dialogového stavu.
Algoritmus bude otestován na datech z reálného dialogového systému Let's Go a porovnán s dalšími systémy, které se účastnily soutěže Dialog State Tracking Challenge 2013.
Důležitou části algoritmů pro inferenci je určení pořadí, v jakém má inference probíhat.
Práce bude obsahovat implementaci několika strategií pro inferenci, které budou umožňovat efektivní inferenci pro různé druhy bayesovských sítí (stromy, dynamické sítě, obecné grafy).

Nakonec se práce bude zabývat učením parametrů sítě a představí algoritmus Expectation Propagation, který umožňuje inferenci v bayesovských sítích se spojitými náhodnými proměnnými.
Algoritmus Expectation Propagation je pouze konceptem, do kterého je třeba zasadit výpočet parametrů pravděpodobnostních rozložení z konkrétního grafického modelu.
Cílem práce není vytvořit obecný systém pro inferenci v libovolném grafickém modelu.
Takový cíl není ani reálný, protože při většině reálných použití dochází k aproximacím a úprávám v závislosti na ontologii problému, aby bylo možné inferenci provádět v reálném čase.
Cílem této práce tedy je vytvořit základní funkcionalitu potřebnou pro inferenci s EP, s kterou bude možné jednoduše vytvořit funkční systém, impolementovaný s doménově specifickými informacemi.
Jako příklad bude ukázán systém pro učení parametrů pravděpodobnostního rozložení pro pozorování v generickém modelu pro reprezentaci dialogového stavu.

% !TEX root = prace.tex
\chapter{Teorie dialogových systémů}

\section{Dialogový systém}

Dialogový systém je počítačový systém, který umožňuje uživatelům používat
hlasových pokynů pro ovládání počítačů a získávání informací. Dialogové systémy
nejsou pouhým hlasovým ovládáním počítače. Obsahují strategii řízení dialogu a
navíc obsahují interní stav dialogu, díky kterému jsou schopny vést dialog s
uživatelem. Dialog pak umožňuje využít složitějších příkazů, pokládat dotazy
nad strukturovanými daty a upřesňovat dotaz.

Dialogové systémy stále mají spoustu problémů k překonání a pro praktické použití je třeba se uchýlit k několika předpokladům a zjednodušením.
Prvním zjednodušením je doménová specializace, v současné době není možné vytvořit dialogový systém, který by se dokázal s uživatelem bavit o libovolném tématu. 
Vždy je potřeba při vývoji dialogového systému vědět, jaké informace má systém poskytovat a o čem se může chtít uživatel bavit.
Další zjednodušení se týkají přímo dialogu. 
Předpokládá se, že dialog probíhá vždy mezi systémem a jedním uživatelem. Navíc se tito pravidelně střídají. Jedna obrátka dialogu je složená z jedné promluvy systému a jedné promluvy uživatele.

Příkladem dialogového systému může být systém pro nalezení spojení pomocí
městské dopravy. Příkladem příkazu od uživatele pak může být např. \uv{Chci jet z
Malostranského náměstí na Anděl}. Dialogový systém se nyní může rozhodnout,
zda-li mu zadané informace stačí pro nalezení spojení. V tomto případě systém
stále neví, kdy chce uživatel jet. Může předpokládat, že uživatel už na
zastávce stojí a tak tedy uživateli nalezne nejbližší spojení.

Důležitou vlastností dialogového systému je robustnost. Pokud budeme používat
dialogový systém v přirozeném prostředí, tak se musíme vyrovnat s tím, že často
nebude uživateli rozumět. Uživatel může chtít jet stejně jako v minulém
případě, ale nyní dialogový systém přeslechne zastávku Anděl. Pro nalezení
spojení už nemá dost informací, může si pořád domyslet, že uživatel chce jet
nyní, ale je mnohem více možností, kterým směrem se vůbec chce uživatel vydat.
Dialogový systém se proto uživatele zeptá na chybějící informace a zároveň může
i implicitně potvrdit, že rozuměl alespoň odchozí zastávku: \uv{Dobře, jedete ze
zastávky Malostranské náměstí, můžete říct na jakou zastávku chcete jet?}.

\subsection{Typy dialogových systémů}

Dialogové systémy můžeme dělit podle několika kritérií. Podle toho, která
strana dialog vede anebo třeba jakým způsobem funguje strategie dialogového
systému.

\subsubsection{Dělení podle iniciativy}

Nejjednodušší dialogové systémy fungují na stejném principu jako automatické
telefonní systémy. Systém je ten, kdo vede dialog, a uživatel pouze odpovídá na
dotazy nebo vybírá z nabídnutých možností. Tento systém je velmi jednoduchý na
implementaci, pokud nutíme uživatel, aby vždy pouze odpověděl na položený
dotaz, tak velmi omezíme jeho vyjadřovací možnosti a díky tomu bude mnohem
jednodušší rozpoznat co řekl. Například pokud se uživatele ptáme jestli musí
nalezené spojení obsahovat pouze bezbariérové dopravní prostředky, tak můžeme
očekávat pouze \uv{Ano} nebo \uv{Ne} a jejich synonyma. Tento přístup ovšem
nebude velmi populární mezi častými uživateli systému, kteří by rádi zrychlili
postup dialogem.

Další možností je tedy povolit uživateli sdělit informace i tehdy, kdy se ho
systém ptá na něco jiného. Systém se tedy stále ptá na informace, které mu
chybějí, ale může dialog začít obecnější otázkou: "Jaké spojení chcete nalézt?"
oproti "Z jaké zastávky chcete jet?". V prvním případě dává systém pokročilému
uživateli možnost popsat celou cestu a velmi rychle se tak dostat ke kýžené
informaci. Pro nové uživatele, ale může být takováto obecně položená otázka
matoucí a tak je většinou třeba přidat i příklad.

Posledním typem je takový, kde uživateli přenecháme iniciativu, příkladem může
být například systém \uv{How may I help you?} \cite{gorin1997may} od
společnosti AT\&T. Ten slouží pro klasifikaci požadavků od uživatele a jejich
přepojení k patřičnému operátorovi.

\subsubsection{Dělení podle strategie řízení}

Důležité rozdělení dialogových systémů je podle způsobu, jakým funguje jejich
strategie řízení. Systém většinou předpokládá, že se v dialogu střídá s
uživatelem. Vždy, když na systém přijde řada, tak se musí rozhodnout, zda-li už
má dost informací k zodpovězení dotazu uživatele, případně zda-li se musí
dotázat na další informace, nechat si nějakou informaci potvrdit od uživatele,
anebo nechat uživatele vybrat si z nabízených možností.

Systém může mít přístup do databáze s informacemi a musí se tedy umět
rozhodnout, zda-li je třeba v ní hledat. Pak musí samozřejmě na nalezená data
reagovat. Může se stát, že je informací, které odpovídají zadání uživatele
příliš velké množství a tak je třeba zpřísnit zadání, anebo naopak nebyla žádná
informace nalezena a tedy je potřeba uživatele o tomto informovat.

Strategii, podle které se systém bude v těchto případech řídit lze
implementovat několika způsoby. Přímým řešením je napsat ji ručně. Pak se
rozhodnutí systému skládá z procházení série podmínek až do nalezení podmínky
odpovídající aktuálnímu stavu a následné provedení akce specifikované touto
podmínkou. Tento způsob může fungovat pro jednoduché systémy a pokud je vhodně
napsaný, tak může fungovat velmi dobře. Dost často ale tento způsob selhává,
pokud je dialogový systém příliš složitý, když se stávají podmínky
neudržitelné, anebo v případě, kdy se nemůžeme spolehnout na rozpoznanou řeč.

Další možností jak vytvořit strategii řízení je použít učení s učitelem.
Předpokládáme, že máme databázi hovorů mezi uživatelem a živým operátorem, z
kterých se dialogový systém může naučit, jak v dialogu postupovat. Tento
přístup se zdá jako dobrý nápad, ovšem naučíme se pouze napodobit operátora,
což nemusí být ten nejlepší způsob jak vést dialog. Navíc je nepřeberné
množství směrů, kterým se může dialog vyvíjet a jakmile se v jedné akci
odchýlíme, tak dostáváme úplně jiný dialog a tedy je velmi těžké z pozorovaných
dialogů zobecňovat. Proto se tento způsob téměř nevyužívá, existují ovšem
pokusy jak jej skloubit s následující metodou učení \cite{levin2000stochastic}.

Nejpoužívanější metodou v současnosti jak učit strategii dialogového systému
je použít zpětnovazební učení \cite{singh1999reinforcement}, \cite{walker2011application}. 
V takovém případě nám stačí pouze zhodnocení výsledku dialogu, které můžeme získat požádáním uživatelů dialogového systému, aby na konci dialogu oznámkovali, jak byli s výsledkem spokojeni. 
Při učení se nesnažíme napodobit existujícího operátora, ale získat co nejlepší hodnocení od uživatelů.
Díky tomu se může dialogový systém naučit pracovat dokonce ještě lépe než operátor.

\section{Součásti dialogového systému}

Dialogový systém se skládá z několika částí, každá se specializuje na jiný úkol a dohromady tvoří spolupracující systém.
Na vstupu je vždy zvukový záznam, o převedení do textu se stará systém rozpoznávání řeči (ASR).
Z textu je potřeba získat sémantické informace pomocí systému porozumění mluvené řeči (SLU).
Nad sémanticky anotovanými informacemi už může pracovat dialogový manager (DM), který rozhoduje o další akci.
Výstupem dialogového manageru jsou informace, které se mají předat uživateli.
O jejich převedení do textu se stará systém generování přirozené řeči (NLG).
Do zvukového záznamu převede text syntetizér řeči (TTS).

\subsection{Systém rozpoznávání řeči (ASR)}

Systém rozpoznávání řeči slouží k převedení mluveného projevu do textové podoby. 
Po získání textové podoby je teprve možné se zabývat významem textu.
Aktuálně nejlepší systémy jsou založené na pravděpodobnostním modelu a využívají skryté Markovské modely (HMM). 
Nejznámější je systém HTK \cite{young2002htk}, dalším otevřeným systémem je např. Kaldi \cite{Povey_ASRU2011}.
Existuje i celá řada komerčního software od firem jako IBM nebo Nuance.

Úspěšnost systému rozpoznání řeči je závislá na obtížnosti úlohy a na počtu trénovacích dat, pocházejících ze stejné domény.
Pro obecnou doménu se problém stává mnohem těžší.
Příkladem obecného rozpoznávače je Google Voice Search.

Systém rozpoznávání řeči může produkovat více hypotéz pro jeden vstup.
Často existuje pro jeden zvukový záznam více možných slovních sekvencí, z kterých by mohl pocházet.
Jejich reprezentace může být seznamem možných hypotéz s věrohodností pro každou hypotézu.
Věrohodnosti jsou skóre přiřazené hypotézám, které určují jakou důvěru má systém rozpoznávání řeči ve správnost dané hypotézy.
V ideálním případě je věrohodnost ekvivalentní aposteriorní pravděpodobnosti sekvence slov, dáno vstupní zvuk.
Ovšem ne všechny rozpoznávače pracují na pravděpodobnostním principu a pak není možné od nich požadovat skutečné pravděpodobnosti.

Další možností jak reprezentovat výstup je použití konfůzní sítě \cite{bertoldi2005new}. 
Konfůzní síť je vážený orientovaný graf obsahující startovní a konečný vrchol a hrany má označené slovy. 
Každá cesta ze startovního do konečného vrcholu vede přes všechny ostatní vrcholy. 
Váhy hran jsou pravděpodobnosti slova přiřazeného dané hraně. 
Hrany mohou obsahovat i prázdné slovo $\epsilon$.
Pravděpodobnost sekvence slov je součinem vah po cestě ze startovního do konečného uzlu.
Výhodou konfůzní sítě je, že umožňuje v komprimované podobě uložit mnohem více hypotéz.

\subsection{Porozumění mluvené řeči (SLU)}

Systém SLU převádí text do dialogových aktů. 
Dialogový akt (DA) je reprezentace promluvy, skládá se z jedné nebo více položek dialogového aktu (DAI), které jsou spojené v konjunkci.
Každá DAI se skládá z typu, názvu slotu a jeho hodnoty. 
Příklad DA z dialogového systému pro hledání restaurací:

\begin{center}
{\tt hello()\&inform(food="chinese")}.
\end{center}

Zde se DA skládá ze dvou položek, první položka má pouze typ {\em hello}, značící pozdrav. 
Druhá položka má typ {\em inform}, tzn. uživatel nás informuje o svém požadavku.
Název slotu je {\em food} a hodnota je \uv{chinese}, tedy uživatel nám říká, že hledá restauraci, kde servírují čínské jídlo.

Typů může být libovolné množství, ale existuje několik základních, jejichž použití je ustálené. 
\begin{itemize}
\item {\em inform} --- sdělujeme informaci, doplňujeme hodnotu do slotu,
\item {\em request} --- požadujeme od protějšku doplnění hodnoty pro dotazovaný slot,
\item {\em confirm} --- chceme potvrdit hodnotu slotu, potvrzení může být implicitní, anebo explicitní. 
	Při explicitním potvrzení očekáváme odpověď buď \uv{Ano} nebo \uv{Ne},
	U implicitního, pokud se nám nedostane odpovědi předpokládáme, že protějšek souhlasí,
\item {\em select} --- žádáme protějšek, aby zvolil z nabízených možností.
\end{itemize}

\subsection{Dialogový manager (DM)}

Dialogový manager tvoří mozek dialogového systému, jeho vstupem jsou dialogové akty z SLU, výstupem opět dialogové akty, pro NLG.
DM se dělí na dvě části. 
První část se stará o udržování dialogového stavu, tzn. informace, které uživatel systému poskytl, historii dialogu, které informace už byly poskytnuty anebo např. potvrzeny či zamítnuty uživatelem.
Dále obsahuje dialogovou strategii, která určuje příští akci v závislosti na dialogovém stavu.

\subsubsection{Reprezentace stavu dialogu}

Markovský rozhodovací proces (MDP) \cite{puterman2009markov} je diskrétní, stochastický a kontrolovaný proces. 
V každém časovém okamžiku se systém nachází v nějakém stavu $s$. 
Uživatel provede akci $a$, dostupnou ve stavu $s$, a ta jej přesune náhodně do nového stavu $s^\prime$ a navíc dostane odměnu $r(s, s^\prime, a)$.
Pravděpodobnost přechodu ze stavu $s$ do stavu $s^\prime$ je dána přechodovou funkcí $p(s^\prime \mid s, a)$.

Zobecněním MDP je částečně pozorovatelný Markovský rozhodovací proces (POMDP). 
Částečně pozorovatelný proto, že na rozdíl od MDP nevíme v jakém stavu se systém nachází.
Naše představa o stavu je dána pouze pravděpodobnostním rozložením přes všechny možné stavy, tzv. belief $b(s)$.
Výsledek provedení akce tedy nezávisí pouze na přechodové pravděpodobnosti, ale také na belief stavu. 
Pro aktualizaci stavu je potřeba vysčítat přes všechny možné stavy
\begin{equation}
b(s^\prime) = \sum_s p(s^\prime \mid s, a) b(s).
\end{equation}

Dialog je tedy popsán pomocí POMDP \cite{williams2007partially}, protože musíme zakomponovat naši neznalost cílů uživatele a problémy s rozpoznáváním. 
Pro reprezentaci stavu dialogu existuje několik různých přístupů.

\section{Grafické modely}


% !TEX root = prace.tex
\chapter{Bayesovské sítě a inference}
\label{ch:kap2}

V této kapitole představíme Bayesovské sítě, grafický model pro efektivní reprezentaci pravděpodobnostních rozdělení a nezávislostí mezi náhodnými proměnnými.
Bayesovská síť zároveň vytváří koncept pro inferenci, tedy zodpovídání dotazů nad proměnnými v síti.
Ukážeme si několik přístupů k inferenci, nejprve naivní výpočet vycházející přímo z definice.
Následně využijeme vlastností sítě a konceptů dynamického programování pro jeho zlepšení.
Analýzou složitosti exaktní inference dojdeme k závěru, že pro větší a komplexnější modely bude třeba se uchýlit k aproximacím.

Nejprve aproximujeme sdruženou pravděpodobnostní distribuci součinem marginálních distribucí a představíme Loopy Belief Propagation algoritmus.
Stále můžeme použít pouze diskrétní pravděpodobnostní distribuce se známými parametry.
Pro učení parametrů lze použít Expectation Maximization metodu, pro dialogové systémy je ovšem těžké získat dostatek učících dat.
Dostaneme se k metodě Expectation Propagation, která je zobecněním LBP a je možné ji použít pro libovolné rozdělení.
Díky ní budeme schopní vytvořit generativní model pro aktualizaci dialogového stavu, který bude pracovat stejně jako LBP, ale bude schopen adaptace za běhu dialogového systému.

\section{Bayesovské sítě}
\label{sec:bn}

Bayesovské sítě jsou pravděpodobnostní grafický model, který využívá podmíněných nezávislostí pro úspornou reprezentaci sdružené pravděpodobnosti.
Bayesovská sít je orientovaný acyklický graf, jeho vrcholy jsou náhodné proměnné a hrany odpovídají přímé závislosti jednoho uzlu na druhý.
Pro každou náhodnou proměnnou v síti platí, že její pravděpodobnost je jednoznačně určena jejími rodiči v grafu.
Podmíněná pravděpodobnostní distribuce (CPD) proměnné $X$ popisuje pravděpodobnost proměnné $X$ dáno její rodiče, $P(X \mid parents(X))$.
Pokud proměnná nemá žádná rodiče, pak její podmíněná pravděpodobnostní distribuce je ekvivalentní marginální pravděpodobnostní distribuci.

Příklad Student~\cite{koller2009probabilistic}: firma zvažuje, zda-li přijme studenta.
Firma chce přijímat chytré studenty, ale nesmí je testovat na inteligenci (I) přímo.
Má však výsledek studentových SAT testů, které ovšem nemusí stačit pro správné zhodnocení inteligence.
Požadují tak tedy i doporučení (D) od jednoho z učitelů.
Učitel studentovi napíše doporučující dopis na základě známky (Z), kterou student získal v jeho předmětu.
Předměty se ovšem liší v obtížnosti (O) a tak je studentova známka v předmětu závislá nejen na jeho inteligenci, ale také na obtížnosti předmětu.
Grafický model reprezentující tento problém je vyobrazen na obrázku~\ref{fig:student}.
\begin{figure}
\begin{center}
\begin{tikzpicture}

\matrix[row sep=0.75cm, column sep=1.2cm]
{
    \node[latent]       (O)     {O};
    && \node[latent]    (I)     {I};
    \\
    &\node[latent]      (Z)     {Z};
    &&\node[latent]     (S)     {SAT};
    \\
    &\node[latent]      (D)     {D};
    \\
};

\edge{O}{Z}
\edge{I}{Z}
\edge{I}{S}
\edge{Z}{D}

\end{tikzpicture}
\end{center}
\label{fig:student}
\caption{Bayesovská síť pro příklad se studentem.}
\end{figure}

V tomto modelu je několik nezávislostí. Obtížnost předmětu a inteligence studenta jsou zjevně nezávislé.
Studentova známka z předmětu je závislá na obtížnosti předmětu a inteligenci studenta, ale je podmíněně nezávislá na jeho výsledku ze SAT, dáno studentova inteligence.
Konečně doporučení, které student obdrží, je podmíněně nezávislé na všech ostatních proměnných, dáno studentova známka.

Sdruženou nezávislost tohoto modelu lze zapsat ve formě podmíněných pravděpodobnostních distribucí s pomocí řetízkového pravidla.
\begin{equation}
P(O, I, Z, S, D) = P(D \mid Z) P(Z \mid O, I) P(SAT \mid I) P(O) P(I)
\end{equation}

Předpokládejme, že obtížnost předmětu, inteligence studenta, doporučující dopis a výsledek SAT jsou binární proměnné.
Známka z předmětu pak je ternární proměnná.
Pokud bychom zapsali sdruženou pravděpodobnost ve formě tabulky, tak se dostaneme k 48 položkám.
Díky rozdělení do podmíněných pravděpodobnostních rozložení, které nám bayesovská síť poskytuje, se dostáváme k $2 + 2 + 12 + 4 + 6 = 26$ položkám.
Tedy i v tomto jednoduchém modelu dochází k značné úspoře.

\section{Inference v Bayesovských sítích}

Bayesovské sítě reprezentují pravděpodobnostní model a umožňují nám nad ním provádět dotazy.
Můžeme se například ptát na marginální pravděpodobnost jednotlivých proměnných. 
Tu získáme marginalizací sdružené pravděpodobnosti, pokud vezmeme příklad se studentem a budeme chtít znát marginální pravděpodobnost známek, musíme vysčítat všechny ostatní proměnné.

\begin{equation}
P(Z) = \sum_{O, I, S, D} P(D \mid Z) P(Z \mid O, I) P(SAT \mid I) P(O) P(I)
\end{equation}

Další a asi nejčastější dotaz nastává, pokud některé náhodné proměnné pozorujeme. 
Pak chceme znát pravděpodobnost jiných proměnných dáno naše pozorování, $P(\vec{X} \mid \vec{E} = \vec{e})$, kde $\vec{X}$ jsou dotazované proměnné, $\vec{E}$ jsou pozorované proměnné a $\vec{e}$ jsou pozorované hodnoty.
Z definice podmíněné pravděpodobnosti dostáváme

\begin{equation}
P(\vec{X} \mid \vec{E} = \vec{e}) = \frac{P(\vec{X}, \vec{E} = \vec{e})}{P(\vec{E} = \vec{e})}
\end{equation}

Každou instanci jmenovatele $P(\vec{X}, \vec{E} = \vec{e})$ jde spočítat sumou sdružených pravděpodobností s ohodnoceními proměnných, které jsou kompatibilní s pozorováním a aktuální instancí. 
Pokud počítáme instanci $P(\vec{X} = \vec{x}, \vec{E} = \vec{e})$, pak získáme výsledek marginalizací všech proměnných, kromě $\vec{X}$ a $\vec{E}$, které jsou fixovány.
Pokud tedy množinu všech proměnných bez $\vec{X}$ a $\vec{E}$ označíme $\mathcal{W} = \mathcal{X} - \vec{X} - \vec{E}$, pak pravděpodobnost dané instance je

\begin{equation}
P(\vec{X} = \vec{x}, \vec{E} = \vec{e}) = \sum_{\vec{w}} P(\vec{x}, \vec{e}, \vec{w})
\end{equation}

Pro výpočet normalizační konstanty $P(\vec{E})$ musíme opět marginalizovat sdruženou pravděpodobnost, anebo si můžeme povšimnout, že platí
\begin{equation}
P(\vec{E} = \vec{e}) = \sum_{\vec{x}} P(\vec{x}, \vec{e})
\end{equation}
a tedy můžeme použít už vypočítané hodnoty.

\subsection{Exaktní inference}

V předchozí části jsme viděli, že pomocí definice podmíněné pravděpodobnosti a marginalizace sdružené pravděpodobnosti lze najít odpověď na libovolný dotaz. 
Nyní si ukážeme algoritmus, který využívá struktury Bayesovské sítě pro inferenci a navíc díky metodám dynamického programování umožňuje samotný výpočet urychlit.
Nakonec ovšem zjistíme, že pro velké sítě, které nás většinou zajímají nejvíce, nám přesná inference nebude stačit a musíme se uchýlit k aproximacím.

Začneme s inferencí v jednoduchém modelu $A \rightarrow B \rightarrow C \rightarrow D$.
Sdružená pravděpodobnost $P(A, B, C, D)$ je součinem jednotlivých podmíněných pravděpodobnostních distribucí
\begin{equation}
P(A, B, C, D) = P(D \mid C) P(C \mid B) P(B \mid A) P(A)
\end{equation}

Pokud nyní budeme chtít spočítat marginální distribuci $D$, tak musíme marginalizovat všechny ostatní proměnné
\begin{equation}
P(D) = \sum_{A, B, C} P(D \mid C) P(C \mid B) P(B \mid A) P(A)
\end{equation}

Můžeme si povšimnout, že spousta členů se bude počítat vícekrát.
Využitím metod dynamického programování a přeuspořádáním sum si můžeme mezivýsledky uložit a použít vícekrát.
\begin{equation}
P(D) = \sum_C P(D \mid C) \sum_B P(C \mid B) \sum_A P(B \mid A) P(A)
\end{equation}

Při výpočtu pak nejprve spočítáme $\psi_1(A, B) =  P(B \mid A) P(A)$,  pak vysčítáme proměnnou $A$ a získáme $\tau_1(B) = \sum_A \psi_1(A, B)$.
Pokračujeme obdobně
\begin{align}
    \psi_2(B, C) &= P(C \mid B) \tau_1(B) \\
    \tau_2(C) &= \sum_B \psi_2(B, C)
\end{align}
A nakonec spočítáme finální marginální pravděpodobnost
\begin{align}
    \psi_3(D, C) &= P(D \mid C) \tau_2(C) \\
    P(D) &= \sum_C \psi_3(D, C)
\end{align}

Nechť $\mathcal{X}$ je množina náhodných proměnných. 
$Val(\mathcal{X})$ definujeme jako množinu možných ohodnocení proměnných $\mathcal{X}$.
Potom definujeme faktor $\phi$ jako zobrazení z $Val(\mathcal{X})$ do $\mathbb{R}$. Faktor je nezáporný, pokud všechny jeho obrazy jsou nezáporné. Množina proměnných $\mathcal{X}$ je doménou faktoru a značíme ji jako $Dom(\phi)$.

Faktor, jehož doménu tvoří diskrétní proměnné, si můžeme představit jako tabulku, která obsahuje jednu hodnotu pro každé možné ohodnocení proměnných z domény.

Nechť $\vec{X}, \vec{Y}, \vec{Z}$ jsou tři disjunktní množiny náhodných proměnných.
Nechť $\phi_1(\vec{X}, \vec{Y})$ a $\phi_2(\vec{Y}, \vec{Z})$ jsou faktory.
Definujeme součin faktorů $\phi_1 \times \phi_2$ jako faktor $\psi: Val(\vec{X}, \vec{Y}, \vec{Z}) \rightarrow \mathbb{R}$ následovně:
\begin{equation*}
\psi(\vec{X}, \vec{Y}, \vec{Z}) = \phi_1(\vec{X}, \vec{Y}) \cdot \phi_2(\vec{Y}, \vec{Z})
\end{equation*}

Násobíme prvky, které mají stejné ohodnocení společných proměnných $\vec{Y}$.
Stejný princip použijeme pro všechny matematické operace.

Nechť $\vec{X}$ je množina náhodných proměnných a $Y \not\in \vec{X}$ náhodná proměnná. 
Nechť $\phi(\vec{X}, Y)$ je faktor.
Definujeme marginalizaci $Y$ v $\phi$, značenou $\sum_Y \phi$, jako faktor $\psi$ s doménou $\vec{X}$ takový, že
\begin{equation*}
    \psi(\vec{X}) = \sum_Y \phi(\vec{X}, Y)
\end{equation*}
Této operaci také říkáme vysčítání $Y$ ve $\phi$.

Faktory se při počítání chovají jako čísla, všechny operace probíhají po prvcích, pro které je ohodnocení náhodných proměnných z průniku domén faktorů stejné.
Proto platí komutativita $\phi_1 \cdot \phi_2 = \phi_2 \cdot \phi_1$ a $\sum_X \sum_Y \phi = \sum_Y \sum_X \phi$.
Dále platí asociativita součinu $(\phi_1 \cdot \phi_2) \cdot \phi_3 = \phi_1 \cdot (\phi_2 \cdot \phi_3)$.
Nakonec můžeme vyměnit sumu a součin, pokud $X \not\in Dom(\phi_1)$, potom $\sum_X (\phi_1 \cdot \phi_2) = \phi_1 \sum_X \phi_2$.

Sdruženou pravděpodobnost z minulého příkladu tedy můžeme přepsat do formy faktorů.
\begin{equation}
    P(A, B, C, D) = \phi_A \cdot \phi_B \cdot \phi_C \cdot \phi_D
\end{equation}

Opět se pokusíme spočítat marginální pravděpodobnost proměnné $D$.
\begin{align}
P(D) &= \sum_C \sum_B \sum_A \phi_A \phi_B \phi_C
\\
&= \sum_C 
	\phi_D \cdot 
	\left( 
		\sum_B 
			\phi_C \cdot 
			\left( 
				\sum_A 
					\phi_A \cdot \phi_B
			\right)
	\right)
\end{align}

Přesuny sum můžeme provést díky doméně jednotlivých faktorů.
Faktory $\phi_C$ a $\phi_D$ neobsahují proměnnou $A$ a tedy je můžeme vytknout před sumu přes $A$.
Stejně tak faktor $\phi_D$ neobsahuje proměnnou $B$ a opět jej můžeme vytknout před sumu přes $B$.
Tyto úpravy můžeme provádět v libovolném pořadí, pokud vždy platí, že vysčítáme proměnnou $X$ až poté, co spolu vynásobíme všechny faktory, které ji obsahují.

V obecnosti vždy počítáme výraz, který je ve tvaru
\[
\sum_X \prod_{\phi \in \Phi} \phi.
\]

Z tohoto také vychází název pro tuto metodu: sum-product.
Jednoduchý algoritmus pro exaktní inferenci využívající tuto metodu se nazývá eliminace proměnných.
Základní myšlenka je, že máme dán seznam náhodných proměnných v pořadí, v jakém se mají eliminovat.
Pro eliminaci proměnné je třeba nejprve vynásobit všechny faktory, které ji obsahují a následně ji vysčítat.
Tak získáme faktor, který už tuto proměnnou neobsahuje, tedy jsme ji eliminovali.
Eliminace proměnných je popsána v algoritmu~\ref{alg:ve}.

\begin{algorithm}[H]
\caption{Eliminace proměnných}
\label{alg:ve}
\begin{algorithmic}
\Function{Sum-Product-VE}{$\Phi$, $\vec{X}$, $\prec$}
\State{$\Phi$ množina všech faktorů.} 
\State{$\vec{X}$ množina náhodných proměnných, které mají být eliminovány.}
\State{$\prec$ pořadí proměnných, v jakém mají být eliminovány.}
\State

\State Nechť $X_1, \dots, X_k$ je seřazení proměnných z $\vec{X}$, t.ž. $X_i \prec X_j \Leftrightarrow i < j$.
\For{$i = 1 \dots k$}
	\State $\Phi \gets$ Sum-Product-Eliminate-Var($\Phi$, $Z_i$)
\EndFor
\State $\phi^* \gets \prod_{\phi \in \Phi} \phi$
\State \Return $\phi^*$
\EndFunction
\State
\Function{Sum-Product-Eliminate-Var}{$\Phi$, $X$}
\State{$\Phi$ množina všech faktorů.}
\State{$X$, proměnná, která má být eliminována.}
\State

\State $\Phi^\prime \gets \{\phi \in \Phi: X \in Dom(\phi)\}$
\State $\Phi^{\prime \prime} \gets \Phi - \Phi^\prime$
\State $\psi \gets \prod_{\phi \in \Phi^\prime} \phi$
\State $\tau \gets \sum_X \psi$
\State \Return $\Phi^{\prime\prime} \bigcup \{\tau\}$
\EndFunction
\end{algorithmic}
\end{algorithm}

Nechť $\vec{X}$ je množina náhodných proměnných, nechť $\Phi$ je množina faktorů, t.ž. pro každé $\phi \in \Phi$, $Dom(\phi) \subseteq \vec{X}$.
Nechť $\vec{Y} \subset \vec{X}$ je množina dotazovaných náhodných proměnných a nechť $\vec{Z} = \vec{X} - \vec{Y}$.
Pak pro každé seřazení $\prec$ nad $\vec{Z}$, Sum-Product-VE($\Phi$, $\vec{Z}$, $\prec$) vrátí faktor $\phi^*(\vec{Y})$ takový, že
\[
\phi^*(\vec{Y}) = \sum_Z \prod_{\phi \in \Phi} \phi
\]

Nyní provedeme analýzu algoritmu eliminace proměnných.
Předpokládejme, že na vstupu je $n$ proměnných.
Bayesovská síť obsahuje pro každou proměnnou jeden faktor.
Pro jednoduchost budeme předpokládat, že algoritmus bude eliminovat všechny proměnné.
Běh algoritmu se skládá z jednotlivých eliminačních kroků, při kterých je vždy eliminována jedna proměnná.

Při jednom eliminačním kroku je vybrána proměnná $X_i$, všechny faktory, které ji obsahují jsou vynásobeny a vytvoří jeden velký faktor $\psi_i$.
Proměnná $X_i$ je pak vysčítána z tohoto faktoru.
Počet operací pro jeden eliminační krok tedy závisí na velikosti faktoru $\psi_i$, označme ji $N_i$.
Maximum z velikostí faktorů označme $N_{max} = max_i N_i$.

Nyní se zaměříme na počet násobení.
Celkem vznikne $n+m$ faktorů, kde $m$ je počet faktorů, které vznikly vysčítáním proměnné.
Každý z těchto faktorů je zahrnut do součinu pouze jednou, pokud je eliminována nějaká proměnná, kterou obsahuje.
Cena násobení faktorů pro vznik $\psi_i$ je nejvýše $N_i$.
Celkový počet násobení tedy bude nejvýše $(n+m)N_{max}$ což je $\mathcal{O}(nN_{max})$.

Pokud $k$ je maximum z velikostí domén proměnných, pak velikost faktoru obsahujícího $n$ proměnných může být až $k^n$.
Složitost eliminace proměnných je tedy dominována velikostí faktorů, které vznikají při výpočtu a je exponenciální.
Navíc bylo dokázáno, že výběr nejlepšího pořadí proměnných pro eliminaci je NP těžký~\cite{arnborg1987complexity}.

\subsection{Posílání zpráv ve faktor grafu}

Algoritmy exaktní inference naráží při reálném použití na příliš velkou složitost násobení faktorů.
Hlavním problémem je velikost sdružené pravděpodobnosti, která roste exponenciálně s počtem náhodných proměnných.
Většinou nás ovšem zajímá marginální pravděpodobnost jedné nebo jen mála proměnných.
Je tedy zbytečné počítat celou sdruženou pravděpodobnost, abychom z ní pak vysčítali většinu proměnných.
Řešení se nabízí ve formě aproximace faktorizací sdružené pravděpodobnosti.

\begin{equation*}
    P(X_1, \dots, X_n) = \prod_i P(X_i)
\end{equation*}

Pro výpočet s faktorizovanou distribucí si zavedeme novou datovou strukturu, tzv. faktor graf.
Faktor graf je bipartitní graf, kde jednu partitu tvoří faktory a druhou partitu tvoří náhodné proměnné.
Hrany ve faktor grafu vedou vždy mezi proměnnou a faktorem, který ji obsahuje.

\begin{figure}[H]
\begin{center}
\begin{tikzpicture}
    
\matrix[row sep=0.75cm, column sep=1.2cm]
{
    \node[latent]   (X1)    {$X_1$};
    &
    \factor         {fa}    {$f_a$}{}{};
    &
    \node[latent]   (X2)    {$X_2$};
    &
    \factor         {fb}    {$f_b$}{}{};
    &
    \node[latent]   (X3)    {$X_3$};
    \\
    &&
    \factor         {fc}    {left:$f_c$}{}{};
    \\
    &&
    \node[latent]   (X4)    {$X_4$}{}{};
    \\
};
\edge[-]{X1,X2} {fa}
\edge[-]{X3,X2} {fb}
\edge[-]{X2,X4} {fc}

\end{tikzpicture}
\end{center}
\caption{Příklad faktor grafu se třemi náhodnými proměnnými $X_1, X_2, X_3$ a třemi faktory $f_a, f_b, f_c$.}
\end{figure}

Nejprve začneme s inferencí na stromech.
Pokud zafixujeme jednu náhodnou proměnnou $X$ ve faktor grafu, pak sdruženou pravděpodobnost můžeme spočítat jako
\begin{equation}
    P(\vec{X}) = \prod_{s \in ne(X)} F_s(X, \vec{X}_s),
\label{eq:fs}
\end{equation}
kde $ne(X)$ jsou faktory obsahující proměnnou $X$ (tedy sousedi ve faktor grafu), $F_s$ je součin všech faktorů v podstromu určeném faktorem $f_s$ a $\vec{X}_s$ je množina všech proměnných v daném podstromu.

Pro výpočet marginální pravděpodobnosti $X$ substituujeme~(\ref{eq:fs}) do výpočtu marginální pravděpodobnosti ze sdružené a po výměně sumy a produktu dostaneme
\begin{align}
P(X) &= \prod_{s \in ne(X_i)} \sum_{\vec{X}_s} F_s(X, \vec{X}_s)
\\
&= \prod_{s \in ne(X)} \mu_{f_s \rightarrow X}(X)
\label{eq:margx}
\end{align}

Zavedli jsme funkce
\begin{equation}
    \mu_{f_s \rightarrow X} \equiv \sum_{\vec{X}_s} F_s(X, \vec{X}_s),
\label{eq:defmsgfx}
\end{equation}
které můžeme nazývat zprávami z faktoru $f_s$ do proměnné $X$.

Každé $F_s(X, \vec{X}_s)$ je popsáno podgrafem faktor grafu a tedy může být znova faktorizováno.
Namísto proměnné $X$ nyní vezmeme faktor $f_s$.
Náhodné proměnné sousedící s faktorem $f_s$ bez $X$ si označíme $X_1, \dots, X_M$.
Součin faktorů v podstromech určených těmito proměnnými označíme $G_i(X_i, \vec{X}_{si})$.
Faktor $F_s(X, \vec{X}_s)$ tedy můžeme přepsat jako
\begin{equation}
F_s(X, \vec{X}_s) = f_s(X, X_1, \dots, X_M) G_1(X_1, \vec{X}_{s1}), \dots, G_M(X_M, \vec{X}_{sM})
\end{equation}

Pokud substitujeme přepsaný faktor $F_s(X, \vec{X}_s)$ do definice zprávy z faktoru, dostaneme
\begin{align}
\mu_{f_s \rightarrow X}(X) &=
    \sum_{X_1} \dots \sum_{X_M}
        f_s(X, X_1, \dots, X_M)
        \prod_{m \in ne(f_s) \backslash X}
            \sum_{\vec{X}_{sm}}
                G_m(X_m, \vec{X}_{sm})
\\
&= \sum_{X_1} \dots \sum_{X_M}
    f_s(X, X_1, \dots, X_M)
    \prod_{m \in ne(f_s) \backslash X}
        \mu_{X_m \rightarrow f_s}(X_m)
\label{eq:mfsx}
\end{align}

Zavedli jsme další funkce 
\begin{equation}
    \mu_{X_m \rightarrow f_s}(X_m) \equiv \sum_{\vec{X}_{sm}} G_m(X_m, \vec{X}_{sm}),
\label{eq:defmsgxf}
\end{equation}
které budeme nazývat zprávami z proměnné $X_m$ do faktoru $f_s$.

Z rovnice~\ref{eq:margx} vidíme, že marginální pravděpodobnost proměnné vypočítáme jako součin zpráv ze všech okolních faktorů.
Každou z těchto zpráv můžeme spočítat jako součin faktoru a zpráv přicházejících z proměnných, které s tímto faktorem sousedí, kromě proměnné, které chceme zprávu posílat. Zbývá nám tedy zjistit v jakém tvaru jsou zprávy z proměnných do faktoru.

Stejně jako u $F_s(X, \vec{X}_s)$ nám tady $G_m(X_m, \vec{X}_{sm})$ definuje podgraf faktor grafu. 
V kořeni tohoto podgrafu leží proměnná $X_m$ a tím se dostáváme na už známý případ.
Budeme ignorovat faktor $f_s$, protože neleží v podgrafu určeném $G_m(X_m, \vec{X}_{sm})$, a pak 
\begin{equation}
    G_m(X_m, \vec{X}_{sm}) = \prod_{l \in ne(X_m) \backslash f_s} F_l(X_m, \vec{X}_{ml})
\end{equation}
Přepsáno ve formě zpráv
\begin{align}
\mu_{X_m \rightarrow f_s}(X_m)
&= \prod_{l \in ne(X_m) \backslash f_s}
    \sum_{\vec{X}_{ml}}
        F_l(X_m, \vec{X}_{ml})
\\
&= \prod_{l \in ne(X_m) \backslash f_s}
    \mu_{f_l \rightarrow X_m}(X_m)
\end{align}
Zpráva z náhodné proměnné je tedy součinem zpráv ze všech ostatních faktorů.

Zprávy tedy dokážeme počítat rekurzivně, chybí nám ovšem pravidla pro zprávy z uzlů, které jsou listy.
V případě faktoru je odchozí zpráva ekvivalentní faktoru.
\begin{equation}
    \mu_{f \rightarrow X}(X) = f(X)
\end{equation}
V případě náhodné proměnné je odchozí zpráva 
\begin{equation}
    \mu_{X \rightarrow f}(X) = 1
\end{equation}
Pro pozorovanou proměnnou je zpráva vždy rovna pozorované hodnotě, bez ohledu na sousední faktory.

\subsection{Belief Propagation}

Algoritmus exaktní inference na stromech s pomocí posílání zpráv se nazývá belief propagation~\cite{pearl1988probabilistic}.
Zpráva z vrcholu faktor grafu může být poslána do jiného vrcholu, až když byly obdrženy zprávy ze všech ostatních vrcholů.
Jakmile vrchol obdržel zprávy ze všech vrcholů, lze spočítat jeho marginální pravděpodobnost.
Strom můžeme zakořenit v libovolné náhodné proměnné a posílat zprávy z listů do kořene.
Kořen pak obdrží všechny zprávy ze všech sousedních faktorů a tedy je možné spočítat jeho marginální pravděpodobnost.
Tomuto se říká dopředný krok propagace.

\begin{figure}[H]
\begin{center}
\begin{tikzpicture}

\matrix[row sep=0.75cm, column sep=1.2cm]
{
    \node[latent]   (X1)    {$X_1$};
    &
    \factor         {fa}    {$f_a$}{}{};
    &
    \node[latent]   (X2)    {$X_2$};
    &
    \factor         {fb}    {$f_b$}{}{};
    &
    \node[latent]   (X3)    {$X_3$};
    \\
    &&
    \factor         {fc}    {left:$f_c$}{}{};
    \\
    &&
    \node[latent]   (X4)    {$X_4$}{}{};
    \\
};
\edge{X1}{fa}
\edge{fa}{X2}
\edge{X4}{fc}
\edge{fc}{X2}
\edge{X2}{fb}
\edge{fb}{X3}

\end{tikzpicture}
\end{center}
\caption{Směr posílání zpráv pro strom zakořeněný ve vrcholu $X_3$ v dopředném kroce.}
\end{figure}

Belief propagation algoritmus tedy umožňuje na stromech spočítat marginální pravděpodobnost jedné proměnné s linárním počtem poslaných zpráv.
Při rozšíření na spočítání marginální pravděpodobnosti všech proměnných není třeba algoritmus pouštět $n$-krát.
Stačí si uvědomit, že všem vrcholům chybí pouze zpráva od faktoru na cestě ke kořenu.
A my můžeme z kořene poslat zprávy zpět k listům hned, jak obdržíme všechny příchozí zprávy.
Této části se říká zpětný krok propagace.

\begin{figure}[H]
\begin{center}
\begin{tikzpicture}

\matrix[row sep=0.75cm, column sep=1.2cm]
{
    \node[latent]   (X1)    {$X_1$};
    &
    \factor         {fa}    {$f_a$}{}{};
    &
    \node[latent]   (X2)    {$X_2$};
    &
    \factor         {fb}    {$f_b$}{}{};
    &
    \node[latent]   (X3)    {$X_3$};
    \\
    &&
    \factor         {fc}    {left:$f_c$}{}{};
    \\
    &&
    \node[latent]   (X4)    {$X_4$}{}{};
    \\
};
\edge{fa}{X1}
\edge{X2}{fa}
\edge{fc}{X4}
\edge{X2}{fc}
\edge{fb}{X2}
\edge{X3}{fb}

\end{tikzpicture}
\end{center}
\caption{Směr posílání zpráv pro strom zakořeněný ve vrcholu $X_3$ ve zpětném kroce.}
\end{figure}

Po propagaci všech zpráv až k listům už každá proměnná získala zprávy od všech sousedních faktorů a tedy je možné spočítat marginální pravděpodobnost všech proměnných.

\subsection{Loopy Belief Propagation}
\label{sec:lbp}

Algoritmus Belief Propagation funguje pro inferenci na stromech.
Hlavní důvod proč nelze použít Belief Propagation na obecných grafech je, že můžeme narazit na cyklus v grafu, což znamená, že žádný z vrcholů v tomto cyklu nebude nikdy mít dostatek příchozích zpráv, aby mohl nějakou zprávu odeslat.

Algoritmus Loopy Belief Propagation řeší tento problém relaxací podmínky na příchozí zprávy.
Pro odeslání zprávy do uzlu není třeba znát zprávy ze všech ostatních sousedních vrcholů.
Chybějící zprávy jsou nastaveny na jedničku.
Loopy Belief Propagation je iterativní algoritmus.
V grafu nemusí existovat vrchol, který může posílat zprávu k nějakému ze svých sousedů.
Je tedy třeba vybrat vrchol podle nějaké strategie a z něj poslat zprávy do všech sousedních vrcholů.
V další iteraci je pak vybrán zase jiný vrchol, který bude posílat zprávy.
Iterace končí ve chvíli, kdy se už nemění marginální pravděpodobnosti proměnných.

Vzhledem k tomu, že posílané zprávy už neodpovídají skutečným faktorům, je důležitá otázka, zda-li Loopy Belief Propagation vůbec nalezne správné marginální pravděpodobnosti.
Existují podmínky za kterých algoritmus bude konvergovat~\cite{tatikonda2002loopy}.
V obecném případě ovšem může algoritmus konvergovat ke špatným pravděpodobnostem, anebo nemusí konvergovat vůbec a pak dochází k oscilacím.
Pokud algoritmus nekonverguje ke správným pravděpodobnostem, tak není ani možné se k nim přiblížit nebo je odvodit z oscilací~\cite{murphy1999loopy}.
V praxi ovšem k problémům většinou nedochází a inference konverguje v rozumném čase.

\begin{algorithm}
\caption{Loopy Belief Propagation}
\label{alg:lbp}
\begin{algorithmic}
\Function{LBP}{$F$, $S$}
\State{$F$ -- faktor graf}
\State{$S$ -- strategie}
\State

\State{$\textsc{Init-Factor-Graph}(F)$}
\Repeat
    \For{vrchol $v \in F$ vybraný podle strategie $S$}
        \For{soused $n$ vrcholu $v$}
            \If{$v$ je faktor}
                \State{$\mu_{v \rightarrow n}(n) \gets \textsc{Message-To-Var}(v, n)$}
            \Else
                \State{$\mu_{v \rightarrow n}(v) \gets \textsc{Message-To-Factor}(v, n)$}
            \EndIf
        \EndFor
    \EndFor
\Until{konvergence}

\For{každá proměnná $X$}
    \State{$P(X) = \prod_{f \in ne(X)} \mu_{f \rightarrow X}(X)$}
\EndFor
\EndFunction

\State

\Function{Init-Factor-Graph}{$F$}
\State{$F$ -- faktor graf}
\State

\For{každou proměnnou $X \in F$}
    \For{sousední faktor $f$}
        \State{$\mu_{X \rightarrow f}(X) \gets 1$}
        \State{$\mu_{f \rightarrow X}(X) \gets 1$}
    \EndFor
\EndFor
\EndFunction

\State

\Function{Message-To-Var}{$f$, $v$}
\State{$f$ -- zdrojový faktor}
\State{$v$ -- cílová proměnná}
\State

\State{$\vec{X} \gets \text{Dom}(f)$}
\State{$\mu_{f \rightarrow v}(v) \gets \sum_{\vec{X} \backslash v} f(\vec{X}) \prod_{u \in ne(f) \backslash v} \mu_{u \rightarrow f}(u)$}
\EndFunction

\State

\Function{Message-To-Factor}{$v$, $f$}
\State{$v$ -- zdrojová proměnná}
\State{$f$ -- cílový faktor}
\State

\State{$\mu_{v \rightarrow f}(v) \gets \prod_{g \in ne(v) \backslash f} \mu_{g \rightarrow v}(v)$ }
\EndFunction
\end{algorithmic}
\end{algorithm}

\subsubsection{Popis LBP}

Popis Loopy Belief Propagation je v algoritmu~\ref{alg:lbp}.
Při inferenci ve faktor grafu je nejprve třeba inicializovat všechny zprávy.
Posílání zpráv budeme iterovat dokud zprávy nedokonvergují, anebo můžeme nastavit pevný počet iterací.
Pak vybíráme vrcholy a z každého pošleme zprávu do všech okolních.
Způsobem výběru vrcholů se budeme zabývat později.
I náhodný výběr vrcholů povede k výsledku, dodání externí informace či chytřejší heuristiky ovšem může výpočet výrazně urychlit.

Poslání zprávy se liší podle toho, zda-li posíláme z proměnné do faktoru, anebo naopak.
V obou případech vynásobíme příchozí zprávy ze všech sousedních vrcholů, kromě toho, do kterého zprávu posíláme.
Při posílání zprávy z vecholu pro faktor zprávy vynásobíme faktorem a marginalizujeme všechny proměnné, kromě té, které zprávu posíláme.

Ve faktor grafu se tedy zprávy šíří pouze přes pravděpodobnostní rozdělení jedné proměnné.
To právě může vést k oscilacím anebo konvergenci ke špatným hodnotám. 
Příkladem je například skrytý markovský model se dvěma skrytými binárními proměnnými a pozorováním pro každou skrytou proměnnou.
Pokud je mezi proměnnými závislost XOR, pak nikdy nebudeme pozorovat $(1, 1)$, ale přesto bude každá proměnná konvergovat k rovnoměrnému rozdělení. A tedy sdružená pravděpodobnost pro pozorování $(1, 1)$ dostane nenulovou pravděpodobnost.

Při počítání zpráv často musíme pronásobit spoustu zpráv z okolních vrcholů.
Přitom se zprávy zas tak často nemění.
Častou optimalizací je pro každý vrchol si pamatovat součin všech jeho zpráv
\begin{equation}
    b(v) = \begin{cases}
        \prod_{u \in ne(v)} \mu_{u \rightarrow v}(v) & \mbox{$v$ je proměnná} \\
        f_v(\vec{X}) \prod_{u \in ne(v)} \mu_{u \rightarrow v}(u) & \mbox{$v$ je faktor} \\
    \end{cases}
\end{equation}
Pak můžeme při odesílání zprávy z vrcholu $v$ do vrcholu $u$ pouze $b(v)$ vydělit zprávou $\mu_{u \rightarrow v}$, pro faktory marginalizovat, a zpráva je připravena k odeslání.
$b(v)$ můžeme aktualizovat po výběru vrcholu $v$ pro posílání zpráv, anebo po každém přijetí zprávy.

Pro pozorované proměnné je třeba metodu \textsc{Message-To-Factor} upravit tak, aby pouze posílala pozorovanou hodnotu.

Implementace Loopy Belief Algoritmu i se strategiemi pro výběr vrcholů bude popsána v kapitole~\ref{ch:kap4}.

\section{Propagace s aproximovanými zprávami}

V předchozí sekci byla faktorizována sdružená pravděpodobnost a tak došlo k zjednodušení inference.
Mezi jednotlivými vrcholy byly stále posílány exaktní zprávy.
V této sekci se zaměříme na inferenci v modelech, kde není možné spočítat zprávy exaktně a je třeba je aproximovat.
Tento problém nastává u některých modelů se spojitými náhodnými proměnnými.

Předpokládáme, že máme pravděpodobnostní grafický model, reprezentující sdruženou pravděpodobnost dat a pozorování pomocí součinu faktorů.
\begin{equation}
P(\vec{X}, \vec{E}) = \prod_i f_i(\vec{X})
\end{equation}
Nás zajímá aposteriorní distribuce $P(\vec{X} \mid \vec{E})$ pro zjištění stavu, stejně tak jako pravděpodobnost pozorování $P(\vec{E})$ pro normalizaci.
Aposteriorní distribuci vyjádříme ze sdružené.
\begin{equation}
P(\vec{X} \mid \vec{E}) = \frac{1}{P(\vec{E})} \prod_i f_i(\vec{X})
\end{equation}
a pravděpodobnost pozorování je dána
\begin{equation}
P(\vec{E}) = \int \prod_i f_i(\vec{X}) \mathrm{d}\vec{X}
\end{equation}

Pro aproximativní inferenci ukážeme algoritmus Expectation propagation (EP)~\cite{minka2001expectation}.
Vytvoříme aproximaci aposteriorní distribuce, která je také dána součinem faktorů
\begin{equation}
q(\vec{X}) = \frac{1}{Z} \prod_i \tilde{f}_i(\vec{X}),
\end{equation}
kde každý faktor $\tilde{f}_i$ je aproximace odpovídající skutečnému faktoru $f_i$ a faktor $\frac{1}{Z}$ je normalizační konstanta.
Aproximované faktory musíme zkonstruovat tak, abychom byli schopni provádět inferenci.

Pro měření vzdálenosti aproximovaného rozdělení od skutečného používáme Kullback-Leiblerovu divergenci (KL)~\cite{kullback1997information}.
KL divergence, také známá jako relativní entropie, mezi dvěma pravděpodobnostními rozděleními $p(x)$ a $q(x)$ je
\begin{equation}
KL(p \| q) = \int p(x) \log \left(\frac{p(x)}{q(x)}\right) \mathrm{d}x.
\end{equation}
Divergence splňuje tři vlastnosti:
\begin{enumerate}
\item $KL(p \| p) = 0$,
\item $KL(p \| q) = 0$ právě tehdy když $p = q$,
\item $KL(p \| q) > 0$ pro všechna $p, q$.
\end{enumerate}

\subsection{Rodina exponenciálních rozdělení}
\label{sec:expfam}

Rodina exponenciálních rozdělení~\cite{bernardo2009bayesian} přes $\vec{X}$, dáno parametry $\vec\eta$ je definována jako množina distribucí ve tvaru
\begin{equation}
P(\vec{X} \mid \vec\eta) = h(\vec{X}) g(\vec\eta) \exp \big(\vec\eta^T u(\vec{X})\big),
\label{eq:gj}
\end{equation}
kde $\vec{X}$ jsou náhodné proměnné, které mohou být diskrétní nebo spojité, $\vec\eta$ jsou nazývány přirozené parametry rozdělení, $u(\vec{X})$ je funkce $\vec{X}$.
Funkce $g(\vec\eta)$ může být interpretována jako koeficient, který zajišťuje, že je distribuce normalizována a tedy splňuje
\begin{equation}
g(\vec\eta) \int h(\vec{X}) \exp \big(\vec\eta^T u(\vec{X})\big) \mathrm{d}\vec{X} = 1
\label{eq:g1}
\end{equation}

Výhody distribucí z exponenciální rodiny si ukážeme na KL divergenci.
Nechť počítáme divergenci $KL(p \| q)$, kde $p(\vec{X})$ je zafixovaná distribuce a $q(\vec{X})$ je distribuce z exponenciální rodiny.
Pokud zapíšeme KL divergenci jako funkci $\vec\eta$, tak dostaneme
\begin{equation}
KL(p \| q) = - \log g(\vec\eta) - \vec\eta^T \mathbb{E}_{p(\vec{X})}[u(\vec{X})] + \mathrm{const}.
\end{equation}
KL divergenci můžeme minimalizovat tak, že nastavíme první derivaci podle $\vec\eta$ rovnou nule, z čehož dostaneme
\begin{equation}
-\nabla \log g(\vec\eta) = \mathbb{E}_{p(\vec{X})}[u(\vec{X})].
\label{eq:kl}
\end{equation}

Nyní se musíme podívat na $- \nabla \log g(\vec\eta)$.
Vezmeme derivaci obou stran (\ref{eq:g1}) podle $\vec\eta$.
Získáme
\begin{align}
& \nabla g(\vec\eta) \int h(\vec{X}) \exp \big(\vec\eta^T u(\vec{X}) \big) \mathrm{d}\vec{X} \\
& + g(\vec\eta) \int h(\vec{X}) \exp \big(\vec\eta^T u(\vec{X})\big) u(\vec{X}) \mathrm{d}\vec{X} = 0 \nonumber
\end{align}
Přeuspořádáním a znova použitím (\ref{eq:gj}) a (\ref{eq:g1}) získáme
\begin{equation}
-\frac{1}{g(\vec\eta)}\nabla g(\vec\eta)
= g(\vec\eta) \int h(\vec{X}) \exp\big(\vec\eta^T u(\vec{X})\big) u(\vec{X}) \mathrm{d}\vec{X} = \mathbb{E}[u(\vec{X})].
\end{equation}
Ve výsledku získáváme 
\begin{equation}
- \nabla \log g(\vec\eta) = \mathbb{E}[u(\vec{X})]
\label {eq:ge}
\end{equation}

Nový poznatek o funkci $g(\vec\eta)$ z (\ref{eq:ge}) můžeme substituovat do (\ref{eq:kl}) a získáme
\begin{equation}
\mathbb{E}_{q(\vec{X})}[u(\vec{X})] = \mathbb{E}_{p(\vec{X})}[u(\vec{X})],
\label{eq:ee}
\end{equation}
z čehož vidíme, že pro minimalizaci KL divergence nám stačí najít parametry rozdělení $q(\vec{X})$ tak, aby mělo stejné očekávané statistiky jako $p(\vec{X})$.

\subsection{Expectation Propagation}
\label{sec:ep}

Expectation propagation je algoritmus založený na aproximaci aposteriorní distribuce součinem aproximovaných faktorů
\begin{equation}
q(\vec{X}) = \frac{1}{Z} \prod_i \tilde{f}_i(\vec{X}).
\end{equation}
Z praktických důvodů uvedených v předchozí sekci jsou aproximované faktory z exponenciální rodiny.
Díky tomu bude jejich součin také z exponenciální rodiny.

Aproximující distribuci bychom chtěli nalézt pomocí minimalizace KL divergence mezi skutečnou a aproximovanou distribucí.
\begin{equation}
KL(p \| q) = KL\bigg(\frac{1}{p(\vec{E})} \prod_i f_i(\vec{X}) \bigg\| \frac{1}{Z} \prod_i \tilde{f}_i(\vec{X} \bigg)
\end{equation}
Pro počítání KL divergence ovšem potřebujeme umět počítat se skutečnou pravděpodobností a když už s ní umíme efektivně počítat, proč ztrácet čas s aproximací.

Další možností je aproximovat jednotlivé faktory, díky tomu bychom v jednom kroku nalezli aproximace pro všechny faktory a měli bychom hotovo.
Naším cílem je ale nalézt nejlepší aproximaci celé aposteriorní pravděpodobnosti a to není v tomto případě zaručeno.

Expectation propagation sice aproximuje jednotlivé faktory, ale vždy v kontextu všech ostatních~\cite{bishop2006pattern}.
Nejprve jsou všechny inicializovány a pak jsou procházeny jeden po druhém a každý je aktualizován.
Tento přístup je podobný metodám, které byly ukázány v předešlých sekcí.
Předpokládejme, že chceme aktualizovat faktor $\tilde{f}_j (\vec{X})$.
Nejprve odstraníme faktor z produktu $\prod_{i \ne j} \tilde{f}_i(\vec{X})$.
Následně nalezneme novou hodnotu pro faktor $\tilde{f}_j$ tak, aby pravděpodobnostní rozložení
\begin{equation}
q^{new}(\vec{X}) \propto \tilde{f}_j(\vec{X}) \prod_{i \ne j} \tilde{f}_i(\vec{X})
\label{eq:qnew}
\end{equation}
bylo co nejblíže 
\begin{equation}
f_j(\vec{X})\prod_{i \ne j} \tilde{f}_i(\vec{X}).
\end{equation}
Všechny faktory $i \ne j$ necháváme zafixované. 
Díky tomu zajistíme, že nová aproximace faktoru je nejpřesnější v oblastech, které mají největší aposteriorní pravděpodobnost definovanou zbývajícími faktory.

Odstranění faktoru $\tilde{f}_j(\vec{X})$ z aktuální aproximace aposteriorní distribuce provedeme vytvořením nenormalizované distribuce
\begin{equation}
q^{\backslash j}(\vec{X}) = \frac{q(\vec{X})}{\tilde{f}_j(\vec{X})}
\end{equation}
Mohli bychom počítat $q^{\backslash j}(\vec{X})$ jako součin všech faktorů kromě $j$, ale v praxi je dělení rychlejší.
Této nenormalizované distribuci se říká cavity distribuce.
Její kombinací se skutečným faktorem $f_j(\vec{X})$ dostaneme distribuci
\begin{equation}
    \frac{1}{Z_j} f_j(\vec{X}) q^{\backslash j}(\vec{X}),
\label{eq:aprox}
\end{equation}
kde $Z_j$ je normalizační konstanta, dána
\begin{equation}
    Z_j = \int f_j(\vec{X}) q^{\backslash j}(\vec{X}) \mathrm{d}\vec{X}.
\label{eq:zj}
\end{equation}

Nyní můžeme nalézt novou hodnotu pro faktor $\tilde{f}_j$ minimalizací KL divergence
\begin{equation}
    KL\bigg(\frac{f_j(\vec{X}) q^{\backslash j}(\vec{X})}{Z_j} \bigg\| q^{new}(\vec{X})\bigg).
\end{equation}
To uděláme jednoduše použitím poznatku z (\ref{eq:ee}), který říká, že nám stačí najít parametry $q^{new}$ tak, aby postačující statistiky odpovídali momentům aproximovaného rozdělení (\ref{eq:aprox}).

Z nalezené aproximace $q^{new}$ pak získáme novou hodnotu faktoru $\tilde{f}_j$ z (\ref{eq:qnew}) vydělením zbývajících faktorů
\begin{equation}
    \tilde{f}_j(\vec{X}) = K \frac{q^{new}(\vec{X})}{q^{\backslash j}(\vec{X})}
\label{eq:newfj}
\end{equation}

Koeficient $K$ získáme tak, že obě strany (\ref{eq:newfj}) vynásobíme $q^{\backslash j}(\vec{X})$ a zintegrováním, čímž získáme
\begin{equation}
    K = \int \tilde{f}_j(\vec{X}) q^{\backslash j}(\vec{X}) \mathrm{d}\vec{X},
\end{equation}
použili jsme navíc toho, že $q^{new}(\vec{X})$ je normalizovaná distribuce.
Hodnota $K$ pak může být nalezena srovnáním
\begin{equation}
    \int \tilde{f}_j (\vec{X}) q^{\backslash j}(\vec{X}) \mathrm{d}\vec{X} = \int f_j(\vec{X}) q^{\backslash j}(\vec{X}) \mathrm{d}\vec{X}
\end{equation}
Z čehož zjistíme, že $K = Z_j$ a tedy může být nalezeno přímo z (\ref{eq:zj}).

Expectation Propagation je popsán v algoritmu~\ref{alg:ep}.
Dokud aproximované rozdělení nekonverguje, vybíráme jednotlivé faktory a aktualizujeme je podle popsaného postupu.

\begin{algorithm}
\caption{Expectation propagation}
\label{alg:ep}
\begin{algorithmic}
\State Inicializujeme všechny aproximativní faktory $\tilde{f}_i(\vec{X})$ na neinformativní.
\State Inicializujeme aposteriorní aproximace nastavením: $$q(\vec{X}) \propto \prod_i \tilde{f}_i(\vec{X})$$
\Repeat
    \State Vyber faktor $\tilde{f}_j(\vec{X})$, který bude aktualizován.
    \State Odeber $\tilde{f}_j(\vec{X})$ z aposteriorní distribuce vydělením: $$q^{\backslash j} (\vec{X}) = \frac{q(\vec{X})}{\tilde{f}_j(\vec{X})}$$
    \State Vypočítej novou aposteriorní distribuci nastavením momentů $q^{new}$ na 
    \State hodnotu momentů $q^{\backslash j}(\vec{X})f_j(\vec{X})$, včetně vyhodnocení normalizační 
    \State konstanty $$Z_j = \int q^{\backslash j}(\vec{X}) f_j(\vec{X}) \mathrm{d}\vec{X}$$
    \State Nastav novou hodnotu faktoru: $$\tilde{f}_j(\vec{X}) = Z_j \frac{q^{new}(\vec{X})}{q^{\backslash j}(\vec{X})}$$
\Until{konvergence}
\State Spočítej aproximaci pozorování: $$p(\vec{E}) \simeq \int \prod_i \tilde{f}_i(\vec{X}) \mathrm{d}\vec{X}$$
\end{algorithmic}
\end{algorithm}

V praxi je třeba provést několik iterací, a v každé aktualizovat všechny faktory.
Stejně jako u LBP zde nemáme žádnou garanci, že bude algoritmus konvergovat.
Pro aproximace $q(\vec{X})$ v exponenciální rodině, pokud iterace konverguje, pak nalezené řešení je stacionární bod specifické potenciální energie~\cite{minka2001expectation}.

% !TEX root = thesis.tex
\chapter{Učení parametrů}
\label{ch:ep}

Představená metoda LBP v~minulé kapitole funguje pro modely, kde máme nastavené parametry faktorů.
V~této kapitole si představíme grafický model pro diskrétní proměnné, u~kterého jsou i parametry faktorů proměnné a je možné pro ně použít Bayesovský přístup.
Díky tomu bude možné nalézt aposteriorní distribuci pro tyto parametry a naučit je podle dat.

Inferenci v~tomto upraveném modelu už není možné dělat pomocí LBP a~budeme muset použít algoritmus Expectation Propagation.

\section{Grafický model}

Model se skládá z~faktorů a proměnných.
Nechť máme vybraný faktor $f_\beta$, spojený s~několika proměnnými
$\vec{x} = (x_0, x_1, \dots, x_{N_x})$
a množinami parametrů
$\vec{\Theta} = (\vec{\theta}_1, \dots, \vec{\theta}_{N_\theta})$.
Tento faktor reprezentuje podmíněnou pravděpodobnost:
$$f_\beta(\vec{x}, \Theta) = p(x_0 | x_1, \dots, x_{N_x}; \Theta).$$

Rodičovské proměnné $x_1, \dots, x_{N_x}$ označujeme jako $\vec{x^\prime}$.
Vektor $\vec{x^\prime}$ určuje, jaká množina parametrů bude použita.
Protože množiny parametrů jsou číslovány $1, \dots, N_\theta$ a rodičovské
proměnné $1, \dots, N_x$, musí být pro vybrání správné množiny parametrů
použito mapování $\rho(\vec{x^\prime})$.
Faktor pak může být zapsán zkráceně:
$$f_\beta(\vec{x}, \Theta) = p(x_0 | x_1, \dots, x_{N_x}; \Theta) =
\theta_{\rho(x^\prime), x_0}$$

Příklad faktoru $f_\beta$ je na obrázku~\ref{ex:factor}, pro ilustraci pouze s~třemi parametry a~čtyřmi proměnnými.
Správně by měl faktor mít tolik parametrů, kolik je možných přiřazení pro rodičovské proměnné.

\begin{figure}[h]
\begin{center}
\begin{tikzpicture}[->, main/.style={circle, draw}, factor/.style={draw}]
    
\node[main] at (-3.7, 2)    (t1) {$\vec\theta_1$};
\node[main] at (-4, 0)      (t2) {$\vec\theta_2$};
\node[main] at (-3.7, -2)   (t3) {$\vec\theta_3$};
\node[factor] at (0, 0)     (f)  {$f_\beta$};
\node[main] at (3.7, 2)     (x1) {$X_1$};
\node[main] at (4, 0)       (x2) {$X_2$};
\node[main] at (3.7, -2)    (x3) {$X_3$};
\node[main] at (0, -4)      (x0) {$X_0$};

\path
	(t1) edge [above right, near start] node {$q^{\backslash \beta}(\vec\theta_1)$} (f)
	(t2) edge [above, near start] node {$q^{\backslash \beta}(\vec\theta_2)$} (f)
	(t3) edge [below right, near start] node {$q^{\backslash \beta}(\vec\theta_3)$} (f)
	(x1) edge [above left, near start] node {$q^{\backslash \beta}(X_1)$} (f)
	(x2) edge [above, near start] node {$q^{\backslash \beta}(X_2)$} (f)
	(x3) edge [below left, near start] node {$q^{\backslash \beta}(X_3)$} (f)
	(f) edge [right, near end] node {$m_{f_\beta \rightarrow X_0}(X_0)$} (x0);


\end{tikzpicture}
\end{center}
\caption{Vybraný faktor $f_\beta$ pro aktualizaci při učení parametrů.}
\label{ex:factor}
\end{figure}

\section{Výpočet marginálních pravděpodobností}

Pro výpočet sdružené pravděpodobnosti používáme plně faktorizovanou distribuci.
Pro každou proměnnou (množinu parametrů) je její marginální pravděpodobnost
rovna součinu zpráv ze sousedních faktorů.
Pro daný faktor $f_\beta$ je neúplná distribuce $q^{\backslash \beta}(x_i)$, popř.
$q^{\backslash \beta}(\vec{\theta}_i)$, rovna součinu zpráv ze všech ostatních faktorů do $x_i$, popř $\vec\theta_i$.
Aproximovaná marginální pravděpodobnost proměnné (množiny parametrů) je pak součinem neúplné
distribuce a zprávy z~faktoru:
$$q(x_i) = q^{\backslash \beta}(x_i) m_{f_\beta \rightarrow x_i}(x_i), 
\quad \quad
q(\vec{\theta}_i) = q^{\backslash \beta}(\vec{\theta}_i) m_{f_\beta \rightarrow \theta}(\vec{\theta}_i)$$

Tyto neúplné distribuce jsou právě zprávy z~proměnné, popř. množiny parametrů, do faktoru:

$$m_{x_i \rightarrow f_\beta} = q^{\backslash \beta}(x_i), 
\quad \quad
m_{\vec\theta_i \rightarrow f_\beta} = q^{\backslash \beta}(\vec\theta_i).$$

\subsection{Marginální pravděpodobnost proměnných}

Pokud chceme aktualizovat hodnotu naší aproximace marginální pravděpodobnosti,
je třeba minimalizovat její vzdálenost od skutečné marginální
pravděpodobnosti, ve tvaru
\begin{align}
p^*(\tilde{x}_j) &=
\sum_{\vec{x} \backslash x_j} \,
	\int_{\vec{\Theta}} \,
    		\prod_i \,
			q^{\backslash \beta}(x_i) \,
		\prod_l \,
			q^{\backslash \beta}(\vec{\theta}_l) \;
		f_\beta(\vec{x};\,
    		  \vec{\Theta})
\label{eq:1} 
\\
&=
\sum_{\vec{x} \backslash x_j} \,
	\prod_i \,
		q^{\backslash \beta}(x_i) \,
    \int_{\vec{\theta_{\rho(\vec{x^\prime})}}} \,
	    q^{\backslash \beta}(\vec{\theta}_{\rho(\vec{x^\prime})})\;
    \theta_{\rho(\vec{x^\prime}), x_0} \label{eq:2} 
\\
&= 
\sum_{\vec{x} \backslash x_j} \,
	\prod_i \,
		q^{\backslash \beta}(x_i)\;
    		\mathbb{E}_{q^{\backslash \beta}} 
			(\theta_{\rho(\vec{x^\prime}), x_0}).
\label{eq:3}
\end{align}

Rovnost~(\ref{eq:1}) vychází z~definice výpočtu marginální pravděpodobnosti ze
sdružené pravděpodobnosti.
V~(\ref{eq:2}) byla použita definice faktoru, z~integrálu byly vytknuty členy neobsahující $\Theta$ a nakonec bylo využito toho, že pro
množiny parametrů, které nejsou spojeny s~faktorem $f_\beta$, je jejich neúplná
distribuce rovna marginální distribuci, a tedy $\int_{\theta_i} q(\theta_i) =
1$. V~(\ref{eq:3}) byla použita definice očekávané hodnoty.

Marginální pravděpodobnost proměnné $x_j$ tedy je
\begin{equation}
p^*(x_j) =
\sum_{\vec{x} \backslash x_j} \;
	\prod_i \;
		q^{\backslash \beta}(x_i)\;
    		\mathbb{E}_{q^{\backslash \beta}} 
			(\theta_{\rho(\vec{x^\prime}), x_0}).
\end{equation}

Tady docházíme k~výsledku, který je velmi podobný výpočtu marginální
pravděpodobnosti v~Loopy Belief Propagation algoritmu.

Zprávu z~faktoru $f_\beta$ do vrcholu $x_j$ pak získáme vydělením zprávy z~$x_j$ do $f_\beta$ z~marginální pravděpodobnosti,

\begin{equation}
m_{f_\beta \rightarrow x_j}(x_j) =
    \sum_{\vec{x} \backslash x_j} \;
        \prod_{i \ne j}\;
            q^{\backslash \beta}(x_i)\;
            \mathbb{E}_{q^{\backslash \beta}}
                (\theta_{\rho(\vec{x^\prime}), x_0}).
\label{eq:msgfromftox}
\end{equation}

\subsection{Marginální pravděpodobnost parametrů}

Pro množiny parametrů se jejich marginální pravděpodobnost spočítá podobně jako
pro proměnné:
\begin{align}
p^*(\tilde{\vec{\theta}}_j) & = 
    \sum_{\vec{x}} \; 
    \int_{\vec{\Theta} \backslash \vec{\theta}_j} \;
         \prod_i \;
             q^{\backslash \beta}(x_i) \;
         \prod_l \;
    q^{\backslash \beta}(\vec{\theta}_l) \;
    f_\beta(\vec{x}; \vec{\Theta}) \label{eq:ep:theta_1}
\\
& = 
    \sum_{l \ne j} \;
        \sum_{\vec{x}: \rho(\vec{x^\prime}) = l} \;
            \prod_i \;
                q^{\backslash \beta}(x_i) \,
                \int_{\vec{\Theta} \backslash \vec{\theta}_j} \;
                    \prod_k \;
                        q^{\backslash \beta}(\vec{\theta}_k) \;
                        \theta_{l, x_0}\; + \label{eq:ep:theta_2}
\\
& + 
    \sum_{\vec{x}: \rho(\vec{x^\prime}) = j} \;
        \prod_i \;
            q^{\backslash \beta}(x_i) \;
            \int_{\vec{\Theta} \backslash \vec{\theta}_j} \;
                \prod_k \;
                    q^{\backslash \beta}(\vec{\theta}_k) \; 
                    \tilde{\theta}_{j, x_0}
\nonumber
\\
& = 
    \left[ \;
        \sum_{l \ne j} \;
            \sum_{\vec{x}: \rho(\vec{x^\prime}) = l} \;
                \prod_i \;
                    q^{\backslash \beta}(x_i) \;
                    \mathbb{E}_{q^{\backslash \beta}(\vec{\theta}_l)} (\theta_{l, x_0}) \;
    \right] \;
    q^{\backslash \beta}(\tilde{\vec{\theta}}_j) \; + \label{eq:ep:theta_3}
\\
& + 
    \sum_{\vec{x}: \rho(\vec{x^\prime}) = j} \;
        \prod_i \;
            q^{\backslash \beta}(x_i) \;
            \tilde{\theta}_{j,x_0} \;
            q^{\backslash \beta}(\tilde{\vec{\theta}}_j) \;
\label{eq:ep:theta_4}
\nonumber
\\
& = w_0 \; q^{\backslash \beta}(\tilde{\vec{\theta}}_j) \, + \, \sum_k \; w_k \;
    \tilde{\theta}_{j,k} \; q^{\backslash \beta}(\tilde{\vec{\theta}}_j),
\end{align}

kde

\begin{align}
w_0 &=
	\sum_{l \ne j} \;
		\sum_{\vec{x}: \rho(\vec{x^\prime}) = l} \;
			\prod_i \;
 				q^{\backslash \beta}(x_i) \;
				\mathbb{E}_{q^{\backslash \beta}(\vec{\theta}_l)} (
					\theta_{l,
					x_0}) 
\\
w_k &=
	\sum_{\vec{x}: \rho(\vec{x^\prime}) = j, x_0 = k} \; 
		\prod_i \;
			q^{\backslash \beta}(x_i)		
\end{align}

Opět vycházíme z~výpočtu marginální pravděpodobnosti podle definice.
V~rovnici \eqref{eq:ep:theta_2} jsme rozdělili sumu přes
$\vec{x}$ na ty, pro které se ve faktoru použije množina parametrů
$\tilde{\vec{\theta}}_j$, a na ty ostatní. Také jsme z~integrálu vytknuli součin
neúplných distribucí pro proměnné. V~dalším kroku \eqref{eq:ep:theta_3} jsme opět
použili toho, že integrál přes $\vec{\Theta}$ je ve skutečnosti několik
integrálů přes jednotlivé množiny parametrů. Díky tomu je můžeme vložit mezi členy produktu neúplných distribucí pro množiny parametrů. Ve výsledku
získáme $q^{\backslash \beta}(\tilde{\vec{\theta}}_j) \, \int_{\vec{\theta}_l} \,
q^{\backslash \beta}(\vec{\theta}_l) \; \theta_{l, x_0}$ a pak zbylé členy, které zmizí.

Docházíme k~vyjádření skutečné marginální pravděpodobnosti, v~níž není třeba
integrovat přes všechny množiny parametrů, ale stačí jen očekávaná hodnota
těchto parametrů.

\section{Aproximace marginálních pravděpodobností}

Při výpočtu aproximující distribuce $q(\vec{\theta}_j)$ se může stát, že skutečná marginální distribuce je v~komplikovaném tvaru a pak je aproximace složitá na výpočet a může se odchylovat od skutečné distribuce.
Například se jedná o~směs distribucí.
Ideální případ je takový, kdy aproximace je stejného typu jako skutečná distribuce.

Pro zjednodušení výpočtů, jak se dále ukáže, zvolíme zprávy z~faktoru do množiny parametrů, $m_{f \rightarrow \vec{\theta}_i}(\vec{\theta}_i)$, ve tvaru Dirichletovské distribuce s~parametry $\vec\alpha_{f \rightarrow \vec\theta_i}$:
\begin{equation}
m_{f \rightarrow \vec{\theta}_i}(\vec{\theta}_i) =
	Dir(\vec{\theta}_i; \vec\alpha_{f \rightarrow \vec\theta_i}) =
        \frac{\Gamma (\sum_j \vec\alpha_{f \rightarrow \vec\theta_i, j})}
             {\prod_j \Gamma(\vec\alpha_{f \rightarrow \vec\theta_i, j})}
        \prod_j \theta_{i,j}^{\vec\alpha_{f \rightarrow \vec\theta_i, j} - 1},
\end{equation}
kde $\Gamma$ je Gamma funkce (zobecnění faktoriálu):
\begin{equation}
    \Gamma(z) = \int_0^\infty \! t^{z-1} \exp(-t) \;\mathrm{d}t.
\end{equation}

Dirichletovská distribuce byla zvolena, protože má důležité vlastnosti pro
součin, které budou využity dále pro výpočet neúplné distribuce a celkové
aproximace. Pokud označíme aproximované faktory indexem $\beta$ a každý bude
mít vlastní parametry $\vec\alpha_{f_\beta \rightarrow \vec\theta_i}$, tak výsledná aproximace bude
tvaru:
\begin{align}
q(\vec{\theta}_i) \,& \, \propto
    \prod_\beta \;
    m_{f_\beta \rightarrow \vec{\theta}_i}(\vec{\theta}_i) \,
 \propto \,
    \prod_\beta \;
        \prod_j \;
            \theta_{i,j}^{\vec\alpha_{f_\beta \rightarrow \vec\theta_i, j} \,-\, 1}
\\
&\propto \,
    Dir(\vec{\theta}_i ;\, \sum_\beta \vec\alpha_{f_\beta \rightarrow \vec\theta_i} - (\,|\beta| - 1\,) \, \vec{1})
\\
&=
    Dir(\vec{\theta}_i;\,
        \vec{\alpha}_i),
\label{eq:ep:aprx_1}
\end{align}
kde $\vec{\alpha}_i = \sum_\beta \vec\alpha_{f_\beta \rightarrow \vec\theta_i} - (\,|\beta| - 1\,)
\vec{1})$.

Při aktualizaci faktoru $\tilde\beta$ tedy neúplná distribuce bude
\begin{align}
q^{\backslash \tilde\beta} (\vec{\theta_i})
&\propto
    \prod_{\beta \ne \tilde\beta}
        m_{f_\beta \rightarrow \vec{\theta}_i}(\vec{\theta}_i)
\\
& \propto
Dir(\vec{\theta}_i;
    \vec{\alpha}_i - \vec\alpha_{f_\beta \rightarrow \vec\theta_i} + \vec{1}).
\end{align}

Naším cílem je nalézt parametry $\vec{\alpha}^*$ aproximované marginální pravděpodobnosti \eqref{eq:ep:aprx_1}, které minimalizují vzdálenost od skutečné marginální pravděpodobnosti \eqref{eq:ep:theta_4}. 

Jednotlivé faktory musíme vždy aproximovat s~ohledem na všechny ostatní.
U~našeho grafického modelu využijeme nezávislostí, a tak nám stačí pouze pracovat se zprávami ze sousedních vrcholů.
Všechny informace se v~modelu šíří přes zprávy mezi vrcholy, a proto je třeba několika iterací pro přesun informace z~jedné části modelu do druhé.

Pro měření vzdálenosti mezi dvěma pravděpodobnostními rozloženími se používá Kullback-Leiblerova divergence:
\begin{equation}
KL(p \| q) =
\int_{-\infty}^{\infty}
    p(x) \log\left(\frac{p(x)}{q(x)}\right) \mathrm{d}x
\end{equation}
Pro nalezení minima použijeme algoritmus srovnání momentů a budeme
minimalizovat $KL(p^*\| q)$.

Pokud se podíváme na skutečnou marginální pravděpodobnost
$p^*(\vec{\theta}_i)$, zjistíme, že můžeme některé její členy upravit.
Využijeme také vlastnosti Gamma funkce $\Gamma(x) = (x-1) \Gamma(x-1)$.

\begin{align}
w_j \, \theta_j \, Dir(\vec{\theta}; \vec{\alpha}) &\propto \,
    w_j \, \theta_j \, \frac{\Gamma (\, \sum_i \, \alpha_{i} \,)}{\prod_i \, \Gamma(\,\alpha_i\,)}\;
    \prod_i \, \theta_{i}^{\alpha_{i} - 1}
\label{eq:4}
\\
&\propto w_j \, \frac{\Gamma (\,\sum_i \, \alpha_{i}\,)} {\prod_i \,
    \Gamma(\,\alpha_{i}\,)} \; \theta_j^{\alpha_j} \prod_{i \ne j} \,
    \theta_{i}^{\alpha_{i} - 1} \\
&\propto w_j \,
    \frac{\Gamma (\,\sum_i \alpha_{i}\,)}
         {\prod_i \, \Gamma(\, \alpha_{i} \,)} \;
    \frac{\Gamma(\, \alpha_{j} + 1 \,) \; \prod_{i \ne j} \, \Gamma(\,\alpha_i\,)}
         {\Gamma (\,1 + \sum_i \alpha_i\,)} \;
    Dir(\vec\theta; \, \vec{\alpha} + \vec{\delta}_j) \\
&\propto w_j \,
    \frac{\Gamma (\,\sum_i \alpha_{i}\,)}
         {\prod_i \, \Gamma(\, \alpha_{i}\,)} \;
    \frac{\alpha_j \, \Gamma(\,\alpha_j\,) \, \prod_{i \ne j} \, \Gamma(\,\alpha_i\,)}
         {(\,\sum_i \, \alpha_i\,) \, \Gamma (\,\sum_i \,\alpha_i\,)}\;
    Dir(\vec\theta; \,\vec{\alpha} + \vec{\delta}_j) \\
&\propto w_j \,
    \frac{\alpha_j}
         {\sum_i \, \alpha_i} \;
    Dir(\vec\theta; \,\vec{\alpha} + \vec{\delta}_j) \\
\end{align}

Zde jsme použili $\vec{\delta}_i$, což je vektor nul, který má pouze na pozici $i$ jedničku.

Díky této úpravě lze $p^*$ vyjádřit jako směs Dirichletovských rozdělení.
\begin{equation}
p^*(\vec{\theta}) =
    w_0^* \; Dir(\vec{\theta};\, \vec{\alpha}) +
    \sum_j \, w^*_j \;
        Dir(\vec{\theta};\, \vec{\alpha} + \vec{\delta}_j),
\label{eq:marginaltheta}
\end{equation}
kde
\begin{align}
    w^*_0 \propto w_0, \quad \quad
    w^*_j \propto w_j \, \frac{\alpha_j}{\sum_i \alpha_i}, \quad \quad
    \sum_{i=0}^k w_i^* = 1.
\end{align}

Jak bylo ukázáno v~sekci~\ref{sec:expfam}, pro minimalizaci KL divergence mezi dvěma rozděleními z~exponenciální rodiny
stačí, pokud se budou rovnat jejich postačující statistiky. Dokážeme jednoduše
spočítat první dva momenty Dirichletovského rozdělení, a proto použijeme
aproximaci a budeme počítat pouze s~nimi a zbylé momenty zanedbáme.
Je tedy třeba nalézt střední hodnotu a rozptyl proměnných z~$p^*(\vec\theta)$:

\begin{align}
\mathbb{E}_{p^*}[\vec\theta] &= \int \vec\theta \, p^*(\vec\theta) ~ \mathrm{d}\vec\theta
\\
&= \int \vec\theta \, (
	w_0^* \, Dir(\vec{\theta};\, \vec{\alpha}) +
	\sum_j \, w^*_j \,
        	Dir(\vec{\theta};\, \vec{\alpha} + \vec{\delta}_j)
    ) ~ \mathrm{d}\vec\theta
\\
&= w_0^* \int \vec\theta \, Dir(\vec{\theta};\, \vec{\alpha}) ~ d\vec\theta +
	\sum_j w_j^* \int \vec\theta \, Dir(\vec{\theta};\, \vec{\alpha} + \vec{\delta}_j)
    ~ \mathrm{d}\vec\theta
\\
&= w_0^* \, \mathbb{E}_{Dir(\vec\alpha)}[\vec\theta] +
	\sum_j \, w_j^* \, \mathbb{E}_{Dir(\vec\alpha + \vec{\delta}_j)}[\vec\theta]
\label{eq:epweimix}
\end{align}
Střední hodnotu proměnných $\vec\theta$ podle rozdělení $p^*$ lze podle (\ref{eq:epweimix}) spočítat
jako vážený součet středních hodnot $\vec\theta$ podle jednotlivých
Dirichletovských distribucí, z~kterých se $p^*$ skládá.

První moment tedy máme spočítaný, pro výpočet rozptylu můžeme využít
definici:
\begin{equation}
	Var[\vec{\theta}] = \mathbb{E}[\vec{\theta}^2] - \mathbb{E}[\vec\theta]^2.
\end{equation}

Chybí nám tedy ještě výpočet střední hodnoty druhé mocniny proměnné
$\vec\theta$ podle $p^*$. Můžeme ji vyjádřit z~definice střední hodnoty:

\begin{align}
\mathbb{E}_{p^*}[\vec\theta^2] &= \int \vec\theta^2 \, p^*(\vec\theta) ~ d\vec\theta
\\
&= w_0^* \int \vec\theta^2 \, Dir(\vec{\theta};\, \vec{\alpha}) ~ d\vec\theta +
	\sum_j w_j^* \int \vec\theta^2 \, Dir(\vec{\theta};\, \vec{\alpha} + \vec{\delta}_j)
	 ~ d\vec\theta
\\
&= w_0^* \, \mathbb{E}_{Dir(\vec\alpha)}[\vec\theta^2] +
	\sum_j w_j^* \, \mathbb{E}_{Dir(\vec\alpha + \vec{\delta}_j)}[\vec\theta^2].
\label{eq:epweimix2}
\end{align}

Opět získáváme vážený součet středních hodnot podle Dirichletovských rozdělení.
Střední hodnotu druhé mocniny proměnné podle Dirichletovského rozdělení
lze opět jednoduše odvodit z~definice:
\begin{align}
\mathbb{E}_{Dir(\vec\alpha)}[x_i^2] &=
	\int x_i^2 \, Dir(\vec x; \, \vec\alpha) \; \mathrm{d}\vec x
\\
&=
	\int x_i^2 \, \frac{\Gamma(\,\alpha_0\,)}
			   {\prod_{j=1}^N \, \Gamma(\,\alpha_j\,)} \;
		\prod_{j=1}^N x_j^{\alpha_j - 1} \; \mathrm{d}\vec x.
\end{align}

Nyní jsme v~podobné situaci jako v~(\ref{eq:4}). Budeme postupovat stejně,
pokusíme se výraz přepsat do formy Dirichletovské distribuce nad $x_i$:

\begin{align}
\mathbb{E}_{Dir(\vec\alpha)}[x_i^2] &=
\int \frac{\Gamma(\,\alpha_0 + 2\,) \; \alpha_i \; (\,\alpha_i + 1\,)}
		{\alpha_0 \; (\,\alpha_0 + 1\,) \; \Gamma(\,\alpha_i + 2\,) \; \prod_{j \ne i} \; \Gamma(\,\alpha_j\,)} \;
    x_i^{\alpha_i + 1} \, \prod_{j \ne i} \, x_j^{\alpha_j - 1} \;
\mathrm{d}\vec x
\\
&= \frac{\alpha_i \; (\,\alpha_i + 1\,)}
		{\alpha_0 \; (\,\alpha_0 + 1\,)} \,
	\int \frac{\Gamma(\,\beta_0\,)}
		{\prod_i \,\Gamma(\,\beta_i\,)} \;
        \prod_i \, x_i^{\beta_i - 1} \;
    \mathrm{d}\vec x
\\
&= \frac{\alpha_i \; (\,\alpha_i + 1\,)}
		{\alpha_0 \; (\,\alpha_0 + 1\,)}
	\int Dir(\vec x;\, \vec\beta)\;
    \mathrm{d}\vec x
\\
&= \frac{\alpha_i \; (\,\alpha_i + 1\,)}
		{\alpha_0 \; (\,\alpha_0 + 1\,)}.
\end{align}

Vyjádřili jsme $\Gamma(\,\alpha_0\,)$ a $\Gamma(\,\alpha_i\,)$ s~pomocí $\Gamma(\,\alpha_0 + 2\,)$ a $\Gamma(\,\alpha_i + 2\,)$:
\begin{align}
\Gamma(\alpha_0) = \frac{\Gamma(\,\alpha_0 + 2\,)}
				{\alpha_0\;(\,\alpha_0 + 1\,)},
\quad \quad
\Gamma(\alpha_i) = \frac{\Gamma(\,\alpha_i + 2\,)}
				{\alpha_i \; (\,\alpha_i + 1\,)}.
\end{align}

Následně jsme vytvořili nové parametry $\vec\beta$:
\begin{align}
\beta_i = \alpha_i + 2, \quad \quad
\beta_{j \ne i} = \alpha_j, \quad \quad
\beta_0 = \sum_i \beta_i.
\end{align}

Nyní tedy dokážeme spočítat $\mathbb{E}_{p^*}[\vec\theta]$ a $\mathbb{E}_{p^*}[\vec\theta^2]$.
Parametry aproximovaného rozdělení nalezneme následovně
\begin{align}
\frac{\mathbb{E}[X_1] - \mathbb{E}[X_1^2]}
     {\mathbb{E}[X_1^2] - \mathbb{E}[X_1]^2} &=
\frac{
	\frac{\alpha_1}
		{\alpha_0}
	- \frac{\alpha_1(\alpha_1 + 1)}
		{\alpha_0(\alpha_0 + 1)}
	}
	{
	\frac{\alpha_1(\alpha_1 + 1)}
		{\alpha_0(\alpha_0 + 1)}
	- \frac{\alpha_1^2}
		{\alpha_0^2}
	}
\label{eq:5}
\\
&=
\frac{
    \frac{\alpha_1 (\alpha_0 + 1) - \alpha_1(\alpha_1 + 1)}
         {\alpha_0 (\alpha_0 + 1)}
}{
    \frac{\alpha_0 \alpha_1 (\alpha_1 + 1) - \alpha_1^2 (\alpha_0 + 1)}
         {\alpha_0^2 (\alpha_0 + 1)}
}
\\
&=
\frac{\alpha_0 \alpha_1 (\alpha_0 - \alpha_1)}
     {\alpha_1 (\alpha_0 \alpha_1 + \alpha_0 - \alpha_0 \alpha_1 - \alpha_1)}
\\
&=
\alpha_0
\\
\alpha_i &= \mathbb{E}[X_i] \;\alpha_0 \label{eq:6}
\end{align}

Z~rovnice~(\ref{eq:5}) vypočítáme sumu všech parametrů $\alpha_0$. 
Protože střední hodnota proměnné z~Dirichletovského rozdělení je právě
$\frac{\alpha_i}{\alpha_0}$, tak jednotlivé parametry získáme z~rovnice~(\ref{eq:6}).

\section{Algoritmus}

V~algoritmu~\ref{alg:ep} je předveden algoritmus Expectation Propagation upravený pro potřeby učení parametrů.

\begin{algorithm}
\caption{Expectation Propagation pro učení parametrů}
\label{alg:ep}
\begin{algorithmic}
\State Parametry zpráv z~faktoru $\beta$ do množiny parametrů
    $\vec{\theta}_i$ označíme $\vec\alpha_{f_\beta \rightarrow \vec\theta_i}$.
\State Parametry zpráv z~množiny parametrů $\vec{\theta}_i$ do faktoru $\beta$
    označíme $\vec\alpha_{ \vec\theta_i \rightarrow f_\beta}$.
\State Parametry marginální distribuce množiny parametrů $\vec{\theta}_i$
    označíme $\vec{\alpha}_i$.
\Init
    \State Nastav zprávy mezi faktory a proměnnými na $1$.
    \State Nastav parametry $\vec\alpha_{f_\beta \rightarrow \vec\theta_i}$ na $1$.
    \State Nastav parametry $\vec{\alpha}_i$ na apriorní hodnotu.
\EndInit
\Repeat
    \State Vyber faktor $f_{\tilde\beta}$, který se bude aktualizovat.

    \State Spočítej všechny zprávy z~parametrů:
    \For{každý parametr $\vec{\theta}_i$ spojený s~faktorem $f_{\tilde\beta}$}
        \State Parametry zprávy z~$\vec{\theta}_i$ do $f_{\tilde\beta}$:
            $\vec\alpha_{ \vec\theta_i \rightarrow f_{\tilde\beta}} = \vec{\alpha}_i -
            \vec\alpha_{f_{\tilde\beta} \rightarrow \vec\theta_i } + 1$.
    \EndFor

    \State Aktualizuj zprávy z~faktoru do proměnných:
    \For{každou proměnnou $X_i$, spojenou s~faktorem $f_{\tilde\beta}$}
    \State Zpráva z~$f_{\tilde\beta}$ do $X_i$ podle (\ref{eq:msgfromftox}):
    \State
        \hskip\algorithmicindent
        $m_{f_{\tilde{\beta}} \rightarrow x_j} =
            \sum_{\vec{x} \backslash x_j}
                \mathbb{E}_{q^{\backslash\tilde\beta}}(\,\theta_{\rho(\vec{x^\prime}), x_0}\,)\,
                \prod_{i \ne j}\,
                    q^{\backslash\tilde\beta}(x_i)$
    \EndFor

    \State Aktualizuj marginální pravděpodobnost parametrů:
    \For{každý parametr $\vec{\theta}_i$, spojený s~faktorem $f_{\tilde\beta}$}
        \State Spočítej parametry $\vec{\alpha}_i^*$ pro Dirichletovské
            rozdělení, které nejlépe
        \State aproximuje cílovou marginální distribuci
            (\ref{eq:marginaltheta}). Metoda popsána
        \State v~předchozí sekci.
        \State Parametry zprávy z~$f_{\tilde\beta}$ do $\vec{\theta}_i$:
        \State
            \hskip\algorithmicindent
            $\vec\alpha_{f_{\tilde\beta} \rightarrow \vec\theta_i} =
                \vec{\alpha}_i^*
                - \vec\alpha_{ \vec\theta_i \rightarrow f_{\tilde\beta}}
                + 1$
        \State Aktualizuj parametry marginální distribuce $q(\vec{\theta}_i)$
        \State
            \hskip\algorithmicindent
            $\vec{\alpha}_i = \vec{\alpha}_i^* =
                \vec\alpha_{f_{\tilde\beta} \rightarrow \vec\theta_i}
                + \vec\alpha_{ \vec\theta_i \rightarrow f_{\tilde\beta}}$
    \EndFor
    \For{každou proměnnou $X_i$, spojenou s~faktorem $f_{\tilde\beta}$}
        \State Aktualizuj zprávy z~proměnných do faktoru:
        \State
            \hskip\algorithmicindent
            $q^{\backslash\tilde\beta}(x_i) =
                \prod_{\beta \ne \tilde\beta} \, m_{f_{\beta} \rightarrow x_i}(x_i)$
    \EndFor
\Until{konvergence}
\end{algorithmic}
\end{algorithm}

% !TEX root = thesis.tex
\chapter{Implementace}
\label{ch:kap4}

V~předchozích kapitolách byly popsány teoretické základy nutné pro implementaci inference v~Bayesovských sítích.
V~této kapitole bude popsána vytvořená knihovna pro dialogové systémy.
Celá knihovna se skládá z~několika vrstev, každá z~nich staví na předchozí a poskytuje větší úroveň abstrakce.
Nejnižší vrstva implementuje efektivní počítání s~faktory.
Nad ní stojí vrstva reprezentující jednotlivé vrcholy ve faktor grafu.
Tato vrstva také obsahuje funkcionalitu pro počítání a~posílání zpráv.
Nejvyšší vrstva se zaměřuje na samotnou inferenci v~grafickém modelu.

Knihovna byla implementována v jazyku Python 2.7 s pomocí knihoven Numpy~\cite{oliphant-2006-guide} a Scipy~\cite{scipy}. 
Jedná se o jediné závislosti, které je třeba mít nainstalované pro funkčnost knihovny.

\section{Existující řešení}

Většina existujících knihoven pro inferenci v~Bayesovských sítích je vytvořena pro obecné grafické modely.
Použití v~dialogových systémech se ovšem liší od předpokládaného využití většiny knihoven.
Při řešení problémů je jen málokdy potřeba získat výsledek v~řádu milisekund, i na úkor kvality aproximace.
Ovšem při hovoru s~uživatelem není možné si dovolit dlouhé pauzy, protože působí rušivě.
Častým problémem pak u~těchto knihoven je nemožnost volby aproximace.

Populární knihovna, která umožňuje použití Bayesovských modelů pro širokou škálu problémů, je Infer.Net, pocházející z~Microsoft Research~\cite{InferNET12}.
Tato knihovna nabízí velkou řadu inferenčních algoritmů, např. Expectation Propagation, Variational Message Passing nebo Gibbsovo vzorkování.
Jak název napovídá, je napsaná v~prostředí .Net, a tedy je nutné ji používat z~jazyka, který běží v~tomto prostředí (C\#, F\# atd.)

Další knihovny jsou například Bayes Net Toolbox pro Matlab (BNL)~\cite{bnl} nebo Probabilistic Network Library (PNL, \url{http://sourceforge.net/projects/openpnl/}).
Většinou obsahují obecné inferenční algoritmy, ale omezený počet podporovaných distribucí a podpora je diskutabilní.
V případě BNT byla poslední verze vydána v roce 2007 a u PNL už neexistuje oficiální web na stránkách firmy Intel, která knihovnu vytvořila.

Dále existují knihovny používající místo variačních metod vzorkování.
Příkladem je knihovna PyMC~\cite{patil2010pymc}. 
Díky využití metody Monte Carlo je jednodušší počítání se složitějšími distribucemi, ovšem na druhou stranu platíme mnohdy pomalou konvergencí.

\section{Diskrétní faktor}

Faktor je základním stavebním kamenem.
Operace s~faktory musí být efektivní, pro výpočet jedné zprávy je potřeba několik násobení faktorů, následně je třeba je marginalizovat atd.
Faktory je třeba úsporně reprezentovat.
Každá zpráva, každé pravděpodobnostní rozdělení je implementováno pomocí faktoru.
Faktory jsou reprezentovány třídou \texttt{Factor} v~modulu \texttt{alex.infer.factor}, která podporuje všechny základní matematické operace a také řadu specifických operací jako je marginalizace, normalizace, nastavení faktoru na pozorovanou hodnotu, vybrání nejpravděpodobnějšího přiřazení atd.

\subsection{Reprezentace faktorů}
\label{sec:repfak}

Každý diskrétní faktor má seznam diskrétních proměnných, které tvoří jeho doménu.
Každá z~těchto proměnných může mít jinou kardinalitu, některé proměnné jsou binární, jiné mají mnohem více hodnot.
Faktor je tedy ve své podstatě multidimenzionální tabulka.
Pro implementaci je tato tabulka zploštěná do jednoduchého pole.

Pole využívají implementaci z~knihovny Numpy, díky ní jsou matematické operace napsané v~rychlejším jazyku než je Python (C, Fortran) a jsou navíc vektorizované.

Jak je tedy možné reprezentovat multidimenzionální tabulku jednodimenzionálním polem?
Proměnné jsou seřazené a pro zjednodušení můžeme předpokládat, že hodnoty proměnných jsou čísla z~$\{0, \dots, n-1\}$, kde $n$ je kardinalita proměnné.
Každá hodnota v~poli je pak hodnotou faktoru pro nějaké přiřazení proměnných, např. $(0, 1, 0)$.
Jednotlivá přiřazení jsou v~poli seřazena lexikograficky.
Pro každou proměnnou si pamatujeme její kardinalitu a také tzv. krok.
Krok určuje o~kolik hodnot v~tabulce se musíme posunout, abychom se dostali na další hodnotu této proměnné se zachováním hodnot všech následujících proměnných.

Příklad faktoru je v~tabulce \ref{tab:stride}.
Doménu faktoru tvoří dvě binární proměnné $X$ a $Y$, jejich kardinalita je tedy 2.
Pro proměnnou $X$ je krok 2, pro proměnnou $Y$ je krok 1.

\begin{table}
\begin{center}
\begin{tabular}{|c|c|r|}
\hline
$X$ & $Y$ & Hodnota \\
\hline
\hline
0 & 0 & 0.2 \\
\hline
0 & 1 & 0.3 \\
\hline
1 & 0 & 0.1 \\
\hline
1 & 1 & 0.4 \\
\hline
\end{tabular}
\end{center}
\caption{Příklad faktoru s~dvěma proměnnými $X$ a $Y$.}
\label{tab:stride}
\end{table}

\subsection{Operace s~faktory}

Implementace diskrétního faktoru obsahuje všechny základní matematické operace i operace určené speciálně pro pravděpodobnostní distribuce, např. marginalizace nebo normalizace.
Každý výpočet zprávy se skládá z několika násobení, dělení, nakonec marginalizace.
Je důležité, aby tyto základní operace byly napsány co nejefektivněji.

Matematické operace s~faktory musí fungovat ve třech různých situacích (příklady uvedeme na násobení):
\begin{enumerate}
    \item násobení faktoru s~faktorem, oba se stejnou doménou,
    \item násobení dvou faktorů, které sdílí jen některé proměnné,
    \item násobení faktoru konstantou.
\end{enumerate}

Násobení faktoru konstantou je triviální, každá položka faktoru bude vynásobena konstantou.
Tato operace může být jednoduše vektorizována.

Při násobení dvou faktorů se stejnou doménou je třeba pronásobit prvky se stejným přiřazením proměnných.
Což znamená pronásobit hodnoty na stejných místech v~poli.
Opět se tedy jedná o~operaci, která je jednoduše vektorizovatelná.

\subsubsection{Operace s~různými doménami}

Poslední možnost je, že se snažíme provést matematickou operaci s~dvěma faktory, které nemají stejnou doménu.
Pak musí být výsledkem  nový faktor, jehož doména je sjednocením domén vstupních faktorů.
Jednotlivé prvky nového faktoru jsou pak výsledkem aplikace operace na prvky ze vstupních faktorů sdílejících přiřazení společných proměnných.
Příklad násobení takových dvou faktorů je v~tabulce~\ref{tab:facmul}.

\begin{table}
\begin{equation*}
\begin{array}{|c|c|r|}
    \hline
    \multicolumn{3}{|c|}{f_1} \\
    \hline
    X & Y & Hodnota \\
    \hline
    \hline
    0 & 0 & 0.2 \\
    \hline
    0 & 1 & 0.3 \\
    \hline
    1 & 0 & 0.1 \\
    \hline
    1 & 1 & 0.4 \\
    \hline
\end{array}
\times
\begin{array}{|c|c|r|}
    \hline
    \multicolumn{3}{|c|}{f_2} \\
    \hline
    Y & Z~& Hodnota \\
    \hline
    \hline
    0 & 0 & 0.2 \\
    \hline
    0 & 1 & 0.2 \\
    \hline
    1 & 0 & 0.2 \\
    \hline
    1 & 1 & 0.4 \\
    \hline
\end{array}
=
\begin{array}{|c|c|c|r|}
    \hline
    \multicolumn{4}{|c|}{f_r} \\
    \hline
    X & Y & Z~& Hodnota \\
    \hline
    \hline
    0 & 0 & 0 & 0.04 \\
    \hline
    0 & 0 & 1 & 0.04 \\
    \hline
    0 & 1 & 0 & 0.06 \\
    \hline
    0 & 1 & 1 & 0.12 \\
    \hline
    1 & 0 & 0 & 0.02 \\
    \hline
    1 & 0 & 1 & 0.02 \\
    \hline
    1 & 1 & 0 & 0.08 \\
    \hline
    1 & 1 & 1 & 0.16\\
    \hline
\end{array}
\end{equation*}
\caption{Příklad násobení dvou faktorů, které sdílí proměnnou.}
\label{tab:facmul}
\end{table}

Výsledek násobení faktorů $f_1$ a $f_2$ je ve faktoru $f_r$.
Faktory sdílí pouze proměnnou $Y$, takže je potřeba pronásobit všechna přiřazení z~$f_1$ se všemi přiřazeními z~$f_2$, která mají stejnou hodnotu $Y$.
Příkladem je přiřazení $(0, 1)$ s~hodnotou $0.3$ vynásobené s~přiřazením $(1, 1)$ s~hodnotou $0.4$.
Výsledek je uložen ve faktoru $f_r$ s~přiřazením $(0, 1, 1)$ a správnou hodnotou $0.3 \cdot 0.4 = 0.12$.

\subsection{Algoritmus pro operace s~různými doménami}

Předvedli jsme možné případy operací s~faktory a ukázali, že dva ze tří jsou triviální na implementaci.
Nyní představíme efektivní implementaci třetí možnosti, tedy aplikaci operace na dva faktory s~různými doménami (algoritmus \ref{alg:apop}).

Ze vstupních faktorů vytvoříme prázdný faktor pro výsledek.
Jeho doména je sjednocením domén vstupních faktorů.
Kardinalita proměnných zůstává stejná.
Krok jednotlivých proměnných je třeba přepočítat.
Spočítáme jej jako součin kardinalit proměnných, které následují po té aktuální.
Velikost pole pro všechny hodnoty je rovna součinu všech kardinalit.

Následně přistoupíme k~vyplňování tabulky.
Pro oba vstupní faktory si budeme udržovat index pozice s~přiřazením, které odpovídá aktuálně vyplňovanému přiřazení ve výsledném faktoru.
Po provedení operace tyto indexy aktualizujeme.

Pokud reprezentujeme přiřazení proměnných jako číslo (kde každá cifra může mít jinou kardinalitu), pak se přesuneme k~dalšímu ohodnocení v~řadě tak, že zvýšíme nejméně signifikantní cifru (proměnnou) o~jedna.
Může se stát, že jsme dosáhli kardinality dané proměnné a pak se musíme vrátit na ohodnocení 0 a~aplikovat přesun na vyšší cifru.
Opakovanou aplikací přesunu můžeme upravit všechny proměnné, příkladem je přechod z~přiřazení $(0, 1, 1)$ na $(1, 0, 0)$, všechny proměnné binární.
Při každé úpravě proměnné také aktualizujeme indexy do vstupních faktorů.

Při přesunu na další hodnotu proměnné stačí k~indexu do faktoru přičíst krok upravené proměnné v~tomto faktoru.
Pokud je třeba se vrátit u~proměnné $v$~na ohodnocení 0, pak od indexu pro vstupní faktor odečteme $(c_v - 1) \cdot s_v$, kde $c_v$ je kardinalita proměnné $v$ a $s_v$ je krok proměnné $v$.
Pro proměnnou, která není obsažena v~doméně faktoru, je její krok roven 0.
Celá metoda pro aplikaci matematické operace je v~algoritmu~\ref{alg:apop}.

\begin{algorithm}
\caption{Aplikace matematické operace na faktory s~různými doménami}
\label{alg:apop}
\begin{algorithmic}
\Function{Apply-Op}{$f_1$, $f_2$, $op$}
\State $f_1, f_2$~-- vstupní faktory
\State $op$~-- operace
\State
\State $f_r \gets$ nový faktor pro výsledek operace na $f_1$ a $f_2$ s~prázdným polem
\State
\State $\text{index}[f_1] \gets 0$
\State $\text{index}[f_2] \gets 0$
\State $\text{přiřazení}[v] \gets 0 \text{ pro každou proměnnou } v \in \text{proměnné}[f_r]$
\State
\For{$i \in \{0, \dots, \textsc{Length}(\text{pole}[f_r])-1\}$}
	\State $\text{pole}[f_r][i] = op(\text{pole}[f_1][\text{index}[f_1]], \text{pole}[f_2][\text{index}[f_2]])$
	\For{$v \in \textsc{Reversed}(\text{proměnné}[f_r])$}
		\State $\text{přiřazení}[v] \mathrel{+}= 1$
		\If{\text{přiřazení}[v] = \text{kardinalita}[v]}
			\State $\text{přiřazení}[v] \gets 0$
			\State $\text{index}[f_1] \mathrel{-}= (\text{kardinalita}[v] - 1) \cdot \text{krok}[f_1][v]$
			\State $\text{index}[f_2] \mathrel{-}= (\text{kardinalita}[v] - 1) \cdot \text{krok}[f_2][v]$
		\Else
			\State $\text{index}[f_1] \mathrel{+}= \text{krok}[f_1][v]$
			\State $\text{index}[f_2] \mathrel{+}= \text{krok}[f_2][v]$
			\State \textbf{break}
		\EndIf
	\EndFor
\EndFor
\State \Return $f_r$
\EndFunction
\end{algorithmic}
\end{algorithm}

Tento algoritmus pracuje přímo nad polem a jeho výhodou je, že nepotřebuje ke svému výpočtu znát jednotlivá přiřazení proměnných.
Na rozdíl třeba od algoritmu, který by generoval všechny možné kombinace přiřazení a pak pomocí nich přistupoval k~odpovídajícím hodnotám ve faktorech, zde pracujeme pouze s~indexy do polí a používáme jednoduché matematické operace.

\subsection{Marginalizace proměnných ve faktoru}

Další důležitou a často používanou operací při počítání s~faktory je marginalizace.
Na vstupu je faktor a podmnožina proměnných z~domény faktoru, která má zůstat po vysčítání zbytku.
Algoritmus je podobný aplikaci matematické operace z~předchozí sekce.
V~tomto případě procházíme původní faktor a každou hodnotu přičteme na správnou pozici v~novém faktoru.
Pro aktualizaci indexu musíme projít všechny zůstávající proměnné a u~každé zkontrolovat, zda-li se v~následujícím kroku změní.
Případně patřičně upravit indexy.
Postup marginalizace je zapsán v~algoritmu~\ref{alg:marg}.

\begin{algorithm}
\caption{Marginalizace faktoru}
\label{alg:marg}
\begin{algorithmic}
\Function{Marginalize}{$f$, $vars$}
\State $f$~-- vstupní faktor
\State $vars$~-- seznam proměnných, které mají zůstat
\State
\State $f_r \gets$ nový faktor, obsahující pouze proměnné z~$vars$
\State $\text{přiřazení}[v] \gets 0 \text{ pro každou proměnnou } v \in vars$
\State $index \gets 0$ \Comment{Index do nového faktoru}
\State
\For{$i \in \{0, \dots, \textsc{Length}(\text{pole}[f])-1\}$}
	\State $\text{pole}[f_r][index] \mathrel{+}= \text{pole}[f][i]$
	\State
	\For{$v \in vars$}
		\If{$(i + 1) \bmod \text{krok}[f][v] = 0$}
			\State $\text{přiřazení}[v] \mathrel{+}= 1$
			\State $index \mathrel{+}= \text{krok}[f_r][v]$
		\EndIf
		\State
		\If{$\text{přiřazení}[v] = \text{kardinalita}[v]$}
			\State $\text{přiřazení}[v] \gets 0$
			\State $index \mathrel{-}= \text{kardinalita}[v] \cdot \text{krok}[f_r][v]$
		\EndIf
	\EndFor
\EndFor
\State \Return $f_r$
\EndFunction
\end{algorithmic}
\end{algorithm}

\section{Vrcholy faktor grafu}

Předchozí sekce hovořila o~práci s~faktory, tato sekce se bude zabývat implementací jednotlivých vrcholů ve faktor grafu.
Vrcholy se dělí na vrcholy pro proměnné a~vrcholy pro faktory.
Zde může být názvosloví matoucí, protože implementace faktorů z~předchozí sekce se používá pro oboje.
Vrchol pro proměnnou reprezentuje marginální distribuci proměnné.
Vrcholy pro faktory reprezentují pouze samotné faktory, ze~kterých se skládá sdružená distribuce.
Vrcholy slouží k~vytvoření reprezentace grafického modelu a obsahují metody pro výpočet a posílání zpráv.
Implementace vrcholů je v~modulu \texttt{alex.infer.node}.

\subsection{Rozhraní vrcholů}

Základní funkcionalita vrcholů je popsána v~abstraktní třídě \texttt{Node}.
Vlastností všech vrcholů je jejich sdružování do sítě.
K~tomu slouží metoda \texttt{connect}, kterou musí obsahovat každá implementace vrcholu, a ta informuje oba vrcholy, že spolu sousedí.
Většina vrcholů si pamatuje své sousedy pro počítání zpráv.
Při připojování proměnných k~faktorům je také možné označit proměnnou za rodiče daného faktoru.
Tato informace je důležitá při normalizaci a učení parametrů.

Dále třída obsahuje metody sloužící pro posílání zpráv.
Nejdůležitější jsou metody \texttt{message\_to} a \texttt{message\_from}.
Vrchol spočítá v~metodě \texttt{message\_to} zprávu pro svého souseda.
Má přístup ke všem ostatním vrcholům, a tak pro něj není problém zprávu spočítat.
Sousední vrchol zprávu přijme, když bude zavolána jeho metoda \texttt{message\_from}.
Díky tomu, že na posílání zpráv se podílí odesílatel i~příjemce, můžeme kombinovat různé vrcholy. 
Stačí, pokud se dohodnou na stejném formátu zprávy.
Dochází tak k~oddělení odesílatele od příjemce.

Další metodou, kterou obsahuje každý vrchol, je \texttt{init\_messages} pro inicializaci zpráv. 
Pro stávající implementace vrcholů není třeba ji volat před posíláním zpráv, protože zprávy mezi sousedními vrcholy jsou inicializovány vždy při zavolání metody \texttt{connect} pro jejich propojení.
Pokud ovšem budeme provádět více výpočtů nad jedním grafickým modelem, je třeba zprávy nastavit na původní hodnoty před dalším výpočtem, jinak by hodnota zpráv z~minulého výpočtu ovlivnila ten následující.

Z~optimalizačních důvodů si každý vrchol udržuje interně součin všech příchozích zpráv.
Při přijímání jsou zprávy pouze uloženy a jejich součin je aktualizován až při zavolání metody \texttt{update}.
Vzhledem k~tomu, že při inferenci je každý faktor vybrán jen jednou a pak jsou z~něj odeslány zprávy do všech ostatních vrcholů, stačí metodu \texttt{update} volat pro každý faktor v~každé iteraci jen jednou, před odesíláním zpráv.
Pro odeslání zpráv do všech vrcholů slouží metoda \texttt{send\_messages}.

Třída \texttt{DiscreteVariableNode} implementuje funkcionalitu vrcholů pro diskrétní proměnné.
Implementace vrcholů pro faktory s~diskrétními proměnnými je ve třídě \texttt{DiscreteFactorNode}. 

\subsection{Rozhraní vrcholů pro proměnné}

Smyslem vrcholů pro proměnné je reprezentovat aposteriorní marginální pravděpodobnost proměnných.
Mají navíc několik metod, které jsou potřeba pro výpočty.
Nejprve je třeba mít možnost nastavit pozorované proměnné.
K~tomu slouží metoda \texttt{observed}.
Tato metoda může být zavolána se slovníkem, kde klíčem jsou pozorovaná přiřazení a hodnotou je pravděpodobnost pozorování.
Zbylá přiřazení dostanou nulovou hodnotu.
S~tímto pozorováním se pak počítá při posílání zpráv.

Po ukončení výpočtu je třeba zjistit, která hodnota nebo více hodnot patří mezi nejpravděpodobnější.
Některé dialogové strategie potřebují jen jednu nejpravděpodobnější hodnotu, ale jiné mohou počítat se seznamem nejpravděpodobnějších možností. 
K~tomu slouží metoda \texttt{most\_probable}.
Vrátí seznam nejpravděpodobnějších přiřazení a jejich pravděpodobností.

\section{Vrcholy pro Dirichletovské parametry}
\label{sec:vrdir}

V~kapitole \ref{ch:ep} jsme ukázali, jakým způsobem je možné pro diskrétní proměnné učit parametry.
V~modulu \texttt{alex.infer.node} jsou implementovány vrcholy pro práci s~Dirichletovskými parametry.
Jedná se o~třídy \texttt{DirichletParameterNode} a \texttt{DirichletFactorNode}.
V~rozhraní se neliší od standardních implementací vrcholů.

Vrchol \texttt{DirichletParameterNode} musí být vždy spojen pouze s~vrcholem \texttt{DirichletFactorNode}, který už ale může být dále propojen s~vrcholy pro diskrétní proměnné.
Tyto vrcholy o~své existenci navzájem ví, a tak je možné použít standardní rozhraní pro posílání zpráv.
Při připojení vrcholu pro parametr k~vrcholu pro faktor tento pozná, že se k~němu připojil speciální vrchol a bude se tak k~němu chovat.
Zprávy od něj bude interpretovat jako parametry pro distribuci zbytku proměnných.
Pro zprávu z~faktoru do parametru je třeba provést aproximaci sdružené pravděpodobnostní distribuce a z~ní odvodit nové parametry (popsáno v~algoritmu~\ref{alg:ep} na straně~\pageref{alg:ep}).

\section{Inferenční algoritmus}

V~této sekci si popíšeme poslední vrstvu implementovanou v~knihovně, a to vrstvu starající se o~inferenci v~grafických modelech.
Předpokládáme, že již máme vytvořený graf z~vrcholů představených v~předchozí sekci.
V~modulu \texttt{alex.infer.lbp} je implementována třída \texttt{LBP}, která implementuje iterativní algoritmus pro aktualizaci grafického modelu a také obsahuje několik implementací strategií pro různé typy grafů.

Pro výpočet v~grafickém modelu je nejprve třeba zaregistrovat všechny vrcholy z~grafu.
Pro to existuje několik metod lišících se podle typu grafu.
Pro Dynamické Bayesovské sítě je možné přidávat vrcholy po vrstvách.
K~tomu slouží metody \texttt{add\_layer} a \texttt{add\_layers}.
První metoda přidá jednu vrstvu na konec sítě.
Druhá metoda umožňuje přidat více vrstev naráz.
Pro obecné grafy slouží metoda \texttt{add\_nodes}, která umožňuje přidat vrcholy bez informace o~jejich organizaci uvnitř sítě.

K~mazání vrcholů ze sítě slouží metody \texttt{clear\_nodes} a \texttt{clear\_layers}.

Typ inference je možné zadat při vytváření třídy, patřičná strategie je pak použita při samotném výpočtu uvnitř metody \texttt{run}. 
Tato metoda pak provádí inferenci, je možné zadat počet iterací, v~případě inference v~dynamické síti i~zadat, od které vrstvy se má inference provádět.

Pro více výpočtů nad jedním grafickým modelem je možné, místo inicializace zpráv pro každý vrchol zvlášť, použít metodu \texttt{init\_messages}, inicializující všechny vrcholy naráz.

\subsection{Strategie výběru vrcholu v~LBP algoritmu}
\label{sec:noch}

Inference v~grafu se může lišit podle typu faktor grafu, ale také podle nároků na výsledek.
Nejjednodušší metodou pro výběr je nechat pořadí na uživateli algoritmu.
Pro stromy máme speciální strategii zaručující konvergenci po jednom dopředném a jednom zpětném kroku propagace.
V~dialogových systémech je často používána Dynamická Bayesovská síť a nás zajímá pravděpodobnost proměnných v~poslední vrstvě sítě.
V~takovém případě můžeme některé zprávy zanedbat, protože už příliš neovlivní proměnné, které nás zajímají.

\subsubsection{Inference na stromě}

Pro efektivní inferenci je třeba si pro každý vrchol pamatovat, kolik mu chybí zpráv, aby mohl jednu sám odeslat.
Pro každý vrchol $v$ s~$k$ sousedy je na začátku počet chybějících zpráv $k-1$.
Strom na alespoň dvou vrcholech obsahuje alespoň dva listy.
Budeme tedy postupně odebírat vrcholy, jejichž počet chybějících zpráv je nulový.
Rozešleme z~nich vždy zprávu do souseda, ze~kterého ještě zpráva nepřišla a snížíme jeho počet chybějících zpráv.
Odebráním listu ze stromu vždy dostaneme opět strom.
Postupně tak graf zmenšujeme, až dostaneme právě jeden vrchol se všemi zprávami.

Po vypočtení marginální pravděpodobnosti ve stromě o~jednom vrcholu můžeme zase přidávat vrcholy v~obráceném pořadí, než v~jakém jsme je odebírali.
Do každého přidaného vrcholu pak můžeme poslat zprávu a spočítat jeho marginální pravděpodobnost.

\subsubsection{Inference v~Dynamické Bayesovské síti}

Dynamická Bayesovská síť je nejčastější reprezentace dialogového stavu.
V~jedné vrstvě sítě je popsána jedna obrátka dialogu.
Vrcholy v~jedné vrstvě většinou závisí pouze na vrcholech ve stejné nebo předchozí vrstvě.
Po každé obrátce nás pro účely komponenty řízení dialogu zajímají hlavně pravděpodobnosti proměnných v~poslední vrstvě, tedy po aktuální obrátce.

Inferenci provádíme způsobem podobným indukci.
Pro první vrstvu provedeme inferenci libovolným způsobem, může dokonce platit, že v~rámci jedné vrstvy se jedná o~strom.
Po přidání $k$-té vrstvy předpokládáme, že byla provedena inference na předchozích $k-1$ vrstvách, a tedy zprávy v~této části grafu jsou správné.
Zprávy ve směru k~nové vrstvě jsou stále správné, neznáme ovšem hodnotu zpráv z~předposlední vrstvy do nové vrstvy a hodnotu zpráv v~nové vrstvě.
Pošleme tedy zprávy z~předposlední vrstvy a provedeme inferenci v~nové vrstvě.

Stále ještě zbývají zprávy z~nové vrstvy zpět v~síti.
Čím dál do historie jdeme, tím menší vliv naše nová pozorování budou mít na vrcholy v~dané vrstvě.
Proto se většinou omezíme jen na posledních několik vrstev (1 až 3).

\section{Příklady}

Po popisu jednotlivých vrstev v~předchozí sekci nyní přejdeme k~příkladům použití knihovny.
Nejprve ukážeme použití jednotlivých tříd a metod.
Nakonec bude představen příklad systému použitého v~Dialog State Tracking Challenge (DSTC)~\cite{dstc2013}.

\subsection{Použití jednotlivých komponent}
\label{sec:usage}

Představíme příklady vytvoření a používání jednotlivých komponent, od implementace faktorů až po inferenci v~grafickém modelu.

\subsubsection{Faktor}

Třída \texttt{Faktor} je popsána v~sekci \ref{sec:repfak}.
Vytvoření faktoru je ukázáno v~příkladu~\ref{lst:crfac}, zde máme jednoduchý generativní model pro výběr restaurace podle typu jídla.
Vytvořený faktor obsahuje dvě proměnné, první je požadovaný typ jídla, druhá je pozorovaný požadavek.
Faktor může pocházet z~generativního modelu, kde pravděpodobnost pozorování závisí na skutečné hodnotě.
Zde vidíme, že je větší pravděpodobnost pozorování typu jídla, které uživatel opravdu chce.
Pro indické jídlo je zde pravděpodobnost větší než pro čínské, v~reálném dialogovém systému na tyto pravděpodobnosti může mít vliv například jazykový model nebo podobnost slov.

\begin{example}
\begin{Verbatim}[
    commandchars=\\\{\},
]

\PY{k+kn}{from} \PY{n+nn}{factor} \PY{k+kn}{import} \PY{n}{Factor}

\PY{n}{factor} \PY{o}{=} \PY{n}{Factor}\PY{p}{(}
    \PY{p}{[}\PY{l+s}{\PYZsq{}}\PY{l+s}{food}\PY{l+s}{\PYZsq{}}\PY{p}{,} \PY{l+s}{\PYZsq{}}\PY{l+s}{food\PYZus{}obs}\PY{l+s}{\PYZsq{}}\PY{p}{]}\PY{p}{,}
    \PY{p}{\PYZob{}}
        \PY{l+s}{\PYZsq{}}\PY{l+s}{food}\PY{l+s}{\PYZsq{}}\PY{p}{:} \PY{p}{[}\PY{l+s}{\PYZsq{}}\PY{l+s}{chinese}\PY{l+s}{\PYZsq{}}\PY{p}{,} \PY{l+s}{\PYZsq{}}\PY{l+s}{indian}\PY{l+s}{\PYZsq{}}\PY{p}{]}\PY{p}{,}
        \PY{l+s}{\PYZsq{}}\PY{l+s}{food\PYZus{}obs}\PY{l+s}{\PYZsq{}}\PY{p}{:} \PY{p}{[}\PY{l+s}{\PYZsq{}}\PY{l+s}{obs\PYZus{}chinese}\PY{l+s}{\PYZsq{}}\PY{p}{,} \PY{l+s}{\PYZsq{}}\PY{l+s}{obs\PYZus{}indian}\PY{l+s}{\PYZsq{}}\PY{p}{]}\PY{p}{,}
    \PY{p}{\PYZcb{}}\PY{p}{,}
    \PY{p}{\PYZob{}}
		\PY{p}{(}\PY{l+s}{\PYZsq{}}\PY{l+s}{indian}\PY{l+s}{\PYZsq{}}\PY{p}{,} \PY{l+s}{\PYZsq{}}\PY{l+s}{obs\PYZus{}indian}\PY{l+s}{\PYZsq{}}\PY{p}{)}\PY{p}{:} \PY{l+m+mf}{0.9}\PY{p}{,}
		\PY{p}{(}\PY{l+s}{\PYZsq{}}\PY{l+s}{indian}\PY{l+s}{\PYZsq{}}\PY{p}{,} \PY{l+s}{\PYZsq{}}\PY{l+s}{obs\PYZus{}chinese}\PY{l+s}{\PYZsq{}}\PY{p}{)}\PY{p}{:} \PY{l+m+mf}{0.1}\PY{p}{,}
		\PY{p}{(}\PY{l+s}{\PYZsq{}}\PY{l+s}{chinese}\PY{l+s}{\PYZsq{}}\PY{p}{,} \PY{l+s}{\PYZsq{}}\PY{l+s}{obs\PYZus{}indian}\PY{l+s}{\PYZsq{}}\PY{p}{)}\PY{p}{:} \PY{l+m+mf}{0.2}\PY{p}{,}
		\PY{p}{(}\PY{l+s}{\PYZsq{}}\PY{l+s}{chinese}\PY{l+s}{\PYZsq{}}\PY{p}{,} \PY{l+s}{\PYZsq{}}\PY{l+s}{obs\PYZus{}chinese}\PY{l+s}{\PYZsq{}}\PY{p}{)}\PY{p}{:} \PY{l+m+mf}{0.8}\PY{p}{,}
    \PY{p}{\PYZcb{}}\PY{p}{)}
\end{Verbatim}
\caption{Vytvoření faktoru}
\label{lst:crfac}
\end{example}

%\begin{lstlisting}[
%	caption={Vytvoření faktoru},
%	label={lst:crfac}
%]
%from factor import Factor
%
%factor = Factor(
%    ['food', 'food_obs'],
%    {
%        'food': ['chinese', 'indian'],
%        'food_obs': ['obs_chinese', 'obs_indian'],
%    },
%    {
%		('indian', 'obs_indian'): 0.9,
%		('indian', 'obs_chinese'): 0.1,
%		('chinese', 'obs_indian'): 0.2,
%		('chinese', 'obs_chinese'): 0.8,
%    })
%\end{lstlisting}

Vytvořený faktor můžeme zobrazit (příklad~\ref{lst:facprt}), stačí použít příkaz \texttt{print}, popřípadě je možné použít metodu \texttt{pretty\_print}, která vrátí řetězec s~formátovaným výpisem faktoru.
Nabízí možnost nastavení šířky tabulky a počtu desetinných míst.

\begin{example}
\begin{Verbatim}[commandchars=\\\{\}]

\PY{g+gp}{\PYZgt{}\PYZgt{}\PYZgt{} }\PY{k}{print} \PY{n}{factor}
\PY{g+go}{\PYZhy{}\PYZhy{}\PYZhy{}\PYZhy{}\PYZhy{}\PYZhy{}\PYZhy{}\PYZhy{}\PYZhy{}\PYZhy{}\PYZhy{}\PYZhy{}\PYZhy{}\PYZhy{}\PYZhy{}\PYZhy{}\PYZhy{}\PYZhy{}\PYZhy{}\PYZhy{}\PYZhy{}\PYZhy{}\PYZhy{}\PYZhy{}\PYZhy{}\PYZhy{}\PYZhy{}\PYZhy{}\PYZhy{}\PYZhy{}\PYZhy{}\PYZhy{}\PYZhy{}\PYZhy{}\PYZhy{}\PYZhy{}\PYZhy{}\PYZhy{}\PYZhy{}\PYZhy{}\PYZhy{}\PYZhy{}\PYZhy{}\PYZhy{}\PYZhy{}\PYZhy{}\PYZhy{}\PYZhy{}\PYZhy{}\PYZhy{}\PYZhy{}\PYZhy{}\PYZhy{}\PYZhy{}\PYZhy{}\PYZhy{}\PYZhy{}\PYZhy{}\PYZhy{}\PYZhy{}\PYZhy{}\PYZhy{}\PYZhy{}\PYZhy{}\PYZhy{}\PYZhy{}\PYZhy{}\PYZhy{}\PYZhy{}\PYZhy{}}
\PY{g+go}{         food                    food\PYZus{}obs                 Value}
\PY{g+go}{\PYZhy{}\PYZhy{}\PYZhy{}\PYZhy{}\PYZhy{}\PYZhy{}\PYZhy{}\PYZhy{}\PYZhy{}\PYZhy{}\PYZhy{}\PYZhy{}\PYZhy{}\PYZhy{}\PYZhy{}\PYZhy{}\PYZhy{}\PYZhy{}\PYZhy{}\PYZhy{}\PYZhy{}\PYZhy{}\PYZhy{}\PYZhy{}\PYZhy{}\PYZhy{}\PYZhy{}\PYZhy{}\PYZhy{}\PYZhy{}\PYZhy{}\PYZhy{}\PYZhy{}\PYZhy{}\PYZhy{}\PYZhy{}\PYZhy{}\PYZhy{}\PYZhy{}\PYZhy{}\PYZhy{}\PYZhy{}\PYZhy{}\PYZhy{}\PYZhy{}\PYZhy{}\PYZhy{}\PYZhy{}\PYZhy{}\PYZhy{}\PYZhy{}\PYZhy{}\PYZhy{}\PYZhy{}\PYZhy{}\PYZhy{}\PYZhy{}\PYZhy{}\PYZhy{}\PYZhy{}\PYZhy{}\PYZhy{}\PYZhy{}\PYZhy{}\PYZhy{}\PYZhy{}\PYZhy{}\PYZhy{}\PYZhy{}\PYZhy{}}
\PY{g+go}{       chinese                 obs\PYZus{}chinese             0.8000000119}
\PY{g+go}{       chinese                  obs\PYZus{}indian             0.1999999881}
\PY{g+go}{        indian                 obs\PYZus{}chinese            0.09999999404}
\PY{g+go}{        indian                  obs\PYZus{}indian             0.8999999762}
\PY{g+go}{\PYZhy{}\PYZhy{}\PYZhy{}\PYZhy{}\PYZhy{}\PYZhy{}\PYZhy{}\PYZhy{}\PYZhy{}\PYZhy{}\PYZhy{}\PYZhy{}\PYZhy{}\PYZhy{}\PYZhy{}\PYZhy{}\PYZhy{}\PYZhy{}\PYZhy{}\PYZhy{}\PYZhy{}\PYZhy{}\PYZhy{}\PYZhy{}\PYZhy{}\PYZhy{}\PYZhy{}\PYZhy{}\PYZhy{}\PYZhy{}\PYZhy{}\PYZhy{}\PYZhy{}\PYZhy{}\PYZhy{}\PYZhy{}\PYZhy{}\PYZhy{}\PYZhy{}\PYZhy{}\PYZhy{}\PYZhy{}\PYZhy{}\PYZhy{}\PYZhy{}\PYZhy{}\PYZhy{}\PYZhy{}\PYZhy{}\PYZhy{}\PYZhy{}\PYZhy{}\PYZhy{}\PYZhy{}\PYZhy{}\PYZhy{}\PYZhy{}\PYZhy{}\PYZhy{}\PYZhy{}\PYZhy{}\PYZhy{}\PYZhy{}\PYZhy{}\PYZhy{}\PYZhy{}\PYZhy{}\PYZhy{}\PYZhy{}\PYZhy{}}
\PY{g+gp}{\PYZgt{}\PYZgt{}\PYZgt{} }\PY{k}{print} \PY{n}{factor}\PY{o}{.}\PY{n}{pretty\PYZus{}print}\PY{p}{(}\PY{n}{width}\PY{o}{=}\PY{l+m+mi}{40}\PY{p}{,} \PY{n}{precision}\PY{o}{=}\PY{l+m+mi}{2}\PY{p}{)}
\PY{g+go}{\PYZhy{}\PYZhy{}\PYZhy{}\PYZhy{}\PYZhy{}\PYZhy{}\PYZhy{}\PYZhy{}\PYZhy{}\PYZhy{}\PYZhy{}\PYZhy{}\PYZhy{}\PYZhy{}\PYZhy{}\PYZhy{}\PYZhy{}\PYZhy{}\PYZhy{}\PYZhy{}\PYZhy{}\PYZhy{}\PYZhy{}\PYZhy{}\PYZhy{}\PYZhy{}\PYZhy{}\PYZhy{}\PYZhy{}\PYZhy{}\PYZhy{}\PYZhy{}\PYZhy{}\PYZhy{}\PYZhy{}\PYZhy{}\PYZhy{}\PYZhy{}\PYZhy{}\PYZhy{}}
\PY{g+go}{    food       food\PYZus{}obs       Value}
\PY{g+go}{\PYZhy{}\PYZhy{}\PYZhy{}\PYZhy{}\PYZhy{}\PYZhy{}\PYZhy{}\PYZhy{}\PYZhy{}\PYZhy{}\PYZhy{}\PYZhy{}\PYZhy{}\PYZhy{}\PYZhy{}\PYZhy{}\PYZhy{}\PYZhy{}\PYZhy{}\PYZhy{}\PYZhy{}\PYZhy{}\PYZhy{}\PYZhy{}\PYZhy{}\PYZhy{}\PYZhy{}\PYZhy{}\PYZhy{}\PYZhy{}\PYZhy{}\PYZhy{}\PYZhy{}\PYZhy{}\PYZhy{}\PYZhy{}\PYZhy{}\PYZhy{}\PYZhy{}\PYZhy{}}
\PY{g+go}{   chinese    obs\PYZus{}chinese      0.8}
\PY{g+go}{   chinese    obs\PYZus{}indian       0.2}
\PY{g+go}{   indian     obs\PYZus{}chinese      0.1}
\PY{g+go}{   indian     obs\PYZus{}indian       0.9}
\PY{g+go}{\PYZhy{}\PYZhy{}\PYZhy{}\PYZhy{}\PYZhy{}\PYZhy{}\PYZhy{}\PYZhy{}\PYZhy{}\PYZhy{}\PYZhy{}\PYZhy{}\PYZhy{}\PYZhy{}\PYZhy{}\PYZhy{}\PYZhy{}\PYZhy{}\PYZhy{}\PYZhy{}\PYZhy{}\PYZhy{}\PYZhy{}\PYZhy{}\PYZhy{}\PYZhy{}\PYZhy{}\PYZhy{}\PYZhy{}\PYZhy{}\PYZhy{}\PYZhy{}\PYZhy{}\PYZhy{}\PYZhy{}\PYZhy{}\PYZhy{}\PYZhy{}\PYZhy{}\PYZhy{}              }
\end{Verbatim}
\caption{Zobrazení faktoru}
\label{lst:facprt}
\end{example}


%\begin{lstlisting}[
%	caption={Zobrazení faktoru},
%	label={lst:facprt},
%	float=h,
%]
%>>> print factor
%--------------------------------------------------------------------------
%           food                    food_obs                   Value
%--------------------------------------------------------------------------
%         chinese                 obs_chinese               0.8000000119
%         chinese                  obs_indian               0.1999999881
%          indian                 obs_chinese              0.09999999404
%          indian                  obs_indian               0.8999999762
%--------------------------------------------------------------------------
%>>> print factor.pretty_print(width=40, precision=2)
%----------------------------------------
%    food       food_obs       Value
%----------------------------------------
%   chinese    obs_chinese      0.8
%   chinese    obs_indian       0.2
%   indian     obs_chinese      0.1
%   indian     obs_indian       0.9
%----------------------------------------
%\end{lstlisting}

Dále obsahuje faktor implementaci matematických operací, je možné používat standardní operátory.
V~příkladu~\ref{lst:facmul} je ukázáno násobení dvou faktorů.
Operace fungují i s~konstantami (příklad~\ref{lst:facmulk}).

\begin{example}
\begin{Verbatim}[commandchars=\\\{\}]

\PY{g+gp}{\PYZgt{}\PYZgt{}\PYZgt{} }\PY{n}{factor\PYZus{}trans} \PY{o}{=} \PY{n}{Factor}\PY{p}{(}
\PY{g+gp}{... }   \PY{p}{[}\PY{l+s}{\PYZsq{}}\PY{l+s}{food}\PY{l+s}{\PYZsq{}}\PY{p}{,} \PY{l+s}{\PYZsq{}}\PY{l+s}{food\PYZus{}next}\PY{l+s}{\PYZsq{}}\PY{p}{]}\PY{p}{,}
\PY{g+gp}{... }   \PY{p}{\PYZob{}}
\PY{g+gp}{... }       \PY{l+s}{\PYZsq{}}\PY{l+s}{food}\PY{l+s}{\PYZsq{}}\PY{p}{:} \PY{p}{[}\PY{l+s}{\PYZsq{}}\PY{l+s}{chinese}\PY{l+s}{\PYZsq{}}\PY{p}{,} \PY{l+s}{\PYZsq{}}\PY{l+s}{indian}\PY{l+s}{\PYZsq{}}\PY{p}{]}\PY{p}{,}
\PY{g+gp}{... }       \PY{l+s}{\PYZsq{}}\PY{l+s}{food\PYZus{}next}\PY{l+s}{\PYZsq{}}\PY{p}{:} \PY{p}{[}\PY{l+s}{\PYZsq{}}\PY{l+s}{chinese}\PY{l+s}{\PYZsq{}}\PY{p}{,} \PY{l+s}{\PYZsq{}}\PY{l+s}{indian}\PY{l+s}{\PYZsq{}}\PY{p}{]}\PY{p}{,}
\PY{g+gp}{... }   \PY{p}{\PYZcb{}}\PY{p}{,}
\PY{g+gp}{... }   \PY{p}{\PYZob{}}
\PY{g+gp}{... }       \PY{p}{(}\PY{l+s}{\PYZsq{}}\PY{l+s}{indian}\PY{l+s}{\PYZsq{}}\PY{p}{,} \PY{l+s}{\PYZsq{}}\PY{l+s}{indian}\PY{l+s}{\PYZsq{}}\PY{p}{)}\PY{p}{:} \PY{l+m+mf}{0.9}\PY{p}{,}
\PY{g+gp}{... }       \PY{p}{(}\PY{l+s}{\PYZsq{}}\PY{l+s}{indian}\PY{l+s}{\PYZsq{}}\PY{p}{,} \PY{l+s}{\PYZsq{}}\PY{l+s}{chinese}\PY{l+s}{\PYZsq{}}\PY{p}{)}\PY{p}{:} \PY{l+m+mf}{0.1}\PY{p}{,}
\PY{g+gp}{... }       \PY{p}{(}\PY{l+s}{\PYZsq{}}\PY{l+s}{chinese}\PY{l+s}{\PYZsq{}}\PY{p}{,} \PY{l+s}{\PYZsq{}}\PY{l+s}{indian}\PY{l+s}{\PYZsq{}}\PY{p}{)}\PY{p}{:} \PY{l+m+mf}{0.1}\PY{p}{,}
\PY{g+gp}{... }       \PY{p}{(}\PY{l+s}{\PYZsq{}}\PY{l+s}{chinese}\PY{l+s}{\PYZsq{}}\PY{p}{,} \PY{l+s}{\PYZsq{}}\PY{l+s}{chinese}\PY{l+s}{\PYZsq{}}\PY{p}{)}\PY{p}{:} \PY{l+m+mf}{0.9}\PY{p}{,}
\PY{g+gp}{... }   \PY{p}{\PYZcb{}}\PY{p}{)}
\PY{g+gp}{\PYZgt{}\PYZgt{}\PYZgt{} }\PY{n}{result} \PY{o}{=} \PY{n}{factor} \PY{o}{*} \PY{n}{factor\PYZus{}trans}
\PY{g+gp}{\PYZgt{}\PYZgt{}\PYZgt{} }\PY{k}{print} \PY{n}{result}\PY{o}{.}\PY{n}{pretty\PYZus{}print}\PY{p}{(}\PY{l+m+mi}{50}\PY{p}{,} \PY{l+m+mi}{2}\PY{p}{)}
\PY{g+go}{\PYZhy{}\PYZhy{}\PYZhy{}\PYZhy{}\PYZhy{}\PYZhy{}\PYZhy{}\PYZhy{}\PYZhy{}\PYZhy{}\PYZhy{}\PYZhy{}\PYZhy{}\PYZhy{}\PYZhy{}\PYZhy{}\PYZhy{}\PYZhy{}\PYZhy{}\PYZhy{}\PYZhy{}\PYZhy{}\PYZhy{}\PYZhy{}\PYZhy{}\PYZhy{}\PYZhy{}\PYZhy{}\PYZhy{}\PYZhy{}\PYZhy{}\PYZhy{}\PYZhy{}\PYZhy{}\PYZhy{}\PYZhy{}\PYZhy{}\PYZhy{}\PYZhy{}\PYZhy{}\PYZhy{}\PYZhy{}\PYZhy{}\PYZhy{}\PYZhy{}\PYZhy{}\PYZhy{}\PYZhy{}\PYZhy{}\PYZhy{}}
\PY{g+go}{    food     food\PYZus{}next    food\PYZus{}obs     Value}
\PY{g+go}{\PYZhy{}\PYZhy{}\PYZhy{}\PYZhy{}\PYZhy{}\PYZhy{}\PYZhy{}\PYZhy{}\PYZhy{}\PYZhy{}\PYZhy{}\PYZhy{}\PYZhy{}\PYZhy{}\PYZhy{}\PYZhy{}\PYZhy{}\PYZhy{}\PYZhy{}\PYZhy{}\PYZhy{}\PYZhy{}\PYZhy{}\PYZhy{}\PYZhy{}\PYZhy{}\PYZhy{}\PYZhy{}\PYZhy{}\PYZhy{}\PYZhy{}\PYZhy{}\PYZhy{}\PYZhy{}\PYZhy{}\PYZhy{}\PYZhy{}\PYZhy{}\PYZhy{}\PYZhy{}\PYZhy{}\PYZhy{}\PYZhy{}\PYZhy{}\PYZhy{}\PYZhy{}\PYZhy{}\PYZhy{}\PYZhy{}\PYZhy{}}
\PY{g+go}{  chinese     chinese   obs\PYZus{}chinese     0.72}
\PY{g+go}{  chinese     chinese    obs\PYZus{}indian     0.18}
\PY{g+go}{  chinese      indian   obs\PYZus{}chinese     0.08}
\PY{g+go}{  chinese      indian    obs\PYZus{}indian     0.02}
\PY{g+go}{   indian     chinese   obs\PYZus{}chinese     0.01}
\PY{g+go}{   indian     chinese    obs\PYZus{}indian     0.09}
\PY{g+go}{   indian      indian   obs\PYZus{}chinese     0.09}
\PY{g+go}{   indian      indian    obs\PYZus{}indian     0.81}
\PY{g+go}{\PYZhy{}\PYZhy{}\PYZhy{}\PYZhy{}\PYZhy{}\PYZhy{}\PYZhy{}\PYZhy{}\PYZhy{}\PYZhy{}\PYZhy{}\PYZhy{}\PYZhy{}\PYZhy{}\PYZhy{}\PYZhy{}\PYZhy{}\PYZhy{}\PYZhy{}\PYZhy{}\PYZhy{}\PYZhy{}\PYZhy{}\PYZhy{}\PYZhy{}\PYZhy{}\PYZhy{}\PYZhy{}\PYZhy{}\PYZhy{}\PYZhy{}\PYZhy{}\PYZhy{}\PYZhy{}\PYZhy{}\PYZhy{}\PYZhy{}\PYZhy{}\PYZhy{}\PYZhy{}\PYZhy{}\PYZhy{}\PYZhy{}\PYZhy{}\PYZhy{}\PYZhy{}\PYZhy{}\PYZhy{}\PYZhy{}\PYZhy{}}

\end{Verbatim}
\caption{Násobení faktorů}
\label{lst:facmul}
\end{example}

%\begin{lstlisting}[
%	caption={Násobení faktorů},
%	label={lst:facmul}
%]
%>>> factor_trans = Factor(
%...    ['food', 'food_next'],
%...    {
%...        'food': ['chinese', 'indian'],
%...        'food_next': ['chinese', 'indian'],
%...    },
%...    {
%...        ('indian', 'indian'): 0.9,
%...        ('indian', 'chinese'): 0.1,
%...        ('chinese', 'indian'): 0.1,
%...        ('chinese', 'chinese'): 0.9,
%...    })
%>>> result = factor * factor_trans
%>>> print result.pretty_print(50, 2)
%--------------------------------------------------
%    food     food_next    food_obs     Value
%--------------------------------------------------
%  chinese     chinese   obs_chinese     0.72
%  chinese     chinese    obs_indian     0.18
%  chinese      indian   obs_chinese     0.08
%  chinese      indian    obs_indian     0.02
%   indian     chinese   obs_chinese     0.01
%   indian     chinese    obs_indian     0.09
%   indian      indian   obs_chinese     0.09
%   indian      indian    obs_indian     0.81
%--------------------------------------------------
%\end{lstlisting}

\begin{example}
\begin{Verbatim}[commandchars=\\\{\}]

\PY{g+gp}{\PYZgt{}\PYZgt{}\PYZgt{} }\PY{n}{result} \PY{o}{=} \PY{n}{factor} \PY{o}{*} \PY{l+m+mf}{0.5}
\PY{g+gp}{\PYZgt{}\PYZgt{}\PYZgt{} }\PY{k}{print} \PY{n}{result}\PY{o}{.}\PY{n}{pretty\PYZus{}print}\PY{p}{(}\PY{l+m+mi}{50}\PY{p}{,} \PY{l+m+mi}{2}\PY{p}{)}
\PY{g+go}{\PYZhy{}\PYZhy{}\PYZhy{}\PYZhy{}\PYZhy{}\PYZhy{}\PYZhy{}\PYZhy{}\PYZhy{}\PYZhy{}\PYZhy{}\PYZhy{}\PYZhy{}\PYZhy{}\PYZhy{}\PYZhy{}\PYZhy{}\PYZhy{}\PYZhy{}\PYZhy{}\PYZhy{}\PYZhy{}\PYZhy{}\PYZhy{}\PYZhy{}\PYZhy{}\PYZhy{}\PYZhy{}\PYZhy{}\PYZhy{}\PYZhy{}\PYZhy{}\PYZhy{}\PYZhy{}\PYZhy{}\PYZhy{}\PYZhy{}\PYZhy{}\PYZhy{}\PYZhy{}\PYZhy{}\PYZhy{}\PYZhy{}\PYZhy{}\PYZhy{}\PYZhy{}\PYZhy{}\PYZhy{}\PYZhy{}\PYZhy{}}
\PY{g+go}{      food          food\PYZus{}obs         Value}
\PY{g+go}{\PYZhy{}\PYZhy{}\PYZhy{}\PYZhy{}\PYZhy{}\PYZhy{}\PYZhy{}\PYZhy{}\PYZhy{}\PYZhy{}\PYZhy{}\PYZhy{}\PYZhy{}\PYZhy{}\PYZhy{}\PYZhy{}\PYZhy{}\PYZhy{}\PYZhy{}\PYZhy{}\PYZhy{}\PYZhy{}\PYZhy{}\PYZhy{}\PYZhy{}\PYZhy{}\PYZhy{}\PYZhy{}\PYZhy{}\PYZhy{}\PYZhy{}\PYZhy{}\PYZhy{}\PYZhy{}\PYZhy{}\PYZhy{}\PYZhy{}\PYZhy{}\PYZhy{}\PYZhy{}\PYZhy{}\PYZhy{}\PYZhy{}\PYZhy{}\PYZhy{}\PYZhy{}\PYZhy{}\PYZhy{}\PYZhy{}\PYZhy{}}
\PY{g+go}{    chinese       obs\PYZus{}chinese         0.4}
\PY{g+go}{    chinese        obs\PYZus{}indian         0.1}
\PY{g+go}{     indian       obs\PYZus{}chinese         0.05}
\PY{g+go}{     indian        obs\PYZus{}indian         0.45}
\PY{g+go}{\PYZhy{}\PYZhy{}\PYZhy{}\PYZhy{}\PYZhy{}\PYZhy{}\PYZhy{}\PYZhy{}\PYZhy{}\PYZhy{}\PYZhy{}\PYZhy{}\PYZhy{}\PYZhy{}\PYZhy{}\PYZhy{}\PYZhy{}\PYZhy{}\PYZhy{}\PYZhy{}\PYZhy{}\PYZhy{}\PYZhy{}\PYZhy{}\PYZhy{}\PYZhy{}\PYZhy{}\PYZhy{}\PYZhy{}\PYZhy{}\PYZhy{}\PYZhy{}\PYZhy{}\PYZhy{}\PYZhy{}\PYZhy{}\PYZhy{}\PYZhy{}\PYZhy{}\PYZhy{}\PYZhy{}\PYZhy{}\PYZhy{}\PYZhy{}\PYZhy{}\PYZhy{}\PYZhy{}\PYZhy{}\PYZhy{}\PYZhy{}}
\end{Verbatim}
\caption{Násobení faktoru konstantou}
\label{lst:facmulk}
\end{example}

%\begin{lstlisting}[
%	caption={Násobení faktorů konstantou},
%	label={lst:facmulk},
%	float=h,
%]
%>>> result = factor * 0.5
%>>> print result.pretty_print(50, 2)
%--------------------------------------------------
%      food          food_obs         Value
%--------------------------------------------------
%    chinese       obs_chinese         0.4
%    chinese        obs_indian         0.1
%     indian       obs_chinese         0.05
%     indian        obs_indian         0.45
%--------------------------------------------------
%\end{lstlisting}

Další důležitou metodou, kterou faktory nabízí, je marginalizace proměnných.
Předvedeme si ji na faktoru pro pozorování (příklad~\ref{lst:facmar}), ze~kterého chceme získat pouze marginální pravděpodobnosti pozorování.

\begin{example}
\begin{Verbatim}[commandchars=\\\{\}]

\PY{g+gp}{\PYZgt{}\PYZgt{}\PYZgt{} }\PY{n}{marginalized} \PY{o}{=} \PY{n}{factor}\PY{o}{.}\PY{n}{marginalize}\PY{p}{(}\PY{p}{[}\PY{l+s}{\PYZsq{}}\PY{l+s}{food\PYZus{}obs}\PY{l+s}{\PYZsq{}}\PY{p}{]}\PY{p}{)}
\PY{g+gp}{\PYZgt{}\PYZgt{}\PYZgt{} }\PY{k}{print} \PY{n}{marginalized}\PY{o}{.}\PY{n}{pretty\PYZus{}print}\PY{p}{(}\PY{l+m+mi}{50}\PY{p}{,} \PY{l+m+mi}{2}\PY{p}{)}
\PY{g+go}{\PYZhy{}\PYZhy{}\PYZhy{}\PYZhy{}\PYZhy{}\PYZhy{}\PYZhy{}\PYZhy{}\PYZhy{}\PYZhy{}\PYZhy{}\PYZhy{}\PYZhy{}\PYZhy{}\PYZhy{}\PYZhy{}\PYZhy{}\PYZhy{}\PYZhy{}\PYZhy{}\PYZhy{}\PYZhy{}\PYZhy{}\PYZhy{}\PYZhy{}\PYZhy{}\PYZhy{}\PYZhy{}\PYZhy{}\PYZhy{}\PYZhy{}\PYZhy{}\PYZhy{}\PYZhy{}\PYZhy{}\PYZhy{}\PYZhy{}\PYZhy{}\PYZhy{}\PYZhy{}\PYZhy{}\PYZhy{}\PYZhy{}\PYZhy{}\PYZhy{}\PYZhy{}\PYZhy{}\PYZhy{}\PYZhy{}\PYZhy{}}
\PY{g+go}{        food\PYZus{}obs                   Value}
\PY{g+go}{\PYZhy{}\PYZhy{}\PYZhy{}\PYZhy{}\PYZhy{}\PYZhy{}\PYZhy{}\PYZhy{}\PYZhy{}\PYZhy{}\PYZhy{}\PYZhy{}\PYZhy{}\PYZhy{}\PYZhy{}\PYZhy{}\PYZhy{}\PYZhy{}\PYZhy{}\PYZhy{}\PYZhy{}\PYZhy{}\PYZhy{}\PYZhy{}\PYZhy{}\PYZhy{}\PYZhy{}\PYZhy{}\PYZhy{}\PYZhy{}\PYZhy{}\PYZhy{}\PYZhy{}\PYZhy{}\PYZhy{}\PYZhy{}\PYZhy{}\PYZhy{}\PYZhy{}\PYZhy{}\PYZhy{}\PYZhy{}\PYZhy{}\PYZhy{}\PYZhy{}\PYZhy{}\PYZhy{}\PYZhy{}\PYZhy{}\PYZhy{}}
\PY{g+go}{       obs\PYZus{}chinese                  0.9}
\PY{g+go}{       obs\PYZus{}indian                   1.1}
\PY{g+go}{\PYZhy{}\PYZhy{}\PYZhy{}\PYZhy{}\PYZhy{}\PYZhy{}\PYZhy{}\PYZhy{}\PYZhy{}\PYZhy{}\PYZhy{}\PYZhy{}\PYZhy{}\PYZhy{}\PYZhy{}\PYZhy{}\PYZhy{}\PYZhy{}\PYZhy{}\PYZhy{}\PYZhy{}\PYZhy{}\PYZhy{}\PYZhy{}\PYZhy{}\PYZhy{}\PYZhy{}\PYZhy{}\PYZhy{}\PYZhy{}\PYZhy{}\PYZhy{}\PYZhy{}\PYZhy{}\PYZhy{}\PYZhy{}\PYZhy{}\PYZhy{}\PYZhy{}\PYZhy{}\PYZhy{}\PYZhy{}\PYZhy{}\PYZhy{}\PYZhy{}\PYZhy{}\PYZhy{}\PYZhy{}\PYZhy{}\PYZhy{}}
\end{Verbatim}
\caption{Marginalizace faktoru}
\label{lst:facmar}
\end{example}

%\begin{lstlisting}[
%	caption={Marginalizace faktoru},
%	label={lst:facmar},
%	float=h,
%]
%>>> marginalized = factor.marginalize(['food_obs'])
%>>> print marginalized.pretty_print(50, 2)
%--------------------------------------------------
%        food_obs                   Value
%--------------------------------------------------
%       obs_chinese                  0.9
%       obs_indian                   1.1
%--------------------------------------------------
%\end{lstlisting}

Pokud je faktor použit pro reprezentaci pravděpodobnostního rozdělení, pak je možné, že při některých úkonech bude výsledkem nenormalizované pravděpodobnostní rozdělení.
Faktor nabízí metodu pro normalizaci hodnot, která navíc bere v~úvahu i podmíněné pravděpodobnosti (příklad~\ref{lst:facnor}).

\begin{example}
\begin{Verbatim}[commandchars=\\\{\}]

\PY{g+gp}{\PYZgt{}\PYZgt{}\PYZgt{} }\PY{n}{uniform} \PY{o}{=} \PY{n}{Factor}\PY{p}{(}
\PY{g+gp}{... }   \PY{p}{[}\PY{l+s}{\PYZsq{}}\PY{l+s}{food}\PY{l+s}{\PYZsq{}}\PY{p}{,} \PY{l+s}{\PYZsq{}}\PY{l+s}{food\PYZus{}next}\PY{l+s}{\PYZsq{}}\PY{p}{]}\PY{p}{,}
\PY{g+gp}{... }   \PY{p}{\PYZob{}}
\PY{g+gp}{... }       \PY{l+s}{\PYZsq{}}\PY{l+s}{food}\PY{l+s}{\PYZsq{}}\PY{p}{:} \PY{p}{[}\PY{l+s}{\PYZsq{}}\PY{l+s}{chinese}\PY{l+s}{\PYZsq{}}\PY{p}{,} \PY{l+s}{\PYZsq{}}\PY{l+s}{indian}\PY{l+s}{\PYZsq{}}\PY{p}{]}\PY{p}{,}
\PY{g+gp}{... }       \PY{l+s}{\PYZsq{}}\PY{l+s}{food\PYZus{}next}\PY{l+s}{\PYZsq{}}\PY{p}{:} \PY{p}{[}\PY{l+s}{\PYZsq{}}\PY{l+s}{chinese}\PY{l+s}{\PYZsq{}}\PY{p}{,} \PY{l+s}{\PYZsq{}}\PY{l+s}{indian}\PY{l+s}{\PYZsq{}}\PY{p}{]}\PY{p}{,}
\PY{g+gp}{... }   \PY{p}{\PYZcb{}}\PY{p}{,}
\PY{g+gp}{... }   \PY{p}{\PYZob{}}
\PY{g+gp}{... }       \PY{p}{(}\PY{l+s}{\PYZsq{}}\PY{l+s}{indian}\PY{l+s}{\PYZsq{}}\PY{p}{,} \PY{l+s}{\PYZsq{}}\PY{l+s}{indian}\PY{l+s}{\PYZsq{}}\PY{p}{)}\PY{p}{:} \PY{l+m+mi}{1}\PY{p}{,}
\PY{g+gp}{... }       \PY{p}{(}\PY{l+s}{\PYZsq{}}\PY{l+s}{indian}\PY{l+s}{\PYZsq{}}\PY{p}{,} \PY{l+s}{\PYZsq{}}\PY{l+s}{chinese}\PY{l+s}{\PYZsq{}}\PY{p}{)}\PY{p}{:} \PY{l+m+mi}{1}\PY{p}{,}
\PY{g+gp}{... }       \PY{p}{(}\PY{l+s}{\PYZsq{}}\PY{l+s}{chinese}\PY{l+s}{\PYZsq{}}\PY{p}{,} \PY{l+s}{\PYZsq{}}\PY{l+s}{indian}\PY{l+s}{\PYZsq{}}\PY{p}{)}\PY{p}{:} \PY{l+m+mi}{2}\PY{p}{,}
\PY{g+gp}{... }       \PY{p}{(}\PY{l+s}{\PYZsq{}}\PY{l+s}{chinese}\PY{l+s}{\PYZsq{}}\PY{p}{,} \PY{l+s}{\PYZsq{}}\PY{l+s}{chinese}\PY{l+s}{\PYZsq{}}\PY{p}{)}\PY{p}{:} \PY{l+m+mi}{2}\PY{p}{,}
\PY{g+gp}{... }   \PY{p}{\PYZcb{}}\PY{p}{)}
\PY{g+gp}{\PYZgt{}\PYZgt{}\PYZgt{} }\PY{n}{uniform}\PY{o}{.}\PY{n}{normalize}\PY{p}{(}\PY{n}{parents}\PY{o}{=}\PY{p}{[}\PY{l+s}{\PYZsq{}}\PY{l+s}{food}\PY{l+s}{\PYZsq{}}\PY{p}{]}\PY{p}{)}
\PY{g+gp}{\PYZgt{}\PYZgt{}\PYZgt{} }\PY{k}{print} \PY{n}{uniform}\PY{o}{.}\PY{n}{pretty\PYZus{}print}\PY{p}{(}\PY{l+m+mi}{50}\PY{p}{,} \PY{l+m+mi}{2}\PY{p}{)}
\PY{g+go}{\PYZhy{}\PYZhy{}\PYZhy{}\PYZhy{}\PYZhy{}\PYZhy{}\PYZhy{}\PYZhy{}\PYZhy{}\PYZhy{}\PYZhy{}\PYZhy{}\PYZhy{}\PYZhy{}\PYZhy{}\PYZhy{}\PYZhy{}\PYZhy{}\PYZhy{}\PYZhy{}\PYZhy{}\PYZhy{}\PYZhy{}\PYZhy{}\PYZhy{}\PYZhy{}\PYZhy{}\PYZhy{}\PYZhy{}\PYZhy{}\PYZhy{}\PYZhy{}\PYZhy{}\PYZhy{}\PYZhy{}\PYZhy{}\PYZhy{}\PYZhy{}\PYZhy{}\PYZhy{}\PYZhy{}\PYZhy{}\PYZhy{}\PYZhy{}\PYZhy{}\PYZhy{}\PYZhy{}\PYZhy{}\PYZhy{}\PYZhy{}}
\PY{g+go}{      food         food\PYZus{}next         Value}
\PY{g+go}{\PYZhy{}\PYZhy{}\PYZhy{}\PYZhy{}\PYZhy{}\PYZhy{}\PYZhy{}\PYZhy{}\PYZhy{}\PYZhy{}\PYZhy{}\PYZhy{}\PYZhy{}\PYZhy{}\PYZhy{}\PYZhy{}\PYZhy{}\PYZhy{}\PYZhy{}\PYZhy{}\PYZhy{}\PYZhy{}\PYZhy{}\PYZhy{}\PYZhy{}\PYZhy{}\PYZhy{}\PYZhy{}\PYZhy{}\PYZhy{}\PYZhy{}\PYZhy{}\PYZhy{}\PYZhy{}\PYZhy{}\PYZhy{}\PYZhy{}\PYZhy{}\PYZhy{}\PYZhy{}\PYZhy{}\PYZhy{}\PYZhy{}\PYZhy{}\PYZhy{}\PYZhy{}\PYZhy{}\PYZhy{}\PYZhy{}\PYZhy{}}
\PY{g+go}{    chinese         chinese           0.5}
\PY{g+go}{    chinese          indian           0.5}
\PY{g+go}{     indian         chinese           0.5}
\PY{g+go}{     indian          indian           0.5}
\PY{g+go}{\PYZhy{}\PYZhy{}\PYZhy{}\PYZhy{}\PYZhy{}\PYZhy{}\PYZhy{}\PYZhy{}\PYZhy{}\PYZhy{}\PYZhy{}\PYZhy{}\PYZhy{}\PYZhy{}\PYZhy{}\PYZhy{}\PYZhy{}\PYZhy{}\PYZhy{}\PYZhy{}\PYZhy{}\PYZhy{}\PYZhy{}\PYZhy{}\PYZhy{}\PYZhy{}\PYZhy{}\PYZhy{}\PYZhy{}\PYZhy{}\PYZhy{}\PYZhy{}\PYZhy{}\PYZhy{}\PYZhy{}\PYZhy{}\PYZhy{}\PYZhy{}\PYZhy{}\PYZhy{}\PYZhy{}\PYZhy{}\PYZhy{}\PYZhy{}\PYZhy{}\PYZhy{}\PYZhy{}\PYZhy{}\PYZhy{}\PYZhy{}}
\end{Verbatim}
\caption{Normalizace faktoru}
\label{lst:facnor}
\end{example}

%\begin{lstlisting}[
%	caption={Normalizace faktoru},
%	label={lst:facnor}
%]
%>>> uniform = Factor(
%...    ['food', 'food_next'],
%...    {
%...        'food': ['chinese', 'indian'],
%...        'food_next': ['chinese', 'indian'],
%...    },
%...    {
%...        ('indian', 'indian'): 1,
%...	       ('indian', 'chinese'): 1,
%...        ('chinese', 'indian'): 2,
%...        ('chinese', 'chinese'): 2,
%...    })
%>>> uniform.normalize(parents=['food'])
%>>> print uniform.pretty_print(50, 2)
%--------------------------------------------------
%      food         food_next         Value
%--------------------------------------------------
%    chinese         chinese           0.5
%    chinese          indian           0.5
%     indian         chinese           0.5
%     indian          indian           0.5
%--------------------------------------------------
%\end{lstlisting}

Nakonec faktor nabízí možnost zjistit, která přiřazení jsou nejpravděpodobnější (příklad~\ref{lst:facmop}).

\begin{example}
\begin{Verbatim}[commandchars=\\\{\}]

\PY{g+gp}{\PYZgt{}\PYZgt{}\PYZgt{} }\PY{n}{food} \PY{o}{=} \PY{n}{Factor}\PY{p}{(}
\PY{g+gp}{... }   \PY{p}{[}\PY{l+s}{\PYZsq{}}\PY{l+s}{food}\PY{l+s}{\PYZsq{}}\PY{p}{]}\PY{p}{,}
\PY{g+gp}{... }   \PY{p}{\PYZob{}}
\PY{g+gp}{... }       \PY{l+s}{\PYZsq{}}\PY{l+s}{food}\PY{l+s}{\PYZsq{}}\PY{p}{:} \PY{p}{[}\PY{l+s}{\PYZsq{}}\PY{l+s}{chinese}\PY{l+s}{\PYZsq{}}\PY{p}{,} \PY{l+s}{\PYZsq{}}\PY{l+s}{indian}\PY{l+s}{\PYZsq{}}\PY{p}{,} \PY{l+s}{\PYZsq{}}\PY{l+s}{french}\PY{l+s}{\PYZsq{}}\PY{p}{]}\PY{p}{,}
\PY{g+gp}{... }   \PY{p}{\PYZcb{}}\PY{p}{,}
\PY{g+gp}{... }   \PY{p}{\PYZob{}}
\PY{g+gp}{... }       \PY{p}{(}\PY{l+s}{\PYZsq{}}\PY{l+s}{chinese}\PY{l+s}{\PYZsq{}}\PY{p}{,}\PY{p}{)}\PY{p}{:} \PY{l+m+mf}{0.3}\PY{p}{,}
\PY{g+gp}{... }       \PY{p}{(}\PY{l+s}{\PYZsq{}}\PY{l+s}{indian}\PY{l+s}{\PYZsq{}}\PY{p}{,}\PY{p}{)}\PY{p}{:} \PY{l+m+mf}{0.6}\PY{p}{,}
\PY{g+gp}{... }       \PY{p}{(}\PY{l+s}{\PYZsq{}}\PY{l+s}{french}\PY{l+s}{\PYZsq{}}\PY{p}{,}\PY{p}{)}\PY{p}{:} \PY{l+m+mf}{0.1}\PY{p}{,}
\PY{g+gp}{... }   \PY{p}{\PYZcb{}}\PY{p}{)}
\PY{g+gp}{\PYZgt{}\PYZgt{}\PYZgt{} }\PY{n}{food}\PY{o}{.}\PY{n}{most\PYZus{}probable}\PY{p}{(}\PY{n}{n}\PY{o}{=}\PY{l+m+mi}{2}\PY{p}{)}
\PY{g+go}{[(\PYZsq{}indian\PYZsq{}, 0.6), (\PYZsq{}chinese\PYZsq{}, 0.3)]}
\end{Verbatim}
\caption{Nejpravděpodobnější přiřazení}
\label{lst:facmop}
\end{example}

%\begin{lstlisting}[
%	caption={Nejpravděpodobnější hodnoty},
%	label={lst:facmop}
%]
%>>> food = Factor(
%...    ['food'],
%...    {
%...        'food': ['chinese', 'indian', 'french'],
%...    },
%...    {
%...        ('chinese',): 0.3,
%...        ('indian',): 0.6,
%...        ('french',): 0.1,
%...    })
%>>> food.most_probable(n=2)
%[('indian', 0.6), ('chinese', 0.3)]
%\end{lstlisting}

\subsubsection{Vrcholy}

V~této sekci ukážeme vytváření jednoduchého grafického modelu a následně poslání zpráv z~jednoho vrcholu do druhého.
V~příkladu~\ref{lst:nodex} vytvoříme jednoduchý Skrytý Markovský model pro výběr typu jídla.
Model bude simulovat dvě obrátky, v~každé bude jedna skrytá a jedna pozorovaná proměnná.
Tento model bude generativní s faktorem pro generování pozorování na základě skutečné hodnoty a~také faktorem pro přechodovou pravděpodobnost z~jedné obrátky do druhé.

\begin{example}
\begin{Verbatim}[commandchars=\\\{\}, fontsize=\relsize{-1}]

\PY{k+kn}{from} \PY{n+nn}{alex.infer.node} \PY{k+kn}{import} \PY{n}{DiscreteVariableNode}\PY{p}{,} \PY{n}{DiscreteFactorNode}
\PY{k+kn}{from} \PY{n+nn}{alex.infer.factor} \PY{k+kn}{import} \PY{n}{Factor}

\PY{n}{hid\PYZus{}1} \PY{o}{=} \PY{n}{DiscreteVariableNode}\PY{p}{(}\PY{l+s}{\PYZsq{}}\PY{l+s}{food\PYZus{}1}\PY{l+s}{\PYZsq{}}\PY{p}{,} \PY{p}{[}\PY{l+s}{\PYZsq{}}\PY{l+s}{chinese}\PY{l+s}{\PYZsq{}}\PY{p}{,} \PY{l+s}{\PYZsq{}}\PY{l+s}{indian}\PY{l+s}{\PYZsq{}}\PY{p}{]}\PY{p}{)}
\PY{n}{obs\PYZus{}1} \PY{o}{=} \PY{n}{DiscreteVariableNode}\PY{p}{(}\PY{l+s}{\PYZsq{}}\PY{l+s}{food\PYZus{}obs\PYZus{}1}\PY{l+s}{\PYZsq{}}\PY{p}{,} \PY{p}{[}\PY{l+s}{\PYZsq{}}\PY{l+s}{chinese}\PY{l+s}{\PYZsq{}}\PY{p}{,} \PY{l+s}{\PYZsq{}}\PY{l+s}{indian}\PY{l+s}{\PYZsq{}}\PY{p}{]}\PY{p}{)}
\PY{n}{obs\PYZus{}factor\PYZus{}1} \PY{o}{=} \PY{n}{DiscreteFactorNode}\PY{p}{(}\PY{l+s}{\PYZsq{}}\PY{l+s}{food\PYZus{}obs\PYZus{}factor\PYZus{}1}\PY{l+s}{\PYZsq{}}\PY{p}{,} \PY{n}{Factor}\PY{p}{(}
    \PY{p}{[}\PY{l+s}{\PYZsq{}}\PY{l+s}{food\PYZus{}1}\PY{l+s}{\PYZsq{}}\PY{p}{,} \PY{l+s}{\PYZsq{}}\PY{l+s}{food\PYZus{}obs\PYZus{}1}\PY{l+s}{\PYZsq{}}\PY{p}{]}\PY{p}{,}
    \PY{p}{\PYZob{}}
        \PY{l+s}{\PYZsq{}}\PY{l+s}{food\PYZus{}1}\PY{l+s}{\PYZsq{}}\PY{p}{:} \PY{p}{[}\PY{l+s}{\PYZsq{}}\PY{l+s}{chinese}\PY{l+s}{\PYZsq{}}\PY{p}{,} \PY{l+s}{\PYZsq{}}\PY{l+s}{indian}\PY{l+s}{\PYZsq{}}\PY{p}{]}\PY{p}{,}
        \PY{l+s}{\PYZsq{}}\PY{l+s}{food\PYZus{}obs\PYZus{}1}\PY{l+s}{\PYZsq{}}\PY{p}{:} \PY{p}{[}\PY{l+s}{\PYZsq{}}\PY{l+s}{chinese}\PY{l+s}{\PYZsq{}}\PY{p}{,} \PY{l+s}{\PYZsq{}}\PY{l+s}{indian}\PY{l+s}{\PYZsq{}}\PY{p}{]}\PY{p}{,}
    \PY{p}{\PYZcb{}}\PY{p}{,}
    \PY{p}{\PYZob{}}
        \PY{p}{(}\PY{l+s}{\PYZsq{}}\PY{l+s}{chinese}\PY{l+s}{\PYZsq{}}\PY{p}{,} \PY{l+s}{\PYZsq{}}\PY{l+s}{chinese}\PY{l+s}{\PYZsq{}}\PY{p}{)}\PY{p}{:} \PY{l+m+mf}{0.9}\PY{p}{,}
        \PY{p}{(}\PY{l+s}{\PYZsq{}}\PY{l+s}{chinese}\PY{l+s}{\PYZsq{}}\PY{p}{,} \PY{l+s}{\PYZsq{}}\PY{l+s}{indian}\PY{l+s}{\PYZsq{}}\PY{p}{)}\PY{p}{:} \PY{l+m+mf}{0.1}\PY{p}{,}
        \PY{p}{(}\PY{l+s}{\PYZsq{}}\PY{l+s}{indian}\PY{l+s}{\PYZsq{}}\PY{p}{,} \PY{l+s}{\PYZsq{}}\PY{l+s}{chinese}\PY{l+s}{\PYZsq{}}\PY{p}{)}\PY{p}{:} \PY{l+m+mf}{0.1}\PY{p}{,}
        \PY{p}{(}\PY{l+s}{\PYZsq{}}\PY{l+s}{indian}\PY{l+s}{\PYZsq{}}\PY{p}{,} \PY{l+s}{\PYZsq{}}\PY{l+s}{indian}\PY{l+s}{\PYZsq{}}\PY{p}{)}\PY{p}{:} \PY{l+m+mf}{0.9}\PY{p}{,}
    \PY{p}{\PYZcb{}}\PY{p}{)}\PY{p}{)}

\PY{n}{hid\PYZus{}2} \PY{o}{=} \PY{n}{DiscreteVariableNode}\PY{p}{(}\PY{l+s}{\PYZsq{}}\PY{l+s}{food\PYZus{}2}\PY{l+s}{\PYZsq{}}\PY{p}{,} \PY{p}{[}\PY{l+s}{\PYZsq{}}\PY{l+s}{chinese}\PY{l+s}{\PYZsq{}}\PY{p}{,} \PY{l+s}{\PYZsq{}}\PY{l+s}{indian}\PY{l+s}{\PYZsq{}}\PY{p}{]}\PY{p}{)}
\PY{n}{obs\PYZus{}2} \PY{o}{=} \PY{n}{DiscreteVariableNode}\PY{p}{(}\PY{l+s}{\PYZsq{}}\PY{l+s}{food\PYZus{}obs\PYZus{}2}\PY{l+s}{\PYZsq{}}\PY{p}{,} \PY{p}{[}\PY{l+s}{\PYZsq{}}\PY{l+s}{chinese}\PY{l+s}{\PYZsq{}}\PY{p}{,} \PY{l+s}{\PYZsq{}}\PY{l+s}{indian}\PY{l+s}{\PYZsq{}}\PY{p}{]}\PY{p}{)}
\PY{n}{obs\PYZus{}factor\PYZus{}2} \PY{o}{=} \PY{n}{DiscreteFactorNode}\PY{p}{(}\PY{l+s}{\PYZsq{}}\PY{l+s}{food\PYZus{}obs\PYZus{}factor\PYZus{}2}\PY{l+s}{\PYZsq{}}\PY{p}{,} \PY{n}{Factor}\PY{p}{(}
    \PY{p}{[}\PY{l+s}{\PYZsq{}}\PY{l+s}{food\PYZus{}2}\PY{l+s}{\PYZsq{}}\PY{p}{,} \PY{l+s}{\PYZsq{}}\PY{l+s}{food\PYZus{}obs\PYZus{}2}\PY{l+s}{\PYZsq{}}\PY{p}{]}\PY{p}{,}
    \PY{p}{\PYZob{}}
        \PY{l+s}{\PYZsq{}}\PY{l+s}{food\PYZus{}2}\PY{l+s}{\PYZsq{}}\PY{p}{:} \PY{p}{[}\PY{l+s}{\PYZsq{}}\PY{l+s}{chinese}\PY{l+s}{\PYZsq{}}\PY{p}{,} \PY{l+s}{\PYZsq{}}\PY{l+s}{indian}\PY{l+s}{\PYZsq{}}\PY{p}{]}\PY{p}{,}
        \PY{l+s}{\PYZsq{}}\PY{l+s}{food\PYZus{}obs\PYZus{}2}\PY{l+s}{\PYZsq{}}\PY{p}{:} \PY{p}{[}\PY{l+s}{\PYZsq{}}\PY{l+s}{chinese}\PY{l+s}{\PYZsq{}}\PY{p}{,} \PY{l+s}{\PYZsq{}}\PY{l+s}{indian}\PY{l+s}{\PYZsq{}}\PY{p}{]}\PY{p}{,}
    \PY{p}{\PYZcb{}}\PY{p}{,}
    \PY{p}{\PYZob{}}
        \PY{p}{(}\PY{l+s}{\PYZsq{}}\PY{l+s}{chinese}\PY{l+s}{\PYZsq{}}\PY{p}{,} \PY{l+s}{\PYZsq{}}\PY{l+s}{chinese}\PY{l+s}{\PYZsq{}}\PY{p}{)}\PY{p}{:} \PY{l+m+mf}{0.9}\PY{p}{,}
        \PY{p}{(}\PY{l+s}{\PYZsq{}}\PY{l+s}{chinese}\PY{l+s}{\PYZsq{}}\PY{p}{,} \PY{l+s}{\PYZsq{}}\PY{l+s}{indian}\PY{l+s}{\PYZsq{}}\PY{p}{)}\PY{p}{:} \PY{l+m+mf}{0.1}\PY{p}{,}
        \PY{p}{(}\PY{l+s}{\PYZsq{}}\PY{l+s}{indian}\PY{l+s}{\PYZsq{}}\PY{p}{,} \PY{l+s}{\PYZsq{}}\PY{l+s}{chinese}\PY{l+s}{\PYZsq{}}\PY{p}{)}\PY{p}{:} \PY{l+m+mf}{0.1}\PY{p}{,}
        \PY{p}{(}\PY{l+s}{\PYZsq{}}\PY{l+s}{indian}\PY{l+s}{\PYZsq{}}\PY{p}{,} \PY{l+s}{\PYZsq{}}\PY{l+s}{indian}\PY{l+s}{\PYZsq{}}\PY{p}{)}\PY{p}{:} \PY{l+m+mf}{0.9}\PY{p}{,}
    \PY{p}{\PYZcb{}}\PY{p}{)}\PY{p}{)}

\PY{n}{trans\PYZus{}factor} \PY{o}{=} \PY{n}{DiscreteFactorNode}\PY{p}{(}\PY{l+s}{\PYZsq{}}\PY{l+s}{food\PYZus{}trans\PYZus{}factor}\PY{l+s}{\PYZsq{}}\PY{p}{,} \PY{n}{Factor}\PY{p}{(}
    \PY{p}{[}\PY{l+s}{\PYZsq{}}\PY{l+s}{food\PYZus{}1}\PY{l+s}{\PYZsq{}}\PY{p}{,} \PY{l+s}{\PYZsq{}}\PY{l+s}{food\PYZus{}2}\PY{l+s}{\PYZsq{}}\PY{p}{]}\PY{p}{,}
    \PY{p}{\PYZob{}}
        \PY{l+s}{\PYZsq{}}\PY{l+s}{food\PYZus{}1}\PY{l+s}{\PYZsq{}}\PY{p}{:} \PY{p}{[}\PY{l+s}{\PYZsq{}}\PY{l+s}{chinese}\PY{l+s}{\PYZsq{}}\PY{p}{,} \PY{l+s}{\PYZsq{}}\PY{l+s}{indian}\PY{l+s}{\PYZsq{}}\PY{p}{]}\PY{p}{,}
        \PY{l+s}{\PYZsq{}}\PY{l+s}{food\PYZus{}2}\PY{l+s}{\PYZsq{}}\PY{p}{:} \PY{p}{[}\PY{l+s}{\PYZsq{}}\PY{l+s}{chinese}\PY{l+s}{\PYZsq{}}\PY{p}{,} \PY{l+s}{\PYZsq{}}\PY{l+s}{indian}\PY{l+s}{\PYZsq{}}\PY{p}{]}\PY{p}{,}
    \PY{p}{\PYZcb{}}\PY{p}{,}
    \PY{p}{\PYZob{}}
        \PY{p}{(}\PY{l+s}{\PYZsq{}}\PY{l+s}{chinese}\PY{l+s}{\PYZsq{}}\PY{p}{,} \PY{l+s}{\PYZsq{}}\PY{l+s}{chinese}\PY{l+s}{\PYZsq{}}\PY{p}{)}\PY{p}{:} \PY{l+m+mf}{0.99}\PY{p}{,}
        \PY{p}{(}\PY{l+s}{\PYZsq{}}\PY{l+s}{chinese}\PY{l+s}{\PYZsq{}}\PY{p}{,} \PY{l+s}{\PYZsq{}}\PY{l+s}{indian}\PY{l+s}{\PYZsq{}}\PY{p}{)}\PY{p}{:} \PY{l+m+mf}{0.01}\PY{p}{,}
        \PY{p}{(}\PY{l+s}{\PYZsq{}}\PY{l+s}{indian}\PY{l+s}{\PYZsq{}}\PY{p}{,} \PY{l+s}{\PYZsq{}}\PY{l+s}{chinese}\PY{l+s}{\PYZsq{}}\PY{p}{)}\PY{p}{:} \PY{l+m+mf}{0.01}\PY{p}{,}
        \PY{p}{(}\PY{l+s}{\PYZsq{}}\PY{l+s}{indian}\PY{l+s}{\PYZsq{}}\PY{p}{,} \PY{l+s}{\PYZsq{}}\PY{l+s}{indian}\PY{l+s}{\PYZsq{}}\PY{p}{)}\PY{p}{:} \PY{l+m+mf}{0.99}\PY{p}{,}
    \PY{p}{\PYZcb{}}\PY{p}{)}\PY{p}{)}

\PY{n}{obs\PYZus{}factor\PYZus{}1}\PY{o}{.}\PY{n}{connect}\PY{p}{(}\PY{n}{hid\PYZus{}1}\PY{p}{,} \PY{n}{parent}\PY{o}{=}\PY{n+nb+bp}{True}\PY{p}{)}
\PY{n}{obs\PYZus{}factor\PYZus{}1}\PY{o}{.}\PY{n}{connect}\PY{p}{(}\PY{n}{obs\PYZus{}1}\PY{p}{,} \PY{n}{parent}\PY{o}{=}\PY{n+nb+bp}{False}\PY{p}{)}

\PY{n}{obs\PYZus{}factor\PYZus{}2}\PY{o}{.}\PY{n}{connect}\PY{p}{(}\PY{n}{hid\PYZus{}2}\PY{p}{,} \PY{n}{parent}\PY{o}{=}\PY{n+nb+bp}{True}\PY{p}{)}
\PY{n}{obs\PYZus{}factor\PYZus{}2}\PY{o}{.}\PY{n}{connect}\PY{p}{(}\PY{n}{obs\PYZus{}2}\PY{p}{,} \PY{n}{parent}\PY{o}{=}\PY{n+nb+bp}{False}\PY{p}{)}

\PY{n}{trans\PYZus{}factor}\PY{o}{.}\PY{n}{connect}\PY{p}{(}\PY{n}{hid\PYZus{}1}\PY{p}{,} \PY{n}{parent}\PY{o}{=}\PY{n+nb+bp}{True}\PY{p}{)}
\PY{n}{trans\PYZus{}factor}\PY{o}{.}\PY{n}{connect}\PY{p}{(}\PY{n}{hid\PYZus{}2}\PY{p}{,} \PY{n}{parent}\PY{o}{=}\PY{n+nb+bp}{False}\PY{p}{)}
\end{Verbatim}
\caption{Jednoduchý generativní model}
\label{lst:nodex}
\end{example}


%\begin{lstlisting}[
%	caption={Jednoduchý generativní model},
%	label={lst:nodex},
%	float=hp
%]
%from alex.ml.bn.node import DiscreteVariableNode, DiscreteFactorNode
%from alex.ml.bn.factor import Factor
%
%hid_1 = DiscreteVariableNode('food_1', ['chinese', 'indian'])
%obs_1 = DiscreteVariableNode('food_obs_1', ['chinese', 'indian'])
%obs_factor_1 = DiscreteFactorNode('food_obs_factor_1', Factor(
%    ['food_1', 'food_obs_1'],
%    {
%        'food_1': ['chinese', 'indian'],
%        'food_obs_1': ['chinese', 'indian'],
%    },
%    {
%        ('chinese', 'chinese'): 0.9,
%        ('chinese', 'indian'): 0.1,
%        ('indian', 'chinese'): 0.1,
%        ('indian', 'indian'): 0.9,
%    }))
%
%hid_2 = DiscreteVariableNode('food_2', ['chinese', 'indian'])
%obs_2 = DiscreteVariableNode('food_obs_2', ['chinese', 'indian'])
%obs_factor_2 = DiscreteFactorNode('food_obs_factor_2', Factor(
%    ['food_2', 'food_obs_2'],
%    {
%        'food_2': ['chinese', 'indian'],
%        'food_obs_2': ['chinese', 'indian'],
%    },
%    {
%        ('chinese', 'chinese'): 0.9,
%        ('chinese', 'indian'): 0.1,
%        ('indian', 'chinese'): 0.1,
%        ('indian', 'indian'): 0.9,
%    }))
%
%trans_factor = DiscreteFactorNode('food_trans_factor', Factor(
%    ['food_1', 'food_2'],
%    {
%        'food_1': ['chinese', 'indian'],
%        'food_2': ['chinese', 'indian'],
%    },
%    {
%        ('chinese', 'chinese'): 0.99,
%        ('chinese', 'indian'): 0.01,
%        ('indian', 'chinese'): 0.01,
%        ('indian', 'indian'): 0.99,
%    }))
%
%obs_factor_1.connect(hid_1, parent=True)
%obs_factor_1.connect(obs_1, parent=False)
%
%obs_factor_2.connect(hid_2, parent=True)
%obs_factor_2.connect(obs_2, parent=False)
%
%trans_factor.connect(hid_1, parent=True)
%trans_factor.connect(hid_2, parent=False)
%\end{lstlisting}

Po vytvoření grafického modelu je dalším krokem nastavení hodnot pozorovaných přiřazení (příklad~\ref{lst:facobs}).
Předpokládejme, že v~první obrátce bylo pozorováno čínské jídlo s~pravděpodobností $0.6$ a indické s~pravděpodobností $0.4$.
V~druhé obrátce bylo pozorováno čínské jídlo s~pravděpodobností $0.5$ a indické s~pravděpodobností $0.5$.
V~první obrátce si myslíme, že uživatel spíše preferuje čínské jídlo.
Z~druhé obrátky víme jen, že uživatel preferuje čínské nebo indické jídlo. 
Podíváme se, co z~těchto informací zjistí náš grafický model.

\begin{example}
\begin{Verbatim}[commandchars=\\\{\}]

\PY{n}{obs\PYZus{}1}\PY{o}{.}\PY{n}{observed}\PY{p}{(}\PY{p}{\PYZob{}}
	\PY{p}{(}\PY{l+s}{\PYZsq{}}\PY{l+s}{chinese}\PY{l+s}{\PYZsq{}}\PY{p}{,}\PY{p}{)}\PY{p}{:} \PY{l+m+mf}{0.6}\PY{p}{,}
	\PY{p}{(}\PY{l+s}{\PYZsq{}}\PY{l+s}{indian}\PY{l+s}{\PYZsq{}}\PY{p}{,}\PY{p}{)}\PY{p}{:} \PY{l+m+mf}{0.4}\PY{p}{,}
\PY{p}{\PYZcb{}}\PY{p}{)}

\PY{n}{obs\PYZus{}2}\PY{o}{.}\PY{n}{observed}\PY{p}{(}\PY{p}{\PYZob{}}
	\PY{p}{(}\PY{l+s}{\PYZsq{}}\PY{l+s}{chinese}\PY{l+s}{\PYZsq{}}\PY{p}{,}\PY{p}{)}\PY{p}{:} \PY{l+m+mf}{0.5}\PY{p}{,}
	\PY{p}{(}\PY{l+s}{\PYZsq{}}\PY{l+s}{indian}\PY{l+s}{\PYZsq{}}\PY{p}{,}\PY{p}{)}\PY{p}{:} \PY{l+m+mf}{0.5}\PY{p}{,}
\PY{p}{\PYZcb{}}\PY{p}{)}
\end{Verbatim}
\caption{Nastavení hodnot pozorovaných přiřazení}
\label{lst:facobs}
\end{example}

%\begin{lstlisting}[
%	caption={Nastavení pozorovaných hodnot},
%	label={lst:facobs},
%    float=h,
%]
%obs_1.observed({
%	('chinese',): 0.6,
%	('indian',): 0.4,
%})
%
%obs_2.observed({
%	('chinese',): 0.5,
%	('indian',): 0.5,
%})
%\end{lstlisting}

Strom zakořeníme ve vrcholu \texttt{trans\_factor} a začneme posílat zprávy od listů ke kořeni a pak zpět (příklad~\ref{lst:facmsg}).
Naším cílem je zjistit aposteriorní marginální pravděpodobnosti skrytých proměnných.
Před odesláním je vždy třeba aktualizovat vnitřní reprezentaci vrcholů.
Nakonec musíme normalizovat vrcholy s~proměnnými.

\begin{example}
\begin{Verbatim}[commandchars=\\\{\}]

\PY{n}{obs\PYZus{}1}\PY{o}{.}\PY{n}{message\PYZus{}to}\PY{p}{(}\PY{n}{obs\PYZus{}factor\PYZus{}1}\PY{p}{)}
\PY{n}{obs\PYZus{}2}\PY{o}{.}\PY{n}{message\PYZus{}to}\PY{p}{(}\PY{n}{obs\PYZus{}factor\PYZus{}2}\PY{p}{)}

\PY{n}{obs\PYZus{}factor\PYZus{}1}\PY{o}{.}\PY{n}{update}\PY{p}{(}\PY{p}{)}
\PY{n}{obs\PYZus{}factor\PYZus{}1}\PY{o}{.}\PY{n}{message\PYZus{}to}\PY{p}{(}\PY{n}{hid\PYZus{}1}\PY{p}{)}

\PY{n}{obs\PYZus{}factor\PYZus{}2}\PY{o}{.}\PY{n}{update}\PY{p}{(}\PY{p}{)}
\PY{n}{obs\PYZus{}factor\PYZus{}2}\PY{o}{.}\PY{n}{message\PYZus{}to}\PY{p}{(}\PY{n}{hid\PYZus{}2}\PY{p}{)}

\PY{n}{hid\PYZus{}1}\PY{o}{.}\PY{n}{update}\PY{p}{(}\PY{p}{)}
\PY{n}{hid\PYZus{}1}\PY{o}{.}\PY{n}{message\PYZus{}to}\PY{p}{(}\PY{n}{trans\PYZus{}factor}\PY{p}{)}

\PY{n}{hid\PYZus{}2}\PY{o}{.}\PY{n}{update}\PY{p}{(}\PY{p}{)}
\PY{n}{hid\PYZus{}2}\PY{o}{.}\PY{n}{message\PYZus{}to}\PY{p}{(}\PY{n}{trans\PYZus{}factor}\PY{p}{)}

\PY{n}{trans\PYZus{}factor}\PY{o}{.}\PY{n}{update}\PY{p}{(}\PY{p}{)}
\PY{n}{trans\PYZus{}factor}\PY{o}{.}\PY{n}{send\PYZus{}messages}\PY{p}{(}\PY{p}{)}

\PY{n}{hid\PYZus{}1}\PY{o}{.}\PY{n}{update}\PY{p}{(}\PY{p}{)}
\PY{n}{hid\PYZus{}1}\PY{o}{.}\PY{n}{normalize}\PY{p}{(}\PY{p}{)}

\PY{n}{hid\PYZus{}2}\PY{o}{.}\PY{n}{update}\PY{p}{(}\PY{p}{)}
\PY{n}{hid\PYZus{}2}\PY{o}{.}\PY{n}{normalize}\PY{p}{(}\PY{p}{)}               
\end{Verbatim}
\caption{Posílání zpráv}
\label{lst:facmsg}
\end{example}

%\begin{lstlisting}[
%	caption={Posílání zpráv},
%	label={lst:facmsg},
%	float=h!
%]
%obs_1.message_to(obs_factor_1)
%obs_2.message_to(obs_factor_2)
%
%obs_factor_1.update()
%obs_factor_1.message_to(hid_1)
%
%obs_factor_2.update()
%obs_factor_2.message_to(hid_2)
%
%hid_1.update()
%hid_1.message_to(trans_factor)
%
%hid_2.update()
%hid_2.message_to(trans_factor)
%
%trans_factor.update()
%trans_factor.send_messages()
%
%hid_1.update()
%hid_1.normalize()
%
%hid_2.update()
%hid_2.normalize()
%\end{lstlisting}

\begin{example}
\begin{Verbatim}[commandchars=\\\{\}]

\PY{g+gp}{\PYZgt{}\PYZgt{}\PYZgt{} }\PY{k}{print} \PY{n}{hid\PYZus{}1}\PY{o}{.}\PY{n}{belief}\PY{o}{.}\PY{n}{pretty\PYZus{}print}\PY{p}{(}\PY{l+m+mi}{50}\PY{p}{,} \PY{l+m+mi}{2}\PY{p}{)}
\PY{g+go}{\PYZhy{}\PYZhy{}\PYZhy{}\PYZhy{}\PYZhy{}\PYZhy{}\PYZhy{}\PYZhy{}\PYZhy{}\PYZhy{}\PYZhy{}\PYZhy{}\PYZhy{}\PYZhy{}\PYZhy{}\PYZhy{}\PYZhy{}\PYZhy{}\PYZhy{}\PYZhy{}\PYZhy{}\PYZhy{}\PYZhy{}\PYZhy{}\PYZhy{}\PYZhy{}\PYZhy{}\PYZhy{}\PYZhy{}\PYZhy{}\PYZhy{}\PYZhy{}\PYZhy{}\PYZhy{}\PYZhy{}\PYZhy{}\PYZhy{}\PYZhy{}\PYZhy{}\PYZhy{}\PYZhy{}\PYZhy{}\PYZhy{}\PYZhy{}\PYZhy{}\PYZhy{}\PYZhy{}\PYZhy{}\PYZhy{}\PYZhy{}}
\PY{g+go}{         food\PYZus{}1                    Value          }
\PY{g+go}{\PYZhy{}\PYZhy{}\PYZhy{}\PYZhy{}\PYZhy{}\PYZhy{}\PYZhy{}\PYZhy{}\PYZhy{}\PYZhy{}\PYZhy{}\PYZhy{}\PYZhy{}\PYZhy{}\PYZhy{}\PYZhy{}\PYZhy{}\PYZhy{}\PYZhy{}\PYZhy{}\PYZhy{}\PYZhy{}\PYZhy{}\PYZhy{}\PYZhy{}\PYZhy{}\PYZhy{}\PYZhy{}\PYZhy{}\PYZhy{}\PYZhy{}\PYZhy{}\PYZhy{}\PYZhy{}\PYZhy{}\PYZhy{}\PYZhy{}\PYZhy{}\PYZhy{}\PYZhy{}\PYZhy{}\PYZhy{}\PYZhy{}\PYZhy{}\PYZhy{}\PYZhy{}\PYZhy{}\PYZhy{}\PYZhy{}\PYZhy{}}
\PY{g+go}{         chinese                   0.58           }
\PY{g+go}{         indian                    0.42           }
\PY{g+go}{\PYZhy{}\PYZhy{}\PYZhy{}\PYZhy{}\PYZhy{}\PYZhy{}\PYZhy{}\PYZhy{}\PYZhy{}\PYZhy{}\PYZhy{}\PYZhy{}\PYZhy{}\PYZhy{}\PYZhy{}\PYZhy{}\PYZhy{}\PYZhy{}\PYZhy{}\PYZhy{}\PYZhy{}\PYZhy{}\PYZhy{}\PYZhy{}\PYZhy{}\PYZhy{}\PYZhy{}\PYZhy{}\PYZhy{}\PYZhy{}\PYZhy{}\PYZhy{}\PYZhy{}\PYZhy{}\PYZhy{}\PYZhy{}\PYZhy{}\PYZhy{}\PYZhy{}\PYZhy{}\PYZhy{}\PYZhy{}\PYZhy{}\PYZhy{}\PYZhy{}\PYZhy{}\PYZhy{}\PYZhy{}\PYZhy{}\PYZhy{}}
\PY{g+gp}{\PYZgt{}\PYZgt{}\PYZgt{} }\PY{k}{print} \PY{n}{hid\PYZus{}2}\PY{o}{.}\PY{n}{belief}\PY{o}{.}\PY{n}{pretty\PYZus{}print}\PY{p}{(}\PY{l+m+mi}{50}\PY{p}{,} \PY{l+m+mi}{2}\PY{p}{)}
\PY{g+go}{\PYZhy{}\PYZhy{}\PYZhy{}\PYZhy{}\PYZhy{}\PYZhy{}\PYZhy{}\PYZhy{}\PYZhy{}\PYZhy{}\PYZhy{}\PYZhy{}\PYZhy{}\PYZhy{}\PYZhy{}\PYZhy{}\PYZhy{}\PYZhy{}\PYZhy{}\PYZhy{}\PYZhy{}\PYZhy{}\PYZhy{}\PYZhy{}\PYZhy{}\PYZhy{}\PYZhy{}\PYZhy{}\PYZhy{}\PYZhy{}\PYZhy{}\PYZhy{}\PYZhy{}\PYZhy{}\PYZhy{}\PYZhy{}\PYZhy{}\PYZhy{}\PYZhy{}\PYZhy{}\PYZhy{}\PYZhy{}\PYZhy{}\PYZhy{}\PYZhy{}\PYZhy{}\PYZhy{}\PYZhy{}\PYZhy{}\PYZhy{}}
\PY{g+go}{         food\PYZus{}2                    Value          }
\PY{g+go}{\PYZhy{}\PYZhy{}\PYZhy{}\PYZhy{}\PYZhy{}\PYZhy{}\PYZhy{}\PYZhy{}\PYZhy{}\PYZhy{}\PYZhy{}\PYZhy{}\PYZhy{}\PYZhy{}\PYZhy{}\PYZhy{}\PYZhy{}\PYZhy{}\PYZhy{}\PYZhy{}\PYZhy{}\PYZhy{}\PYZhy{}\PYZhy{}\PYZhy{}\PYZhy{}\PYZhy{}\PYZhy{}\PYZhy{}\PYZhy{}\PYZhy{}\PYZhy{}\PYZhy{}\PYZhy{}\PYZhy{}\PYZhy{}\PYZhy{}\PYZhy{}\PYZhy{}\PYZhy{}\PYZhy{}\PYZhy{}\PYZhy{}\PYZhy{}\PYZhy{}\PYZhy{}\PYZhy{}\PYZhy{}\PYZhy{}\PYZhy{}}
\PY{g+go}{         chinese                   0.58           }
\PY{g+go}{         indian                    0.42           }
\PY{g+go}{\PYZhy{}\PYZhy{}\PYZhy{}\PYZhy{}\PYZhy{}\PYZhy{}\PYZhy{}\PYZhy{}\PYZhy{}\PYZhy{}\PYZhy{}\PYZhy{}\PYZhy{}\PYZhy{}\PYZhy{}\PYZhy{}\PYZhy{}\PYZhy{}\PYZhy{}\PYZhy{}\PYZhy{}\PYZhy{}\PYZhy{}\PYZhy{}\PYZhy{}\PYZhy{}\PYZhy{}\PYZhy{}\PYZhy{}\PYZhy{}\PYZhy{}\PYZhy{}\PYZhy{}\PYZhy{}\PYZhy{}\PYZhy{}\PYZhy{}\PYZhy{}\PYZhy{}\PYZhy{}\PYZhy{}\PYZhy{}\PYZhy{}\PYZhy{}\PYZhy{}\PYZhy{}\PYZhy{}\PYZhy{}\PYZhy{}\PYZhy{}}
\end{Verbatim}
\caption{Výsledek inference}
\label{lst:facres}
\end{example}

%\begin{lstlisting}[
%	caption={Výsledek inference},
%	label={lst:facres},
%	float=h!
%]
%>>> print hid_1.belief.pretty_print(50, 2)
%--------------------------------------------------
%         food_1                    Value          
%--------------------------------------------------
%         chinese                   0.58           
%         indian                    0.42           
%--------------------------------------------------
%>>> print hid_2.belief.pretty_print(50, 2)
%--------------------------------------------------
%         food_2                    Value          
%--------------------------------------------------
%         chinese                   0.58           
%         indian                    0.42           
%--------------------------------------------------
%\end{lstlisting}

Ve výsledku (příklad ~\ref{lst:facres}) vidíme, že ve~druhé obrátce z~pozorování sice nezískáme žádnou informaci, ale díky vysoké pravděpodobnosti přechodu zůstala pravděpodobnost čínského jídla stále vysoká.
Toto pozorování zároveň snížilo pravděpodobnost v~první obrátce.

\subsubsection{Inference}

Nyní se zbavíme manuálního posílání zpráv z~příkladu~\ref{lst:facmsg} a nahradíme jej použitím třídy pro Loopy Belief Propagation (příklad~\ref{lst:lbpex}). 
Zvolíme strategii pro stromy a výsledek bude stejný jako v~případě s~manuálním posíláním zpráv.

\begin{example}
\begin{Verbatim}[commandchars=\\\{\}]

\PY{k+kn}{from} \PY{n+nn}{alex.infer.lbp} \PY{k+kn}{import} \PY{n}{LBP}

\PY{n}{lbp} \PY{o}{=} \PY{n}{LBP}\PY{p}{(}\PY{n}{strategy}\PY{o}{=}\PY{l+s}{\PYZsq{}}\PY{l+s}{tree}\PY{l+s}{\PYZsq{}}\PY{p}{)}
\PY{n}{lbp}\PY{o}{.}\PY{n}{add\PYZus{}nodes}\PY{p}{(}\PY{p}{[}\PY{n}{obs\PYZus{}1}\PY{p}{,} \PY{n}{obs\PYZus{}2}\PY{p}{,} \PY{n}{hid\PYZus{}1}\PY{p}{,} \PY{n}{hid\PYZus{}2}\PY{p}{,} \PY{n}{obs\PYZus{}factor\PYZus{}1}\PY{p}{,} \PY{n}{obs\PYZus{}factor\PYZus{}2}\PY{p}{,}
               \PY{n}{trans\PYZus{}factor}\PY{p}{]}\PY{p}{)}
\PY{n}{lbp}\PY{o}{.}\PY{n}{run}\PY{p}{(}\PY{p}{)}
\end{Verbatim}
\caption{Inference s~pomocí LBP}
\label{lst:lbpex}
\end{example}

%\begin{lstlisting}[
%	caption={Inference s pomocí LBP},
%	label={lst:lbpex},
%	float=h
%]
%from alex.infer.lbp import LBP
%
%lbp = LBP(strategy='tree')
%lbp.add_nodes([obs_1, obs_2, hid_1, hid_2, obs_factor_1, obs_factor_2,
%               trans_factor])
%lbp.run()
%\end{lstlisting}

\subsection{Učení Dirichletovských parametrů}

Nyní grafický model z~minulé sekce upravíme tak, aby parametry přechodové pravděpodobnosti pocházely z~Dirichletovského rozdělení.
Pro přechodovou pravděpodobnost použijeme faktor \texttt{DirichletFactorNode} a k~němu navíc připojíme nový vrchol \texttt{DirichletParameterNode}, který bude reprezentovat parametry.
Vytvoření grafického modelu je v~příkladu~\ref{lst:crtep}.

Druhý parametr konstruktoru \texttt{DirichletParameterNode} nyní neurčuje pravděpodobnosti, ale parametry $\alpha$ Dirichletovských rozdělení.
Pro každou kombinaci hodnot rodičů existuje jedno Dirichletovské rozdělení popisující apriorní rozdělení nad parametry faktoru.
V~tomto případě je to skrytá proměnná v~druhé obrátce.
Parametry mohou být buď nastaveny bez apriorní znalosti, pak jsou všechny rovny jedné, anebo můžeme nějakou naši apriorní znalost přidat.
V~tomto případě budeme předpokládat, že se cíl spíše nemění.

Pozorování budou nastavena tak, že je zřejmé, který typ jídla uživatel požaduje.
Pokud dostaneme kvalitní vstup od uživatele, můžeme dostat velmi dobré hypotézy.
Inferenci provedeme naprosto stejně jako v~předchozím případě (příklad~\ref{lst:lbpep}).
Zajímavý je výsledek inference, protože se nezměnily pouze hodnoty skrytých proměnných, ale také Dirichletovský parametr.
V~příkladu~\ref{lst:epres} je vidět, že vypočítané hodnoty skrytých proměnných jsou blízké pozorovaným hodnotám v~každé obrátce.
To je způsobeno tím, že parametry pro přechodovou pravděpodobnost jsme nenastavili příliš informativně, a tak nejsou skryté proměnné mezi sebou příliš svázané.
Zároveň je vidět, že byly parametry modifikovány.
Pokud uživatel v~první obrátce chce čínské jídlo, pak je větší pravděpodobnost, že v~druhé obrátce jej bude chtít také.

Nyní ještě ukážeme, jak je možné použít stejný model vícekrát, s~různými pozorováními.
Představme si, že máme další hovor o~dvou obrátkách, avšak s~jinými pozorovanými hodnotami.
Stačí nastavit nová pozorování, znovu inicializovat zprávy a spustit inferenci.
Výhodou je, že není třeba znovu vytvářet síť.
Navíc počítáme už s~upraveným Dirichletovským parametrem pro přechodovou pravděpodobnost.
Výsledek je v~příkladu~\ref{lst:epint}.

\begin{example}[H]
\caption{Vytvoření grafického modelu s~Dirichletovskými parametry}
\label{lst:crtep}
\begin{Verbatim}[commandchars=\\\{\}, fontsize=\relsize{-1}]

\PY{k+kn}{from} \PY{n+nn}{alex.infer.node} \PY{k+kn}{import} \PY{p}{(}\PY{n}{DiscreteVariableNode}\PY{p}{,} \PY{n}{DiscreteFactorNode}\PY{p}{,}
                             \PY{n}{DirichletParameterNode}\PY{p}{,}
                             \PY{n}{DirichletFactorNode}\PY{p}{)}
\PY{k+kn}{from} \PY{n+nn}{alex.infer.factor} \PY{k+kn}{import} \PY{n}{Factor}
\PY{k+kn}{from} \PY{n+nn}{alex.infer.lbp} \PY{k+kn}{import} \PY{n}{LBP}

\PY{n}{obs\PYZus{}probability} \PY{o}{=} \PY{p}{\PYZob{}}
    \PY{p}{(}\PY{l+s}{\PYZsq{}}\PY{l+s}{chinese}\PY{l+s}{\PYZsq{}}\PY{p}{,} \PY{l+s}{\PYZsq{}}\PY{l+s}{chinese}\PY{l+s}{\PYZsq{}}\PY{p}{)}\PY{p}{:} \PY{l+m+mf}{0.9}\PY{p}{,}
    \PY{p}{(}\PY{l+s}{\PYZsq{}}\PY{l+s}{chinese}\PY{l+s}{\PYZsq{}}\PY{p}{,} \PY{l+s}{\PYZsq{}}\PY{l+s}{indian}\PY{l+s}{\PYZsq{}}\PY{p}{)}\PY{p}{:} \PY{l+m+mf}{0.1}\PY{p}{,}
    \PY{p}{(}\PY{l+s}{\PYZsq{}}\PY{l+s}{indian}\PY{l+s}{\PYZsq{}}\PY{p}{,} \PY{l+s}{\PYZsq{}}\PY{l+s}{chinese}\PY{l+s}{\PYZsq{}}\PY{p}{)}\PY{p}{:} \PY{l+m+mf}{0.1}\PY{p}{,}
    \PY{p}{(}\PY{l+s}{\PYZsq{}}\PY{l+s}{indian}\PY{l+s}{\PYZsq{}}\PY{p}{,} \PY{l+s}{\PYZsq{}}\PY{l+s}{indian}\PY{l+s}{\PYZsq{}}\PY{p}{)}\PY{p}{:} \PY{l+m+mf}{0.9}\PY{p}{,}
\PY{p}{\PYZcb{}}

\PY{n}{hid\PYZus{}1} \PY{o}{=} \PY{n}{DiscreteVariableNode}\PY{p}{(}\PY{l+s}{\PYZsq{}}\PY{l+s}{food\PYZus{}1}\PY{l+s}{\PYZsq{}}\PY{p}{,} \PY{p}{[}\PY{l+s}{\PYZsq{}}\PY{l+s}{chinese}\PY{l+s}{\PYZsq{}}\PY{p}{,} \PY{l+s}{\PYZsq{}}\PY{l+s}{indian}\PY{l+s}{\PYZsq{}}\PY{p}{]}\PY{p}{)}
\PY{n}{obs\PYZus{}1} \PY{o}{=} \PY{n}{DiscreteVariableNode}\PY{p}{(}\PY{l+s}{\PYZsq{}}\PY{l+s}{food\PYZus{}obs\PYZus{}1}\PY{l+s}{\PYZsq{}}\PY{p}{,} \PY{p}{[}\PY{l+s}{\PYZsq{}}\PY{l+s}{chinese}\PY{l+s}{\PYZsq{}}\PY{p}{,} \PY{l+s}{\PYZsq{}}\PY{l+s}{indian}\PY{l+s}{\PYZsq{}}\PY{p}{]}\PY{p}{)}
\PY{n}{obs\PYZus{}factor\PYZus{}1} \PY{o}{=} \PY{n}{DiscreteFactorNode}\PY{p}{(}\PY{l+s}{\PYZsq{}}\PY{l+s}{food\PYZus{}obs\PYZus{}factor\PYZus{}1}\PY{l+s}{\PYZsq{}}\PY{p}{,} \PY{n}{Factor}\PY{p}{(}
    \PY{p}{[}\PY{l+s}{\PYZsq{}}\PY{l+s}{food\PYZus{}1}\PY{l+s}{\PYZsq{}}\PY{p}{,} \PY{l+s}{\PYZsq{}}\PY{l+s}{food\PYZus{}obs\PYZus{}1}\PY{l+s}{\PYZsq{}}\PY{p}{]}\PY{p}{,}
    \PY{p}{\PYZob{}}
        \PY{l+s}{\PYZsq{}}\PY{l+s}{food\PYZus{}1}\PY{l+s}{\PYZsq{}}\PY{p}{:} \PY{p}{[}\PY{l+s}{\PYZsq{}}\PY{l+s}{chinese}\PY{l+s}{\PYZsq{}}\PY{p}{,} \PY{l+s}{\PYZsq{}}\PY{l+s}{indian}\PY{l+s}{\PYZsq{}}\PY{p}{]}\PY{p}{,}
        \PY{l+s}{\PYZsq{}}\PY{l+s}{food\PYZus{}obs\PYZus{}1}\PY{l+s}{\PYZsq{}}\PY{p}{:} \PY{p}{[}\PY{l+s}{\PYZsq{}}\PY{l+s}{chinese}\PY{l+s}{\PYZsq{}}\PY{p}{,} \PY{l+s}{\PYZsq{}}\PY{l+s}{indian}\PY{l+s}{\PYZsq{}}\PY{p}{]}\PY{p}{,}
    \PY{p}{\PYZcb{}}\PY{p}{,}
    \PY{n}{obs\PYZus{}probability}\PY{p}{)}\PY{p}{)}

\PY{n}{hid\PYZus{}2} \PY{o}{=} \PY{n}{DiscreteVariableNode}\PY{p}{(}\PY{l+s}{\PYZsq{}}\PY{l+s}{food\PYZus{}2}\PY{l+s}{\PYZsq{}}\PY{p}{,} \PY{p}{[}\PY{l+s}{\PYZsq{}}\PY{l+s}{chinese}\PY{l+s}{\PYZsq{}}\PY{p}{,} \PY{l+s}{\PYZsq{}}\PY{l+s}{indian}\PY{l+s}{\PYZsq{}}\PY{p}{]}\PY{p}{)}
\PY{n}{obs\PYZus{}2} \PY{o}{=} \PY{n}{DiscreteVariableNode}\PY{p}{(}\PY{l+s}{\PYZsq{}}\PY{l+s}{food\PYZus{}obs\PYZus{}2}\PY{l+s}{\PYZsq{}}\PY{p}{,} \PY{p}{[}\PY{l+s}{\PYZsq{}}\PY{l+s}{chinese}\PY{l+s}{\PYZsq{}}\PY{p}{,} \PY{l+s}{\PYZsq{}}\PY{l+s}{indian}\PY{l+s}{\PYZsq{}}\PY{p}{]}\PY{p}{)}
\PY{n}{obs\PYZus{}factor\PYZus{}2} \PY{o}{=} \PY{n}{DiscreteFactorNode}\PY{p}{(}\PY{l+s}{\PYZsq{}}\PY{l+s}{food\PYZus{}obs\PYZus{}factor\PYZus{}2}\PY{l+s}{\PYZsq{}}\PY{p}{,} \PY{n}{Factor}\PY{p}{(}
    \PY{p}{[}\PY{l+s}{\PYZsq{}}\PY{l+s}{food\PYZus{}2}\PY{l+s}{\PYZsq{}}\PY{p}{,} \PY{l+s}{\PYZsq{}}\PY{l+s}{food\PYZus{}obs\PYZus{}2}\PY{l+s}{\PYZsq{}}\PY{p}{]}\PY{p}{,}
    \PY{p}{\PYZob{}}
        \PY{l+s}{\PYZsq{}}\PY{l+s}{food\PYZus{}2}\PY{l+s}{\PYZsq{}}\PY{p}{:} \PY{p}{[}\PY{l+s}{\PYZsq{}}\PY{l+s}{chinese}\PY{l+s}{\PYZsq{}}\PY{p}{,} \PY{l+s}{\PYZsq{}}\PY{l+s}{indian}\PY{l+s}{\PYZsq{}}\PY{p}{]}\PY{p}{,}
        \PY{l+s}{\PYZsq{}}\PY{l+s}{food\PYZus{}obs\PYZus{}2}\PY{l+s}{\PYZsq{}}\PY{p}{:} \PY{p}{[}\PY{l+s}{\PYZsq{}}\PY{l+s}{chinese}\PY{l+s}{\PYZsq{}}\PY{p}{,} \PY{l+s}{\PYZsq{}}\PY{l+s}{indian}\PY{l+s}{\PYZsq{}}\PY{p}{]}\PY{p}{,}
    \PY{p}{\PYZcb{}}\PY{p}{,}
    \PY{n}{obs\PYZus{}probability}\PY{p}{)}\PY{p}{)}

\PY{n}{trans\PYZus{}factor} \PY{o}{=} \PY{n}{DirichletFactorNode}\PY{p}{(}\PY{l+s}{\PYZsq{}}\PY{l+s}{food\PYZus{}trans\PYZus{}factor}\PY{l+s}{\PYZsq{}}\PY{p}{)}
\PY{n}{trans\PYZus{}param} \PY{o}{=} \PY{n}{DirichletParameterNode}\PY{p}{(}\PY{l+s}{\PYZsq{}}\PY{l+s}{food\PYZus{}trans\PYZus{}param}\PY{l+s}{\PYZsq{}}\PY{p}{,} \PY{n}{Factor}\PY{p}{(}
    \PY{p}{[}\PY{l+s}{\PYZsq{}}\PY{l+s}{food\PYZus{}1}\PY{l+s}{\PYZsq{}}\PY{p}{,} \PY{l+s}{\PYZsq{}}\PY{l+s}{food\PYZus{}2}\PY{l+s}{\PYZsq{}}\PY{p}{]}\PY{p}{,}
    \PY{p}{\PYZob{}}
        \PY{l+s}{\PYZsq{}}\PY{l+s}{food\PYZus{}1}\PY{l+s}{\PYZsq{}}\PY{p}{:} \PY{p}{[}\PY{l+s}{\PYZsq{}}\PY{l+s}{chinese}\PY{l+s}{\PYZsq{}}\PY{p}{,} \PY{l+s}{\PYZsq{}}\PY{l+s}{indian}\PY{l+s}{\PYZsq{}}\PY{p}{]}\PY{p}{,}
        \PY{l+s}{\PYZsq{}}\PY{l+s}{food\PYZus{}2}\PY{l+s}{\PYZsq{}}\PY{p}{:} \PY{p}{[}\PY{l+s}{\PYZsq{}}\PY{l+s}{chinese}\PY{l+s}{\PYZsq{}}\PY{p}{,} \PY{l+s}{\PYZsq{}}\PY{l+s}{indian}\PY{l+s}{\PYZsq{}}\PY{p}{]}\PY{p}{,}
    \PY{p}{\PYZcb{}}\PY{p}{,}
    \PY{p}{\PYZob{}}
        \PY{p}{(}\PY{l+s}{\PYZsq{}}\PY{l+s}{chinese}\PY{l+s}{\PYZsq{}}\PY{p}{,} \PY{l+s}{\PYZsq{}}\PY{l+s}{chinese}\PY{l+s}{\PYZsq{}}\PY{p}{)}\PY{p}{:} \PY{l+m+mi}{2}\PY{p}{,}
        \PY{p}{(}\PY{l+s}{\PYZsq{}}\PY{l+s}{chinese}\PY{l+s}{\PYZsq{}}\PY{p}{,} \PY{l+s}{\PYZsq{}}\PY{l+s}{indian}\PY{l+s}{\PYZsq{}}\PY{p}{)}\PY{p}{:} \PY{l+m+mi}{1}\PY{p}{,}
        \PY{p}{(}\PY{l+s}{\PYZsq{}}\PY{l+s}{indian}\PY{l+s}{\PYZsq{}}\PY{p}{,} \PY{l+s}{\PYZsq{}}\PY{l+s}{chinese}\PY{l+s}{\PYZsq{}}\PY{p}{)}\PY{p}{:} \PY{l+m+mi}{1}\PY{p}{,}
        \PY{p}{(}\PY{l+s}{\PYZsq{}}\PY{l+s}{indian}\PY{l+s}{\PYZsq{}}\PY{p}{,} \PY{l+s}{\PYZsq{}}\PY{l+s}{indian}\PY{l+s}{\PYZsq{}}\PY{p}{)}\PY{p}{:} \PY{l+m+mi}{2}\PY{p}{,}
    \PY{p}{\PYZcb{}}\PY{p}{)}\PY{p}{)}

\PY{n}{obs\PYZus{}factor\PYZus{}1}\PY{o}{.}\PY{n}{connect}\PY{p}{(}\PY{n}{hid\PYZus{}1}\PY{p}{,} \PY{n}{parent}\PY{o}{=}\PY{n+nb+bp}{True}\PY{p}{)}
\PY{n}{obs\PYZus{}factor\PYZus{}1}\PY{o}{.}\PY{n}{connect}\PY{p}{(}\PY{n}{obs\PYZus{}1}\PY{p}{,} \PY{n}{parent}\PY{o}{=}\PY{n+nb+bp}{False}\PY{p}{)}

\PY{n}{obs\PYZus{}factor\PYZus{}2}\PY{o}{.}\PY{n}{connect}\PY{p}{(}\PY{n}{hid\PYZus{}2}\PY{p}{,} \PY{n}{parent}\PY{o}{=}\PY{n+nb+bp}{True}\PY{p}{)}
\PY{n}{obs\PYZus{}factor\PYZus{}2}\PY{o}{.}\PY{n}{connect}\PY{p}{(}\PY{n}{obs\PYZus{}2}\PY{p}{,} \PY{n}{parent}\PY{o}{=}\PY{n+nb+bp}{False}\PY{p}{)}

\PY{n}{trans\PYZus{}factor}\PY{o}{.}\PY{n}{connect}\PY{p}{(}\PY{n}{hid\PYZus{}1}\PY{p}{,} \PY{n}{parent}\PY{o}{=}\PY{n+nb+bp}{True}\PY{p}{)}
\PY{n}{trans\PYZus{}factor}\PY{o}{.}\PY{n}{connect}\PY{p}{(}\PY{n}{hid\PYZus{}2}\PY{p}{,} \PY{n}{parent}\PY{o}{=}\PY{n+nb+bp}{False}\PY{p}{)}
\PY{n}{trans\PYZus{}factor}\PY{o}{.}\PY{n}{connect}\PY{p}{(}\PY{n}{trans\PYZus{}param}\PY{p}{)}
\end{Verbatim}
\end{example}

\begin{example}[H]
\caption{Nastavení pozorovaných proměnných a inference}
\label{lst:lbpep}
\begin{Verbatim}[commandchars=\\\{\}]

\PY{n}{obs\PYZus{}1}\PY{o}{.}\PY{n}{observed}\PY{p}{(}\PY{p}{\PYZob{}}
    \PY{p}{(}\PY{l+s}{\PYZsq{}}\PY{l+s}{chinese}\PY{l+s}{\PYZsq{}}\PY{p}{,}\PY{p}{)}\PY{p}{:} \PY{l+m+mf}{0.9}\PY{p}{,}
    \PY{p}{(}\PY{l+s}{\PYZsq{}}\PY{l+s}{indian}\PY{l+s}{\PYZsq{}}\PY{p}{,}\PY{p}{)}\PY{p}{:} \PY{l+m+mf}{0.1}\PY{p}{,}
\PY{p}{\PYZcb{}}\PY{p}{)}

\PY{n}{obs\PYZus{}2}\PY{o}{.}\PY{n}{observed}\PY{p}{(}\PY{p}{\PYZob{}}
    \PY{p}{(}\PY{l+s}{\PYZsq{}}\PY{l+s}{chinese}\PY{l+s}{\PYZsq{}}\PY{p}{,}\PY{p}{)}\PY{p}{:} \PY{l+m+mf}{0.8}\PY{p}{,}
    \PY{p}{(}\PY{l+s}{\PYZsq{}}\PY{l+s}{indian}\PY{l+s}{\PYZsq{}}\PY{p}{,}\PY{p}{)}\PY{p}{:} \PY{l+m+mf}{0.2}\PY{p}{,}
\PY{p}{\PYZcb{}}\PY{p}{)}

\PY{n}{lbp} \PY{o}{=} \PY{n}{LBP}\PY{p}{(}\PY{n}{strategy}\PY{o}{=}\PY{l+s}{\PYZsq{}}\PY{l+s}{tree}\PY{l+s}{\PYZsq{}}\PY{p}{)}
\PY{n}{lbp}\PY{o}{.}\PY{n}{add\PYZus{}nodes}\PY{p}{(}\PY{p}{[}\PY{n}{obs\PYZus{}1}\PY{p}{,} \PY{n}{obs\PYZus{}2}\PY{p}{,} \PY{n}{hid\PYZus{}1}\PY{p}{,} \PY{n}{hid\PYZus{}2}\PY{p}{,} \PY{n}{obs\PYZus{}factor\PYZus{}1}\PY{p}{,} \PY{n}{obs\PYZus{}factor\PYZus{}2}\PY{p}{,}
               \PY{n}{trans\PYZus{}factor}\PY{p}{,} \PY{n}{trans\PYZus{}param}\PY{p}{]}\PY{p}{)}
\PY{n}{lbp}\PY{o}{.}\PY{n}{run}\PY{p}{(}\PY{p}{)}
\end{Verbatim}
\end{example}


\begin{example}[H]
\caption{Vypočtené hodnoty skrytých proměnných}
\label{lst:epres}
\begin{Verbatim}[commandchars=\\\{\}]

\PY{g+gp}{\PYZgt{}\PYZgt{}\PYZgt{} }\PY{k}{print} \PY{n}{hid\PYZus{}1}\PY{o}{.}\PY{n}{belief}\PY{o}{.}\PY{n}{pretty\PYZus{}print}\PY{p}{(}\PY{l+m+mi}{50}\PY{p}{,} \PY{l+m+mi}{2}\PY{p}{)}
\PY{g+go}{\PYZhy{}\PYZhy{}\PYZhy{}\PYZhy{}\PYZhy{}\PYZhy{}\PYZhy{}\PYZhy{}\PYZhy{}\PYZhy{}\PYZhy{}\PYZhy{}\PYZhy{}\PYZhy{}\PYZhy{}\PYZhy{}\PYZhy{}\PYZhy{}\PYZhy{}\PYZhy{}\PYZhy{}\PYZhy{}\PYZhy{}\PYZhy{}\PYZhy{}\PYZhy{}\PYZhy{}\PYZhy{}\PYZhy{}\PYZhy{}\PYZhy{}\PYZhy{}\PYZhy{}\PYZhy{}\PYZhy{}\PYZhy{}\PYZhy{}\PYZhy{}\PYZhy{}\PYZhy{}\PYZhy{}\PYZhy{}\PYZhy{}\PYZhy{}\PYZhy{}\PYZhy{}\PYZhy{}\PYZhy{}\PYZhy{}\PYZhy{}}
\PY{g+go}{         food\PYZus{}1                    Value          }
\PY{g+go}{\PYZhy{}\PYZhy{}\PYZhy{}\PYZhy{}\PYZhy{}\PYZhy{}\PYZhy{}\PYZhy{}\PYZhy{}\PYZhy{}\PYZhy{}\PYZhy{}\PYZhy{}\PYZhy{}\PYZhy{}\PYZhy{}\PYZhy{}\PYZhy{}\PYZhy{}\PYZhy{}\PYZhy{}\PYZhy{}\PYZhy{}\PYZhy{}\PYZhy{}\PYZhy{}\PYZhy{}\PYZhy{}\PYZhy{}\PYZhy{}\PYZhy{}\PYZhy{}\PYZhy{}\PYZhy{}\PYZhy{}\PYZhy{}\PYZhy{}\PYZhy{}\PYZhy{}\PYZhy{}\PYZhy{}\PYZhy{}\PYZhy{}\PYZhy{}\PYZhy{}\PYZhy{}\PYZhy{}\PYZhy{}\PYZhy{}\PYZhy{}}
\PY{g+go}{         chinese                   0.86           }
\PY{g+go}{         indian                    0.14           }
\PY{g+go}{\PYZhy{}\PYZhy{}\PYZhy{}\PYZhy{}\PYZhy{}\PYZhy{}\PYZhy{}\PYZhy{}\PYZhy{}\PYZhy{}\PYZhy{}\PYZhy{}\PYZhy{}\PYZhy{}\PYZhy{}\PYZhy{}\PYZhy{}\PYZhy{}\PYZhy{}\PYZhy{}\PYZhy{}\PYZhy{}\PYZhy{}\PYZhy{}\PYZhy{}\PYZhy{}\PYZhy{}\PYZhy{}\PYZhy{}\PYZhy{}\PYZhy{}\PYZhy{}\PYZhy{}\PYZhy{}\PYZhy{}\PYZhy{}\PYZhy{}\PYZhy{}\PYZhy{}\PYZhy{}\PYZhy{}\PYZhy{}\PYZhy{}\PYZhy{}\PYZhy{}\PYZhy{}\PYZhy{}\PYZhy{}\PYZhy{}\PYZhy{}}

\PY{g+gp}{\PYZgt{}\PYZgt{}\PYZgt{} }\PY{k}{print} \PY{n}{hid\PYZus{}2}\PY{o}{.}\PY{n}{belief}\PY{o}{.}\PY{n}{pretty\PYZus{}print}\PY{p}{(}\PY{l+m+mi}{50}\PY{p}{,} \PY{l+m+mi}{2}\PY{p}{)}
\PY{g+go}{\PYZhy{}\PYZhy{}\PYZhy{}\PYZhy{}\PYZhy{}\PYZhy{}\PYZhy{}\PYZhy{}\PYZhy{}\PYZhy{}\PYZhy{}\PYZhy{}\PYZhy{}\PYZhy{}\PYZhy{}\PYZhy{}\PYZhy{}\PYZhy{}\PYZhy{}\PYZhy{}\PYZhy{}\PYZhy{}\PYZhy{}\PYZhy{}\PYZhy{}\PYZhy{}\PYZhy{}\PYZhy{}\PYZhy{}\PYZhy{}\PYZhy{}\PYZhy{}\PYZhy{}\PYZhy{}\PYZhy{}\PYZhy{}\PYZhy{}\PYZhy{}\PYZhy{}\PYZhy{}\PYZhy{}\PYZhy{}\PYZhy{}\PYZhy{}\PYZhy{}\PYZhy{}\PYZhy{}\PYZhy{}\PYZhy{}\PYZhy{}}
\PY{g+go}{         food\PYZus{}2                    Value          }
\PY{g+go}{\PYZhy{}\PYZhy{}\PYZhy{}\PYZhy{}\PYZhy{}\PYZhy{}\PYZhy{}\PYZhy{}\PYZhy{}\PYZhy{}\PYZhy{}\PYZhy{}\PYZhy{}\PYZhy{}\PYZhy{}\PYZhy{}\PYZhy{}\PYZhy{}\PYZhy{}\PYZhy{}\PYZhy{}\PYZhy{}\PYZhy{}\PYZhy{}\PYZhy{}\PYZhy{}\PYZhy{}\PYZhy{}\PYZhy{}\PYZhy{}\PYZhy{}\PYZhy{}\PYZhy{}\PYZhy{}\PYZhy{}\PYZhy{}\PYZhy{}\PYZhy{}\PYZhy{}\PYZhy{}\PYZhy{}\PYZhy{}\PYZhy{}\PYZhy{}\PYZhy{}\PYZhy{}\PYZhy{}\PYZhy{}\PYZhy{}\PYZhy{}}
\PY{g+go}{         chinese                   0.81           }
\PY{g+go}{         indian                    0.19           }
\PY{g+go}{\PYZhy{}\PYZhy{}\PYZhy{}\PYZhy{}\PYZhy{}\PYZhy{}\PYZhy{}\PYZhy{}\PYZhy{}\PYZhy{}\PYZhy{}\PYZhy{}\PYZhy{}\PYZhy{}\PYZhy{}\PYZhy{}\PYZhy{}\PYZhy{}\PYZhy{}\PYZhy{}\PYZhy{}\PYZhy{}\PYZhy{}\PYZhy{}\PYZhy{}\PYZhy{}\PYZhy{}\PYZhy{}\PYZhy{}\PYZhy{}\PYZhy{}\PYZhy{}\PYZhy{}\PYZhy{}\PYZhy{}\PYZhy{}\PYZhy{}\PYZhy{}\PYZhy{}\PYZhy{}\PYZhy{}\PYZhy{}\PYZhy{}\PYZhy{}\PYZhy{}\PYZhy{}\PYZhy{}\PYZhy{}\PYZhy{}\PYZhy{}}

\PY{g+gp}{\PYZgt{}\PYZgt{}\PYZgt{} }\PY{k}{print} \PY{n}{trans\PYZus{}param}\PY{o}{.}\PY{n}{alpha}\PY{o}{.}\PY{n}{pretty\PYZus{}print}\PY{p}{(}\PY{l+m+mi}{50}\PY{p}{,} \PY{l+m+mi}{2}\PY{p}{)}
\PY{g+go}{\PYZhy{}\PYZhy{}\PYZhy{}\PYZhy{}\PYZhy{}\PYZhy{}\PYZhy{}\PYZhy{}\PYZhy{}\PYZhy{}\PYZhy{}\PYZhy{}\PYZhy{}\PYZhy{}\PYZhy{}\PYZhy{}\PYZhy{}\PYZhy{}\PYZhy{}\PYZhy{}\PYZhy{}\PYZhy{}\PYZhy{}\PYZhy{}\PYZhy{}\PYZhy{}\PYZhy{}\PYZhy{}\PYZhy{}\PYZhy{}\PYZhy{}\PYZhy{}\PYZhy{}\PYZhy{}\PYZhy{}\PYZhy{}\PYZhy{}\PYZhy{}\PYZhy{}\PYZhy{}\PYZhy{}\PYZhy{}\PYZhy{}\PYZhy{}\PYZhy{}\PYZhy{}\PYZhy{}\PYZhy{}\PYZhy{}\PYZhy{}}
\PY{g+go}{     food\PYZus{}1          food\PYZus{}2          Value      }
\PY{g+go}{\PYZhy{}\PYZhy{}\PYZhy{}\PYZhy{}\PYZhy{}\PYZhy{}\PYZhy{}\PYZhy{}\PYZhy{}\PYZhy{}\PYZhy{}\PYZhy{}\PYZhy{}\PYZhy{}\PYZhy{}\PYZhy{}\PYZhy{}\PYZhy{}\PYZhy{}\PYZhy{}\PYZhy{}\PYZhy{}\PYZhy{}\PYZhy{}\PYZhy{}\PYZhy{}\PYZhy{}\PYZhy{}\PYZhy{}\PYZhy{}\PYZhy{}\PYZhy{}\PYZhy{}\PYZhy{}\PYZhy{}\PYZhy{}\PYZhy{}\PYZhy{}\PYZhy{}\PYZhy{}\PYZhy{}\PYZhy{}\PYZhy{}\PYZhy{}\PYZhy{}\PYZhy{}\PYZhy{}\PYZhy{}\PYZhy{}\PYZhy{}}
\PY{g+go}{    chinese         chinese           2.3       }
\PY{g+go}{    chinese          indian           0.96      }
\PY{g+go}{     indian         chinese           1.0       }
\PY{g+go}{     indian          indian           2.0       }
\PY{g+go}{\PYZhy{}\PYZhy{}\PYZhy{}\PYZhy{}\PYZhy{}\PYZhy{}\PYZhy{}\PYZhy{}\PYZhy{}\PYZhy{}\PYZhy{}\PYZhy{}\PYZhy{}\PYZhy{}\PYZhy{}\PYZhy{}\PYZhy{}\PYZhy{}\PYZhy{}\PYZhy{}\PYZhy{}\PYZhy{}\PYZhy{}\PYZhy{}\PYZhy{}\PYZhy{}\PYZhy{}\PYZhy{}\PYZhy{}\PYZhy{}\PYZhy{}\PYZhy{}\PYZhy{}\PYZhy{}\PYZhy{}\PYZhy{}\PYZhy{}\PYZhy{}\PYZhy{}\PYZhy{}\PYZhy{}\PYZhy{}\PYZhy{}\PYZhy{}\PYZhy{}\PYZhy{}\PYZhy{}\PYZhy{}\PYZhy{}\PYZhy{}}
\end{Verbatim}
\end{example}

\begin{example}
\caption{Využití modelu pro nová pozorování}
\label{lst:epint}
\begin{Verbatim}[commandchars=\\\{\}]
\PY{g+gp}{\PYZgt{}\PYZgt{}\PYZgt{} }\PY{n}{obs\PYZus{}1}\PY{o}{.}\PY{n}{observed}\PY{p}{(}\PY{p}{\PYZob{}}
\PY{g+gp}{... }    \PY{p}{(}\PY{l+s}{\PYZsq{}}\PY{l+s}{chinese}\PY{l+s}{\PYZsq{}}\PY{p}{,}\PY{p}{)}\PY{p}{:} \PY{l+m+mf}{0.1}\PY{p}{,}
\PY{g+gp}{... }    \PY{p}{(}\PY{l+s}{\PYZsq{}}\PY{l+s}{indian}\PY{l+s}{\PYZsq{}}\PY{p}{,}\PY{p}{)}\PY{p}{:} \PY{l+m+mf}{0.9}\PY{p}{,}
\PY{g+gp}{... }\PY{p}{\PYZcb{}}\PY{p}{)}      
\PY{g+gp}{\PYZgt{}\PYZgt{}\PYZgt{} }\PY{n}{obs\PYZus{}2}\PY{o}{.}\PY{n}{observed}\PY{p}{(}\PY{p}{\PYZob{}}
\PY{g+gp}{... }    \PY{p}{(}\PY{l+s}{\PYZsq{}}\PY{l+s}{chinese}\PY{l+s}{\PYZsq{}}\PY{p}{,}\PY{p}{)}\PY{p}{:} \PY{l+m+mf}{0.2}\PY{p}{,}
\PY{g+gp}{... }    \PY{p}{(}\PY{l+s}{\PYZsq{}}\PY{l+s}{indian}\PY{l+s}{\PYZsq{}}\PY{p}{,}\PY{p}{)}\PY{p}{:} \PY{l+m+mf}{0.8}\PY{p}{,}
\PY{g+gp}{... }\PY{p}{\PYZcb{}}\PY{p}{)}   
\PY{g+gp}{\PYZgt{}\PYZgt{}\PYZgt{} }\PY{n}{lbp}\PY{o}{.}\PY{n}{init\PYZus{}messages}\PY{p}{(}\PY{p}{)}
\PY{g+gp}{\PYZgt{}\PYZgt{}\PYZgt{} }\PY{n}{lbp}\PY{o}{.}\PY{n}{run}\PY{p}{(}\PY{p}{)}
\PY{g+gp}{\PYZgt{}\PYZgt{}\PYZgt{} }\PY{k}{print} \PY{n}{trans\PYZus{}param}\PY{o}{.}\PY{n}{alpha}\PY{o}{.}\PY{n}{pretty\PYZus{}print}\PY{p}{(}\PY{l+m+mi}{50}\PY{p}{,} \PY{l+m+mi}{2}\PY{p}{)}        
\PY{g+go}{\PYZhy{}\PYZhy{}\PYZhy{}\PYZhy{}\PYZhy{}\PYZhy{}\PYZhy{}\PYZhy{}\PYZhy{}\PYZhy{}\PYZhy{}\PYZhy{}\PYZhy{}\PYZhy{}\PYZhy{}\PYZhy{}\PYZhy{}\PYZhy{}\PYZhy{}\PYZhy{}\PYZhy{}\PYZhy{}\PYZhy{}\PYZhy{}\PYZhy{}\PYZhy{}\PYZhy{}\PYZhy{}\PYZhy{}\PYZhy{}\PYZhy{}\PYZhy{}\PYZhy{}\PYZhy{}\PYZhy{}\PYZhy{}\PYZhy{}\PYZhy{}\PYZhy{}\PYZhy{}\PYZhy{}\PYZhy{}\PYZhy{}\PYZhy{}\PYZhy{}\PYZhy{}\PYZhy{}\PYZhy{}\PYZhy{}\PYZhy{}}
\PY{g+go}{     food\PYZus{}1          food\PYZus{}2          Value      }
\PY{g+go}{\PYZhy{}\PYZhy{}\PYZhy{}\PYZhy{}\PYZhy{}\PYZhy{}\PYZhy{}\PYZhy{}\PYZhy{}\PYZhy{}\PYZhy{}\PYZhy{}\PYZhy{}\PYZhy{}\PYZhy{}\PYZhy{}\PYZhy{}\PYZhy{}\PYZhy{}\PYZhy{}\PYZhy{}\PYZhy{}\PYZhy{}\PYZhy{}\PYZhy{}\PYZhy{}\PYZhy{}\PYZhy{}\PYZhy{}\PYZhy{}\PYZhy{}\PYZhy{}\PYZhy{}\PYZhy{}\PYZhy{}\PYZhy{}\PYZhy{}\PYZhy{}\PYZhy{}\PYZhy{}\PYZhy{}\PYZhy{}\PYZhy{}\PYZhy{}\PYZhy{}\PYZhy{}\PYZhy{}\PYZhy{}\PYZhy{}\PYZhy{}}
\PY{g+go}{    chinese         chinese           2.3       }
\PY{g+go}{    chinese          indian           0.97      }
\PY{g+go}{     indian         chinese           0.97      }
\PY{g+go}{     indian          indian           2.3       }
\PY{g+go}{\PYZhy{}\PYZhy{}\PYZhy{}\PYZhy{}\PYZhy{}\PYZhy{}\PYZhy{}\PYZhy{}\PYZhy{}\PYZhy{}\PYZhy{}\PYZhy{}\PYZhy{}\PYZhy{}\PYZhy{}\PYZhy{}\PYZhy{}\PYZhy{}\PYZhy{}\PYZhy{}\PYZhy{}\PYZhy{}\PYZhy{}\PYZhy{}\PYZhy{}\PYZhy{}\PYZhy{}\PYZhy{}\PYZhy{}\PYZhy{}\PYZhy{}\PYZhy{}\PYZhy{}\PYZhy{}\PYZhy{}\PYZhy{}\PYZhy{}\PYZhy{}\PYZhy{}\PYZhy{}\PYZhy{}\PYZhy{}\PYZhy{}\PYZhy{}\PYZhy{}\PYZhy{}\PYZhy{}\PYZhy{}\PYZhy{}\PYZhy{}}
\end{Verbatim}
\end{example}


\section{Dialog State Tracking Challenge}
\label{sec:dstc}

Knihovna pro inferenci byla použita v~Dialog State Tracking Challenge (DSTC) 2013~\cite{dstc2013}.
Nejprve popíšeme problém, následně představíme použitý model a nakonec porovnáme výsledky modelu s~ostatními zúčastněnými systémy.

Cílem DSTC bylo vytvořit prostředky pro porovnání různých přístupů k~inferenci dialogového stavu a vyhodnotit jejich úspěšnost pomocí společné množiny metrik.
Organizátoři poskytli několik anotovaných množin dat, které pocházejí z~reálného použití tří různých dialogových systémů pro úlohu Let's Go!~\cite{raux2005let}.
Zároveň poskytli i baseline dialogový systém.
Cílem účastníků soutěže bylo vytvořit systém pro odhad dialového stavu, který správně predikuje dialogový stav na základě vstupu od uživatele a minulé akce systému.

\subsection{Let's Go!}

Úlohou dialogového systému Let's Go! je poskytnout telefonní službu pro zjišťování autobusového spojení v~Pittsburghu.
Systém rozpoznává devět různých slotů, některé z~nich se sestávají z~podslotů.
Tyto sloty jsou
\begin{itemize}
\item \emph{route}~-- linka,
\item \emph{time}~-- čas odjezdu nebo příjezdu,
\item \emph{date}~-- datum odjezdu,
\item \emph{from.description}~-- popis místa odjezdu,
\item \emph{to.description}~-- popis místa příjezdu,
\item \emph{from.monument}~-- významný monument v~místě odjezdu,
\item \emph{to.monument}~-- významný monument v~místě příjezdu,
\item \emph{from.neighborhood}~-- část města (sousedství) z~které chce uživatel jet,
\item \emph{to.neighborhood}~-- část města (sousedství) kam chce uživatel jet.
\end{itemize}
Navíc se sloty \emph{time} a \emph{date} skládají z~dalších pěti, respektive čtyř slotů (např. \emph{time.hour}, \emph{time.minute}, \emph{date.day}).
Systém Let's Go! celkem rozpoznával více než 5000 míst a zhruba 150 různých linek.

Jednotlivé dialogy se skládají z~obrátek, ve kterých se střídá uživatel a systém.
Příklad rozhovoru je v~tabulce~\ref{tab:exampledialog}.
Uživatel může informovat systém o~hodnotě libovolného slotu.
Systém může uživatele požádat o~hodnotu libovolného slotu, popř. o~potvrzení hodnoty slotu.
Vstup a výstup je reprezentován dialogovými akty.
Používány jsou pouze dialogové akty \emph{inform}, \emph{deny}, \emph{affirm} a \emph{negate}.
Zbytek dialogových aktů je ignorován, protože nemění cíle uživatele.

\begin{table}
\catcode`\-=12
\begin{center}
    \begin{tabular}{|c|l|p{10cm}|}
\hline
Obrátka & Aktér & Přepis \\
\hline
1   & Systém   & \textit{Welcome to the Lets Go! bus information system for Pittsburgh.
              Where would you like to leave from?}\\
              & dact  & hello()\&request(from.desc) \\
\cline{2-3}
    & Uživatel  & \textit{The South Side.} \\
    & dact  & inform(from.desc=\uv{the south side}) \\
\hline
2   & Systém   & \textit{Leaving from the south side neighborhood. Did I get that
              right?}\\
    & dact  & confirm(from.desc=\uv{the south side}) \\
\cline{2-3}
    & Uživatel  & \textit{Correct.} \\
    & dact  & affirm() \\
\hline
3   & Systém   & \textit{Where do you wanna go}? \\
    & dact  & request(to.desc) \\
\cline{2-3}
    & Uživatel  & \textit{Oakland.} \\
    & dact  & inform(to.desc=\uv{Oakland}) \\
\hline
4   & Systém   & \textit{When do you wanna travel?} \\
    & dact  & request(time) \\
\cline{2-3}
    & Uživatel  & \textit{Next}. \\
    & dact  & inform(time=next) \\
\hline
5   & Systém   & \textit{The next bus is the 54C, leaving from south eighteenth street at
              east Carson in the south side at 10:25pm. It arrives at fifth
              avenue at Bigelow in Oakland at 10:35pm. You may now say. next
              bus, previous bus, restart, or goodbye.}\\
    & dact  & ... \\
\hline
\end{tabular}
\end{center}
\caption{Příklad dialogu mezi uživatelem a dialogovým systémem Let's Go!, \textit{dact} znamená dialogový akt.}
\label{tab:exampledialog}
\end{table}

Organizátoři poskytli čtyři datové sady pocházející ze tří různých dialogových systémů.
Každá datová sada obsahovala hypotézy z~živého běhu dialogového systému (live data).
Dva dialogové systémy produkovaly seznam $n$ nejlepších hypotéz.
Jeden systém produkoval pouze nejlepší hypotézu, záznamy hovorů proto byly zpracovány zpětně pro vygenerování seznamů $n$ nejlepších hypotéz (batch data).
Vygenerované hypotézy ovšem obsahovaly pouze skóre, které muselo být přepočítáno pro získání pravděpodobností.
Organizátoři nevydali žádné informace o~povaze skóre.

Data obsahovala zjevně chybné hypotézy, např. v~jedné hypotéze se objevily dvě různé hodnoty pro jeden slot.
Tyto hypotézy byly z~dat odstraněny ještě před samotnou inferencí.

\subsection{Popis systému}

Dialogový stav byl v~použitém systému modelován jednoduchým generativním modelem, kde pro každý slot v~jedné obrátce existují dva vrcholy, skutečná hodnota slotu $s_t$ a pozorovaná hodnota $o_t$.
Stav důvěry skutečné hodnoty slotu $b(s_t)$ závisí pouze na hodnotě v~minulé obrátce $s_{t-1}$ a na poslední systémové akci $a_{t-1}$.
Pozorovaná hodnota $o_t$ závisí pouze na skutečné hodnotě $s_t$.

Stav důvěry můžeme vyjádřit ze sdružené pravděpodobnosti generativního modelu:
\begin{equation}
b(s_t) = \sum_{s_{t-1}, o_t} p(s_t \mid a_{t-1}, s_{t-1}) \; p(o_t \mid s_t) \; b(s_{t-1}).
\end{equation}

Pro zrychlení výpočtu bylo využito stažených parametrů s~manuálně nastavenými hodnotami~\cite{thomson2010bayesian}:
\begin{equation}
p(s_t \mid a_{t-1}, s_{t-1}) = \begin{cases}
\theta_t & \text{pokud } s_t = s_{t-1} \\
\frac{1 - \theta_t}{|hodnoty| - 1} & \text{jinak}
\end{cases},
\end{equation}
kde $\theta_t$ je pravděpodobnost, že hodnota slotu se nemění a $|hodnoty|$ je počet hodnot pro slot.

Pro model pozorování byly také použity stažené pravděpodobnosti:
\begin{equation}
p(o_t \mid s_t) = \begin{cases}
\theta_o & \text{pokud } o_t = s_t \\
\frac{1 - \theta_o}{|hodnoty| - 1} & \text{jinak}
\end{cases},
\end{equation}
kde $\theta_o$ je pravděpodobnost, že pozorování bude odpovídat skutečné hodnotě slotu.

Parametr $\theta_t$ určuje, jak moc bude systém zapomínat, v~případě hodnoty blízké~1 systém téměř nezapomíná a jakmile nějakou hodnotu pozoruje, bude jej těžké přesvědčit o~čemkoliv jiném.
V~opačném případě systém zapomíná a~i~z~téměř jisté důvěry v~přiřazení stavu může během pár obrátek dojít do stavu, kdy mají všechny možné přiřazení stejnou pravděpodobnost.
Parametr $\theta_o$ vyjadřuje důvěru systému v~SLU.
Pokud je jeho hodnota vysoká, předpokládáme, že SLU téměř nedělá chyby.
V~opačném případě je systém tolerantní k~chybám SLU.
Na základě výsledků dialogového systému na trénovacích datech byly manuálně zvoleny vhodné hodnoty parametrů: $\theta_t = 0.8$ a $\theta_o = 0.8$.

Představený model pro pozorování platí pro \emph{inform} dialogové akty, dialogové sloty \emph{affirm} a \emph{negate} byly převedeny na \emph{inform}.

Pro inferenci bylo použito LBP, a protože možných přiřazení slotů bylo v~řádu stovek, všechny nepozorované hodnoty byly staženy do jedné speciální hodnoty a~jejich pravděpodobnost byla počítána pouze dohromady.

Popis použitého systému a jeho srovnání s diskriminativním systémem bylo publikováno v samostatném článku~\cite{zilka2013}.

\subsection{Evaluace}

Použitý model byl testován na čtyřech testovacích datových sadách s~live daty a na dvou testovacích datových sadách s~batch daty.
Výsledky byly porovnány s~baseline dialogovým systémem vytvořeným organizátory.
Použité metriky byly přesnost (accuracy) dialogového systému a Brier score.
Přesnost je podíl obrátek, ve kterých byla nejlepší hypotéza systému pro odhad stavu správná oproti celkovému počtu obrátek.
Brier score měří přesnost prediktivní pravděpodobnostní distribuce systému pro odhad stavu~\cite{brier1950verification} (čím menší, tím lepší).

Výsledky ukazují (viz tabulka~\ref{t:all:datasets}), že generativní model překonává baseline na všech datových sadách, kromě batch dat.
Manuální inspekcí bylo zjištěno, že generativní model je velmi citlivý na vstupní pravděpodobnosti.
Pravděpodobnosti hypotéz z~batch dat, zřejmě kvůli způsobu přepočítání skóre na pravděpodobnosti, jsou velmi blízko sebe, a tedy nedávají dostatek informací pro pravděpodobnostní model.

\begin{table}[htb]
\begin{center}
\begin{tabular}{|l|l|l|l|l|l|}
  \hline
live data & metric      & BT & GT \\
\hline
test1 & přesnost    & 0.77 & 0.88 \\
        & Brier score & 0.29 & 0.21 \\
\hline
test2 & přesnost    & 0.79 & 0.85 \\
        & Brier score & 0.27 & 0.23 \\
\hline
test3  & přesnost    & 0.92 & 0.93 \\
        & Brier score & 0.14 & 0.16 \\
\hline
test4  & přesnost    & 0.82 & 0.87 \\
        & Brier score & 0.24 & 0.20 \\
\hline
\textsc{All}    & přesnost    & 0.83 & 0.88 \\
        & Brier score & 0.24 & 0.20 \\
\hline
\hline
batch data & metric      & BT & GT \\
\hline
test1 & přesnost    & 0.75 & 0.74 \\
        & Brier score & 0.35 & 0.39 \\
\hline
test2 & přesnost    & 0.79 & 0.77 \\
        & Brier score & 0.30 & 0.33 \\
\hline
\textsc{All}    & přesnost    & 0.77 & 0.76 \\
        & Brier score & 0.32 & 0.36 \\
\hline
\end{tabular}
\end{center}
\caption{Přesnost systému pro odhad stavu na live a batch testovacích datech, kde BT je baseline systém a GT je generativní systém. \textsc{All} značí průměrné skóre přes všechny testovací sady.}
\label{t:all:datasets}
\end{table}

Generativní systém byl také porovnán s~ostatními dialogovými systémy přihlášenými do DSTC.
Zde byl rovnocenným soupeřem většině ostatních systémů (viz tabulka~\ref{t:DSTC:ranking}).
V oficiálních výsledcích bude generativní systém veden jako druhý systém týmu číslo 3.
Je třeba ovšem zdůraznit, že v systému byla chyba, která vedla k~velkým výkyvům v kvalitě výstupu systému napříč datovými sadami.
Tato chyba byla odhalena až po odeslání řešení.
Výsledky v této práci pocházejí z~opraveného systému.

\begin{table}[H]
\begin{center}
\begin{tabular}{|l|l|l|l|l|l|}
  \hline
tým / systém & přesnost & Brier score \\
\hline
BT - C         & 0.81 & 0.27\\
BT          & 0.83 & 0.24\\
GT          & 0.88 & 0.20\\
team1       & 0.88 & 0.23\\
team2       & 0.88 & 0.21\\
team4         & 0.81 & 0.28\\
team5         & 0.88 & 0.21\\
team6         & \textbf{0.91} & \textbf{0.18}\\
team7         & 0.85 & 0.23\\
team8         & 0.83 & 0.24\\
team9         & 0.89 & 0.20\\
\hline
\end{tabular}
\end{center}
\caption{Přesnost systémů přihlášených do DSTC. BT - C značí baseline systém bez mazání chybných hypotéz, BT je baseline systém, GT je generativní systém a dále následují systémy ostatních týmů. Skóre jsou zprůměrována přes všechny 4 testovací sady dat.}
\label{t:DSTC:ranking}
\end{table}

Tabulky výsledků a příklad dialogu byly převzaty z~článku popisujícího přihlášené systémy do DSTC a porovnávajícího generativní model s~jednoduchým diskriminativním modelem~\cite{zilka2013}.

Popis jednotlivých součástí systému a jeho používání je v~příloze~\ref{ap:dstc}.



% Ukázka použití některých konstrukcí LateXu (odkomentujte, chcete-li)

% !TEX root = prace.tex
\chapter*{Závěr}
\addcontentsline{toc}{chapter}{Závěr}

Práce splnila všechny vytyčené cíle.
Výsledkem je implementace metod pro odhad stavu a parametrů v~dialogových systémech.
Byly představeny metody pro inferenci v~grafických modelech založené na posílání zpráv, navíc byl odvozen algoritmus pro inferenci parametrů dirichletovského rozdělení srovnáním momentů.
Práce obsahuje příklady použití vytvořené knihovny a otestování na reálných datech v~rámci soutěže Dialog State Tracking Challenge, kde generativní systém pro odhad stavu používající implementovanou knihovnu dosáhl výsledků srovnatelných s~ostatními přihlášenými týmy, které se dají považovat za state of the art.
Bližší detaily podle cílů vytyčených v~úvodu práce:
\begin{enumerate}

\item V~kapitole \ref{ch:kap1} byly představeny dialogové systémy a jejich jednotlivé součásti.
V~sekci \ref{sec:bn} pak byly představeny bayesovské sítě, které se nabízí jako ideální model pro reprezentaci dialogového stavu.

\item V~kapitole \ref{ch:kap2} jsou prezentovány možné inferenční algoritmy pro odhad stavu v~bayesovských sítích.
V~sekci \ref{sec:lbp} je představen algoritmus Loopy Belief Propagation.
Jeho použití je ukázáno v~sekci \ref{sec:usage}.

\item Implementace LBP byla použita pro DSTC a popis celé soutěže i s~výsledky je v~sekci \ref{sec:dstc}.
Implementovaný systém pro odhad stavu byl lepší než baseline systém poskytnutý organizátory (tabulka \ref{t:all:datasets}) a systém se navíc umístil příznivě i vzhledem k~ostatním přihlášeným týmům (tabulka \ref{t:DSTC:ranking}).

\item Implementované strategie pro výběr vrcholu v~LBP algoritmu jsou popsány v~sekci \ref{sec:noch}.

\item Algoritmus Expectation Propagation pro inferenci v~obecných grafických modelech je představen v~sekci \ref{sec:ep}.

\item Knihovna a jednotlivé její součásti jsou popsány v~kapitole \ref{ch:kap4}.
Příklad vrcholů vytvořených speciálně pro použití v~EP je v~sekci \ref{sec:vrdir}.
Učení parametrů pro odhad stavu v~dialogových systémech je odvozeno v~kapitole \ref{ch:ep}.
\end{enumerate}


\appendix
% !TEX root = prace.tex
\chapter{Zpracování dialogů z DSTC}
\label{ap:dstc}

\section{Soubory}

Předvedeme použití vytvořené knihovny na data z Dialog State Tracking Challenge.
Všechny potřebné data a zdrojové kódy jsou na přiloženém cd.
V adresáři \texttt{src/dstc} je velké množství modulů a dalších adresářů:
\begin{itemize}
\item \texttt{common.py}, \texttt{dstcutils.py}, \texttt{jsonpath.py}, \texttt{ontology.py}, \texttt{pddiscrete.py}, \texttt{preprocessing.py} -- pomocné moduly pro práci s daty, vytvořené členy výzkumné skupiny pracující na projektu Vystadial,

\item \texttt{loader.py} -- modul pro načítání dat z formátu JSON~\cite{crockford2006application} do interní reprezentace, 

\item \texttt{system\_dm.py} -- skript pro zpracování dat používající vytvořenou knihovnu pro odhad stavu v jednotlivých dialozích.
Použitá datová sada, vstupní a výstupní soubory, parametry grafického modelu a další parametry nutné pro výpočet je možné skriptu zadat buď pomocí přepínačů na příkazové řádce, anebo v konfiguračním souboru.

\item \texttt{test\_run.sh} -- bashový skript pro spuštění celé pipeline od zpracování dat až po zobrazení výsledků,

\item \texttt{test1.conf}, \texttt{test2.conf}, \texttt{test3.conf}, \texttt{test4.conf} -- konfigurační soubory pro spuštění skriptu \texttt{system\_dm.py} na testovací datové sady ze všech čtyř dialogových systémů.

\item V adresáři \texttt{src/dstc/bin/} je baseline dialogový systém a skripty pro zpracování výsledků od organizátorů soutěže:
	\begin{itemize}
	\item \texttt{baseline} -- baseline dialogový systém,
	\item \texttt{score} -- skript pro vyhodnocení výstupu soutěžního systému,
	\item \texttt{report} -- skript pro zobrazení výsledků soutěžního systému v čitelné podobě,
	\item \texttt{validate} -- skript pro validaci korektního výstupu soutěžního systému,
	\item Pomocné moduly pro tyto skripty jsou v adresáři \texttt{src/dstc/lib}.
	\end{itemize}

\item V adresáři \texttt{src/dstc/data/} jsou všechny 4 testovací datové sady, každá ve vlastním adresáři (\texttt{test1/}, \texttt{test2/}, \texttt{test3/}, \texttt{test4/}).
Každá datová sada se pak skládá z podadresářů pro jednotlivé hovory.
Každý hovor se skládá ze dvou souborů:
	\begin{itemize}
	\item \texttt{dstc.log.json} -- záznam z dialogového systému obsahující všechny informace o dialogu ve formátu JSON,
	\item \texttt{dstc.labels.json} -- soubor obsahující správné hypotézy, tedy řešení úlohy, tento soubor samozřejmě před skončením soutěže v datech  chyběl.
	\end{itemize}

\item V adresáři \texttt{src/dstc/config/} jsou soubory obsahující cesty k jednotlivým dialogům pro každou testovací sadu.
\end{itemize}

\section{Konfigurace}

Skript \texttt{system\_dm.py} očekává řadu parametrů pro správné fungování.
Parametry je možné zadat na příkazové řádce, například:
\begin{verbatim}
$ python system_dm.py \
    --dataroot=data/ \
    --dataset=config/test1.sessions \
    --scorefile=track_test1.json
\end{verbatim}

Popřípadě je možné vytvořit konfigurační soubor obsahující pouze parametry:
\begin{verbatim}
--dataroot=data/
--dataset=config/test1.sessions
--scorefile=track_test1.json
\end{verbatim}
a tento soubor pak předat jako jediný parametr:
\begin{verbatim}
$ python system_dm.py --flagfile test1.conf
\end{verbatim}

Většina vlastností systému pro odhad stavu je konfigurovatelná, ale jen malý počet parametrů je povinný.
Představíme si ty nejdůležitější parametry:
\begin{itemize}
\item \emph{dataroot} -- povinný, adresář obsahující datové sady,
\item \emph{dataset} -- povinný, konfigurační soubor obsahující seznam cest k jednotlivým dialogům,
\item \emph{scorefile} -- povinný, název souboru, do kterého má být zapsán výsledek zpracování dialogů,
\item \emph{slu} (live / batch) -- přepínač, zda-li pracovat s hypotézami z live nebo batch dat,
\item \emph{remove\_inconsistent} -- mazat nekonzistentní hypotézy (např. v jedné hypotéze více hodnot pro jeden slot),
\item \emph{processes} -- systém podporuje zpracování na více jádrech, pomocí tohoto přepínače je možné nastavit počet jader, pro jedno jádro systém přepne do ladícího módu, kdy si udržuje v paměti všechny dialogy; nehodí se pro velké datové sady,
\item \emph{prob\_same} -- hodnota parametru $\theta_o$, tedy pravděpodobnost, že pozorování odpovídá skryté proměnné,
\item \emph{prob\_trans} -- hodnota parametru $\theta_t$, tedy pravděpodobnost, že se hodnota skryté proměnné nezmění v další obrátce.
\end{itemize}

\section{Výpočet}

V předchozí sekci bylo ukázáno jakým způsobem lze nastavit a spustit skript \texttt{system\_dm.py}, který provede odhad dialogového stavu.
Výsledkem je soubor ve formátu JSON, kde pro každý dialog a každou jeho obrátku jsou vyznačeny hypotézy pro každý slot.
V příkladu~\ref{ex:res} je ukázáno jak vypadá výstup pro jeden slot, v tomto případě se jedná o počáteční zastávku, kde hodnota \uv{shadyside} má pravděpodobnost přibližně $0.56$.

\begin{example}
\begin{Verbatim}[commandchars=\\\{\}]

\PY{n+nt}{\PYZdq{}from.neighborhood\PYZdq{}}\PY{p}{:} \PY{p}{\PYZob{}}
    \PY{n+nt}{\PYZdq{}hyps\PYZdq{}}\PY{p}{:} \PY{p}{[}
        \PY{p}{\PYZob{}}
	    \PY{n+nt}{\PYZdq{}score\PYZdq{}}\PY{p}{:} \PY{l+m+mf}{0.5637651968841264}\PY{p}{,}
            \PY{n+nt}{\PYZdq{}slots\PYZdq{}}\PY{p}{:} \PY{p}{\PYZob{}}
	        \PY{n+nt}{\PYZdq{}from.neighborhood\PYZdq{}}\PY{p}{:} \PY{l+s+s2}{\PYZdq{}shadyside\PYZdq{}}
            \PY{p}{\PYZcb{}}
        \PY{p}{\PYZcb{}}
    \PY{p}{]}
\PY{p}{\PYZcb{}}\PY{p}{,}
\end{Verbatim}
\caption{Část výsledku systému pro odhad stavu.}
\label{ex:res}
\end{example}

Nyní by bylo třeba výsledky zpracovat a následně zobrazit.
Všechny tři kroky je možné provést naráz použitím skriptu \texttt{test\_run.sh}:

\begin{verbatim}
$ ./test_run.sh test1.conf
\end{verbatim}

Ukázka výsledků je v příklad~\ref{ex:res1}~a~\ref{ex:res2}. 
Ve výsledku jsou zobrazeny tři tabulky s výsledky. 
V první tabulce (\uv{schedule1}) jsou statistiky počítány po každé obrátce.
V druhé tabulce (\uv{schedule2}) jsou statistiky počítány pro slot pouze pokud se o něm v dané obrátce mluvilo.
Nakonec ve třetí tabulce (\uv{schedule3}) jsou statistiky počítány pouze na konci dialogu.

Jednotlivé sloupce reprezentují různé sloty, navíc je zde sloupec \emph{joint}, který hodnotí sdruženou informaci o všech slotech v dané obrátce, a sloupec \emph{all}, který obsahuje průměrné skóre.

Jednotlivé řádky obsahují statistiky o slotech, zvlášť zajímavé jsou:
\begin{itemize}
\item accuracy -- přesnost,
\item avgp -- průměrné skóre přiřazené správné hypotéze,
\item l2 -- Brier skóre.
\end{itemize}

Nakonec jsou zobrazeny i základní informace o celém výpočtu, například jaká datová sada byla zpracována, počet hovorů, počet obrátek, celkový čas a průměrný čas na jednu obrátku.

{
\footnotesize
\begin{example}
\begin{lstlisting}[language=c,frame=none]
----------------------------------------------------------------------------
                                    schedule1
----------------------------------------------------------------------------
               route from.d from.m from.n to.des to.mon to.nei  joint    all
            N  10085  10085  10085  10085  10085  10085  10085  10085  90765
     accuracy 0.8817 0.6969 0.9841 0.7771 0.8256 0.9845 0.7327 0.2642 0.8534
         avgp 0.8794 0.6779 0.9811 0.6765 0.8235 0.9828 0.7185 0.2465 0.8370
     hypcount 0.0773 0.2038 0.0103 0.2442 0.0435 0.0045 0.0353 0.4943 0.0729
           l2 0.1703 0.4542 0.0268 0.4575 0.2495 0.0243 0.3981 1.0655 0.2304
          mrr 0.8977 0.7239 0.9850 0.7838 0.8329 0.9846 0.7333 0.2888 0.8606
     nonempty 0.0682 0.1660 0.0103 0.2442 0.0367 0.0045 0.0348 0.4943 0.0669
  roc.v1_ca05 0.0000 0.0000 0.9841 0.0000 0.0000 0.9845 0.0000 0.0000 0.0000
  roc.v1_ca10 0.8617 0.0000 0.9841 0.0000 0.0000 0.9845 0.0000 0.0000 0.0000
  roc.v1_ca20 0.8817 0.0000 0.9841 0.0000 0.8256 0.9845 0.0000 0.0000 0.8534
   roc.v1_eer 0.8817 0.6969 0.9841 0.4288 0.8256 0.9845 0.7327 0.2642 0.8534
  roc.v2_ca05 0.0000 0.0000 0.0000 0.0000 0.0000 0.0000 0.0000 0.0000 0.0000
  roc.v2_ca10 0.0000 0.0000 0.0000 0.0000 0.0000 0.0000 0.0000 0.0000 0.0000
  roc.v2_ca20 0.0000 0.0000 0.0000 0.0000 0.0000 0.0000 0.0000 0.0000 0.0000


----------------------------------------------------------------------------
                                    schedule2
----------------------------------------------------------------------------
               route from.d from.m from.n to.des to.mon to.nei  joint    all
            N   6100   4837    926   2706   4070    904   2283   9974  27536
     accuracy 0.8590 0.6729 0.9471 0.6936 0.6784 0.9436 0.4717 0.2669 0.7309
         avgp 0.8589 0.6685 0.9449 0.6414 0.6773 0.9351 0.4496 0.2490 0.7212
     hypcount 0.0807 0.1784 0.0248 0.1778 0.0744 0.0243 0.0604 0.4927 0.0934
           l2 0.1994 0.4670 0.0780 0.5071 0.4562 0.0918 0.7784 1.0621 0.3939
          mrr 0.8769 0.6997 0.9514 0.7051 0.6907 0.9441 0.4733 0.2916 0.7443
     nonempty 0.0723 0.1333 0.0248 0.1778 0.0585 0.0243 0.0600 0.4927 0.0812
  roc.v1_ca05 0.0000 0.0000 0.9471 0.0000 0.0000 0.0000 0.0000 0.0000 0.0000
  roc.v1_ca10 0.0000 0.0000 0.9471 0.0000 0.0000 0.9436 0.0000 0.0000 0.0000
  roc.v1_ca20 0.8590 0.0000 0.9471 0.0000 0.0000 0.9436 0.0000 0.0000 0.0000
   roc.v1_eer 0.8590 0.6729 0.9471 0.6936 0.6784 0.9436 0.4717 0.2669 0.7309
  roc.v2_ca05 0.0000 0.0000 0.0000 0.0000 0.0000 0.0000 0.0000 0.0000 0.0000
  roc.v2_ca10 0.0000 0.0000 0.0000 0.0000 0.0000 0.0000 0.0000 0.0000 0.0000
  roc.v2_ca20 0.0000 0.0000 0.0000 0.0000 0.0000 0.0000 0.0000 0.0000 0.0000
\end{lstlisting}
\caption{Výsledek systému pro odhad stavu podle prvních dvou typů hodnocení.}
\label{ex:res1}
\end{example}

\begin{example}
\begin{lstlisting}[language=c,frame=none]
----------------------------------------------------------------------------
                                    schedule3
----------------------------------------------------------------------------
               route from.d from.m from.n to.des to.mon to.nei  joint    all
            N    897    874    422    767    767    354    620    901   5784
     accuracy 0.8863 0.6739 0.9645 0.7471 0.7249 0.9520 0.4048 0.1576 0.7076
         avgp 0.8827 0.6368 0.9496 0.5949 0.7159 0.9426 0.3581 0.1526 0.6692
     hypcount 0.0836 0.2574 0.0355 0.3494 0.0704 0.0198 0.1145 0.6071 0.1376
           l2 0.1656 0.5109 0.0713 0.5730 0.4012 0.0811 0.9077 1.1984 0.4673
          mrr 0.9036 0.7008 0.9656 0.7562 0.7347 0.9520 0.4073 0.1770 0.7184
     nonempty 0.0758 0.2082 0.0355 0.3494 0.0600 0.0198 0.1129 0.6071 0.1272
  roc.v1_ca05 0.0000 0.0000 0.9645 0.0000 0.0000 0.9520 0.0000 0.0000 0.0000
  roc.v1_ca10 0.8740 0.0000 0.9645 0.0000 0.0000 0.9520 0.0000 0.0000 0.0000
  roc.v1_ca20 0.8863 0.0000 0.9645 0.0000 0.0000 0.9520 0.0000 0.0000 0.0000
   roc.v1_eer 0.8863 0.6739 0.0664 0.5176 0.7249 0.9520 0.4048 0.1576 0.7076
  roc.v2_ca05 0.0000 0.0000 0.0000 0.0000 0.0000 0.0000 0.0000 0.0000 0.0000
  roc.v2_ca10 0.0000 0.0000 0.0000 0.0000 0.0000 0.0000 0.0000 0.0000 0.0000
  roc.v2_ca20 0.0000 0.0000 0.0000 0.0000 0.0000 0.0000 0.0000 0.0000 0.0000


----------------------------------------------------------------------------
                                    basic stats
----------------------------------------------------------------------------
             dataset : config/test1.sessions
      scorer_version : 1.7
            sessions : 715
     total_wall_time : 114.580698967
               turns : 10085
  wall_time_per_turn : 0.0113614971707
\end{lstlisting}
\caption{Výsledek systému pro odhad stavu podle třetího typu hodnocení a základní statistiky.}
\label{ex:res2}
\end{example}
}
% !TEX root = prace.tex

\chapter{Dokumentace tříd}
\label{ap:doc}

\section{Balíček \texttt{alex.infer}}
\label{alex.infer::doc}\label{alex.infer:balicek-infer}

\subsection{Modul \texttt{factor}}
\label{alex.infer:modul-factor}\label{alex.infer:module-alex.infer.factor}\index{alex.infer.factor (module)}
Modul {\hyperref[alex.infer:module-alex.infer.factor]{\code{factor}}} obsahuje implementaci faktorů pro počítání s
diskrétními pravděpodobnostními rozděleními.
\index{FactorError}

\begin{fulllineitems}
\phantomsection\label{alex.infer:alex.infer.factor.FactorError}\pysigline{\strong{exception }\code{alex.infer.factor.}\bfcode{FactorError}}
Výjimka značící chybu při manipulaci s faktorem.

\end{fulllineitems}

\index{Factor() (třída v alex.infer.factor)}

\begin{fulllineitems}
\phantomsection\label{alex.infer:alex.infer.factor.Factor}\pysigline{\strong{class }\code{alex.infer.factor.}\bfcode{Factor}}
Třída reprezentující diskrétní faktor s definovanými operacemi.

Je možné použít všechny standardní matematické operátory (+, -, *, /,
**). Operátor je možné aplikovat na dva faktory, anebo na faktor a
skalární hodnotu.
\index{\_\_init\_\_() (metoda alex.infer.factor.Factor)}

\begin{fulllineitems}
\phantomsection\label{alex.infer:alex.infer.factor.Factor.__init__}\pysiglinewithargsret{\bfcode{\_\_init\_\_}}{\emph{variables}, \emph{variable\_values}, \emph{prob\_table}\optional{, \emph{logarithmetic=True}}}{}
Vytvoření nového faktoru.

\begin{Verbatim}[commandchars=\\\{\}]
\PY{n}{f} \PY{o}{=} \PY{n}{Factor}\PY{p}{(}\PY{p}{[}\PY{l+s}{\PYZsq{}}\PY{l+s}{A}\PY{l+s}{\PYZsq{}}\PY{p}{,} \PY{l+s}{\PYZsq{}}\PY{l+s}{B}\PY{l+s}{\PYZsq{}}\PY{p}{]}\PY{p}{,}
    \PY{p}{\PYZob{}}\PY{l+s}{\PYZsq{}}\PY{l+s}{A}\PY{l+s}{\PYZsq{}}\PY{p}{:} \PY{p}{[}\PY{l+s}{\PYZsq{}}\PY{l+s}{a1}\PY{l+s}{\PYZsq{}}\PY{p}{,} \PY{l+s}{\PYZsq{}}\PY{l+s}{a2}\PY{l+s}{\PYZsq{}}\PY{p}{]}\PY{p}{,} \PY{l+s}{\PYZsq{}}\PY{l+s}{B}\PY{l+s}{\PYZsq{}}\PY{p}{:} \PY{p}{[}\PY{l+s}{\PYZsq{}}\PY{l+s}{b1}\PY{l+s}{\PYZsq{}}\PY{p}{,} \PY{l+s}{\PYZsq{}}\PY{l+s}{b2}\PY{l+s}{\PYZsq{}}\PY{p}{]}\PY{p}{\PYZcb{}}\PY{p}{,}
    \PY{p}{\PYZob{}}
        \PY{p}{(}\PY{l+s}{\PYZsq{}}\PY{l+s}{a1}\PY{l+s}{\PYZsq{}}\PY{p}{,} \PY{l+s}{\PYZsq{}}\PY{l+s}{b1}\PY{l+s}{\PYZsq{}}\PY{p}{)}\PY{p}{:} \PY{l+m+mf}{0.8}\PY{p}{,}
        \PY{p}{(}\PY{l+s}{\PYZsq{}}\PY{l+s}{a1}\PY{l+s}{\PYZsq{}}\PY{p}{,} \PY{l+s}{\PYZsq{}}\PY{l+s}{b2}\PY{l+s}{\PYZsq{}}\PY{p}{)}\PY{p}{:} \PY{l+m+mf}{0.2}\PY{p}{,}
        \PY{p}{(}\PY{l+s}{\PYZsq{}}\PY{l+s}{a2}\PY{l+s}{\PYZsq{}}\PY{p}{,} \PY{l+s}{\PYZsq{}}\PY{l+s}{b1}\PY{l+s}{\PYZsq{}}\PY{p}{)}\PY{p}{:} \PY{l+m+mf}{0.3}\PY{p}{,}
        \PY{p}{(}\PY{l+s}{\PYZsq{}}\PY{l+s}{a2}\PY{l+s}{\PYZsq{}}\PY{p}{,} \PY{l+s}{\PYZsq{}}\PY{l+s}{b2}\PY{l+s}{\PYZsq{}}\PY{p}{)}\PY{p}{:} \PY{l+m+mf}{0.7}
    \PY{p}{\PYZcb{}}\PY{p}{)}
\end{Verbatim}
\begin{quote}\begin{description}
\item[{Parametry}] \leavevmode\begin{itemize}
\item {} 
\textbf{variables} -- seznam proměnných ve faktoru.

\item {} 
\textbf{variable\_values} -- slovník obsahující pro každou proměnnou seznam možných hodnot.

\item {} 
\textbf{prob\_table} -- slovník možných přiřazení proměnných a jejich ohodnocení.

\item {} 
\textbf{logarithmetic} -- přepínač určující, zda-li použít aritmetiku nad zlogaritmovanými hodnotami.

\end{itemize}

\end{description}\end{quote}

\end{fulllineitems}

\index{\_\_getitem\_\_() (metoda alex.infer.factor.Factor)}

\begin{fulllineitems}
\phantomsection\label{alex.infer:alex.infer.factor.Factor.__getitem__}\pysiglinewithargsret{\bfcode{\_\_getitem\_\_}}{\emph{assignment}}{}
Získání ohodnocení pro dané přiřazení.

\begin{Verbatim}[commandchars=\\\{\}]
\PY{g+gp}{\PYZgt{}\PYZgt{}\PYZgt{} }\PY{n}{f} \PY{o}{=} \PY{n}{Factor}\PY{p}{(}\PY{p}{[}\PY{l+s}{\PYZsq{}}\PY{l+s}{A}\PY{l+s}{\PYZsq{}}\PY{p}{,} \PY{l+s}{\PYZsq{}}\PY{l+s}{B}\PY{l+s}{\PYZsq{}}\PY{p}{]}\PY{p}{,}
\PY{g+gp}{... }    \PY{p}{\PYZob{}}\PY{l+s}{\PYZsq{}}\PY{l+s}{A}\PY{l+s}{\PYZsq{}}\PY{p}{:} \PY{p}{[}\PY{l+s}{\PYZsq{}}\PY{l+s}{a1}\PY{l+s}{\PYZsq{}}\PY{p}{,} \PY{l+s}{\PYZsq{}}\PY{l+s}{a2}\PY{l+s}{\PYZsq{}}\PY{p}{]}\PY{p}{,} \PY{l+s}{\PYZsq{}}\PY{l+s}{B}\PY{l+s}{\PYZsq{}}\PY{p}{:} \PY{p}{[}\PY{l+s}{\PYZsq{}}\PY{l+s}{b1}\PY{l+s}{\PYZsq{}}\PY{p}{,} \PY{l+s}{\PYZsq{}}\PY{l+s}{b2}\PY{l+s}{\PYZsq{}}\PY{p}{]}\PY{p}{\PYZcb{}}\PY{p}{,}
\PY{g+gp}{... }    \PY{p}{\PYZob{}}
\PY{g+gp}{... }        \PY{p}{(}\PY{l+s}{\PYZsq{}}\PY{l+s}{a1}\PY{l+s}{\PYZsq{}}\PY{p}{,} \PY{l+s}{\PYZsq{}}\PY{l+s}{b1}\PY{l+s}{\PYZsq{}}\PY{p}{)}\PY{p}{:} \PY{l+m+mf}{0.8}\PY{p}{,}
\PY{g+gp}{... }        \PY{p}{(}\PY{l+s}{\PYZsq{}}\PY{l+s}{a1}\PY{l+s}{\PYZsq{}}\PY{p}{,} \PY{l+s}{\PYZsq{}}\PY{l+s}{b2}\PY{l+s}{\PYZsq{}}\PY{p}{)}\PY{p}{:} \PY{l+m+mf}{0.2}\PY{p}{,}
\PY{g+gp}{... }        \PY{p}{(}\PY{l+s}{\PYZsq{}}\PY{l+s}{a2}\PY{l+s}{\PYZsq{}}\PY{p}{,} \PY{l+s}{\PYZsq{}}\PY{l+s}{b1}\PY{l+s}{\PYZsq{}}\PY{p}{)}\PY{p}{:} \PY{l+m+mf}{0.3}\PY{p}{,}
\PY{g+gp}{... }        \PY{p}{(}\PY{l+s}{\PYZsq{}}\PY{l+s}{a2}\PY{l+s}{\PYZsq{}}\PY{p}{,} \PY{l+s}{\PYZsq{}}\PY{l+s}{b2}\PY{l+s}{\PYZsq{}}\PY{p}{)}\PY{p}{:} \PY{l+m+mf}{0.7}
\PY{g+gp}{... }    \PY{p}{\PYZcb{}}\PY{p}{)}
\PY{g+gp}{\PYZgt{}\PYZgt{}\PYZgt{} }\PY{n}{f}\PY{p}{[}\PY{p}{(}\PY{l+s}{\PYZsq{}}\PY{l+s}{a2}\PY{l+s}{\PYZsq{}}\PY{p}{,} \PY{l+s}{\PYZsq{}}\PY{l+s}{b1}\PY{l+s}{\PYZsq{}}\PY{p}{)}\PY{p}{]}
\PY{g+go}{0.3}
\end{Verbatim}
\begin{quote}\begin{description}
\item[{Parametr}] \leavevmode
\textbf{assignment} -- přiřazení proměnných.

\item[{Vrací}] \leavevmode
hodnotu dotazovaného přiřazení.

\end{description}\end{quote}

\end{fulllineitems}

\index{\_\_setitem\_\_() (metoda alex.infer.factor.Factor)}

\begin{fulllineitems}
\phantomsection\label{alex.infer:alex.infer.factor.Factor.__setitem__}\pysiglinewithargsret{\bfcode{\_\_setitem\_\_}}{\emph{assignment}, \emph{value}}{}
Nastavení hodnoty pro přiřazení proměnných.

\begin{Verbatim}[commandchars=\\\{\}]
\PY{g+gp}{\PYZgt{}\PYZgt{}\PYZgt{} }\PY{n}{f} \PY{o}{=} \PY{n}{Factor}\PY{p}{(}\PY{p}{[}\PY{l+s}{\PYZsq{}}\PY{l+s}{A}\PY{l+s}{\PYZsq{}}\PY{p}{,} \PY{l+s}{\PYZsq{}}\PY{l+s}{B}\PY{l+s}{\PYZsq{}}\PY{p}{]}\PY{p}{,}
\PY{g+gp}{... }    \PY{p}{\PYZob{}}\PY{l+s}{\PYZsq{}}\PY{l+s}{A}\PY{l+s}{\PYZsq{}}\PY{p}{:} \PY{p}{[}\PY{l+s}{\PYZsq{}}\PY{l+s}{a1}\PY{l+s}{\PYZsq{}}\PY{p}{,} \PY{l+s}{\PYZsq{}}\PY{l+s}{a2}\PY{l+s}{\PYZsq{}}\PY{p}{]}\PY{p}{,} \PY{l+s}{\PYZsq{}}\PY{l+s}{B}\PY{l+s}{\PYZsq{}}\PY{p}{:} \PY{p}{[}\PY{l+s}{\PYZsq{}}\PY{l+s}{b1}\PY{l+s}{\PYZsq{}}\PY{p}{,} \PY{l+s}{\PYZsq{}}\PY{l+s}{b2}\PY{l+s}{\PYZsq{}}\PY{p}{]}\PY{p}{\PYZcb{}}\PY{p}{,}
\PY{g+gp}{... }    \PY{p}{\PYZob{}}
\PY{g+gp}{... }        \PY{p}{(}\PY{l+s}{\PYZsq{}}\PY{l+s}{a1}\PY{l+s}{\PYZsq{}}\PY{p}{,} \PY{l+s}{\PYZsq{}}\PY{l+s}{b1}\PY{l+s}{\PYZsq{}}\PY{p}{)}\PY{p}{:} \PY{l+m+mf}{0.8}\PY{p}{,}
\PY{g+gp}{... }        \PY{p}{(}\PY{l+s}{\PYZsq{}}\PY{l+s}{a1}\PY{l+s}{\PYZsq{}}\PY{p}{,} \PY{l+s}{\PYZsq{}}\PY{l+s}{b2}\PY{l+s}{\PYZsq{}}\PY{p}{)}\PY{p}{:} \PY{l+m+mf}{0.2}\PY{p}{,}
\PY{g+gp}{... }        \PY{p}{(}\PY{l+s}{\PYZsq{}}\PY{l+s}{a2}\PY{l+s}{\PYZsq{}}\PY{p}{,} \PY{l+s}{\PYZsq{}}\PY{l+s}{b1}\PY{l+s}{\PYZsq{}}\PY{p}{)}\PY{p}{:} \PY{l+m+mf}{0.3}\PY{p}{,}
\PY{g+gp}{... }        \PY{p}{(}\PY{l+s}{\PYZsq{}}\PY{l+s}{a2}\PY{l+s}{\PYZsq{}}\PY{p}{,} \PY{l+s}{\PYZsq{}}\PY{l+s}{b2}\PY{l+s}{\PYZsq{}}\PY{p}{)}\PY{p}{:} \PY{l+m+mf}{0.7}
\PY{g+gp}{... }    \PY{p}{\PYZcb{}}\PY{p}{)}
\PY{g+gp}{\PYZgt{}\PYZgt{}\PYZgt{} }\PY{n}{f}\PY{p}{[}\PY{p}{(}\PY{l+s}{\PYZsq{}}\PY{l+s}{a2}\PY{l+s}{\PYZsq{}}\PY{p}{,} \PY{l+s}{\PYZsq{}}\PY{l+s}{b1}\PY{l+s}{\PYZsq{}}\PY{p}{)}\PY{p}{]} \PY{o}{=} \PY{l+m+mf}{0.4}
\PY{g+gp}{\PYZgt{}\PYZgt{}\PYZgt{} }\PY{n}{f}\PY{p}{[}\PY{p}{(}\PY{l+s}{\PYZsq{}}\PY{l+s}{a2}\PY{l+s}{\PYZsq{}}\PY{p}{,} \PY{l+s}{\PYZsq{}}\PY{l+s}{b1}\PY{l+s}{\PYZsq{}}\PY{p}{)}\PY{p}{]}
\PY{g+go}{0.4}
\end{Verbatim}
\begin{quote}\begin{description}
\item[{Parametry}] \leavevmode\begin{itemize}
\item {} 
\textbf{assignment} -- přiřazení proměnných.

\item {} 
\textbf{value} -- nová hodnota.

\end{itemize}

\end{description}\end{quote}

\end{fulllineitems}

\index{pretty\_print() (metoda alex.infer.factor.Factor)}

\begin{fulllineitems}
\phantomsection\label{alex.infer:alex.infer.factor.Factor.pretty_print}\pysiglinewithargsret{\bfcode{pretty\_print}}{\emph{width=79}, \emph{precision=10}}{}
Vypsání faktoru ve formě tabulky.

\begin{Verbatim}[commandchars=\\\{\}]
\PY{g+gp}{\PYZgt{}\PYZgt{}\PYZgt{} }\PY{n}{f} \PY{o}{=} \PY{n}{Factor}\PY{p}{(}\PY{p}{[}\PY{l+s}{\PYZsq{}}\PY{l+s}{A}\PY{l+s}{\PYZsq{}}\PY{p}{,} \PY{l+s}{\PYZsq{}}\PY{l+s}{B}\PY{l+s}{\PYZsq{}}\PY{p}{]}\PY{p}{,}
\PY{g+gp}{... }    \PY{p}{\PYZob{}}\PY{l+s}{\PYZsq{}}\PY{l+s}{A}\PY{l+s}{\PYZsq{}}\PY{p}{:} \PY{p}{[}\PY{l+s}{\PYZsq{}}\PY{l+s}{a1}\PY{l+s}{\PYZsq{}}\PY{p}{,} \PY{l+s}{\PYZsq{}}\PY{l+s}{a2}\PY{l+s}{\PYZsq{}}\PY{p}{]}\PY{p}{,} \PY{l+s}{\PYZsq{}}\PY{l+s}{B}\PY{l+s}{\PYZsq{}}\PY{p}{:} \PY{p}{[}\PY{l+s}{\PYZsq{}}\PY{l+s}{b1}\PY{l+s}{\PYZsq{}}\PY{p}{,} \PY{l+s}{\PYZsq{}}\PY{l+s}{b2}\PY{l+s}{\PYZsq{}}\PY{p}{]}\PY{p}{\PYZcb{}}\PY{p}{,}
\PY{g+gp}{... }    \PY{p}{\PYZob{}}
\PY{g+gp}{... }        \PY{p}{(}\PY{l+s}{\PYZsq{}}\PY{l+s}{a1}\PY{l+s}{\PYZsq{}}\PY{p}{,} \PY{l+s}{\PYZsq{}}\PY{l+s}{b1}\PY{l+s}{\PYZsq{}}\PY{p}{)}\PY{p}{:} \PY{l+m+mf}{0.8}\PY{p}{,}
\PY{g+gp}{... }        \PY{p}{(}\PY{l+s}{\PYZsq{}}\PY{l+s}{a1}\PY{l+s}{\PYZsq{}}\PY{p}{,} \PY{l+s}{\PYZsq{}}\PY{l+s}{b2}\PY{l+s}{\PYZsq{}}\PY{p}{)}\PY{p}{:} \PY{l+m+mf}{0.2}\PY{p}{,}
\PY{g+gp}{... }        \PY{p}{(}\PY{l+s}{\PYZsq{}}\PY{l+s}{a2}\PY{l+s}{\PYZsq{}}\PY{p}{,} \PY{l+s}{\PYZsq{}}\PY{l+s}{b1}\PY{l+s}{\PYZsq{}}\PY{p}{)}\PY{p}{:} \PY{l+m+mf}{0.3}\PY{p}{,}
\PY{g+gp}{... }        \PY{p}{(}\PY{l+s}{\PYZsq{}}\PY{l+s}{a2}\PY{l+s}{\PYZsq{}}\PY{p}{,} \PY{l+s}{\PYZsq{}}\PY{l+s}{b2}\PY{l+s}{\PYZsq{}}\PY{p}{)}\PY{p}{:} \PY{l+m+mf}{0.7}
\PY{g+gp}{... }    \PY{p}{\PYZcb{}}\PY{p}{)}
\PY{g+gp}{\PYZgt{}\PYZgt{}\PYZgt{} }\PY{k}{print} \PY{n}{f}\PY{o}{.}\PY{n}{pretty\PYZus{}print}\PY{p}{(}\PY{l+m+mi}{50}\PY{p}{,} \PY{l+m+mi}{2}\PY{p}{)}
\PY{g+go}{\PYZhy{}\PYZhy{}\PYZhy{}\PYZhy{}\PYZhy{}\PYZhy{}\PYZhy{}\PYZhy{}\PYZhy{}\PYZhy{}\PYZhy{}\PYZhy{}\PYZhy{}\PYZhy{}\PYZhy{}\PYZhy{}\PYZhy{}\PYZhy{}\PYZhy{}\PYZhy{}\PYZhy{}\PYZhy{}\PYZhy{}\PYZhy{}\PYZhy{}\PYZhy{}\PYZhy{}\PYZhy{}\PYZhy{}\PYZhy{}\PYZhy{}\PYZhy{}\PYZhy{}\PYZhy{}\PYZhy{}\PYZhy{}\PYZhy{}\PYZhy{}\PYZhy{}\PYZhy{}\PYZhy{}\PYZhy{}\PYZhy{}\PYZhy{}\PYZhy{}\PYZhy{}\PYZhy{}\PYZhy{}\PYZhy{}\PYZhy{}}
\PY{g+go}{    A               B             Value}
\PY{g+go}{\PYZhy{}\PYZhy{}\PYZhy{}\PYZhy{}\PYZhy{}\PYZhy{}\PYZhy{}\PYZhy{}\PYZhy{}\PYZhy{}\PYZhy{}\PYZhy{}\PYZhy{}\PYZhy{}\PYZhy{}\PYZhy{}\PYZhy{}\PYZhy{}\PYZhy{}\PYZhy{}\PYZhy{}\PYZhy{}\PYZhy{}\PYZhy{}\PYZhy{}\PYZhy{}\PYZhy{}\PYZhy{}\PYZhy{}\PYZhy{}\PYZhy{}\PYZhy{}\PYZhy{}\PYZhy{}\PYZhy{}\PYZhy{}\PYZhy{}\PYZhy{}\PYZhy{}\PYZhy{}\PYZhy{}\PYZhy{}\PYZhy{}\PYZhy{}\PYZhy{}\PYZhy{}\PYZhy{}\PYZhy{}\PYZhy{}\PYZhy{}}
\PY{g+go}{    a1              b1             0.8}
\PY{g+go}{    a1              b2             0.2}
\PY{g+go}{    a2              b1             0.3}
\PY{g+go}{    a2              b2             0.7}
\PY{g+go}{\PYZhy{}\PYZhy{}\PYZhy{}\PYZhy{}\PYZhy{}\PYZhy{}\PYZhy{}\PYZhy{}\PYZhy{}\PYZhy{}\PYZhy{}\PYZhy{}\PYZhy{}\PYZhy{}\PYZhy{}\PYZhy{}\PYZhy{}\PYZhy{}\PYZhy{}\PYZhy{}\PYZhy{}\PYZhy{}\PYZhy{}\PYZhy{}\PYZhy{}\PYZhy{}\PYZhy{}\PYZhy{}\PYZhy{}\PYZhy{}\PYZhy{}\PYZhy{}\PYZhy{}\PYZhy{}\PYZhy{}\PYZhy{}\PYZhy{}\PYZhy{}\PYZhy{}\PYZhy{}\PYZhy{}\PYZhy{}\PYZhy{}\PYZhy{}\PYZhy{}\PYZhy{}\PYZhy{}\PYZhy{}\PYZhy{}\PYZhy{}}
\end{Verbatim}
\begin{quote}\begin{description}
\item[{Parametry}] \leavevmode\begin{itemize}
\item {} 
\textbf{width} -- šířka tabulky.

\item {} 
\textbf{precision} -- počet desetinných míst.

\end{itemize}

\item[{Vrací}] \leavevmode
řetězec s naformátovanou tabulkou.

\end{description}\end{quote}

\end{fulllineitems}

\index{\_\_iter\_\_() (metoda alex.infer.factor.Factor)}

\begin{fulllineitems}
\phantomsection\label{alex.infer:alex.infer.factor.Factor.__iter__}\pysiglinewithargsret{\bfcode{\_\_iter\_\_}}{}{}
Iterace přes jednotlivé hodnoty přiřazení.

\begin{Verbatim}[commandchars=\\\{\}]
\PY{g+gp}{\PYZgt{}\PYZgt{}\PYZgt{} }\PY{n}{f} \PY{o}{=} \PY{n}{Factor}\PY{p}{(}\PY{p}{[}\PY{l+s}{\PYZsq{}}\PY{l+s}{A}\PY{l+s}{\PYZsq{}}\PY{p}{,} \PY{l+s}{\PYZsq{}}\PY{l+s}{B}\PY{l+s}{\PYZsq{}}\PY{p}{]}\PY{p}{,}
\PY{g+gp}{... }    \PY{p}{\PYZob{}}\PY{l+s}{\PYZsq{}}\PY{l+s}{A}\PY{l+s}{\PYZsq{}}\PY{p}{:} \PY{p}{[}\PY{l+s}{\PYZsq{}}\PY{l+s}{a1}\PY{l+s}{\PYZsq{}}\PY{p}{,} \PY{l+s}{\PYZsq{}}\PY{l+s}{a2}\PY{l+s}{\PYZsq{}}\PY{p}{]}\PY{p}{,} \PY{l+s}{\PYZsq{}}\PY{l+s}{B}\PY{l+s}{\PYZsq{}}\PY{p}{:} \PY{p}{[}\PY{l+s}{\PYZsq{}}\PY{l+s}{b1}\PY{l+s}{\PYZsq{}}\PY{p}{,} \PY{l+s}{\PYZsq{}}\PY{l+s}{b2}\PY{l+s}{\PYZsq{}}\PY{p}{]}\PY{p}{\PYZcb{}}\PY{p}{,}
\PY{g+gp}{... }    \PY{p}{\PYZob{}}
\PY{g+gp}{... }        \PY{p}{(}\PY{l+s}{\PYZsq{}}\PY{l+s}{a1}\PY{l+s}{\PYZsq{}}\PY{p}{,} \PY{l+s}{\PYZsq{}}\PY{l+s}{b1}\PY{l+s}{\PYZsq{}}\PY{p}{)}\PY{p}{:} \PY{l+m+mf}{0.8}\PY{p}{,}
\PY{g+gp}{... }        \PY{p}{(}\PY{l+s}{\PYZsq{}}\PY{l+s}{a1}\PY{l+s}{\PYZsq{}}\PY{p}{,} \PY{l+s}{\PYZsq{}}\PY{l+s}{b2}\PY{l+s}{\PYZsq{}}\PY{p}{)}\PY{p}{:} \PY{l+m+mf}{0.2}\PY{p}{,}
\PY{g+gp}{... }        \PY{p}{(}\PY{l+s}{\PYZsq{}}\PY{l+s}{a2}\PY{l+s}{\PYZsq{}}\PY{p}{,} \PY{l+s}{\PYZsq{}}\PY{l+s}{b1}\PY{l+s}{\PYZsq{}}\PY{p}{)}\PY{p}{:} \PY{l+m+mf}{0.3}\PY{p}{,}
\PY{g+gp}{... }        \PY{p}{(}\PY{l+s}{\PYZsq{}}\PY{l+s}{a2}\PY{l+s}{\PYZsq{}}\PY{p}{,} \PY{l+s}{\PYZsq{}}\PY{l+s}{b2}\PY{l+s}{\PYZsq{}}\PY{p}{)}\PY{p}{:} \PY{l+m+mf}{0.7}
\PY{g+gp}{... }    \PY{p}{\PYZcb{}}\PY{p}{)}
\PY{g+gp}{\PYZgt{}\PYZgt{}\PYZgt{} }\PY{k}{for} \PY{p}{(}\PY{n}{assignment}\PY{p}{,} \PY{n}{value}\PY{p}{)} \PY{o+ow}{in} \PY{n}{f}\PY{p}{:}
\PY{g+gp}{... }    \PY{k}{print} \PY{n}{assignment}\PY{p}{,} \PY{l+s}{\PYZsq{}}\PY{l+s}{\PYZhy{}\PYZgt{}}\PY{l+s}{\PYZsq{}}\PY{p}{,} \PY{n}{value}
\PY{g+go}{(\PYZsq{}a1\PYZsq{}, \PYZsq{}b1\PYZsq{}) \PYZhy{}\PYZgt{} 0.8}
\PY{g+go}{(\PYZsq{}a1\PYZsq{}, \PYZsq{}b2\PYZsq{}) \PYZhy{}\PYZgt{} 0.2}
\PY{g+go}{(\PYZsq{}a2\PYZsq{}, \PYZsq{}b1\PYZsq{}) \PYZhy{}\PYZgt{} 0.3}
\PY{g+go}{(\PYZsq{}a2\PYZsq{}, \PYZsq{}b2\PYZsq{}) \PYZhy{}\PYZgt{} 0.7}
\end{Verbatim}

\end{fulllineitems}

\index{observed() (metoda alex.infer.factor.Factor)}

\begin{fulllineitems}
\phantomsection\label{alex.infer:alex.infer.factor.Factor.observed}\pysiglinewithargsret{\bfcode{observed}}{\emph{assignment\_dict}}{}
Nastavení pozorovaných hodnot.

\begin{Verbatim}[commandchars=\\\{\}]
\PY{g+gp}{\PYZgt{}\PYZgt{}\PYZgt{} }\PY{n}{f} \PY{o}{=} \PY{n}{Factor}\PY{p}{(}\PY{p}{[}\PY{l+s}{\PYZsq{}}\PY{l+s}{A}\PY{l+s}{\PYZsq{}}\PY{p}{,} \PY{l+s}{\PYZsq{}}\PY{l+s}{B}\PY{l+s}{\PYZsq{}}\PY{p}{]}\PY{p}{,}
\PY{g+gp}{... }    \PY{p}{\PYZob{}}\PY{l+s}{\PYZsq{}}\PY{l+s}{A}\PY{l+s}{\PYZsq{}}\PY{p}{:} \PY{p}{[}\PY{l+s}{\PYZsq{}}\PY{l+s}{a1}\PY{l+s}{\PYZsq{}}\PY{p}{,} \PY{l+s}{\PYZsq{}}\PY{l+s}{a2}\PY{l+s}{\PYZsq{}}\PY{p}{]}\PY{p}{,} \PY{l+s}{\PYZsq{}}\PY{l+s}{B}\PY{l+s}{\PYZsq{}}\PY{p}{:} \PY{p}{[}\PY{l+s}{\PYZsq{}}\PY{l+s}{b1}\PY{l+s}{\PYZsq{}}\PY{p}{,} \PY{l+s}{\PYZsq{}}\PY{l+s}{b2}\PY{l+s}{\PYZsq{}}\PY{p}{]}\PY{p}{\PYZcb{}}\PY{p}{,}
\PY{g+gp}{... }    \PY{p}{\PYZob{}}
\PY{g+gp}{... }        \PY{p}{(}\PY{l+s}{\PYZsq{}}\PY{l+s}{a1}\PY{l+s}{\PYZsq{}}\PY{p}{,} \PY{l+s}{\PYZsq{}}\PY{l+s}{b1}\PY{l+s}{\PYZsq{}}\PY{p}{)}\PY{p}{:} \PY{l+m+mf}{0.8}\PY{p}{,}
\PY{g+gp}{... }        \PY{p}{(}\PY{l+s}{\PYZsq{}}\PY{l+s}{a1}\PY{l+s}{\PYZsq{}}\PY{p}{,} \PY{l+s}{\PYZsq{}}\PY{l+s}{b2}\PY{l+s}{\PYZsq{}}\PY{p}{)}\PY{p}{:} \PY{l+m+mf}{0.2}\PY{p}{,}
\PY{g+gp}{... }        \PY{p}{(}\PY{l+s}{\PYZsq{}}\PY{l+s}{a2}\PY{l+s}{\PYZsq{}}\PY{p}{,} \PY{l+s}{\PYZsq{}}\PY{l+s}{b1}\PY{l+s}{\PYZsq{}}\PY{p}{)}\PY{p}{:} \PY{l+m+mf}{0.3}\PY{p}{,}
\PY{g+gp}{... }        \PY{p}{(}\PY{l+s}{\PYZsq{}}\PY{l+s}{a2}\PY{l+s}{\PYZsq{}}\PY{p}{,} \PY{l+s}{\PYZsq{}}\PY{l+s}{b2}\PY{l+s}{\PYZsq{}}\PY{p}{)}\PY{p}{:} \PY{l+m+mf}{0.7}
\PY{g+gp}{... }    \PY{p}{\PYZcb{}}\PY{p}{)}
\PY{g+gp}{\PYZgt{}\PYZgt{}\PYZgt{} }\PY{n}{f}\PY{o}{.}\PY{n}{observed}\PY{p}{(}\PY{p}{\PYZob{}}
\PY{g+gp}{... }    \PY{p}{(}\PY{l+s}{\PYZsq{}}\PY{l+s}{a1}\PY{l+s}{\PYZsq{}}\PY{p}{,} \PY{l+s}{\PYZsq{}}\PY{l+s}{b1}\PY{l+s}{\PYZsq{}}\PY{p}{)}\PY{p}{:} \PY{l+m+mf}{0.6}\PY{p}{,}
\PY{g+gp}{... }    \PY{p}{(}\PY{l+s}{\PYZsq{}}\PY{l+s}{a1}\PY{l+s}{\PYZsq{}}\PY{p}{,} \PY{l+s}{\PYZsq{}}\PY{l+s}{b2}\PY{l+s}{\PYZsq{}}\PY{p}{)}\PY{p}{:} \PY{l+m+mf}{0.4}\PY{p}{,}
\PY{g+gp}{... }\PY{p}{\PYZcb{}}\PY{p}{)}
\PY{g+gp}{\PYZgt{}\PYZgt{}\PYZgt{} }\PY{k}{print} \PY{n}{f}\PY{o}{.}\PY{n}{pretty\PYZus{}print}\PY{p}{(}\PY{l+m+mi}{50}\PY{p}{,} \PY{l+m+mi}{2}\PY{p}{)}
\PY{g+go}{\PYZhy{}\PYZhy{}\PYZhy{}\PYZhy{}\PYZhy{}\PYZhy{}\PYZhy{}\PYZhy{}\PYZhy{}\PYZhy{}\PYZhy{}\PYZhy{}\PYZhy{}\PYZhy{}\PYZhy{}\PYZhy{}\PYZhy{}\PYZhy{}\PYZhy{}\PYZhy{}\PYZhy{}\PYZhy{}\PYZhy{}\PYZhy{}\PYZhy{}\PYZhy{}\PYZhy{}\PYZhy{}\PYZhy{}\PYZhy{}\PYZhy{}\PYZhy{}\PYZhy{}\PYZhy{}\PYZhy{}\PYZhy{}\PYZhy{}\PYZhy{}\PYZhy{}\PYZhy{}\PYZhy{}\PYZhy{}\PYZhy{}\PYZhy{}\PYZhy{}\PYZhy{}\PYZhy{}\PYZhy{}\PYZhy{}\PYZhy{}}
\PY{g+go}{    A               B             Value}
\PY{g+go}{\PYZhy{}\PYZhy{}\PYZhy{}\PYZhy{}\PYZhy{}\PYZhy{}\PYZhy{}\PYZhy{}\PYZhy{}\PYZhy{}\PYZhy{}\PYZhy{}\PYZhy{}\PYZhy{}\PYZhy{}\PYZhy{}\PYZhy{}\PYZhy{}\PYZhy{}\PYZhy{}\PYZhy{}\PYZhy{}\PYZhy{}\PYZhy{}\PYZhy{}\PYZhy{}\PYZhy{}\PYZhy{}\PYZhy{}\PYZhy{}\PYZhy{}\PYZhy{}\PYZhy{}\PYZhy{}\PYZhy{}\PYZhy{}\PYZhy{}\PYZhy{}\PYZhy{}\PYZhy{}\PYZhy{}\PYZhy{}\PYZhy{}\PYZhy{}\PYZhy{}\PYZhy{}\PYZhy{}\PYZhy{}\PYZhy{}\PYZhy{}}
\PY{g+go}{    a1              b1             0.6}
\PY{g+go}{    a1              b2             0.4}
\PY{g+go}{    a2              b1             0.0}
\PY{g+go}{    a2              b2             0.0}
\PY{g+go}{\PYZhy{}\PYZhy{}\PYZhy{}\PYZhy{}\PYZhy{}\PYZhy{}\PYZhy{}\PYZhy{}\PYZhy{}\PYZhy{}\PYZhy{}\PYZhy{}\PYZhy{}\PYZhy{}\PYZhy{}\PYZhy{}\PYZhy{}\PYZhy{}\PYZhy{}\PYZhy{}\PYZhy{}\PYZhy{}\PYZhy{}\PYZhy{}\PYZhy{}\PYZhy{}\PYZhy{}\PYZhy{}\PYZhy{}\PYZhy{}\PYZhy{}\PYZhy{}\PYZhy{}\PYZhy{}\PYZhy{}\PYZhy{}\PYZhy{}\PYZhy{}\PYZhy{}\PYZhy{}\PYZhy{}\PYZhy{}\PYZhy{}\PYZhy{}\PYZhy{}\PYZhy{}\PYZhy{}\PYZhy{}\PYZhy{}\PYZhy{}}
\end{Verbatim}
\begin{quote}\begin{description}
\item[{Parametr}] \leavevmode
\textbf{assignment\_dict} -- slovník s pozorovanými přiřazeními a jejich ohodnocením.

\end{description}\end{quote}

\end{fulllineitems}

\index{normalize() (metoda alex.infer.factor.Factor)}

\begin{fulllineitems}
\phantomsection\label{alex.infer:alex.infer.factor.Factor.normalize}\pysiglinewithargsret{\bfcode{normalize}}{\optional{\emph{parents}}}{}
Normalizace hodnot ve faktoru, aby jejich součet byl roven 1.

Pokud je specifikován parametr \code{parents}, pak musí být součet všech
proměnných, které mají stejné přiřazení pro proměnné z \code{parents}, roven
1.

\begin{Verbatim}[commandchars=\\\{\}]
\PY{g+gp}{\PYZgt{}\PYZgt{}\PYZgt{} }\PY{n}{f} \PY{o}{=} \PY{n}{Factor}\PY{p}{(}\PY{p}{[}\PY{l+s}{\PYZsq{}}\PY{l+s}{A}\PY{l+s}{\PYZsq{}}\PY{p}{,} \PY{l+s}{\PYZsq{}}\PY{l+s}{B}\PY{l+s}{\PYZsq{}}\PY{p}{]}\PY{p}{,}
\PY{g+gp}{... }    \PY{p}{\PYZob{}}\PY{l+s}{\PYZsq{}}\PY{l+s}{A}\PY{l+s}{\PYZsq{}}\PY{p}{:} \PY{p}{[}\PY{l+s}{\PYZsq{}}\PY{l+s}{a1}\PY{l+s}{\PYZsq{}}\PY{p}{,} \PY{l+s}{\PYZsq{}}\PY{l+s}{a2}\PY{l+s}{\PYZsq{}}\PY{p}{]}\PY{p}{,} \PY{l+s}{\PYZsq{}}\PY{l+s}{B}\PY{l+s}{\PYZsq{}}\PY{p}{:} \PY{p}{[}\PY{l+s}{\PYZsq{}}\PY{l+s}{b1}\PY{l+s}{\PYZsq{}}\PY{p}{,} \PY{l+s}{\PYZsq{}}\PY{l+s}{b2}\PY{l+s}{\PYZsq{}}\PY{p}{]}\PY{p}{\PYZcb{}}\PY{p}{,}
\PY{g+gp}{... }    \PY{p}{\PYZob{}}
\PY{g+gp}{... }        \PY{p}{(}\PY{l+s}{\PYZsq{}}\PY{l+s}{a1}\PY{l+s}{\PYZsq{}}\PY{p}{,} \PY{l+s}{\PYZsq{}}\PY{l+s}{b1}\PY{l+s}{\PYZsq{}}\PY{p}{)}\PY{p}{:} \PY{l+m+mi}{1}\PY{p}{,}
\PY{g+gp}{... }        \PY{p}{(}\PY{l+s}{\PYZsq{}}\PY{l+s}{a1}\PY{l+s}{\PYZsq{}}\PY{p}{,} \PY{l+s}{\PYZsq{}}\PY{l+s}{b2}\PY{l+s}{\PYZsq{}}\PY{p}{)}\PY{p}{:} \PY{l+m+mi}{3}\PY{p}{,}
\PY{g+gp}{... }        \PY{p}{(}\PY{l+s}{\PYZsq{}}\PY{l+s}{a2}\PY{l+s}{\PYZsq{}}\PY{p}{,} \PY{l+s}{\PYZsq{}}\PY{l+s}{b1}\PY{l+s}{\PYZsq{}}\PY{p}{)}\PY{p}{:} \PY{l+m+mi}{2}\PY{p}{,}
\PY{g+gp}{... }        \PY{p}{(}\PY{l+s}{\PYZsq{}}\PY{l+s}{a2}\PY{l+s}{\PYZsq{}}\PY{p}{,} \PY{l+s}{\PYZsq{}}\PY{l+s}{b2}\PY{l+s}{\PYZsq{}}\PY{p}{)}\PY{p}{:} \PY{l+m+mi}{2}\PY{p}{,}
\PY{g+gp}{... }    \PY{p}{\PYZcb{}}\PY{p}{)}
\PY{g+gp}{\PYZgt{}\PYZgt{}\PYZgt{} }\PY{n}{f}\PY{o}{.}\PY{n}{normalize}\PY{p}{(}\PY{n}{parents}\PY{o}{=}\PY{p}{[}\PY{l+s}{\PYZsq{}}\PY{l+s}{A}\PY{l+s}{\PYZsq{}}\PY{p}{]}\PY{p}{)}
\PY{g+gp}{\PYZgt{}\PYZgt{}\PYZgt{} }\PY{k}{print} \PY{n}{f}\PY{o}{.}\PY{n}{pretty\PYZus{}print}\PY{p}{(}\PY{l+m+mi}{50}\PY{p}{,} \PY{l+m+mi}{2}\PY{p}{)}
\PY{g+go}{\PYZhy{}\PYZhy{}\PYZhy{}\PYZhy{}\PYZhy{}\PYZhy{}\PYZhy{}\PYZhy{}\PYZhy{}\PYZhy{}\PYZhy{}\PYZhy{}\PYZhy{}\PYZhy{}\PYZhy{}\PYZhy{}\PYZhy{}\PYZhy{}\PYZhy{}\PYZhy{}\PYZhy{}\PYZhy{}\PYZhy{}\PYZhy{}\PYZhy{}\PYZhy{}\PYZhy{}\PYZhy{}\PYZhy{}\PYZhy{}\PYZhy{}\PYZhy{}\PYZhy{}\PYZhy{}\PYZhy{}\PYZhy{}\PYZhy{}\PYZhy{}\PYZhy{}\PYZhy{}\PYZhy{}\PYZhy{}\PYZhy{}\PYZhy{}\PYZhy{}\PYZhy{}\PYZhy{}\PYZhy{}\PYZhy{}\PYZhy{}}
\PY{g+go}{    A               B             Value}
\PY{g+go}{\PYZhy{}\PYZhy{}\PYZhy{}\PYZhy{}\PYZhy{}\PYZhy{}\PYZhy{}\PYZhy{}\PYZhy{}\PYZhy{}\PYZhy{}\PYZhy{}\PYZhy{}\PYZhy{}\PYZhy{}\PYZhy{}\PYZhy{}\PYZhy{}\PYZhy{}\PYZhy{}\PYZhy{}\PYZhy{}\PYZhy{}\PYZhy{}\PYZhy{}\PYZhy{}\PYZhy{}\PYZhy{}\PYZhy{}\PYZhy{}\PYZhy{}\PYZhy{}\PYZhy{}\PYZhy{}\PYZhy{}\PYZhy{}\PYZhy{}\PYZhy{}\PYZhy{}\PYZhy{}\PYZhy{}\PYZhy{}\PYZhy{}\PYZhy{}\PYZhy{}\PYZhy{}\PYZhy{}\PYZhy{}\PYZhy{}\PYZhy{}}
\PY{g+go}{    a1              b1             0.25}
\PY{g+go}{    a1              b2             0.75}
\PY{g+go}{    a2              b1             0.5}
\PY{g+go}{    a2              b2             0.5}
\PY{g+go}{\PYZhy{}\PYZhy{}\PYZhy{}\PYZhy{}\PYZhy{}\PYZhy{}\PYZhy{}\PYZhy{}\PYZhy{}\PYZhy{}\PYZhy{}\PYZhy{}\PYZhy{}\PYZhy{}\PYZhy{}\PYZhy{}\PYZhy{}\PYZhy{}\PYZhy{}\PYZhy{}\PYZhy{}\PYZhy{}\PYZhy{}\PYZhy{}\PYZhy{}\PYZhy{}\PYZhy{}\PYZhy{}\PYZhy{}\PYZhy{}\PYZhy{}\PYZhy{}\PYZhy{}\PYZhy{}\PYZhy{}\PYZhy{}\PYZhy{}\PYZhy{}\PYZhy{}\PYZhy{}\PYZhy{}\PYZhy{}\PYZhy{}\PYZhy{}\PYZhy{}\PYZhy{}\PYZhy{}\PYZhy{}\PYZhy{}\PYZhy{}}
\end{Verbatim}
\begin{quote}\begin{description}
\item[{Parametr}] \leavevmode
\textbf{parents} -- seznam rodičovských proměnných.

\end{description}\end{quote}

\end{fulllineitems}

\index{marginalize() (metoda alex.infer.factor.Factor)}

\begin{fulllineitems}
\phantomsection\label{alex.infer:alex.infer.factor.Factor.marginalize}\pysiglinewithargsret{\bfcode{marginalize}}{\emph{keep}}{}
Vysčítání proměnných tak, aby zůstaly pouze proměnné z \code{keep}.

\begin{Verbatim}[commandchars=\\\{\}]
\PY{g+gp}{\PYZgt{}\PYZgt{}\PYZgt{} }\PY{n}{f} \PY{o}{=} \PY{n}{Factor}\PY{p}{(}\PY{p}{[}\PY{l+s}{\PYZsq{}}\PY{l+s}{A}\PY{l+s}{\PYZsq{}}\PY{p}{,} \PY{l+s}{\PYZsq{}}\PY{l+s}{B}\PY{l+s}{\PYZsq{}}\PY{p}{]}\PY{p}{,}
\PY{g+gp}{... }    \PY{p}{\PYZob{}}\PY{l+s}{\PYZsq{}}\PY{l+s}{A}\PY{l+s}{\PYZsq{}}\PY{p}{:} \PY{p}{[}\PY{l+s}{\PYZsq{}}\PY{l+s}{a1}\PY{l+s}{\PYZsq{}}\PY{p}{,} \PY{l+s}{\PYZsq{}}\PY{l+s}{a2}\PY{l+s}{\PYZsq{}}\PY{p}{]}\PY{p}{,} \PY{l+s}{\PYZsq{}}\PY{l+s}{B}\PY{l+s}{\PYZsq{}}\PY{p}{:} \PY{p}{[}\PY{l+s}{\PYZsq{}}\PY{l+s}{b1}\PY{l+s}{\PYZsq{}}\PY{p}{,} \PY{l+s}{\PYZsq{}}\PY{l+s}{b2}\PY{l+s}{\PYZsq{}}\PY{p}{]}\PY{p}{\PYZcb{}}\PY{p}{,}
\PY{g+gp}{... }    \PY{p}{\PYZob{}}
\PY{g+gp}{... }        \PY{p}{(}\PY{l+s}{\PYZsq{}}\PY{l+s}{a1}\PY{l+s}{\PYZsq{}}\PY{p}{,} \PY{l+s}{\PYZsq{}}\PY{l+s}{b1}\PY{l+s}{\PYZsq{}}\PY{p}{)}\PY{p}{:} \PY{l+m+mf}{0.8}\PY{p}{,}
\PY{g+gp}{... }        \PY{p}{(}\PY{l+s}{\PYZsq{}}\PY{l+s}{a1}\PY{l+s}{\PYZsq{}}\PY{p}{,} \PY{l+s}{\PYZsq{}}\PY{l+s}{b2}\PY{l+s}{\PYZsq{}}\PY{p}{)}\PY{p}{:} \PY{l+m+mf}{0.2}\PY{p}{,}
\PY{g+gp}{... }        \PY{p}{(}\PY{l+s}{\PYZsq{}}\PY{l+s}{a2}\PY{l+s}{\PYZsq{}}\PY{p}{,} \PY{l+s}{\PYZsq{}}\PY{l+s}{b1}\PY{l+s}{\PYZsq{}}\PY{p}{)}\PY{p}{:} \PY{l+m+mf}{0.3}\PY{p}{,}
\PY{g+gp}{... }        \PY{p}{(}\PY{l+s}{\PYZsq{}}\PY{l+s}{a2}\PY{l+s}{\PYZsq{}}\PY{p}{,} \PY{l+s}{\PYZsq{}}\PY{l+s}{b2}\PY{l+s}{\PYZsq{}}\PY{p}{)}\PY{p}{:} \PY{l+m+mf}{0.7}\PY{p}{,}
\PY{g+gp}{... }    \PY{p}{\PYZcb{}}\PY{p}{)}
\PY{g+gp}{\PYZgt{}\PYZgt{}\PYZgt{} }\PY{n}{g} \PY{o}{=} \PY{n}{f}\PY{o}{.}\PY{n}{marginalize}\PY{p}{(}\PY{n}{keep}\PY{o}{=}\PY{p}{[}\PY{l+s}{\PYZsq{}}\PY{l+s}{A}\PY{l+s}{\PYZsq{}}\PY{p}{]}\PY{p}{)}
\PY{g+gp}{\PYZgt{}\PYZgt{}\PYZgt{} }\PY{k}{print} \PY{n}{g}\PY{o}{.}\PY{n}{pretty\PYZus{}print}\PY{p}{(}\PY{l+m+mi}{50}\PY{p}{,} \PY{l+m+mi}{2}\PY{p}{)}
\PY{g+go}{\PYZhy{}\PYZhy{}\PYZhy{}\PYZhy{}\PYZhy{}\PYZhy{}\PYZhy{}\PYZhy{}\PYZhy{}\PYZhy{}\PYZhy{}\PYZhy{}\PYZhy{}\PYZhy{}\PYZhy{}\PYZhy{}\PYZhy{}\PYZhy{}\PYZhy{}\PYZhy{}\PYZhy{}\PYZhy{}\PYZhy{}\PYZhy{}\PYZhy{}\PYZhy{}\PYZhy{}\PYZhy{}\PYZhy{}\PYZhy{}\PYZhy{}\PYZhy{}\PYZhy{}\PYZhy{}\PYZhy{}\PYZhy{}\PYZhy{}\PYZhy{}\PYZhy{}\PYZhy{}\PYZhy{}\PYZhy{}\PYZhy{}\PYZhy{}\PYZhy{}\PYZhy{}\PYZhy{}\PYZhy{}\PYZhy{}\PYZhy{}}
\PY{g+go}{            A                      Value}
\PY{g+go}{\PYZhy{}\PYZhy{}\PYZhy{}\PYZhy{}\PYZhy{}\PYZhy{}\PYZhy{}\PYZhy{}\PYZhy{}\PYZhy{}\PYZhy{}\PYZhy{}\PYZhy{}\PYZhy{}\PYZhy{}\PYZhy{}\PYZhy{}\PYZhy{}\PYZhy{}\PYZhy{}\PYZhy{}\PYZhy{}\PYZhy{}\PYZhy{}\PYZhy{}\PYZhy{}\PYZhy{}\PYZhy{}\PYZhy{}\PYZhy{}\PYZhy{}\PYZhy{}\PYZhy{}\PYZhy{}\PYZhy{}\PYZhy{}\PYZhy{}\PYZhy{}\PYZhy{}\PYZhy{}\PYZhy{}\PYZhy{}\PYZhy{}\PYZhy{}\PYZhy{}\PYZhy{}\PYZhy{}\PYZhy{}\PYZhy{}\PYZhy{}}
\PY{g+go}{           a1                       1.0}
\PY{g+go}{           a2                       1.0}
\PY{g+go}{\PYZhy{}\PYZhy{}\PYZhy{}\PYZhy{}\PYZhy{}\PYZhy{}\PYZhy{}\PYZhy{}\PYZhy{}\PYZhy{}\PYZhy{}\PYZhy{}\PYZhy{}\PYZhy{}\PYZhy{}\PYZhy{}\PYZhy{}\PYZhy{}\PYZhy{}\PYZhy{}\PYZhy{}\PYZhy{}\PYZhy{}\PYZhy{}\PYZhy{}\PYZhy{}\PYZhy{}\PYZhy{}\PYZhy{}\PYZhy{}\PYZhy{}\PYZhy{}\PYZhy{}\PYZhy{}\PYZhy{}\PYZhy{}\PYZhy{}\PYZhy{}\PYZhy{}\PYZhy{}\PYZhy{}\PYZhy{}\PYZhy{}\PYZhy{}\PYZhy{}\PYZhy{}\PYZhy{}\PYZhy{}\PYZhy{}\PYZhy{}}
\end{Verbatim}
\begin{quote}\begin{description}
\item[{Parametr}] \leavevmode
\textbf{keep} -- seznam proměnných, které mají zůstat.

\end{description}\end{quote}

\end{fulllineitems}

\index{most\_probable() (metoda alex.infer.factor.Factor)}

\begin{fulllineitems}
\phantomsection\label{alex.infer:alex.infer.factor.Factor.most_probable}\pysiglinewithargsret{\bfcode{most\_probable}}{\optional{\emph{n}}}{}
Nalezení nejpravděpodobnějších přiřazení proměnných.

\begin{Verbatim}[commandchars=\\\{\}]
\PY{g+gp}{\PYZgt{}\PYZgt{}\PYZgt{} }\PY{n}{f} \PY{o}{=} \PY{n}{Factor}\PY{p}{(}\PY{p}{[}\PY{l+s}{\PYZsq{}}\PY{l+s}{A}\PY{l+s}{\PYZsq{}}\PY{p}{]}\PY{p}{,}
\PY{g+gp}{... }    \PY{p}{\PYZob{}}\PY{l+s}{\PYZsq{}}\PY{l+s}{A}\PY{l+s}{\PYZsq{}}\PY{p}{:} \PY{p}{[}\PY{l+s}{\PYZsq{}}\PY{l+s}{a1}\PY{l+s}{\PYZsq{}}\PY{p}{,} \PY{l+s}{\PYZsq{}}\PY{l+s}{a2}\PY{l+s}{\PYZsq{}}\PY{p}{]}\PY{p}{\PYZcb{}}\PY{p}{,}
\PY{g+gp}{... }    \PY{p}{\PYZob{}}
\PY{g+gp}{... }        \PY{p}{(}\PY{l+s}{\PYZsq{}}\PY{l+s}{a1}\PY{l+s}{\PYZsq{}}\PY{p}{,}\PY{p}{)}\PY{p}{:} \PY{l+m+mf}{0.8}\PY{p}{,}
\PY{g+gp}{... }        \PY{p}{(}\PY{l+s}{\PYZsq{}}\PY{l+s}{a2}\PY{l+s}{\PYZsq{}}\PY{p}{,}\PY{p}{)}\PY{p}{:} \PY{l+m+mf}{0.2}\PY{p}{,}
\PY{g+gp}{... }    \PY{p}{\PYZcb{}}\PY{p}{)}
\PY{g+gp}{\PYZgt{}\PYZgt{}\PYZgt{} }\PY{n}{f}\PY{o}{.}\PY{n}{most\PYZus{}probable}\PY{p}{(}\PY{l+m+mi}{2}\PY{p}{)}
\PY{g+go}{[(\PYZsq{}a1\PYZsq{}, 0.8), (\PYZsq{}a2\PYZsq{}, 0.2)]}
\end{Verbatim}
\begin{quote}\begin{description}
\item[{Parametr}] \leavevmode
\textbf{n} -- počet požadovaných nejpravděpodobnějších přiřazení.

\item[{Vrací}] \leavevmode
seznam nejpravděpodobnějších přiřazení a jejich hodnot.

\end{description}\end{quote}

\end{fulllineitems}

\index{rename\_variables() (metoda alex.infer.factor.Factor)}

\begin{fulllineitems}
\phantomsection\label{alex.infer:alex.infer.factor.Factor.rename_variables}\pysiglinewithargsret{\bfcode{rename\_variables}}{\emph{mapping}}{}
Přejmenování proměnných.

Může se stát, že je potřeba přejmenovat proměnné, aby odpovídaly proměnným v jiném faktoru.

\begin{Verbatim}[commandchars=\\\{\}]
\PY{g+gp}{\PYZgt{}\PYZgt{}\PYZgt{} }\PY{n}{f} \PY{o}{=} \PY{n}{Factor}\PY{p}{(}\PY{p}{[}\PY{l+s}{\PYZsq{}}\PY{l+s}{A}\PY{l+s}{\PYZsq{}}\PY{p}{,} \PY{l+s}{\PYZsq{}}\PY{l+s}{B}\PY{l+s}{\PYZsq{}}\PY{p}{]}\PY{p}{,}
\PY{g+gp}{... }    \PY{p}{\PYZob{}}\PY{l+s}{\PYZsq{}}\PY{l+s}{A}\PY{l+s}{\PYZsq{}}\PY{p}{:} \PY{p}{[}\PY{l+s}{\PYZsq{}}\PY{l+s}{a1}\PY{l+s}{\PYZsq{}}\PY{p}{,} \PY{l+s}{\PYZsq{}}\PY{l+s}{a2}\PY{l+s}{\PYZsq{}}\PY{p}{]}\PY{p}{,} \PY{l+s}{\PYZsq{}}\PY{l+s}{B}\PY{l+s}{\PYZsq{}}\PY{p}{:} \PY{p}{[}\PY{l+s}{\PYZsq{}}\PY{l+s}{b1}\PY{l+s}{\PYZsq{}}\PY{p}{,} \PY{l+s}{\PYZsq{}}\PY{l+s}{b2}\PY{l+s}{\PYZsq{}}\PY{p}{]}\PY{p}{\PYZcb{}}\PY{p}{,}
\PY{g+gp}{... }    \PY{p}{\PYZob{}}
\PY{g+gp}{... }        \PY{p}{(}\PY{l+s}{\PYZsq{}}\PY{l+s}{a1}\PY{l+s}{\PYZsq{}}\PY{p}{,} \PY{l+s}{\PYZsq{}}\PY{l+s}{b1}\PY{l+s}{\PYZsq{}}\PY{p}{)}\PY{p}{:} \PY{l+m+mf}{0.8}\PY{p}{,}
\PY{g+gp}{... }        \PY{p}{(}\PY{l+s}{\PYZsq{}}\PY{l+s}{a1}\PY{l+s}{\PYZsq{}}\PY{p}{,} \PY{l+s}{\PYZsq{}}\PY{l+s}{b2}\PY{l+s}{\PYZsq{}}\PY{p}{)}\PY{p}{:} \PY{l+m+mf}{0.2}\PY{p}{,}
\PY{g+gp}{... }        \PY{p}{(}\PY{l+s}{\PYZsq{}}\PY{l+s}{a2}\PY{l+s}{\PYZsq{}}\PY{p}{,} \PY{l+s}{\PYZsq{}}\PY{l+s}{b1}\PY{l+s}{\PYZsq{}}\PY{p}{)}\PY{p}{:} \PY{l+m+mf}{0.3}\PY{p}{,}
\PY{g+gp}{... }        \PY{p}{(}\PY{l+s}{\PYZsq{}}\PY{l+s}{a2}\PY{l+s}{\PYZsq{}}\PY{p}{,} \PY{l+s}{\PYZsq{}}\PY{l+s}{b2}\PY{l+s}{\PYZsq{}}\PY{p}{)}\PY{p}{:} \PY{l+m+mf}{0.7}\PY{p}{,}
\PY{g+gp}{... }    \PY{p}{\PYZcb{}}\PY{p}{)}
\PY{g+gp}{\PYZgt{}\PYZgt{}\PYZgt{} }\PY{n}{f}\PY{o}{.}\PY{n}{rename\PYZus{}variables}\PY{p}{(}\PY{p}{\PYZob{}}\PY{l+s}{\PYZsq{}}\PY{l+s}{A}\PY{l+s}{\PYZsq{}}\PY{p}{:} \PY{l+s}{\PYZsq{}}\PY{l+s}{X}\PY{l+s}{\PYZsq{}}\PY{p}{,} \PY{l+s}{\PYZsq{}}\PY{l+s}{B}\PY{l+s}{\PYZsq{}}\PY{p}{:} \PY{l+s}{\PYZsq{}}\PY{l+s}{Y}\PY{l+s}{\PYZsq{}}\PY{p}{\PYZcb{}}\PY{p}{)}
\PY{g+gp}{\PYZgt{}\PYZgt{}\PYZgt{} }\PY{k}{print} \PY{n}{f}\PY{o}{.}\PY{n}{pretty\PYZus{}print}\PY{p}{(}\PY{l+m+mi}{50}\PY{p}{,} \PY{l+m+mi}{2}\PY{p}{)}
\PY{g+go}{\PYZhy{}\PYZhy{}\PYZhy{}\PYZhy{}\PYZhy{}\PYZhy{}\PYZhy{}\PYZhy{}\PYZhy{}\PYZhy{}\PYZhy{}\PYZhy{}\PYZhy{}\PYZhy{}\PYZhy{}\PYZhy{}\PYZhy{}\PYZhy{}\PYZhy{}\PYZhy{}\PYZhy{}\PYZhy{}\PYZhy{}\PYZhy{}\PYZhy{}\PYZhy{}\PYZhy{}\PYZhy{}\PYZhy{}\PYZhy{}\PYZhy{}\PYZhy{}\PYZhy{}\PYZhy{}\PYZhy{}\PYZhy{}\PYZhy{}\PYZhy{}\PYZhy{}\PYZhy{}\PYZhy{}\PYZhy{}\PYZhy{}\PYZhy{}\PYZhy{}\PYZhy{}\PYZhy{}\PYZhy{}\PYZhy{}\PYZhy{}}
\PY{g+go}{    X               Y             Value}
\PY{g+go}{\PYZhy{}\PYZhy{}\PYZhy{}\PYZhy{}\PYZhy{}\PYZhy{}\PYZhy{}\PYZhy{}\PYZhy{}\PYZhy{}\PYZhy{}\PYZhy{}\PYZhy{}\PYZhy{}\PYZhy{}\PYZhy{}\PYZhy{}\PYZhy{}\PYZhy{}\PYZhy{}\PYZhy{}\PYZhy{}\PYZhy{}\PYZhy{}\PYZhy{}\PYZhy{}\PYZhy{}\PYZhy{}\PYZhy{}\PYZhy{}\PYZhy{}\PYZhy{}\PYZhy{}\PYZhy{}\PYZhy{}\PYZhy{}\PYZhy{}\PYZhy{}\PYZhy{}\PYZhy{}\PYZhy{}\PYZhy{}\PYZhy{}\PYZhy{}\PYZhy{}\PYZhy{}\PYZhy{}\PYZhy{}\PYZhy{}\PYZhy{}}
\PY{g+go}{    a1              b1             0.8}
\PY{g+go}{    a1              b2             0.2}
\PY{g+go}{    a2              b1             0.3}
\PY{g+go}{    a2              b2             0.7}
\PY{g+go}{\PYZhy{}\PYZhy{}\PYZhy{}\PYZhy{}\PYZhy{}\PYZhy{}\PYZhy{}\PYZhy{}\PYZhy{}\PYZhy{}\PYZhy{}\PYZhy{}\PYZhy{}\PYZhy{}\PYZhy{}\PYZhy{}\PYZhy{}\PYZhy{}\PYZhy{}\PYZhy{}\PYZhy{}\PYZhy{}\PYZhy{}\PYZhy{}\PYZhy{}\PYZhy{}\PYZhy{}\PYZhy{}\PYZhy{}\PYZhy{}\PYZhy{}\PYZhy{}\PYZhy{}\PYZhy{}\PYZhy{}\PYZhy{}\PYZhy{}\PYZhy{}\PYZhy{}\PYZhy{}\PYZhy{}\PYZhy{}\PYZhy{}\PYZhy{}\PYZhy{}\PYZhy{}\PYZhy{}\PYZhy{}\PYZhy{}\PYZhy{}}
\end{Verbatim}
\begin{quote}\begin{description}
\item[{Parametr}] \leavevmode
\textbf{mapping} -- slovník se zobrazením ze stávajících názvů proměnných do nových názvů.

\end{description}\end{quote}

\end{fulllineitems}


\textbf{Příklady matematických operací}

Násobení:

\begin{Verbatim}[commandchars=\\\{\}]
\PY{g+gp}{\PYZgt{}\PYZgt{}\PYZgt{} }\PY{n}{f1} \PY{o}{=} \PY{n}{Factor}\PY{p}{(}\PY{p}{[}\PY{l+s}{\PYZsq{}}\PY{l+s}{A}\PY{l+s}{\PYZsq{}}\PY{p}{,} \PY{l+s}{\PYZsq{}}\PY{l+s}{B}\PY{l+s}{\PYZsq{}}\PY{p}{]}\PY{p}{,}
\PY{g+gp}{... }    \PY{p}{\PYZob{}}\PY{l+s}{\PYZsq{}}\PY{l+s}{A}\PY{l+s}{\PYZsq{}}\PY{p}{:} \PY{p}{[}\PY{l+s}{\PYZsq{}}\PY{l+s}{a1}\PY{l+s}{\PYZsq{}}\PY{p}{,} \PY{l+s}{\PYZsq{}}\PY{l+s}{a2}\PY{l+s}{\PYZsq{}}\PY{p}{]}\PY{p}{,} \PY{l+s}{\PYZsq{}}\PY{l+s}{B}\PY{l+s}{\PYZsq{}}\PY{p}{:} \PY{p}{[}\PY{l+s}{\PYZsq{}}\PY{l+s}{b1}\PY{l+s}{\PYZsq{}}\PY{p}{,} \PY{l+s}{\PYZsq{}}\PY{l+s}{b2}\PY{l+s}{\PYZsq{}}\PY{p}{]}\PY{p}{\PYZcb{}}\PY{p}{,}
\PY{g+gp}{... }    \PY{p}{\PYZob{}}
\PY{g+gp}{... }        \PY{p}{(}\PY{l+s}{\PYZsq{}}\PY{l+s}{a1}\PY{l+s}{\PYZsq{}}\PY{p}{,} \PY{l+s}{\PYZsq{}}\PY{l+s}{b1}\PY{l+s}{\PYZsq{}}\PY{p}{)}\PY{p}{:} \PY{l+m+mf}{0.8}\PY{p}{,}
\PY{g+gp}{... }        \PY{p}{(}\PY{l+s}{\PYZsq{}}\PY{l+s}{a1}\PY{l+s}{\PYZsq{}}\PY{p}{,} \PY{l+s}{\PYZsq{}}\PY{l+s}{b2}\PY{l+s}{\PYZsq{}}\PY{p}{)}\PY{p}{:} \PY{l+m+mf}{0.2}\PY{p}{,}
\PY{g+gp}{... }        \PY{p}{(}\PY{l+s}{\PYZsq{}}\PY{l+s}{a2}\PY{l+s}{\PYZsq{}}\PY{p}{,} \PY{l+s}{\PYZsq{}}\PY{l+s}{b1}\PY{l+s}{\PYZsq{}}\PY{p}{)}\PY{p}{:} \PY{l+m+mf}{0.3}\PY{p}{,}
\PY{g+gp}{... }        \PY{p}{(}\PY{l+s}{\PYZsq{}}\PY{l+s}{a2}\PY{l+s}{\PYZsq{}}\PY{p}{,} \PY{l+s}{\PYZsq{}}\PY{l+s}{b2}\PY{l+s}{\PYZsq{}}\PY{p}{)}\PY{p}{:} \PY{l+m+mf}{0.7}
\PY{g+gp}{... }    \PY{p}{\PYZcb{}}\PY{p}{)}
\PY{g+gp}{\PYZgt{}\PYZgt{}\PYZgt{} }\PY{n}{f2} \PY{o}{=} \PY{n}{Factor}\PY{p}{(}\PY{p}{[}\PY{l+s}{\PYZsq{}}\PY{l+s}{B}\PY{l+s}{\PYZsq{}}\PY{p}{,} \PY{l+s}{\PYZsq{}}\PY{l+s}{C}\PY{l+s}{\PYZsq{}}\PY{p}{]}\PY{p}{,}
\PY{g+gp}{... }    \PY{p}{\PYZob{}}\PY{l+s}{\PYZsq{}}\PY{l+s}{B}\PY{l+s}{\PYZsq{}}\PY{p}{:} \PY{p}{[}\PY{l+s}{\PYZsq{}}\PY{l+s}{b1}\PY{l+s}{\PYZsq{}}\PY{p}{,} \PY{l+s}{\PYZsq{}}\PY{l+s}{b2}\PY{l+s}{\PYZsq{}}\PY{p}{]}\PY{p}{,} \PY{l+s}{\PYZsq{}}\PY{l+s}{C}\PY{l+s}{\PYZsq{}}\PY{p}{:} \PY{p}{[}\PY{l+s}{\PYZsq{}}\PY{l+s}{c1}\PY{l+s}{\PYZsq{}}\PY{p}{,} \PY{l+s}{\PYZsq{}}\PY{l+s}{c2}\PY{l+s}{\PYZsq{}}\PY{p}{]}\PY{p}{\PYZcb{}}\PY{p}{,}
\PY{g+gp}{... }    \PY{p}{\PYZob{}}
\PY{g+gp}{... }        \PY{p}{(}\PY{l+s}{\PYZsq{}}\PY{l+s}{b1}\PY{l+s}{\PYZsq{}}\PY{p}{,}\PY{l+s}{\PYZsq{}}\PY{l+s}{c1}\PY{l+s}{\PYZsq{}}\PY{p}{)}\PY{p}{:} \PY{l+m+mf}{0.3}\PY{p}{,}
\PY{g+gp}{... }        \PY{p}{(}\PY{l+s}{\PYZsq{}}\PY{l+s}{b1}\PY{l+s}{\PYZsq{}}\PY{p}{,}\PY{l+s}{\PYZsq{}}\PY{l+s}{c2}\PY{l+s}{\PYZsq{}}\PY{p}{)}\PY{p}{:} \PY{l+m+mf}{0.7}\PY{p}{,}
\PY{g+gp}{... }        \PY{p}{(}\PY{l+s}{\PYZsq{}}\PY{l+s}{b2}\PY{l+s}{\PYZsq{}}\PY{p}{,}\PY{l+s}{\PYZsq{}}\PY{l+s}{c1}\PY{l+s}{\PYZsq{}}\PY{p}{)}\PY{p}{:} \PY{l+m+mf}{0.5}\PY{p}{,}
\PY{g+gp}{... }        \PY{p}{(}\PY{l+s}{\PYZsq{}}\PY{l+s}{b2}\PY{l+s}{\PYZsq{}}\PY{p}{,}\PY{l+s}{\PYZsq{}}\PY{l+s}{c2}\PY{l+s}{\PYZsq{}}\PY{p}{)}\PY{p}{:} \PY{l+m+mf}{0.5}\PY{p}{,}
\PY{g+gp}{... }    \PY{p}{\PYZcb{}}\PY{p}{)}
\PY{g+gp}{\PYZgt{}\PYZgt{}\PYZgt{} }\PY{n}{fr} \PY{o}{=} \PY{n}{f1} \PY{o}{*} \PY{n}{f2}
\PY{g+gp}{\PYZgt{}\PYZgt{}\PYZgt{} }\PY{k}{print} \PY{n}{fr}\PY{o}{.}\PY{n}{pretty\PYZus{}print}\PY{p}{(}\PY{l+m+mi}{50}\PY{p}{,} \PY{l+m+mi}{2}\PY{p}{)}
\PY{g+go}{\PYZhy{}\PYZhy{}\PYZhy{}\PYZhy{}\PYZhy{}\PYZhy{}\PYZhy{}\PYZhy{}\PYZhy{}\PYZhy{}\PYZhy{}\PYZhy{}\PYZhy{}\PYZhy{}\PYZhy{}\PYZhy{}\PYZhy{}\PYZhy{}\PYZhy{}\PYZhy{}\PYZhy{}\PYZhy{}\PYZhy{}\PYZhy{}\PYZhy{}\PYZhy{}\PYZhy{}\PYZhy{}\PYZhy{}\PYZhy{}\PYZhy{}\PYZhy{}\PYZhy{}\PYZhy{}\PYZhy{}\PYZhy{}\PYZhy{}\PYZhy{}\PYZhy{}\PYZhy{}\PYZhy{}\PYZhy{}\PYZhy{}\PYZhy{}\PYZhy{}\PYZhy{}\PYZhy{}\PYZhy{}\PYZhy{}\PYZhy{}}
\PY{g+go}{    A           B           C         Value}
\PY{g+go}{\PYZhy{}\PYZhy{}\PYZhy{}\PYZhy{}\PYZhy{}\PYZhy{}\PYZhy{}\PYZhy{}\PYZhy{}\PYZhy{}\PYZhy{}\PYZhy{}\PYZhy{}\PYZhy{}\PYZhy{}\PYZhy{}\PYZhy{}\PYZhy{}\PYZhy{}\PYZhy{}\PYZhy{}\PYZhy{}\PYZhy{}\PYZhy{}\PYZhy{}\PYZhy{}\PYZhy{}\PYZhy{}\PYZhy{}\PYZhy{}\PYZhy{}\PYZhy{}\PYZhy{}\PYZhy{}\PYZhy{}\PYZhy{}\PYZhy{}\PYZhy{}\PYZhy{}\PYZhy{}\PYZhy{}\PYZhy{}\PYZhy{}\PYZhy{}\PYZhy{}\PYZhy{}\PYZhy{}\PYZhy{}\PYZhy{}\PYZhy{}}
\PY{g+go}{    a1          b1          c1         0.24}
\PY{g+go}{    a1          b1          c2         0.56}
\PY{g+go}{    a1          b2          c1         0.10}
\PY{g+go}{    a1          b2          c2         0.10}
\PY{g+go}{    a2          b1          c1         0.09}
\PY{g+go}{    a2          b1          c2         0.21}
\PY{g+go}{    a2          b2          c1         0.35}
\PY{g+go}{    a2          b2          c2         0.35}
\PY{g+go}{\PYZhy{}\PYZhy{}\PYZhy{}\PYZhy{}\PYZhy{}\PYZhy{}\PYZhy{}\PYZhy{}\PYZhy{}\PYZhy{}\PYZhy{}\PYZhy{}\PYZhy{}\PYZhy{}\PYZhy{}\PYZhy{}\PYZhy{}\PYZhy{}\PYZhy{}\PYZhy{}\PYZhy{}\PYZhy{}\PYZhy{}\PYZhy{}\PYZhy{}\PYZhy{}\PYZhy{}\PYZhy{}\PYZhy{}\PYZhy{}\PYZhy{}\PYZhy{}\PYZhy{}\PYZhy{}\PYZhy{}\PYZhy{}\PYZhy{}\PYZhy{}\PYZhy{}\PYZhy{}\PYZhy{}\PYZhy{}\PYZhy{}\PYZhy{}\PYZhy{}\PYZhy{}\PYZhy{}\PYZhy{}\PYZhy{}\PYZhy{}}
\end{Verbatim}

Umocnění:

\begin{Verbatim}[commandchars=\\\{\}]
\PY{g+gp}{\PYZgt{}\PYZgt{}\PYZgt{} }\PY{n}{f1} \PY{o}{=} \PY{n}{Factor}\PY{p}{(}\PY{p}{[}\PY{l+s}{\PYZsq{}}\PY{l+s}{A}\PY{l+s}{\PYZsq{}}\PY{p}{,} \PY{l+s}{\PYZsq{}}\PY{l+s}{B}\PY{l+s}{\PYZsq{}}\PY{p}{]}\PY{p}{,}
\PY{g+gp}{... }    \PY{p}{\PYZob{}}\PY{l+s}{\PYZsq{}}\PY{l+s}{A}\PY{l+s}{\PYZsq{}}\PY{p}{:} \PY{p}{[}\PY{l+s}{\PYZsq{}}\PY{l+s}{a1}\PY{l+s}{\PYZsq{}}\PY{p}{,} \PY{l+s}{\PYZsq{}}\PY{l+s}{a2}\PY{l+s}{\PYZsq{}}\PY{p}{]}\PY{p}{,} \PY{l+s}{\PYZsq{}}\PY{l+s}{B}\PY{l+s}{\PYZsq{}}\PY{p}{:} \PY{p}{[}\PY{l+s}{\PYZsq{}}\PY{l+s}{b1}\PY{l+s}{\PYZsq{}}\PY{p}{,} \PY{l+s}{\PYZsq{}}\PY{l+s}{b2}\PY{l+s}{\PYZsq{}}\PY{p}{]}\PY{p}{\PYZcb{}}\PY{p}{,}
\PY{g+gp}{... }    \PY{p}{\PYZob{}}
\PY{g+gp}{... }        \PY{p}{(}\PY{l+s}{\PYZsq{}}\PY{l+s}{a1}\PY{l+s}{\PYZsq{}}\PY{p}{,} \PY{l+s}{\PYZsq{}}\PY{l+s}{b1}\PY{l+s}{\PYZsq{}}\PY{p}{)}\PY{p}{:} \PY{l+m+mf}{0.8}\PY{p}{,}
\PY{g+gp}{... }        \PY{p}{(}\PY{l+s}{\PYZsq{}}\PY{l+s}{a1}\PY{l+s}{\PYZsq{}}\PY{p}{,} \PY{l+s}{\PYZsq{}}\PY{l+s}{b2}\PY{l+s}{\PYZsq{}}\PY{p}{)}\PY{p}{:} \PY{l+m+mf}{0.2}\PY{p}{,}
\PY{g+gp}{... }        \PY{p}{(}\PY{l+s}{\PYZsq{}}\PY{l+s}{a2}\PY{l+s}{\PYZsq{}}\PY{p}{,} \PY{l+s}{\PYZsq{}}\PY{l+s}{b1}\PY{l+s}{\PYZsq{}}\PY{p}{)}\PY{p}{:} \PY{l+m+mf}{0.3}\PY{p}{,}
\PY{g+gp}{... }        \PY{p}{(}\PY{l+s}{\PYZsq{}}\PY{l+s}{a2}\PY{l+s}{\PYZsq{}}\PY{p}{,} \PY{l+s}{\PYZsq{}}\PY{l+s}{b2}\PY{l+s}{\PYZsq{}}\PY{p}{)}\PY{p}{:} \PY{l+m+mf}{0.7}
\PY{g+gp}{... }    \PY{p}{\PYZcb{}}\PY{p}{)}
\PY{g+gp}{\PYZgt{}\PYZgt{}\PYZgt{} }\PY{n}{fr} \PY{o}{=} \PY{n}{f1} \PY{o}{*}\PY{o}{*} \PY{l+m+mi}{2}
\PY{g+gp}{\PYZgt{}\PYZgt{}\PYZgt{} }\PY{k}{print} \PY{n}{fr}\PY{o}{.}\PY{n}{pretty\PYZus{}print}\PY{p}{(}\PY{l+m+mi}{50}\PY{p}{,} \PY{l+m+mi}{2}\PY{p}{)}
\PY{g+go}{\PYZhy{}\PYZhy{}\PYZhy{}\PYZhy{}\PYZhy{}\PYZhy{}\PYZhy{}\PYZhy{}\PYZhy{}\PYZhy{}\PYZhy{}\PYZhy{}\PYZhy{}\PYZhy{}\PYZhy{}\PYZhy{}\PYZhy{}\PYZhy{}\PYZhy{}\PYZhy{}\PYZhy{}\PYZhy{}\PYZhy{}\PYZhy{}\PYZhy{}\PYZhy{}\PYZhy{}\PYZhy{}\PYZhy{}\PYZhy{}\PYZhy{}\PYZhy{}\PYZhy{}\PYZhy{}\PYZhy{}\PYZhy{}\PYZhy{}\PYZhy{}\PYZhy{}\PYZhy{}\PYZhy{}\PYZhy{}\PYZhy{}\PYZhy{}\PYZhy{}\PYZhy{}\PYZhy{}\PYZhy{}\PYZhy{}\PYZhy{}}
\PY{g+go}{    A               B             Value}
\PY{g+go}{\PYZhy{}\PYZhy{}\PYZhy{}\PYZhy{}\PYZhy{}\PYZhy{}\PYZhy{}\PYZhy{}\PYZhy{}\PYZhy{}\PYZhy{}\PYZhy{}\PYZhy{}\PYZhy{}\PYZhy{}\PYZhy{}\PYZhy{}\PYZhy{}\PYZhy{}\PYZhy{}\PYZhy{}\PYZhy{}\PYZhy{}\PYZhy{}\PYZhy{}\PYZhy{}\PYZhy{}\PYZhy{}\PYZhy{}\PYZhy{}\PYZhy{}\PYZhy{}\PYZhy{}\PYZhy{}\PYZhy{}\PYZhy{}\PYZhy{}\PYZhy{}\PYZhy{}\PYZhy{}\PYZhy{}\PYZhy{}\PYZhy{}\PYZhy{}\PYZhy{}\PYZhy{}\PYZhy{}\PYZhy{}\PYZhy{}\PYZhy{}}
\PY{g+go}{    a1              b1             0.64}
\PY{g+go}{    a1              b2             0.04}
\PY{g+go}{    a2              b1             0.09}
\PY{g+go}{    a2              b2             0.49}
\PY{g+go}{\PYZhy{}\PYZhy{}\PYZhy{}\PYZhy{}\PYZhy{}\PYZhy{}\PYZhy{}\PYZhy{}\PYZhy{}\PYZhy{}\PYZhy{}\PYZhy{}\PYZhy{}\PYZhy{}\PYZhy{}\PYZhy{}\PYZhy{}\PYZhy{}\PYZhy{}\PYZhy{}\PYZhy{}\PYZhy{}\PYZhy{}\PYZhy{}\PYZhy{}\PYZhy{}\PYZhy{}\PYZhy{}\PYZhy{}\PYZhy{}\PYZhy{}\PYZhy{}\PYZhy{}\PYZhy{}\PYZhy{}\PYZhy{}\PYZhy{}\PYZhy{}\PYZhy{}\PYZhy{}\PYZhy{}\PYZhy{}\PYZhy{}\PYZhy{}\PYZhy{}\PYZhy{}\PYZhy{}\PYZhy{}\PYZhy{}\PYZhy{}}
\end{Verbatim}

\end{fulllineitems}



\subsection{Modul \texttt{node}}
\label{alex.infer:modul-node}
Modul obsahuje implementace vrcholů tvořících grafický model.
Pro výpočty jsou používány třídy z modulu {\hyperref[alex.infer:module-alex.infer.factor]{\code{factor}}}.
\index{alex.infer.node.Node() (třída v alex.infer.factor)}

\hskip\linewidth
\begin{fulllineitems}
\phantomsection\label{alex.infer:alex.infer.factor.alex.infer.node.Node}\pysigline{\strong{class }\code{alex.infer.node.}\bfcode{Node}}
Abstraktní třída pro reprezentaci vrcholů.

Obsahuje definice základních metod, které musí obsahovat každá implementace vrcholu.
\index{\_\_init\_\_() (metoda alex.infer.factor.alex.infer.node.Node)}

\begin{fulllineitems}
\phantomsection\label{alex.infer:alex.infer.factor.alex.infer.node.Node.__init__}\pysiglinewithargsret{\bfcode{\_\_init\_\_}}{\emph{name}\optional{, \emph{aliases}}}{}
Vytvoření nového vrcholu.

Při vytvoření vrcholu je třeba zadat jeho jméno, pod kterým bude možné
jej nalézt při inferenci a je možné předat aliasy pro proměnné. Aliasy pro
proměnné slouží k přejmenování proměnných ve zprávách, které přicházejí
do vrcholu.
\begin{quote}\begin{description}
\item[{Parametry}] \leavevmode\begin{itemize}
\item {} 
\textbf{name} -- jméno vrcholu.

\item {} 
\textbf{aliases} -- zobrazení ze jmen proměnných, které se mohou objevit ve zprávě, do proměnných v parametru.

\end{itemize}

\end{description}\end{quote}

\end{fulllineitems}

\index{add\_neighbor() (metoda alex.infer.factor.alex.infer.node.Node)}

\begin{fulllineitems}
\phantomsection\label{alex.infer:alex.infer.factor.alex.infer.node.Node.add_neighbor}\pysiglinewithargsret{\bfcode{add\_neighbor}}{\emph{node}, \emph{**options}}{}
Přidání souseda.

Pro většinu vrcholů bude stačit pouze vrchol, zůstává tu ovšem možnost přidat implementačně závislé parametry.
\begin{quote}\begin{description}
\item[{Parametry}] \leavevmode\begin{itemize}
\item {} 
\textbf{node} -- sousední vrchol.

\item {} 
\textbf{options} -- implementačně závislé parametry.

\end{itemize}

\end{description}\end{quote}

\end{fulllineitems}

\index{connect() (metoda alex.infer.factor.alex.infer.node.Node)}

\begin{fulllineitems}
\phantomsection\label{alex.infer:alex.infer.factor.alex.infer.node.Node.connect}\pysiglinewithargsret{\bfcode{connect}}{\emph{node}, \emph{**options}}{}
Propojení dvou vrcholů.

Propojení dvou vrcholů, stačí zavolat pouze jednou pro dvojici vrcholů.
Interně zavolá metodu {\hyperref[alex.infer:alex.infer.factor.alex.infer.node.Node.add_neighbor]{\code{add\_neighbor()}}} a předá jí parametry.
\begin{quote}\begin{description}
\item[{Parametry}] \leavevmode\begin{itemize}
\item {} 
\textbf{node} -- sousední vrchol.

\item {} 
\textbf{options} -- implementačně závislé parametry.

\end{itemize}

\end{description}\end{quote}

\end{fulllineitems}

\index{init\_messages() (metoda alex.infer.factor.alex.infer.node.Node)}

\begin{fulllineitems}
\phantomsection\label{alex.infer:alex.infer.factor.alex.infer.node.Node.init_messages}\pysiglinewithargsret{\bfcode{init\_messages}}{}{}
Inicializace všech zpráv mezi vrcholem a jeho sousedy.

\end{fulllineitems}

\index{message\_from() (metoda alex.infer.factor.alex.infer.node.Node)}

\begin{fulllineitems}
\phantomsection\label{alex.infer:alex.infer.factor.alex.infer.node.Node.message_from}\pysiglinewithargsret{\bfcode{message\_from}}{\emph{node}, \emph{message}}{}
Přijmutí zprávy ze sousedního vrcholu.

Zpráva je uložena a použita při dalších výpočtech.
\begin{quote}\begin{description}
\item[{Parametry}] \leavevmode\begin{itemize}
\item {} 
\textbf{node} -- sousední vrchol.

\item {} 
\textbf{message} -- zpráva ze sousedního vrcholu.

\end{itemize}

\end{description}\end{quote}

\end{fulllineitems}

\index{message\_to() (metoda alex.infer.factor.alex.infer.node.Node)}

\begin{fulllineitems}
\phantomsection\label{alex.infer:alex.infer.factor.alex.infer.node.Node.message_to}\pysiglinewithargsret{\bfcode{message\_to}}{\emph{node}}{}
Poslání zprávy do sousedního vrcholu.
\begin{quote}\begin{description}
\item[{Parametr}] \leavevmode
\textbf{node} -- sousední vrchol, kterému má být zpráva zaslána.

\end{description}\end{quote}

\end{fulllineitems}

\index{normalize() (metoda alex.infer.factor.alex.infer.node.Node)}

\begin{fulllineitems}
\phantomsection\label{alex.infer:alex.infer.factor.alex.infer.node.Node.normalize}\pysiglinewithargsret{\bfcode{normalize}}{\optional{\emph{parents}}}{}
Normalizace pravděpodobnostního rozdělení.

Při počítání nejsou zprávy normalizované, proto ani výsledek nemusí být
normalizovaný a je tedy třeba jej normalizovat před jakýmkoliv dalším
použitím. Pro normalizaci je možné zadat, které proměnné jsou rodiče,
pak musí být suma ohodnocení všech přiřazení se stejnými rodiči rovna
jedné.
\begin{quote}\begin{description}
\item[{Parametr}] \leavevmode
\textbf{parents} -- seznam proměnných, které jsou rodiči vrcholu.

\end{description}\end{quote}

\end{fulllineitems}

\index{update() (metoda alex.infer.factor.alex.infer.node.Node)}

\begin{fulllineitems}
\phantomsection\label{alex.infer:alex.infer.factor.alex.infer.node.Node.update}\pysiglinewithargsret{\bfcode{update}}{}{}
Aktualizace interní pravděpodobnostní distribuce.

Příchozí zprávy jsou pouze uloženy a je potřeba aktualizovat interní
reprezentaci jejich součinu před počítáním odchozích zpráv.

\end{fulllineitems}

\end{fulllineitems}
\hskip\linewidth
\index{alex.infer.node.FactorNode() (třída v alex.infer.factor)}

\begin{fulllineitems}
\phantomsection\label{alex.infer:alex.infer.factor.alex.infer.node.FactorNode}\pysigline{\strong{class }\code{alex.infer.node.}\bfcode{FactorNode}}
Předek: {\hyperref[alex.infer:alex.infer.factor.alex.infer.node.Node]{\code{alex.infer.node.Node}}}

Abstraktní třída reprezentující vrcholy pro faktory.
\end{fulllineitems}

\hskip\linewidth

\index{alex.infer.node.VariableNode() (třída v alex.infer.factor)}

\begin{fulllineitems}
\phantomsection\label{alex.infer:alex.infer.factor.alex.infer.node.VariableNode}\pysigline{\strong{class }\code{alex.infer.node.}\bfcode{VariableNode}}
Předek: {\hyperref[alex.infer:alex.infer.factor.alex.infer.node.Node]{\code{alex.infer.node.Node}}}

Abstraktní třída reprezentující vrcholy pro proměnné.

\end{fulllineitems}



\subsubsection{Vrcholy pro diskrétní vrcholy}
\label{alex.infer:vrcholy-pro-diskretni-vrcholy}\index{alex.infer.node.DiscreteVariableNode() (třída v alex.infer.factor)}

\begin{fulllineitems}
\phantomsection\label{alex.infer:alex.infer.factor.alex.infer.node.DiscreteVariableNode}\pysigline{\strong{class }\code{alex.infer.node.}\bfcode{DiscreteVariableNode}}
Předek: {\hyperref[alex.infer:alex.infer.factor.alex.infer.node.VariableNode]{\code{alex.infer.node.VariableNode}}}

Reprezentace diskrétních proměnných.
\index{belief() (atribut alex.infer.factor.alex.infer.node.DiscreteVariableNode)}

\begin{fulllineitems}
\phantomsection\label{alex.infer:alex.infer.factor.alex.infer.node.DiscreteVariableNode.belief}\pysigline{\bfcode{belief}}
Marginální pravděpodobnostní rozdělení pro proměnnou. Typu {\hyperref[alex.infer:alex.infer.factor.Factor]{\code{Factor}}}.

\end{fulllineitems}

\index{\_\_init\_\_() (metoda alex.infer.factor.alex.infer.node.DiscreteVariableNode)}

\begin{fulllineitems}
\phantomsection\label{alex.infer:alex.infer.factor.alex.infer.node.DiscreteVariableNode.__init__}\pysiglinewithargsret{\bfcode{\_\_init\_\_}}{\emph{name}, \emph{values}, \emph{logarithmetic=True}}{}
Vytvoření nového vrcholu pro diskrétní proměnné.

Pro vytvoření je třeba jméno proměnné a hodnoty, kterých může nabývat.

Také je možné zadat, jestli používat při výpočtech trik s aritmetikou nad zlogaritmovanými čísly.
Je třeba, aby všechny vrcholy v grafickém modelu byly nastaveny stejně.
\begin{quote}\begin{description}
\item[{Parametry}] \leavevmode\begin{itemize}
\item {} 
\textbf{name} -- název proměnné.

\item {} 
\textbf{values} -- seznam možných hodnot.

\item {} 
\textbf{logarithmetic} -- zda-li používat trik s logaritmy.

\end{itemize}

\end{description}\end{quote}

\end{fulllineitems}

\index{add\_neighbor() (metoda alex.infer.factor.alex.infer.node.DiscreteVariableNode)}

\begin{fulllineitems}
\phantomsection\label{alex.infer:alex.infer.factor.alex.infer.node.DiscreteVariableNode.add_neighbor}\pysiglinewithargsret{\bfcode{add\_neighbor}}{\emph{node}}{}
Přidání souseda.
\begin{quote}\begin{description}
\item[{Parametr}] \leavevmode
\textbf{node} -- sousední vrchol.

\end{description}\end{quote}

\end{fulllineitems}

\index{connect() (metoda alex.infer.factor.alex.infer.node.DiscreteVariableNode)}

\begin{fulllineitems}
\phantomsection\label{alex.infer:alex.infer.factor.alex.infer.node.DiscreteVariableNode.connect}\pysiglinewithargsret{\bfcode{connect}}{\emph{node}, \emph{**options}}{}
Propojení dvou vrcholů.

Propojení dvou vrcholů, stačí zavolat pouze jednou pro dvojici vrcholů.
Interně zavolá metodu {\hyperref[alex.infer:alex.infer.factor.alex.infer.node.DiscreteVariableNode.add_neighbor]{\code{add\_neighbor()}}} a předá jí parametry.
\begin{quote}\begin{description}
\item[{Parametr}] \leavevmode
\textbf{node} -- sousední vrchol.

\end{description}\end{quote}

\end{fulllineitems}

\index{init\_messages() (metoda alex.infer.factor.alex.infer.node.DiscreteVariableNode)}

\begin{fulllineitems}
\phantomsection\label{alex.infer:alex.infer.factor.alex.infer.node.DiscreteVariableNode.init_messages}\pysiglinewithargsret{\bfcode{init\_messages}}{}{}
Inicializace všech zpráv mezi vrcholem a jeho sousedy.

\end{fulllineitems}

\index{message\_from() (metoda alex.infer.factor.alex.infer.node.DiscreteVariableNode)}

\begin{fulllineitems}
\phantomsection\label{alex.infer:alex.infer.factor.alex.infer.node.DiscreteVariableNode.message_from}\pysiglinewithargsret{\bfcode{message\_from}}{\emph{node}, \emph{message}}{}
Přijmutí zprávy ze sousedního vrcholu.

Zpráva je uložena a použita při dalších výpočtech.
\begin{quote}\begin{description}
\item[{Parametry}] \leavevmode\begin{itemize}
\item {} 
\textbf{node} -- sousední vrchol.

\item {} 
\textbf{message} -- zpráva ze sousedního vrcholu.

\end{itemize}

\end{description}\end{quote}

\end{fulllineitems}

\index{message\_to() (metoda alex.infer.factor.alex.infer.node.DiscreteVariableNode)}

\begin{fulllineitems}
\phantomsection\label{alex.infer:alex.infer.factor.alex.infer.node.DiscreteVariableNode.message_to}\pysiglinewithargsret{\bfcode{message\_to}}{\emph{node}}{}
Poslání zprávy do sousedního vrcholu.
\begin{quote}\begin{description}
\item[{Parametr}] \leavevmode
\textbf{node} -- sousední vrchol, kterému má být zpráva zaslána.

\end{description}\end{quote}

\end{fulllineitems}

\index{normalize() (metoda alex.infer.factor.alex.infer.node.DiscreteVariableNode)}

\begin{fulllineitems}
\phantomsection\label{alex.infer:alex.infer.factor.alex.infer.node.DiscreteVariableNode.normalize}\pysiglinewithargsret{\bfcode{normalize}}{}{}
Normalizace pravděpodobnostního rozdělení.

Při počítání nejsou zprávy normalizované, proto ani výsledek nemusí být
normalizovaný a je tedy jej třeba normalizovat před jakýmkoliv dalším
použitím.

\end{fulllineitems}

\index{update() (metoda alex.infer.factor.alex.infer.node.DiscreteVariableNode)}

\begin{fulllineitems}
\phantomsection\label{alex.infer:alex.infer.factor.alex.infer.node.DiscreteVariableNode.update}\pysiglinewithargsret{\bfcode{update}}{}{}
Aktualizace marginální pravděpodobnostní distribuce.

Příchozí zprávy jsou pouze uloženy a je potřeba aktualizovat interní
reprezentaci jejich součinu před počítáním odchozích zpráv.

\end{fulllineitems}

\index{observed() (metoda alex.infer.factor.alex.infer.node.DiscreteVariableNode)}

\begin{fulllineitems}
\phantomsection\label{alex.infer:alex.infer.factor.alex.infer.node.DiscreteVariableNode.observed}\pysiglinewithargsret{\bfcode{observed}}{\emph{assignment\_dict}}{}
Nastavení pozorovaných hodnot.

Proměnná je nastavená na pozorované rozdělení. Pravděpodobnost nepozorovaných hodnot
je nulová.
\begin{quote}\begin{description}
\item[{Parametr}] \leavevmode
\textbf{assignment\_dict} -- slovník se zobrazením z hodnot do jejich pravděpodobností.

\end{description}\end{quote}

\end{fulllineitems}

\index{most\_probable() (metoda alex.infer.factor.alex.infer.node.DiscreteVariableNode)}

\begin{fulllineitems}
\phantomsection\label{alex.infer:alex.infer.factor.alex.infer.node.DiscreteVariableNode.most_probable}\pysiglinewithargsret{\bfcode{most\_probable}}{\emph{n}}{}
Získání nejpravděpodobnějších hodnot.
\begin{quote}\begin{description}
\item[{Parametr}] \leavevmode
\textbf{n} -- počet nejpravděpodobnějších hodnot, které mají být vráceny.

\item[{Vrací}] \leavevmode
seznam dvojic ve tvaru (hodnota, její pravděpodobnost) seřazený podle pravděpodobnosti.

\end{description}\end{quote}

\end{fulllineitems}


\end{fulllineitems}

\hskip\linewidth
\index{alex.infer.node.DiscreteFactorNode() (třída v alex.infer.factor)}

\begin{fulllineitems}
\phantomsection\label{alex.infer:alex.infer.factor.alex.infer.node.DiscreteFactorNode}\pysigline{\strong{class }\code{alex.infer.node.}\bfcode{DiscreteFactorNode}}
Předek: {\hyperref[alex.infer:alex.infer.factor.alex.infer.node.FactorNode]{\code{alex.infer.node.FactorNode}}}

Vrchol pro diskrétní faktor, reprezentující diskrétní pravděpodobnostní rozdělení.
\index{\_\_init\_\_() (metoda alex.infer.factor.alex.infer.node.DiscreteFactorNode)}

\begin{fulllineitems}
\phantomsection\label{alex.infer:alex.infer.factor.alex.infer.node.DiscreteFactorNode.__init__}\pysiglinewithargsret{\bfcode{\_\_init\_\_}}{\emph{name}, \emph{factor}}{}
Vytvoření nového vrcholu pro diskrétní faktor.

Pro vytvoření je třeba jméno faktoru a samotný faktor.

Také je možné zadat, jestli používat při výpočtech trik s aritmetikou nad zlogaritmovanými čísly.
Je třeba, aby všechny vrcholy v grafickém modelu byly nastaveny stejně.
\begin{quote}\begin{description}
\item[{Parametry}] \leavevmode\begin{itemize}
\item {} 
\textbf{name} -- název proměnné.

\item {} 
\textbf{factor} ({\hyperref[alex.infer:alex.infer.factor.Factor]{\code{Factor}}}) -- faktor reprezentující diskrétní pravděpodobnostní rozdělení.

\end{itemize}

\end{description}\end{quote}

\end{fulllineitems}

\index{add\_neighbor() (metoda alex.infer.factor.alex.infer.node.DiscreteFactorNode)}

\begin{fulllineitems}
\phantomsection\label{alex.infer:alex.infer.factor.alex.infer.node.DiscreteFactorNode.add_neighbor}\pysiglinewithargsret{\bfcode{add\_neighbor}}{\emph{node}}{}
Přidání souseda.
\begin{quote}\begin{description}
\item[{Parametr}] \leavevmode
\textbf{node} -- sousední vrchol.

\end{description}\end{quote}

\end{fulllineitems}

\index{connect() (metoda alex.infer.factor.alex.infer.node.DiscreteFactorNode)}

\begin{fulllineitems}
\phantomsection\label{alex.infer:alex.infer.factor.alex.infer.node.DiscreteFactorNode.connect}\pysiglinewithargsret{\bfcode{connect}}{\emph{node}}{}
Propojení dvou vrcholů.

Propojení dvou vrcholů, stačí zavolat pouze jednou pro dvojici vrcholů.
Interně zavolá metodu {\hyperref[alex.infer:alex.infer.factor.alex.infer.node.DiscreteFactorNode.add_neighbor]{\code{add\_neighbor()}}} obou vrcholů a předá jí parametry.
\begin{quote}\begin{description}
\item[{Parametr}] \leavevmode
\textbf{node} -- sousední vrchol.

\end{description}\end{quote}

\end{fulllineitems}

\index{init\_messages() (metoda alex.infer.factor.alex.infer.node.DiscreteFactorNode)}

\begin{fulllineitems}
\phantomsection\label{alex.infer:alex.infer.factor.alex.infer.node.DiscreteFactorNode.init_messages}\pysiglinewithargsret{\bfcode{init\_messages}}{}{}
Inicializace všech zpráv mezi vrcholem a jeho sousedy.

\end{fulllineitems}

\index{message\_from() (metoda alex.infer.factor.alex.infer.node.DiscreteFactorNode)}

\begin{fulllineitems}
\phantomsection\label{alex.infer:alex.infer.factor.alex.infer.node.DiscreteFactorNode.message_from}\pysiglinewithargsret{\bfcode{message\_from}}{\emph{node}, \emph{message}}{}
Přijmutí zprávy ze sousedního vrcholu.

Zpráva je uložena a použita při dalších výpočtech.
\begin{quote}\begin{description}
\item[{Parametry}] \leavevmode\begin{itemize}
\item {} 
\textbf{node} -- sousední vrchol.

\item {} 
\textbf{message} -- zpráva ze sousedního vrcholu.

\end{itemize}

\end{description}\end{quote}

\end{fulllineitems}

\index{message\_to() (metoda alex.infer.factor.alex.infer.node.DiscreteFactorNode)}

\begin{fulllineitems}
\phantomsection\label{alex.infer:alex.infer.factor.alex.infer.node.DiscreteFactorNode.message_to}\pysiglinewithargsret{\bfcode{message\_to}}{\emph{node}}{}
Poslání zprávy do sousedního vrcholu.
\begin{quote}\begin{description}
\item[{Parametr}] \leavevmode
\textbf{node} -- sousední vrchol, kterému má být zpráva zaslána.

\end{description}\end{quote}

\end{fulllineitems}

\index{normalize() (metoda alex.infer.factor.alex.infer.node.DiscreteFactorNode)}

\begin{fulllineitems}
\phantomsection\label{alex.infer:alex.infer.factor.alex.infer.node.DiscreteFactorNode.normalize}\pysiglinewithargsret{\bfcode{normalize}}{\optional{\emph{parents}}}{}
Normalizace pravděpodobnostního rozdělení.

Při počítání nejsou zprávy normalizované, proto ani výsledek nemusí být
normalizovaný a je tedy jej třeba normalizovat před jakýmkoliv dalším
použitím. Pro normalizaci je možné zadat, které proměnné jsou rodiče,
pak musí být suma ohodnocení všech přiřazení se stejnými rodiči rovna
jedné.
\begin{quote}\begin{description}
\item[{Parametr}] \leavevmode
\textbf{parents} -- seznam proměnných, které jsou rodiči vrcholu.

\end{description}\end{quote}

\end{fulllineitems}

\index{update() (metoda alex.infer.factor.alex.infer.node.DiscreteFactorNode)}

\begin{fulllineitems}
\phantomsection\label{alex.infer:alex.infer.factor.alex.infer.node.DiscreteFactorNode.update}\pysiglinewithargsret{\bfcode{update}}{}{}
Aktualizace interní reprezentace.

Příchozí zprávy jsou pouze uloženy a je potřeba aktualizovat interní
reprezentaci jejich součinu před počítáním odchozích zpráv.

\end{fulllineitems}


\end{fulllineitems}



\subsubsection{Vrcholy pro Dirichletovské parametry}
\label{alex.infer:vrcholy-pro-dirichletovske-parametry}\index{alex.infer.node.DirichletParameterNode() (třída v alex.infer.factor)}

\begin{fulllineitems}
\phantomsection\label{alex.infer:alex.infer.factor.alex.infer.node.DirichletParameterNode}\pysigline{\strong{class }\code{alex.infer.node.}\bfcode{DirichletParameterNode}}
Předek: {\hyperref[alex.infer:alex.infer.factor.alex.infer.node.VariableNode]{\code{alex.infer.node.VariableNode}}}

Reprezentace parametrů Dirichletovského rozdělení.
\index{alpha() (atribut alex.infer.factor.alex.infer.node.DirichletParameterNode)}

\begin{fulllineitems}
\phantomsection\label{alex.infer:alex.infer.factor.alex.infer.node.DirichletParameterNode.alpha}\pysigline{\bfcode{alpha}}
Parametry dirichletovského rozdělení. Typu {\hyperref[alex.infer:alex.infer.factor.Factor]{\code{Factor}}}.

\end{fulllineitems}

\index{\_\_init\_\_() (metoda alex.infer.factor.alex.infer.node.DirichletParameterNode)}

\begin{fulllineitems}
\phantomsection\label{alex.infer:alex.infer.factor.alex.infer.node.DirichletParameterNode.__init__}\pysiglinewithargsret{\bfcode{\_\_init\_\_}}{\emph{name}, \emph{alpha}\optional{, \emph{aliases}}}{}
Vytvoření nového vrcholu pro parametry Dirichletovského rozdělení.

Pro vytvoření je třeba jméno proměnné a parametry.
Aliasy pro proměnné slouží k přejmenování proměnných ve zprávách, které
přicházejí do vrcholu.
\begin{quote}\begin{description}
\item[{Parametry}] \leavevmode\begin{itemize}
\item {} 
\textbf{name} -- název proměnné.

\item {} 
\textbf{alpha} ({\hyperref[alex.infer:alex.infer.factor.Factor]{\code{Factor}}}) -- parametry Dirichletovského rozdělení.

\item {} 
\textbf{aliases} -- zobrazení ze jmen proměnných, které se mohou objevit ve zprávě, do proměnných v parametru.

\end{itemize}

\end{description}\end{quote}

\end{fulllineitems}

\index{add\_neighbor() (metoda alex.infer.factor.alex.infer.node.DirichletParameterNode)}

\begin{fulllineitems}
\phantomsection\label{alex.infer:alex.infer.factor.alex.infer.node.DirichletParameterNode.add_neighbor}\pysiglinewithargsret{\bfcode{add\_neighbor}}{\emph{node}}{}
Přidání souseda.
\begin{quote}\begin{description}
\item[{Parametr}] \leavevmode
\textbf{node} ({\hyperref[alex.infer:alex.infer.factor.alex.infer.node.DirichletFactorNode]{\code{DirichletFactorNode}}}) -- sousední vrchol.

\end{description}\end{quote}

\end{fulllineitems}

\index{connect() (metoda alex.infer.factor.alex.infer.node.DirichletParameterNode)}

\begin{fulllineitems}
\phantomsection\label{alex.infer:alex.infer.factor.alex.infer.node.DirichletParameterNode.connect}\pysiglinewithargsret{\bfcode{connect}}{\emph{node}, \emph{**options}}{}
Propojení dvou vrcholů.

Propojení dvou vrcholů, stačí zavolat pouze jednou pro dvojici vrcholů.
Interně zavolá metodu {\hyperref[alex.infer:alex.infer.factor.alex.infer.node.DirichletParameterNode.add_neighbor]{\code{add\_neighbor()}}} a předá jí parametry.
\begin{quote}\begin{description}
\item[{Parametr}] \leavevmode
\textbf{node} ({\hyperref[alex.infer:alex.infer.factor.alex.infer.node.DirichletFactorNode]{\code{DirichletFactorNode}}}) -- sousední vrchol.

\end{description}\end{quote}

\end{fulllineitems}

\index{init\_messages() (metoda alex.infer.factor.alex.infer.node.DirichletParameterNode)}

\begin{fulllineitems}
\phantomsection\label{alex.infer:alex.infer.factor.alex.infer.node.DirichletParameterNode.init_messages}\pysiglinewithargsret{\bfcode{init\_messages}}{}{}
Inicializace všech zpráv mezi vrcholem a jeho sousedy.

\end{fulllineitems}

\index{message\_from() (metoda alex.infer.factor.alex.infer.node.DirichletParameterNode)}

\begin{fulllineitems}
\phantomsection\label{alex.infer:alex.infer.factor.alex.infer.node.DirichletParameterNode.message_from}\pysiglinewithargsret{\bfcode{message\_from}}{\emph{node}, \emph{message}}{}
Přijmutí zprávy ze sousedního vrcholu.

Zpráva je uložena a použita při dalších výpočtech.
\begin{quote}\begin{description}
\item[{Parametry}] \leavevmode\begin{itemize}
\item {} 
\textbf{node} -- sousední vrchol.

\item {} 
\textbf{message} -- zpráva ze sousedního vrcholu.

\end{itemize}

\end{description}\end{quote}

\end{fulllineitems}

\index{message\_to() (metoda alex.infer.factor.alex.infer.node.DirichletParameterNode)}

\begin{fulllineitems}
\phantomsection\label{alex.infer:alex.infer.factor.alex.infer.node.DirichletParameterNode.message_to}\pysiglinewithargsret{\bfcode{message\_to}}{\emph{node}}{}
Poslání zprávy do sousedního vrcholu.
\begin{quote}\begin{description}
\item[{Parametr}] \leavevmode
\textbf{node} -- sousední vrchol, kterému má být zpráva zaslána.

\end{description}\end{quote}

\end{fulllineitems}


\end{fulllineitems}

\index{alex.infer.node.DirichletFactorNode() (třída v alex.infer.factor)}
\hskip\linewidth
\begin{fulllineitems}
\phantomsection\label{alex.infer:alex.infer.factor.alex.infer.node.DirichletFactorNode}\pysigline{\strong{class }\code{alex.infer.node.}\bfcode{DirichletFactorNode}}
Předek: {\hyperref[alex.infer:alex.infer.factor.alex.infer.node.FactorNode]{\code{alex.infer.node.FactorNode}}}

    Diskrétní faktor, parametry jeho rozdělení ovšem nejsou pevně dány, ale pochází z vrcholu typu
{\hyperref[alex.infer:alex.infer.factor.alex.infer.node.DirichletParameterNode]{\code{DirichletParameterNode}}}.
\index{\_\_init\_\_() (metoda alex.infer.factor.alex.infer.node.DirichletFactorNode)}

\begin{fulllineitems}
\phantomsection\label{alex.infer:alex.infer.factor.alex.infer.node.DirichletFactorNode.__init__}\pysiglinewithargsret{\bfcode{\_\_init\_\_}}{\emph{name}\optional{, \emph{aliases}}}{}
Vytvoření nového vrcholu pro faktor.

Pro vytvoření je třeba jméno faktoru.

Aliasy pro proměnné slouží k přejmenování proměnných ve zprávách, které
přicházejí do vrcholu.
\begin{quote}\begin{description}
\item[{Parametry}] \leavevmode\begin{itemize}
\item {} 
\textbf{name} -- název faktoru.

\item {} 
\textbf{aliases} -- zobrazení ze jmen proměnných, které se mohou objevit ve zprávě do proměnných ve faktoru.

\end{itemize}

\end{description}\end{quote}

\end{fulllineitems}

\index{add\_neighbor() (metoda alex.infer.factor.alex.infer.node.DirichletFactorNode)}

\begin{fulllineitems}
\phantomsection\label{alex.infer:alex.infer.factor.alex.infer.node.DirichletFactorNode.add_neighbor}\pysiglinewithargsret{\bfcode{add\_neighbor}}{\emph{node}}{}
Přidání souseda.
\begin{quote}\begin{description}
\item[{Parametr}] \leavevmode
\textbf{node} -- sousední vrchol.

\end{description}\end{quote}

\end{fulllineitems}

\index{connect() (metoda alex.infer.factor.alex.infer.node.DirichletFactorNode)}

\begin{fulllineitems}
\phantomsection\label{alex.infer:alex.infer.factor.alex.infer.node.DirichletFactorNode.connect}\pysiglinewithargsret{\bfcode{connect}}{\emph{node}}{}
Propojení dvou vrcholů.

Propojení dvou vrcholů, stačí zavolat pouze jednou pro dvojici vrcholů.
Interně zavolá metodu {\hyperref[alex.infer:alex.infer.factor.alex.infer.node.DirichletFactorNode.add_neighbor]{\code{add\_neighbor()}}} obou vrcholů a předá jí parametry.
\begin{quote}\begin{description}
\item[{Parametr}] \leavevmode
\textbf{node} -- sousední vrchol.

\end{description}\end{quote}

\end{fulllineitems}

\index{init\_messages() (metoda alex.infer.factor.alex.infer.node.DirichletFactorNode)}

\begin{fulllineitems}
\phantomsection\label{alex.infer:alex.infer.factor.alex.infer.node.DirichletFactorNode.init_messages}\pysiglinewithargsret{\bfcode{init\_messages}}{}{}
Inicializace všech zpráv mezi vrcholem a jeho sousedy.

\end{fulllineitems}

\index{message\_from() (metoda alex.infer.factor.alex.infer.node.DirichletFactorNode)}

\begin{fulllineitems}
\phantomsection\label{alex.infer:alex.infer.factor.alex.infer.node.DirichletFactorNode.message_from}\pysiglinewithargsret{\bfcode{message\_from}}{\emph{node}, \emph{message}}{}
Přijmutí zprávy ze sousedního vrcholu.

Zpráva je uložena a použita při dalších výpočtech.
\begin{quote}\begin{description}
\item[{Parametry}] \leavevmode\begin{itemize}
\item {} 
\textbf{node} -- sousední vrchol.

\item {} 
\textbf{message} -- zpráva ze sousedního vrcholu.

\end{itemize}

\end{description}\end{quote}

\end{fulllineitems}

\index{message\_to() (metoda alex.infer.factor.alex.infer.node.DirichletFactorNode)}

\begin{fulllineitems}
\phantomsection\label{alex.infer:alex.infer.factor.alex.infer.node.DirichletFactorNode.message_to}\pysiglinewithargsret{\bfcode{message\_to}}{\emph{node}}{}
Poslání zprávy do sousedního vrcholu.
\begin{quote}\begin{description}
\item[{Parametr}] \leavevmode
\textbf{node} -- sousední vrchol, kterému má být zpráva zaslána.

\end{description}\end{quote}

\end{fulllineitems}

\index{update() (metoda alex.infer.factor.alex.infer.node.DirichletFactorNode)}

\begin{fulllineitems}
\phantomsection\label{alex.infer:alex.infer.factor.alex.infer.node.DirichletFactorNode.update}\pysiglinewithargsret{\bfcode{update}}{}{}
Aktualizace interní reprezentace.

Příchozí zprávy jsou pouze uloženy a je potřeba aktualizovat interní
reprezentaci jejich součinu před počítáním odchozích zpráv.

\end{fulllineitems}


\end{fulllineitems}



\subsection{Modul \texttt{lbp}}
\label{alex.infer:modul-lbp}\index{alex.infer.lbp.LBPError}

\begin{fulllineitems}
\phantomsection\label{alex.infer:alex.infer.factor.alex.infer.lbp.LBPError}\pysigline{\strong{exception }\code{alex.infer.lbp.}\bfcode{LBPError}}
Výjimka pro chyby při inferenci.

\end{fulllineitems}

\index{alex.infer.lbp.LBP() (třída v alex.infer.factor)}
\hskip\linewidth

\begin{fulllineitems}
\phantomsection\label{alex.infer:alex.infer.factor.alex.infer.lbp.LBP}\pysigline{\strong{class }\code{alex.infer.lbp.}\bfcode{LBP}}
Inference v grafickém modelu pomocí Loopy Belief Propagation.
\index{\_\_init\_\_() (metoda alex.infer.factor.alex.infer.lbp.LBP)}

\begin{fulllineitems}
\phantomsection\label{alex.infer:alex.infer.factor.alex.infer.lbp.LBP.__init__}\pysiglinewithargsret{\bfcode{\_\_init\_\_}}{\emph{strategy='sequential'}}{}
Vytvoření nové instance LBP algoritmu.

Je možné nastavit strategii výběru vrcholů:
\begin{itemize}
\item {} 
\emph{sequential} -- vrcholy jsou vybírány v takovém pořadí, v jakém byly vloženy,

\item {} 
\emph{tree} -- grafický model je ve formě stromu,

\item {} 
\emph{layers} -- inference na Dynamické Bayesovské síti.

\end{itemize}
\begin{quote}\begin{description}
\item[{Parametr}] \leavevmode
\textbf{strategy} -- strategie, která má být použita pro výběr vrcholů.

\end{description}\end{quote}

\end{fulllineitems}

\index{add\_nodes() (metoda alex.infer.factor.alex.infer.lbp.LBP)}

\begin{fulllineitems}
\phantomsection\label{alex.infer:alex.infer.factor.alex.infer.lbp.LBP.add_nodes}\pysiglinewithargsret{\bfcode{add\_nodes}}{\emph{nodes}}{}
Přidání vrcholů k inferenci.
\begin{quote}\begin{description}
\item[{Parametr}] \leavevmode
\textbf{nodes} -- seznam vrcholů.

\end{description}\end{quote}

\end{fulllineitems}

\index{add\_layer() (metoda alex.infer.factor.alex.infer.lbp.LBP)}

\begin{fulllineitems}
\phantomsection\label{alex.infer:alex.infer.factor.alex.infer.lbp.LBP.add_layer}\pysiglinewithargsret{\bfcode{add\_layer}}{\emph{layer}}{}
Přidání vrstvy na konec grafického modelu.
\begin{quote}\begin{description}
\item[{Parametr}] \leavevmode
\textbf{layer} -- seznam vrcholů tvořících jednu vrstvu.

\end{description}\end{quote}

\end{fulllineitems}

\index{add\_layers() (metoda alex.infer.factor.alex.infer.lbp.LBP)}

\begin{fulllineitems}
\phantomsection\label{alex.infer:alex.infer.factor.alex.infer.lbp.LBP.add_layers}\pysiglinewithargsret{\bfcode{add\_layers}}{\emph{layers}}{}
Přidání vrstev na konec grafického modelu.
\begin{quote}\begin{description}
\item[{Parametr}] \leavevmode
\textbf{layers} -- seznam seznamů s vrcholy.

\end{description}\end{quote}

\end{fulllineitems}

\index{clear\_nodes() (metoda alex.infer.factor.alex.infer.lbp.LBP)}

\begin{fulllineitems}
\phantomsection\label{alex.infer:alex.infer.factor.alex.infer.lbp.LBP.clear_nodes}\pysiglinewithargsret{\bfcode{clear\_nodes}}{}{}
Smazání všech vrcholů.

\end{fulllineitems}

\index{clear\_layers() (metoda alex.infer.factor.alex.infer.lbp.LBP)}

\begin{fulllineitems}
\phantomsection\label{alex.infer:alex.infer.factor.alex.infer.lbp.LBP.clear_layers}\pysiglinewithargsret{\bfcode{clear\_layers}}{}{}
Smazání všech vrstev.

\end{fulllineitems}

\index{init\_messages() (metoda alex.infer.factor.alex.infer.lbp.LBP)}

\begin{fulllineitems}
\phantomsection\label{alex.infer:alex.infer.factor.alex.infer.lbp.LBP.init_messages}\pysiglinewithargsret{\bfcode{init\_messages}}{}{}
Inicializace zpráv pro všechny vrcholy.

\end{fulllineitems}

\index{run() (metoda alex.infer.factor.alex.infer.lbp.LBP)}

\begin{fulllineitems}
\phantomsection\label{alex.infer:alex.infer.factor.alex.infer.lbp.LBP.run}\pysiglinewithargsret{\bfcode{run}}{\emph{n\_iterations=1, from\_layer{]}}}{}
Spustit inferenci v grafickém modelu.

Je možné nastavit počet iterací, po kterém se má inference zastavit.
U stromu není třeba zadávat počet iterací, bude provedena pouze jedna.

U vrstev je možné zadat od které vrstvy inferenci provádět.
Popřípadě budou zprávy propagovány od začátku.
\begin{quote}\begin{description}
\item[{Parametry}] \leavevmode\begin{itemize}
\item {} 
\textbf{n\_iterations} -- počet iterací.

\item {} 
\textbf{from\_layer} -- vrstva, od které se mají zprávy propagovat.

\end{itemize}

\end{description}\end{quote}

\end{fulllineitems}


\end{fulllineitems}


\singlespacing

%%% Seznam použité literatury
\addcontentsline{toc}{chapter}{Seznam použité literatury}
\bibliography{citace}{}
\bibliographystyle{czech}

%%% Tabulky v diplomové práci, existují-li.
%\chapwithtoc{Seznam tabulek}

%%% Použité zkratky v diplomové práci, existují-li, včetně jejich vysvětlení.
%\chapwithtoc{Seznam použitých zkratek}

%%% Přílohy k diplomové práci, existují-li (různé dodatky jako výpisy programů,
%%% diagramy apod.). Každá příloha musí být alespoň jednou odkazována z vlastního
%%% textu práce. Přílohy se číslují.
\chapwithtoc{Přílohy na CD}

Přiložené CD obsahuje:

\newcommand{\priloha}[2]{\item \texttt{#1} -- #2}
\begin{itemize}
\priloha{thesis.pdf}{text diplomové práce,}
\priloha{thesis/}{zdrojové kódy diplomové práce,}
\priloha{src/}{zdrojové kódy implementované knihovny.
	\begin{itemize}
	\priloha{alex/infer/}{balíček \texttt{alex.infer} obsahující všechny moduly,}
	\priloha{doc/}{adresář se zdrojovými kódy pro vytvoření dokumentace,}
	\priloha{LICENSE.txt}{licence zdrojových kódů,}
	\priloha{setup.py}{skript pro instalaci knihovny.}
	\end{itemize}
}
\priloha{docs/}{dokumentace knihovny ve formátu HTML,}
\priloha{dstc/}{adresář se soubory pro Dialog State Tracking Challenge. Blíže popsány jsou v sekci~\ref{sec:dstcfiles}.
}
\end{itemize}

\openright
\end{document}
